\Needspace{8\baselineskip}\originalchapter{2006 作业八: 冷波传播图}
题源: 21 Nov 2006, due 30 Nov. 图形由实际冷流体方程采样, 不作稳定性证明.
\begin{exercise}\label{ex:06-8-1}
原题1. 用绘图软件准确画传播/截止区域, 坐标log10 nref\ensuremath{{}^{2}}对log10(\ensuremath{\omega}/\ensuremath{\Omega}ci), 两套C-Mod参数: (a)ne=$2\times10^{20}$m$^{-3}$,B=5.5T; (b)ne=$5\times10^{20}$m$^{-3}$,B=5.5T; 氘离子. 横范围\ensuremath{-}2..4, 纵\ensuremath{-}2..3.
\end{exercise}
\begin{solution}
原题未指定传播角, 折射率并非单值. 本解分别画平行R/L及垂直O/X四支, 假定完全电离D、ni=ne, mi\ensuremath{\simeq}2mp. 按章\ref{ch:cold}的带符号\ensuremath{\Omega}求R/L/P/S, O=P,X=RL/S. 只画正n\ensuremath{{}^{2}}且在题给纵范围内的段, 负段为空而不取abs, 共振和截止不连线. 普通支截止\ensuremath{\omega}\ensuremath{{}^{2}}=\ensuremath{\omega}pe\ensuremath{{}^{2}}+\ensuremath{\omega}pi\ensuremath{{}^{2}}; R/L零点和S零点分别定X截止/共振. 若要斜传播, 需另给角后解介电行列式.
\begin{center}\begin{tikzpicture}\begin{axis}[width=.82\textwidth,height=5cm,xlabel={$\log_{10}(\omega/\Omega_{ci})$},ylabel={$\log_{10}n_{\rm ref}^2$},xmin=-2,xmax=4,ymin=-2,ymax=3,unbounded coords=jump,title={$n_e=2\times10^{20}\,\mathrm m^{-3}$},legend style={font=\small,at={(1.02,1)},anchor=north west}]
\addplot[thick,structurecolor] coordinates {(-2.00000,nan) (-1.99750,nan) (-1.99500,nan) (-1.99250,nan) (-1.99000,nan) (-1.98750,nan) (-1.98500,nan) (-1.98250,nan) (-1.98000,nan) (-1.97750,nan) (-1.97500,nan) (-1.97250,nan) (-1.97000,nan) (-1.96750,nan) (-1.96500,nan) (-1.96250,nan) (-1.96000,nan) (-1.95750,nan) (-1.95500,nan) (-1.95250,nan) (-1.95000,nan) (-1.94750,nan) (-1.94500,nan) (-1.94250,nan) (-1.94000,nan) (-1.93750,nan) (-1.93500,nan) (-1.93250,nan) (-1.93000,nan) (-1.92750,nan) (-1.92500,nan) (-1.92250,nan) (-1.92000,nan) (-1.91750,nan) (-1.91500,nan) (-1.91250,nan) (-1.91000,nan) (-1.90750,nan) (-1.90500,nan) (-1.90250,nan) (-1.90000,nan) (-1.89750,nan) (-1.89500,nan) (-1.89250,nan) (-1.89000,nan) (-1.88750,nan) (-1.88500,nan) (-1.88250,nan) (-1.88000,nan) (-1.87750,nan) (-1.87500,nan) (-1.87250,nan) (-1.87000,nan) (-1.86750,nan) (-1.86500,nan) (-1.86250,nan) (-1.86000,nan) (-1.85750,nan) (-1.85500,nan) (-1.85250,nan) (-1.85000,nan) (-1.84750,nan) (-1.84500,nan) (-1.84250,nan) (-1.84000,nan) (-1.83750,nan) (-1.83500,nan) (-1.83250,nan) (-1.83000,nan) (-1.82750,nan) (-1.82500,nan) (-1.82250,nan) (-1.82000,nan) (-1.81750,nan) (-1.81500,nan) (-1.81250,nan) (-1.81000,nan) (-1.80750,nan) (-1.80500,nan) (-1.80250,nan) (-1.80000,nan) (-1.79750,nan) (-1.79500,nan) (-1.79250,nan) (-1.79000,nan) (-1.78750,nan) (-1.78500,nan) (-1.78250,nan) (-1.78000,nan) (-1.77750,nan) (-1.77500,nan) (-1.77250,nan) (-1.77000,nan) (-1.76750,nan) (-1.76500,nan) (-1.76250,nan) (-1.76000,nan) (-1.75750,nan) (-1.75500,nan) (-1.75250,nan) (-1.75000,nan) (-1.74750,nan) (-1.74500,nan) (-1.74250,nan) (-1.74000,nan) (-1.73750,nan) (-1.73500,nan) (-1.73250,nan) (-1.73000,nan) (-1.72750,nan) (-1.72500,nan) (-1.72250,nan) (-1.72000,nan) (-1.71750,nan) (-1.71500,nan) (-1.71250,nan) (-1.71000,nan) (-1.70750,nan) (-1.70500,nan) (-1.70250,nan) (-1.70000,nan) (-1.69750,nan) (-1.69500,nan) (-1.69250,nan) (-1.69000,nan) (-1.68750,nan) (-1.68500,nan) (-1.68250,nan) (-1.68000,nan) (-1.67750,nan) (-1.67500,nan) (-1.67250,nan) (-1.67000,nan) (-1.66750,nan) (-1.66500,nan) (-1.66250,nan) (-1.66000,nan) (-1.65750,nan) (-1.65500,nan) (-1.65250,nan) (-1.65000,nan) (-1.64750,nan) (-1.64500,nan) (-1.64250,nan) (-1.64000,nan) (-1.63750,nan) (-1.63500,nan) (-1.63250,nan) (-1.63000,nan) (-1.62750,nan) (-1.62500,nan) (-1.62250,nan) (-1.62000,nan) (-1.61750,nan) (-1.61500,nan) (-1.61250,nan) (-1.61000,nan) (-1.60750,nan) (-1.60500,nan) (-1.60250,nan) (-1.60000,nan) (-1.59750,nan) (-1.59500,nan) (-1.59250,nan) (-1.59000,nan) (-1.58750,nan) (-1.58500,nan) (-1.58250,nan) (-1.58000,nan) (-1.57750,nan) (-1.57500,nan) (-1.57250,nan) (-1.57000,nan) (-1.56750,nan) (-1.56500,nan) (-1.56250,nan) (-1.56000,nan) (-1.55750,nan) (-1.55500,nan) (-1.55250,nan) (-1.55000,nan) (-1.54750,nan) (-1.54500,nan) (-1.54250,nan) (-1.54000,nan) (-1.53750,nan) (-1.53500,nan) (-1.53250,nan) (-1.53000,nan) (-1.52750,nan) (-1.52500,nan) (-1.52250,nan) (-1.52000,nan) (-1.51750,nan) (-1.51500,nan) (-1.51250,nan) (-1.51000,nan) (-1.50750,nan) (-1.50500,nan) (-1.50250,nan) (-1.50000,nan) (-1.49750,nan) (-1.49500,nan) (-1.49250,nan) (-1.49000,nan) (-1.48750,nan) (-1.48500,nan) (-1.48250,nan) (-1.48000,nan) (-1.47750,nan) (-1.47500,nan) (-1.47250,nan) (-1.47000,nan) (-1.46750,nan) (-1.46500,nan) (-1.46250,nan) (-1.46000,nan) (-1.45750,nan) (-1.45500,nan) (-1.45250,nan) (-1.45000,nan) (-1.44750,nan) (-1.44500,nan) (-1.44250,nan) (-1.44000,nan) (-1.43750,nan) (-1.43500,nan) (-1.43250,nan) (-1.43000,nan) (-1.42750,nan) (-1.42500,nan) (-1.42250,nan) (-1.42000,nan) (-1.41750,nan) (-1.41500,nan) (-1.41250,nan) (-1.41000,nan) (-1.40750,nan) (-1.40500,nan) (-1.40250,nan) (-1.40000,nan) (-1.39750,nan) (-1.39500,nan) (-1.39250,nan) (-1.39000,nan) (-1.38750,nan) (-1.38500,nan) (-1.38250,nan) (-1.38000,nan) (-1.37750,nan) (-1.37500,nan) (-1.37250,nan) (-1.37000,nan) (-1.36750,nan) (-1.36500,nan) (-1.36250,nan) (-1.36000,nan) (-1.35750,nan) (-1.35500,nan) (-1.35250,nan) (-1.35000,nan) (-1.34750,nan) (-1.34500,nan) (-1.34250,nan) (-1.34000,nan) (-1.33750,nan) (-1.33500,nan) (-1.33250,nan) (-1.33000,nan) (-1.32750,nan) (-1.32500,nan) (-1.32250,nan) (-1.32000,nan) (-1.31750,nan) (-1.31500,nan) (-1.31250,nan) (-1.31000,nan) (-1.30750,nan) (-1.30500,nan) (-1.30250,nan) (-1.30000,nan) (-1.29750,nan) (-1.29500,nan) (-1.29250,nan) (-1.29000,nan) (-1.28750,nan) (-1.28500,nan) (-1.28250,nan) (-1.28000,nan) (-1.27750,nan) (-1.27500,nan) (-1.27250,nan) (-1.27000,nan) (-1.26750,nan) (-1.26500,nan) (-1.26250,nan) (-1.26000,nan) (-1.25750,nan) (-1.25500,nan) (-1.25250,nan) (-1.25000,nan) (-1.24750,nan) (-1.24500,nan) (-1.24250,nan) (-1.24000,nan) (-1.23750,nan) (-1.23500,nan) (-1.23250,nan) (-1.23000,nan) (-1.22750,nan) (-1.22500,nan) (-1.22250,nan) (-1.22000,nan) (-1.21750,nan) (-1.21500,nan) (-1.21250,nan) (-1.21000,nan) (-1.20750,nan) (-1.20500,nan) (-1.20250,nan) (-1.20000,nan) (-1.19750,nan) (-1.19500,nan) (-1.19250,nan) (-1.19000,nan) (-1.18750,nan) (-1.18500,nan) (-1.18250,nan) (-1.18000,nan) (-1.17750,nan) (-1.17500,nan) (-1.17250,nan) (-1.17000,nan) (-1.16750,nan) (-1.16500,nan) (-1.16250,nan) (-1.16000,nan) (-1.15750,nan) (-1.15500,nan) (-1.15250,nan) (-1.15000,nan) (-1.14750,nan) (-1.14500,nan) (-1.14250,nan) (-1.14000,nan) (-1.13750,nan) (-1.13500,nan) (-1.13250,nan) (-1.13000,nan) (-1.12750,nan) (-1.12500,nan) (-1.12250,nan) (-1.12000,nan) (-1.11750,nan) (-1.11500,nan) (-1.11250,nan) (-1.11000,nan) (-1.10750,nan) (-1.10500,nan) (-1.10250,nan) (-1.10000,nan) (-1.09750,nan) (-1.09500,nan) (-1.09250,nan) (-1.09000,nan) (-1.08750,nan) (-1.08500,nan) (-1.08250,nan) (-1.08000,nan) (-1.07750,nan) (-1.07500,nan) (-1.07250,nan) (-1.07000,nan) (-1.06750,nan) (-1.06500,nan) (-1.06250,nan) (-1.06000,nan) (-1.05750,nan) (-1.05500,nan) (-1.05250,nan) (-1.05000,nan) (-1.04750,nan) (-1.04500,nan) (-1.04250,nan) (-1.04000,nan) (-1.03750,nan) (-1.03500,nan) (-1.03250,nan) (-1.03000,nan) (-1.02750,nan) (-1.02500,nan) (-1.02250,nan) (-1.02000,nan) (-1.01750,nan) (-1.01500,nan) (-1.01250,nan) (-1.01000,nan) (-1.00750,nan) (-1.00500,nan) (-1.00250,nan) (-1.00000,nan) (-0.99750,nan) (-0.99500,nan) (-0.99250,nan) (-0.99000,nan) (-0.98750,nan) (-0.98500,nan) (-0.98250,nan) (-0.98000,nan) (-0.97750,nan) (-0.97500,nan) (-0.97250,nan) (-0.97000,nan) (-0.96750,nan) (-0.96500,nan) (-0.96250,nan) (-0.96000,nan) (-0.95750,nan) (-0.95500,nan) (-0.95250,nan) (-0.95000,nan) (-0.94750,nan) (-0.94500,nan) (-0.94250,nan) (-0.94000,nan) (-0.93750,nan) (-0.93500,nan) (-0.93250,nan) (-0.93000,nan) (-0.92750,nan) (-0.92500,nan) (-0.92250,nan) (-0.92000,nan) (-0.91750,nan) (-0.91500,nan) (-0.91250,nan) (-0.91000,nan) (-0.90750,nan) (-0.90500,nan) (-0.90250,nan) (-0.90000,nan) (-0.89750,nan) (-0.89500,nan) (-0.89250,nan) (-0.89000,nan) (-0.88750,nan) (-0.88500,nan) (-0.88250,nan) (-0.88000,nan) (-0.87750,nan) (-0.87500,nan) (-0.87250,nan) (-0.87000,nan) (-0.86750,nan) (-0.86500,nan) (-0.86250,nan) (-0.86000,nan) (-0.85750,nan) (-0.85500,nan) (-0.85250,nan) (-0.85000,nan) (-0.84750,nan) (-0.84500,nan) (-0.84250,nan) (-0.84000,nan) (-0.83750,nan) (-0.83500,nan) (-0.83250,nan) (-0.83000,nan) (-0.82750,nan) (-0.82500,nan) (-0.82250,nan) (-0.82000,nan) (-0.81750,nan) (-0.81500,nan) (-0.81250,nan) (-0.81000,nan) (-0.80750,nan) (-0.80500,nan) (-0.80250,nan) (-0.80000,nan) (-0.79750,nan) (-0.79500,nan) (-0.79250,nan) (-0.79000,nan) (-0.78750,nan) (-0.78500,nan) (-0.78250,nan) (-0.78000,nan) (-0.77750,nan) (-0.77500,nan) (-0.77250,nan) (-0.77000,nan) (-0.76750,nan) (-0.76500,nan) (-0.76250,nan) (-0.76000,nan) (-0.75750,nan) (-0.75500,nan) (-0.75250,nan) (-0.75000,nan) (-0.74750,nan) (-0.74500,nan) (-0.74250,nan) (-0.74000,nan) (-0.73750,nan) (-0.73500,nan) (-0.73250,nan) (-0.73000,nan) (-0.72750,nan) (-0.72500,nan) (-0.72250,nan) (-0.72000,nan) (-0.71750,nan) (-0.71500,nan) (-0.71250,nan) (-0.71000,nan) (-0.70750,nan) (-0.70500,nan) (-0.70250,nan) (-0.70000,nan) (-0.69750,nan) (-0.69500,nan) (-0.69250,nan) (-0.69000,nan) (-0.68750,nan) (-0.68500,nan) (-0.68250,nan) (-0.68000,nan) (-0.67750,nan) (-0.67500,nan) (-0.67250,nan) (-0.67000,nan) (-0.66750,nan) (-0.66500,nan) (-0.66250,nan) (-0.66000,nan) (-0.65750,nan) (-0.65500,nan) (-0.65250,nan) (-0.65000,nan) (-0.64750,nan) (-0.64500,nan) (-0.64250,nan) (-0.64000,nan) (-0.63750,nan) (-0.63500,nan) (-0.63250,nan) (-0.63000,nan) (-0.62750,nan) (-0.62500,nan) (-0.62250,nan) (-0.62000,nan) (-0.61750,nan) (-0.61500,nan) (-0.61250,nan) (-0.61000,nan) (-0.60750,nan) (-0.60500,nan) (-0.60250,nan) (-0.60000,nan) (-0.59750,nan) (-0.59500,nan) (-0.59250,nan) (-0.59000,nan) (-0.58750,nan) (-0.58500,nan) (-0.58250,nan) (-0.58000,nan) (-0.57750,nan) (-0.57500,nan) (-0.57250,nan) (-0.57000,nan) (-0.56750,nan) (-0.56500,nan) (-0.56250,nan) (-0.56000,nan) (-0.55750,nan) (-0.55500,nan) (-0.55250,nan) (-0.55000,nan) (-0.54750,nan) (-0.54500,nan) (-0.54250,nan) (-0.54000,nan) (-0.53750,nan) (-0.53500,nan) (-0.53250,nan) (-0.53000,nan) (-0.52750,nan) (-0.52500,nan) (-0.52250,nan) (-0.52000,nan) (-0.51750,nan) (-0.51500,nan) (-0.51250,nan) (-0.51000,nan) (-0.50750,nan) (-0.50500,nan) (-0.50250,nan) (-0.50000,nan) (-0.49750,nan) (-0.49500,nan) (-0.49250,nan) (-0.49000,nan) (-0.48750,nan) (-0.48500,nan) (-0.48250,nan) (-0.48000,nan) (-0.47750,nan) (-0.47500,nan) (-0.47250,nan) (-0.47000,nan) (-0.46750,nan) (-0.46500,nan) (-0.46250,nan) (-0.46000,nan) (-0.45750,nan) (-0.45500,nan) (-0.45250,nan) (-0.45000,nan) (-0.44750,nan) (-0.44500,nan) (-0.44250,nan) (-0.44000,nan) (-0.43750,nan) (-0.43500,nan) (-0.43250,nan) (-0.43000,nan) (-0.42750,nan) (-0.42500,nan) (-0.42250,nan) (-0.42000,nan) (-0.41750,nan) (-0.41500,nan) (-0.41250,nan) (-0.41000,nan) (-0.40750,nan) (-0.40500,nan) (-0.40250,nan) (-0.40000,nan) (-0.39750,nan) (-0.39500,nan) (-0.39250,nan) (-0.39000,nan) (-0.38750,nan) (-0.38500,nan) (-0.38250,nan) (-0.38000,nan) (-0.37750,nan) (-0.37500,nan) (-0.37250,nan) (-0.37000,nan) (-0.36750,nan) (-0.36500,nan) (-0.36250,nan) (-0.36000,nan) (-0.35750,nan) (-0.35500,nan) (-0.35250,nan) (-0.35000,nan) (-0.34750,nan) (-0.34500,nan) (-0.34250,nan) (-0.34000,nan) (-0.33750,nan) (-0.33500,nan) (-0.33250,nan) (-0.33000,nan) (-0.32750,nan) (-0.32500,nan) (-0.32250,nan) (-0.32000,nan) (-0.31750,nan) (-0.31500,nan) (-0.31250,nan) (-0.31000,nan) (-0.30750,nan) (-0.30500,nan) (-0.30250,nan) (-0.30000,nan) (-0.29750,nan) (-0.29500,nan) (-0.29250,nan) (-0.29000,nan) (-0.28750,nan) (-0.28500,nan) (-0.28250,nan) (-0.28000,nan) (-0.27750,nan) (-0.27500,nan) (-0.27250,nan) (-0.27000,nan) (-0.26750,nan) (-0.26500,nan) (-0.26250,nan) (-0.26000,nan) (-0.25750,nan) (-0.25500,nan) (-0.25250,nan) (-0.25000,nan) (-0.24750,nan) (-0.24500,nan) (-0.24250,nan) (-0.24000,nan) (-0.23750,nan) (-0.23500,nan) (-0.23250,nan) (-0.23000,nan) (-0.22750,nan) (-0.22500,nan) (-0.22250,nan) (-0.22000,nan) (-0.21750,nan) (-0.21500,nan) (-0.21250,nan) (-0.21000,nan) (-0.20750,nan) (-0.20500,nan) (-0.20250,nan) (-0.20000,nan) (-0.19750,nan) (-0.19500,nan) (-0.19250,nan) (-0.19000,nan) (-0.18750,nan) (-0.18500,nan) (-0.18250,nan) (-0.18000,nan) (-0.17750,nan) (-0.17500,nan) (-0.17250,nan) (-0.17000,nan) (-0.16750,nan) (-0.16500,nan) (-0.16250,nan) (-0.16000,nan) (-0.15750,nan) (-0.15500,nan) (-0.15250,nan) (-0.15000,nan) (-0.14750,nan) (-0.14500,nan) (-0.14250,nan) (-0.14000,nan) (-0.13750,nan) (-0.13500,nan) (-0.13250,nan) (-0.13000,nan) (-0.12750,nan) (-0.12500,nan) (-0.12250,nan) (-0.12000,nan) (-0.11750,nan) (-0.11500,nan) (-0.11250,nan) (-0.11000,nan) (-0.10750,nan) (-0.10500,nan) (-0.10250,nan) (-0.10000,nan) (-0.09750,nan) (-0.09500,nan) (-0.09250,nan) (-0.09000,nan) (-0.08750,nan) (-0.08500,nan) (-0.08250,nan) (-0.08000,nan) (-0.07750,nan) (-0.07500,nan) (-0.07250,nan) (-0.07000,nan) (-0.06750,nan) (-0.06500,nan) (-0.06250,nan) (-0.06000,nan) (-0.05750,nan) (-0.05500,nan) (-0.05250,nan) (-0.05000,nan) (-0.04750,nan) (-0.04500,nan) (-0.04250,nan) (-0.04000,nan) (-0.03750,nan) (-0.03500,nan) (-0.03250,nan) (-0.03000,nan) (-0.02750,nan) (-0.02500,nan) (-0.02250,nan) (-0.02000,nan) (-0.01750,nan) (-0.01500,nan) (-0.01250,nan) (-0.01000,nan) (-0.00750,nan) (-0.00500,nan) (-0.00250,nan) (0.00000,nan) (0.00250,nan) (0.00500,nan) (0.00750,nan) (0.01000,nan) (0.01250,nan) (0.01500,nan) (0.01750,nan) (0.02000,nan) (0.02250,nan) (0.02500,nan) (0.02750,nan) (0.03000,nan) (0.03250,nan) (0.03500,nan) (0.03750,nan) (0.04000,nan) (0.04250,nan) (0.04500,nan) (0.04750,nan) (0.05000,nan) (0.05250,nan) (0.05500,nan) (0.05750,nan) (0.06000,nan) (0.06250,nan) (0.06500,nan) (0.06750,nan) (0.07000,nan) (0.07250,nan) (0.07500,nan) (0.07750,nan) (0.08000,nan) (0.08250,nan) (0.08500,nan) (0.08750,nan) (0.09000,nan) (0.09250,nan) (0.09500,nan) (0.09750,nan) (0.10000,nan) (0.10250,nan) (0.10500,nan) (0.10750,nan) (0.11000,nan) (0.11250,nan) (0.11500,nan) (0.11750,nan) (0.12000,nan) (0.12250,nan) (0.12500,nan) (0.12750,nan) (0.13000,nan) (0.13250,nan) (0.13500,nan) (0.13750,nan) (0.14000,nan) (0.14250,nan) (0.14500,nan) (0.14750,nan) (0.15000,nan) (0.15250,nan) (0.15500,nan) (0.15750,nan) (0.16000,nan) (0.16250,nan) (0.16500,nan) (0.16750,nan) (0.17000,nan) (0.17250,nan) (0.17500,nan) (0.17750,2.9995291) (0.18000,2.9980280) (0.18250,2.9965234) (0.18500,2.9950154) (0.18750,2.9935040) (0.19000,2.9919891) (0.19250,2.9904709) (0.19500,2.9889492) (0.19750,2.9874241) (0.20000,2.9858956) (0.20250,2.9843638) (0.20500,2.9828285) (0.20750,2.9812898) (0.21000,2.9797478) (0.21250,2.9782024) (0.21500,2.9766536) (0.21750,2.9751014) (0.22000,2.9735459) (0.22250,2.9719870) (0.22500,2.9704247) (0.22750,2.9688591) (0.23000,2.9672902) (0.23250,2.9657179) (0.23500,2.9641422) (0.23750,2.9625633) (0.24000,2.9609810) (0.24250,2.9593954) (0.24500,2.9578065) (0.24750,2.9562143) (0.25000,2.9546187) (0.25250,2.9530199) (0.25500,2.9514178) (0.25750,2.9498124) (0.26000,2.9482037) (0.26250,2.9465918) (0.26500,2.9449765) (0.26750,2.9433580) (0.27000,2.9417363) (0.27250,2.9401113) (0.27500,2.9384830) (0.27750,2.9368515) (0.28000,2.9352168) (0.28250,2.9335788) (0.28500,2.9319376) (0.28750,2.9302932) (0.29000,2.9286456) (0.29250,2.9269948) (0.29500,2.9253407) (0.29750,2.9236835) (0.30000,2.9220231) (0.30250,2.9203595) (0.30500,2.9186927) (0.30750,2.9170228) (0.31000,2.9153497) (0.31250,2.9136735) (0.31500,2.9119941) (0.31750,2.9103115) (0.32000,2.9086259) (0.32250,2.9069370) (0.32500,2.9052451) (0.32750,2.9035501) (0.33000,2.9018519) (0.33250,2.9001506) (0.33500,2.8984463) (0.33750,2.8967388) (0.34000,2.8950283) (0.34250,2.8933147) (0.34500,2.8915980) (0.34750,2.8898782) (0.35000,2.8881554) (0.35250,2.8864296) (0.35500,2.8847007) (0.35750,2.8829687) (0.36000,2.8812338) (0.36250,2.8794958) (0.36500,2.8777548) (0.36750,2.8760108) (0.37000,2.8742638) (0.37250,2.8725138) (0.37500,2.8707608) (0.37750,2.8690048) (0.38000,2.8672459) (0.38250,2.8654839) (0.38500,2.8637191) (0.38750,2.8619513) (0.39000,2.8601805) (0.39250,2.8584068) (0.39500,2.8566302) (0.39750,2.8548506) (0.40000,2.8530682) (0.40250,2.8512828) (0.40500,2.8494945) (0.40750,2.8477034) (0.41000,2.8459093) (0.41250,2.8441124) (0.41500,2.8423126) (0.41750,2.8405100) (0.42000,2.8387045) (0.42250,2.8368961) (0.42500,2.8350849) (0.42750,2.8332709) (0.43000,2.8314541) (0.43250,2.8296344) (0.43500,2.8278119) (0.43750,2.8259867) (0.44000,2.8241586) (0.44250,2.8223277) (0.44500,2.8204941) (0.44750,2.8186577) (0.45000,2.8168185) (0.45250,2.8149766) (0.45500,2.8131320) (0.45750,2.8112846) (0.46000,2.8094344) (0.46250,2.8075816) (0.46500,2.8057260) (0.46750,2.8038677) (0.47000,2.8020067) (0.47250,2.8001430) (0.47500,2.7982767) (0.47750,2.7964077) (0.48000,2.7945360) (0.48250,2.7926616) (0.48500,2.7907846) (0.48750,2.7889049) (0.49000,2.7870226) (0.49250,2.7851377) (0.49500,2.7832502) (0.49750,2.7813600) (0.50000,2.7794672) (0.50250,2.7775719) (0.50500,2.7756739) (0.50750,2.7737734) (0.51000,2.7718703) (0.51250,2.7699647) (0.51500,2.7680564) (0.51750,2.7661457) (0.52000,2.7642324) (0.52250,2.7623165) (0.52500,2.7603982) (0.52750,2.7584773) (0.53000,2.7565539) (0.53250,2.7546280) (0.53500,2.7526996) (0.53750,2.7507688) (0.54000,2.7488354) (0.54250,2.7468996) (0.54500,2.7449614) (0.54750,2.7430207) (0.55000,2.7410775) (0.55250,2.7391319) (0.55500,2.7371839) (0.55750,2.7352335) (0.56000,2.7332807) (0.56250,2.7313254) (0.56500,2.7293678) (0.56750,2.7274078) (0.57000,2.7254454) (0.57250,2.7234806) (0.57500,2.7215135) (0.57750,2.7195440) (0.58000,2.7175722) (0.58250,2.7155980) (0.58500,2.7136215) (0.58750,2.7116427) (0.59000,2.7096616) (0.59250,2.7076782) (0.59500,2.7056925) (0.59750,2.7037044) (0.60000,2.7017142) (0.60250,2.6997216) (0.60500,2.6977268) (0.60750,2.6957297) (0.61000,2.6937304) (0.61250,2.6917288) (0.61500,2.6897251) (0.61750,2.6877190) (0.62000,2.6857108) (0.62250,2.6837004) (0.62500,2.6816878) (0.62750,2.6796730) (0.63000,2.6776560) (0.63250,2.6756368) (0.63500,2.6736155) (0.63750,2.6715920) (0.64000,2.6695663) (0.64250,2.6675386) (0.64500,2.6655087) (0.64750,2.6634766) (0.65000,2.6614425) (0.65250,2.6594062) (0.65500,2.6573679) (0.65750,2.6553274) (0.66000,2.6532849) (0.66250,2.6512403) (0.66500,2.6491936) (0.66750,2.6471449) (0.67000,2.6450941) (0.67250,2.6430412) (0.67500,2.6409864) (0.67750,2.6389295) (0.68000,2.6368706) (0.68250,2.6348096) (0.68500,2.6327467) (0.68750,2.6306818) (0.69000,2.6286149) (0.69250,2.6265460) (0.69500,2.6244751) (0.69750,2.6224023) (0.70000,2.6203275) (0.70250,2.6182508) (0.70500,2.6161722) (0.70750,2.6140916) (0.71000,2.6120090) (0.71250,2.6099246) (0.71500,2.6078383) (0.71750,2.6057500) (0.72000,2.6036599) (0.72250,2.6015679) (0.72500,2.5994740) (0.72750,2.5973782) (0.73000,2.5952806) (0.73250,2.5931811) (0.73500,2.5910798) (0.73750,2.5889766) (0.74000,2.5868717) (0.74250,2.5847648) (0.74500,2.5826562) (0.74750,2.5805458) (0.75000,2.5784336) (0.75250,2.5763195) (0.75500,2.5742037) (0.75750,2.5720862) (0.76000,2.5699668) (0.76250,2.5678457) (0.76500,2.5657229) (0.76750,2.5635983) (0.77000,2.5614719) (0.77250,2.5593439) (0.77500,2.5572141) (0.77750,2.5550826) (0.78000,2.5529494) (0.78250,2.5508145) (0.78500,2.5486779) (0.78750,2.5465396) (0.79000,2.5443996) (0.79250,2.5422580) (0.79500,2.5401147) (0.79750,2.5379698) (0.80000,2.5358232) (0.80250,2.5336749) (0.80500,2.5315251) (0.80750,2.5293736) (0.81000,2.5272205) (0.81250,2.5250658) (0.81500,2.5229095) (0.81750,2.5207516) (0.82000,2.5185921) (0.82250,2.5164310) (0.82500,2.5142684) (0.82750,2.5121042) (0.83000,2.5099384) (0.83250,2.5077711) (0.83500,2.5056022) (0.83750,2.5034318) (0.84000,2.5012599) (0.84250,2.4990865) (0.84500,2.4969115) (0.84750,2.4947351) (0.85000,2.4925571) (0.85250,2.4903776) (0.85500,2.4881967) (0.85750,2.4860143) (0.86000,2.4838304) (0.86250,2.4816450) (0.86500,2.4794582) (0.86750,2.4772699) (0.87000,2.4750802) (0.87250,2.4728891) (0.87500,2.4706965) (0.87750,2.4685025) (0.88000,2.4663071) (0.88250,2.4641103) (0.88500,2.4619121) (0.88750,2.4597124) (0.89000,2.4575114) (0.89250,2.4553091) (0.89500,2.4531053) (0.89750,2.4509002) (0.90000,2.4486937) (0.90250,2.4464859) (0.90500,2.4442767) (0.90750,2.4420662) (0.91000,2.4398543) (0.91250,2.4376411) (0.91500,2.4354267) (0.91750,2.4332108) (0.92000,2.4309937) (0.92250,2.4287753) (0.92500,2.4265556) (0.92750,2.4243346) (0.93000,2.4221124) (0.93250,2.4198888) (0.93500,2.4176640) (0.93750,2.4154379) (0.94000,2.4132106) (0.94250,2.4109821) (0.94500,2.4087522) (0.94750,2.4065212) (0.95000,2.4042889) (0.95250,2.4020555) (0.95500,2.3998208) (0.95750,2.3975849) (0.96000,2.3953478) (0.96250,2.3931095) (0.96500,2.3908700) (0.96750,2.3886293) (0.97000,2.3863875) (0.97250,2.3841445) (0.97500,2.3819003) (0.97750,2.3796550) (0.98000,2.3774085) (0.98250,2.3751609) (0.98500,2.3729121) (0.98750,2.3706622) (0.99000,2.3684112) (0.99250,2.3661591) (0.99500,2.3639058) (0.99750,2.3616515) (1.00000,2.3593960) (1.00250,2.3571395) (1.00500,2.3548819) (1.00750,2.3526231) (1.01000,2.3503634) (1.01250,2.3481025) (1.01500,2.3458406) (1.01750,2.3435776) (1.02000,2.3413136) (1.02250,2.3390485) (1.02500,2.3367824) (1.02750,2.3345153) (1.03000,2.3322471) (1.03250,2.3299779) (1.03500,2.3277077) (1.03750,2.3254365) (1.04000,2.3231643) (1.04250,2.3208911) (1.04500,2.3186169) (1.04750,2.3163417) (1.05000,2.3140655) (1.05250,2.3117884) (1.05500,2.3095102) (1.05750,2.3072312) (1.06000,2.3049511) (1.06250,2.3026702) (1.06500,2.3003882) (1.06750,2.2981054) (1.07000,2.2958216) (1.07250,2.2935368) (1.07500,2.2912512) (1.07750,2.2889646) (1.08000,2.2866771) (1.08250,2.2843888) (1.08500,2.2820995) (1.08750,2.2798093) (1.09000,2.2775182) (1.09250,2.2752263) (1.09500,2.2729335) (1.09750,2.2706398) (1.10000,2.2683452) (1.10250,2.2660498) (1.10500,2.2637535) (1.10750,2.2614564) (1.11000,2.2591584) (1.11250,2.2568596) (1.11500,2.2545600) (1.11750,2.2522595) (1.12000,2.2499582) (1.12250,2.2476561) (1.12500,2.2453532) (1.12750,2.2430494) (1.13000,2.2407449) (1.13250,2.2384396) (1.13500,2.2361335) (1.13750,2.2338266) (1.14000,2.2315189) (1.14250,2.2292104) (1.14500,2.2269012) (1.14750,2.2245912) (1.15000,2.2222805) (1.15250,2.2199690) (1.15500,2.2176567) (1.15750,2.2153437) (1.16000,2.2130300) (1.16250,2.2107155) (1.16500,2.2084003) (1.16750,2.2060844) (1.17000,2.2037678) (1.17250,2.2014505) (1.17500,2.1991324) (1.17750,2.1968137) (1.18000,2.1944942) (1.18250,2.1921741) (1.18500,2.1898532) (1.18750,2.1875317) (1.19000,2.1852095) (1.19250,2.1828867) (1.19500,2.1805631) (1.19750,2.1782389) (1.20000,2.1759141) (1.20250,2.1735886) (1.20500,2.1712625) (1.20750,2.1689357) (1.21000,2.1666082) (1.21250,2.1642802) (1.21500,2.1619515) (1.21750,2.1596222) (1.22000,2.1572922) (1.22250,2.1549617) (1.22500,2.1526306) (1.22750,2.1502988) (1.23000,2.1479665) (1.23250,2.1456335) (1.23500,2.1433000) (1.23750,2.1409658) (1.24000,2.1386311) (1.24250,2.1362959) (1.24500,2.1339600) (1.24750,2.1316236) (1.25000,2.1292866) (1.25250,2.1269491) (1.25500,2.1246110) (1.25750,2.1222724) (1.26000,2.1199332) (1.26250,2.1175935) (1.26500,2.1152533) (1.26750,2.1129125) (1.27000,2.1105713) (1.27250,2.1082294) (1.27500,2.1058871) (1.27750,2.1035443) (1.28000,2.1012010) (1.28250,2.0988571) (1.28500,2.0965128) (1.28750,2.0941680) (1.29000,2.0918226) (1.29250,2.0894768) (1.29500,2.0871306) (1.29750,2.0847838) (1.30000,2.0824366) (1.30250,2.0800889) (1.30500,2.0777408) (1.30750,2.0753922) (1.31000,2.0730432) (1.31250,2.0706937) (1.31500,2.0683437) (1.31750,2.0659934) (1.32000,2.0636426) (1.32250,2.0612913) (1.32500,2.0589397) (1.32750,2.0565876) (1.33000,2.0542351) (1.33250,2.0518822) (1.33500,2.0495289) (1.33750,2.0471751) (1.34000,2.0448210) (1.34250,2.0424665) (1.34500,2.0401116) (1.34750,2.0377563) (1.35000,2.0354007) (1.35250,2.0330446) (1.35500,2.0306882) (1.35750,2.0283314) (1.36000,2.0259743) (1.36250,2.0236168) (1.36500,2.0212589) (1.36750,2.0189007) (1.37000,2.0165422) (1.37250,2.0141833) (1.37500,2.0118240) (1.37750,2.0094644) (1.38000,2.0071045) (1.38250,2.0047443) (1.38500,2.0023838) (1.38750,2.0000229) (1.39000,1.9976617) (1.39250,1.9953002) (1.39500,1.9929384) (1.39750,1.9905763) (1.40000,1.9882139) (1.40250,1.9858513) (1.40500,1.9834883) (1.40750,1.9811250) (1.41000,1.9787615) (1.41250,1.9763977) (1.41500,1.9740336) (1.41750,1.9716693) (1.42000,1.9693046) (1.42250,1.9669398) (1.42500,1.9645746) (1.42750,1.9622093) (1.43000,1.9598436) (1.43250,1.9574778) (1.43500,1.9551117) (1.43750,1.9527453) (1.44000,1.9503787) (1.44250,1.9480119) (1.44500,1.9456449) (1.44750,1.9432777) (1.45000,1.9409102) (1.45250,1.9385426) (1.45500,1.9361747) (1.45750,1.9338066) (1.46000,1.9314383) (1.46250,1.9290699) (1.46500,1.9267012) (1.46750,1.9243324) (1.47000,1.9219634) (1.47250,1.9195942) (1.47500,1.9172248) (1.47750,1.9148552) (1.48000,1.9124855) (1.48250,1.9101156) (1.48500,1.9077456) (1.48750,1.9053754) (1.49000,1.9030051) (1.49250,1.9006346) (1.49500,1.8982640) (1.49750,1.8958932) (1.50000,1.8935223) (1.50250,1.8911513) (1.50500,1.8887801) (1.50750,1.8864089) (1.51000,1.8840375) (1.51250,1.8816660) (1.51500,1.8792944) (1.51750,1.8769226) (1.52000,1.8745508) (1.52250,1.8721789) (1.52500,1.8698069) (1.52750,1.8674348) (1.53000,1.8650626) (1.53250,1.8626903) (1.53500,1.8603179) (1.53750,1.8579455) (1.54000,1.8555730) (1.54250,1.8532004) (1.54500,1.8508278) (1.54750,1.8484551) (1.55000,1.8460823) (1.55250,1.8437095) (1.55500,1.8413367) (1.55750,1.8389638) (1.56000,1.8365909) (1.56250,1.8342179) (1.56500,1.8318449) (1.56750,1.8294719) (1.57000,1.8270989) (1.57250,1.8247258) (1.57500,1.8223527) (1.57750,1.8199796) (1.58000,1.8176065) (1.58250,1.8152334) (1.58500,1.8128603) (1.58750,1.8104872) (1.59000,1.8081141) (1.59250,1.8057411) (1.59500,1.8033680) (1.59750,1.8009950) (1.60000,1.7986220) (1.60250,1.7962490) (1.60500,1.7938760) (1.60750,1.7915031) (1.61000,1.7891302) (1.61250,1.7867574) (1.61500,1.7843846) (1.61750,1.7820118) (1.62000,1.7796392) (1.62250,1.7772665) (1.62500,1.7748940) (1.62750,1.7725215) (1.63000,1.7701491) (1.63250,1.7677767) (1.63500,1.7654044) (1.63750,1.7630323) (1.64000,1.7606602) (1.64250,1.7582882) (1.64500,1.7559162) (1.64750,1.7535444) (1.65000,1.7511727) (1.65250,1.7488011) (1.65500,1.7464296) (1.65750,1.7440583) (1.66000,1.7416870) (1.66250,1.7393159) (1.66500,1.7369449) (1.66750,1.7345740) (1.67000,1.7322033) (1.67250,1.7298327) (1.67500,1.7274622) (1.67750,1.7250919) (1.68000,1.7227218) (1.68250,1.7203518) (1.68500,1.7179819) (1.68750,1.7156122) (1.69000,1.7132427) (1.69250,1.7108734) (1.69500,1.7085042) (1.69750,1.7061353) (1.70000,1.7037665) (1.70250,1.7013979) (1.70500,1.6990294) (1.70750,1.6966612) (1.71000,1.6942932) (1.71250,1.6919254) (1.71500,1.6895578) (1.71750,1.6871904) (1.72000,1.6848232) (1.72250,1.6824562) (1.72500,1.6800895) (1.72750,1.6777230) (1.73000,1.6753567) (1.73250,1.6729907) (1.73500,1.6706249) (1.73750,1.6682594) (1.74000,1.6658941) (1.74250,1.6635290) (1.74500,1.6611643) (1.74750,1.6587997) (1.75000,1.6564355) (1.75250,1.6540715) (1.75500,1.6517078) (1.75750,1.6493444) (1.76000,1.6469812) (1.76250,1.6446184) (1.76500,1.6422558) (1.76750,1.6398935) (1.77000,1.6375316) (1.77250,1.6351699) (1.77500,1.6328086) (1.77750,1.6304475) (1.78000,1.6280868) (1.78250,1.6257264) (1.78500,1.6233663) (1.78750,1.6210066) (1.79000,1.6186472) (1.79250,1.6162881) (1.79500,1.6139294) (1.79750,1.6115710) (1.80000,1.6092130) (1.80250,1.6068553) (1.80500,1.6044980) (1.80750,1.6021411) (1.81000,1.5997845) (1.81250,1.5974283) (1.81500,1.5950725) (1.81750,1.5927171) (1.82000,1.5903620) (1.82250,1.5880074) (1.82500,1.5856531) (1.82750,1.5832992) (1.83000,1.5809458) (1.83250,1.5785927) (1.83500,1.5762401) (1.83750,1.5738879) (1.84000,1.5715361) (1.84250,1.5691848) (1.84500,1.5668339) (1.84750,1.5644834) (1.85000,1.5621333) (1.85250,1.5597837) (1.85500,1.5574346) (1.85750,1.5550859) (1.86000,1.5527377) (1.86250,1.5503899) (1.86500,1.5480426) (1.86750,1.5456958) (1.87000,1.5433495) (1.87250,1.5410036) (1.87500,1.5386582) (1.87750,1.5363134) (1.88000,1.5339690) (1.88250,1.5316251) (1.88500,1.5292818) (1.88750,1.5269389) (1.89000,1.5245966) (1.89250,1.5222548) (1.89500,1.5199135) (1.89750,1.5175727) (1.90000,1.5152325) (1.90250,1.5128928) (1.90500,1.5105537) (1.90750,1.5082151) (1.91000,1.5058771) (1.91250,1.5035396) (1.91500,1.5012027) (1.91750,1.4988664) (1.92000,1.4965306) (1.92250,1.4941954) (1.92500,1.4918608) (1.92750,1.4895268) (1.93000,1.4871934) (1.93250,1.4848606) (1.93500,1.4825284) (1.93750,1.4801968) (1.94000,1.4778659) (1.94250,1.4755355) (1.94500,1.4732058) (1.94750,1.4708767) (1.95000,1.4685482) (1.95250,1.4662204) (1.95500,1.4638932) (1.95750,1.4615667) (1.96000,1.4592409) (1.96250,1.4569157) (1.96500,1.4545911) (1.96750,1.4522673) (1.97000,1.4499441) (1.97250,1.4476216) (1.97500,1.4452997) (1.97750,1.4429786) (1.98000,1.4406582) (1.98250,1.4383385) (1.98500,1.4360195) (1.98750,1.4337012) (1.99000,1.4313836) (1.99250,1.4290668) (1.99500,1.4267506) (1.99750,1.4244353) (2.00000,1.4221206) (2.00250,1.4198067) (2.00500,1.4174936) (2.00750,1.4151812) (2.01000,1.4128696) (2.01250,1.4105588) (2.01500,1.4082487) (2.01750,1.4059394) (2.02000,1.4036309) (2.02250,1.4013232) (2.02500,1.3990163) (2.02750,1.3967102) (2.03000,1.3944049) (2.03250,1.3921004) (2.03500,1.3897967) (2.03750,1.3874939) (2.04000,1.3851919) (2.04250,1.3828907) (2.04500,1.3805904) (2.04750,1.3782910) (2.05000,1.3759923) (2.05250,1.3736946) (2.05500,1.3713977) (2.05750,1.3691017) (2.06000,1.3668066) (2.06250,1.3645123) (2.06500,1.3622189) (2.06750,1.3599265) (2.07000,1.3576349) (2.07250,1.3553443) (2.07500,1.3530546) (2.07750,1.3507657) (2.08000,1.3484779) (2.08250,1.3461909) (2.08500,1.3439049) (2.08750,1.3416199) (2.09000,1.3393357) (2.09250,1.3370526) (2.09500,1.3347704) (2.09750,1.3324892) (2.10000,1.3302090) (2.10250,1.3279297) (2.10500,1.3256515) (2.10750,1.3233742) (2.11000,1.3210979) (2.11250,1.3188227) (2.11500,1.3165484) (2.11750,1.3142752) (2.12000,1.3120030) (2.12250,1.3097319) (2.12500,1.3074618) (2.12750,1.3051927) (2.13000,1.3029247) (2.13250,1.3006578) (2.13500,1.2983919) (2.13750,1.2961271) (2.14000,1.2938634) (2.14250,1.2916008) (2.14500,1.2893392) (2.14750,1.2870788) (2.15000,1.2848195) (2.15250,1.2825613) (2.15500,1.2803042) (2.15750,1.2780482) (2.16000,1.2757934) (2.16250,1.2735397) (2.16500,1.2712872) (2.16750,1.2690358) (2.17000,1.2667856) (2.17250,1.2645366) (2.17500,1.2622887) (2.17750,1.2600421) (2.18000,1.2577966) (2.18250,1.2555523) (2.18500,1.2533092) (2.18750,1.2510674) (2.19000,1.2488267) (2.19250,1.2465873) (2.19500,1.2443491) (2.19750,1.2421122) (2.20000,1.2398765) (2.20250,1.2376421) (2.20500,1.2354089) (2.20750,1.2331771) (2.21000,1.2309464) (2.21250,1.2287171) (2.21500,1.2264891) (2.21750,1.2242624) (2.22000,1.2220370) (2.22250,1.2198129) (2.22500,1.2175901) (2.22750,1.2153686) (2.23000,1.2131485) (2.23250,1.2109298) (2.23500,1.2087124) (2.23750,1.2064964) (2.24000,1.2042817) (2.24250,1.2020684) (2.24500,1.1998565) (2.24750,1.1976460) (2.25000,1.1954369) (2.25250,1.1932292) (2.25500,1.1910229) (2.25750,1.1888181) (2.26000,1.1866147) (2.26250,1.1844127) (2.26500,1.1822122) (2.26750,1.1800131) (2.27000,1.1778155) (2.27250,1.1756194) (2.27500,1.1734248) (2.27750,1.1712316) (2.28000,1.1690400) (2.28250,1.1668499) (2.28500,1.1646612) (2.28750,1.1624742) (2.29000,1.1602886) (2.29250,1.1581046) (2.29500,1.1559221) (2.29750,1.1537412) (2.30000,1.1515618) (2.30250,1.1493841) (2.30500,1.1472079) (2.30750,1.1450333) (2.31000,1.1428603) (2.31250,1.1406889) (2.31500,1.1385192) (2.31750,1.1363510) (2.32000,1.1341845) (2.32250,1.1320197) (2.32500,1.1298565) (2.32750,1.1276950) (2.33000,1.1255351) (2.33250,1.1233769) (2.33500,1.1212204) (2.33750,1.1190656) (2.34000,1.1169126) (2.34250,1.1147612) (2.34500,1.1126116) (2.34750,1.1104637) (2.35000,1.1083175) (2.35250,1.1061731) (2.35500,1.1040305) (2.35750,1.1018896) (2.36000,1.0997505) (2.36250,1.0976132) (2.36500,1.0954777) (2.36750,1.0933440) (2.37000,1.0912122) (2.37250,1.0890822) (2.37500,1.0869540) (2.37750,1.0848276) (2.38000,1.0827031) (2.38250,1.0805805) (2.38500,1.0784598) (2.38750,1.0763409) (2.39000,1.0742240) (2.39250,1.0721089) (2.39500,1.0699958) (2.39750,1.0678846) (2.40000,1.0657753) (2.40250,1.0636680) (2.40500,1.0615627) (2.40750,1.0594593) (2.41000,1.0573578) (2.41250,1.0552584) (2.41500,1.0531610) (2.41750,1.0510656) (2.42000,1.0489722) (2.42250,1.0468808) (2.42500,1.0447915) (2.42750,1.0427042) (2.43000,1.0406189) (2.43250,1.0385358) (2.43500,1.0364547) (2.43750,1.0343757) (2.44000,1.0322988) (2.44250,1.0302241) (2.44500,1.0281514) (2.44750,1.0260809) (2.45000,1.0240125) (2.45250,1.0219463) (2.45500,1.0198823) (2.45750,1.0178204) (2.46000,1.0157607) (2.46250,1.0137033) (2.46500,1.0116480) (2.46750,1.0095949) (2.47000,1.0075441) (2.47250,1.0054955) (2.47500,1.0034492) (2.47750,1.0014052) (2.48000,0.9993634) (2.48250,0.9973239) (2.48500,0.9952867) (2.48750,0.9932519) (2.49000,0.9912193) (2.49250,0.9891891) (2.49500,0.9871612) (2.49750,0.9851357) (2.50000,0.9831126) (2.50250,0.9810918) (2.50500,0.9790734) (2.50750,0.9770575) (2.51000,0.9750439) (2.51250,0.9730328) (2.51500,0.9710241) (2.51750,0.9690179) (2.52000,0.9670142) (2.52250,0.9650129) (2.52500,0.9630141) (2.52750,0.9610178) (2.53000,0.9590240) (2.53250,0.9570328) (2.53500,0.9550440) (2.53750,0.9530579) (2.54000,0.9510743) (2.54250,0.9490932) (2.54500,0.9471148) (2.54750,0.9451390) (2.55000,0.9431657) (2.55250,0.9411951) (2.55500,0.9392272) (2.55750,0.9372619) (2.56000,0.9352992) (2.56250,0.9333392) (2.56500,0.9313820) (2.56750,0.9294274) (2.57000,0.9274755) (2.57250,0.9255264) (2.57500,0.9235800) (2.57750,0.9216364) (2.58000,0.9196955) (2.58250,0.9177574) (2.58500,0.9158221) (2.58750,0.9138896) (2.59000,0.9119600) (2.59250,0.9100332) (2.59500,0.9081092) (2.59750,0.9061881) (2.60000,0.9042699) (2.60250,0.9023545) (2.60500,0.9004421) (2.60750,0.8985326) (2.61000,0.8966260) (2.61250,0.8947224) (2.61500,0.8928217) (2.61750,0.8909240) (2.62000,0.8890293) (2.62250,0.8871376) (2.62500,0.8852489) (2.62750,0.8833633) (2.63000,0.8814806) (2.63250,0.8796011) (2.63500,0.8777246) (2.63750,0.8758513) (2.64000,0.8739810) (2.64250,0.8721139) (2.64500,0.8702498) (2.64750,0.8683890) (2.65000,0.8665313) (2.65250,0.8646768) (2.65500,0.8628255) (2.65750,0.8609774) (2.66000,0.8591325) (2.66250,0.8572908) (2.66500,0.8554525) (2.66750,0.8536174) (2.67000,0.8517855) (2.67250,0.8499570) (2.67500,0.8481318) (2.67750,0.8463100) (2.68000,0.8444915) (2.68250,0.8426764) (2.68500,0.8408646) (2.68750,0.8390563) (2.69000,0.8372513) (2.69250,0.8354498) (2.69500,0.8336518) (2.69750,0.8318572) (2.70000,0.8300661) (2.70250,0.8282785) (2.70500,0.8264944) (2.70750,0.8247139) (2.71000,0.8229369) (2.71250,0.8211634) (2.71500,0.8193936) (2.71750,0.8176273) (2.72000,0.8158647) (2.72250,0.8141057) (2.72500,0.8123504) (2.72750,0.8105987) (2.73000,0.8088507) (2.73250,0.8071065) (2.73500,0.8053659) (2.73750,0.8036291) (2.74000,0.8018961) (2.74250,0.8001668) (2.74500,0.7984414) (2.74750,0.7967197) (2.75000,0.7950019) (2.75250,0.7932879) (2.75500,0.7915778) (2.75750,0.7898717) (2.76000,0.7881694) (2.76250,0.7864710) (2.76500,0.7847766) (2.76750,0.7830861) (2.77000,0.7813997) (2.77250,0.7797172) (2.77500,0.7780388) (2.77750,0.7763644) (2.78000,0.7746941) (2.78250,0.7730279) (2.78500,0.7713658) (2.78750,0.7697078) (2.79000,0.7680539) (2.79250,0.7664043) (2.79500,0.7647588) (2.79750,0.7631175) (2.80000,0.7614804) (2.80250,0.7598476) (2.80500,0.7582191) (2.80750,0.7565948) (2.81000,0.7549749) (2.81250,0.7533593) (2.81500,0.7517481) (2.81750,0.7501413) (2.82000,0.7485388) (2.82250,0.7469408) (2.82500,0.7453472) (2.82750,0.7437582) (2.83000,0.7421735) (2.83250,0.7405935) (2.83500,0.7390179) (2.83750,0.7374469) (2.84000,0.7358805) (2.84250,0.7343188) (2.84500,0.7327616) (2.84750,0.7312091) (2.85000,0.7296613) (2.85250,0.7281182) (2.85500,0.7265798) (2.85750,0.7250462) (2.86000,0.7235174) (2.86250,0.7219934) (2.86500,0.7204742) (2.86750,0.7189598) (2.87000,0.7174504) (2.87250,0.7159458) (2.87500,0.7144462) (2.87750,0.7129516) (2.88000,0.7114620) (2.88250,0.7099773) (2.88500,0.7084977) (2.88750,0.7070232) (2.89000,0.7055538) (2.89250,0.7040895) (2.89500,0.7026304) (2.89750,0.7011765) (2.90000,0.6997277) (2.90250,0.6982843) (2.90500,0.6968461) (2.90750,0.6954131) (2.91000,0.6939856) (2.91250,0.6925633) (2.91500,0.6911465) (2.91750,0.6897351) (2.92000,0.6883292) (2.92250,0.6869287) (2.92500,0.6855338) (2.92750,0.6841444) (2.93000,0.6827606) (2.93250,0.6813824) (2.93500,0.6800099) (2.93750,0.6786430) (2.94000,0.6772819) (2.94250,0.6759265) (2.94500,0.6745769) (2.94750,0.6732331) (2.95000,0.6718951) (2.95250,0.6705631) (2.95500,0.6692370) (2.95750,0.6679168) (2.96000,0.6666027) (2.96250,0.6652946) (2.96500,0.6639925) (2.96750,0.6626966) (2.97000,0.6614068) (2.97250,0.6601232) (2.97500,0.6588459) (2.97750,0.6575748) (2.98000,0.6563100) (2.98250,0.6550516) (2.98500,0.6537996) (2.98750,0.6525540) (2.99000,0.6513149) (2.99250,0.6500823) (2.99500,0.6488563) (2.99750,0.6476369) (3.00000,0.6464242) (3.00250,0.6452181) (3.00500,0.6440188) (3.00750,0.6428263) (3.01000,0.6416407) (3.01250,0.6404619) (3.01500,0.6392900) (3.01750,0.6381252) (3.02000,0.6369674) (3.02250,0.6358166) (3.02500,0.6346730) (3.02750,0.6335366) (3.03000,0.6324075) (3.03250,0.6312856) (3.03500,0.6301711) (3.03750,0.6290640) (3.04000,0.6279644) (3.04250,0.6268722) (3.04500,0.6257877) (3.04750,0.6247108) (3.05000,0.6236416) (3.05250,0.6225801) (3.05500,0.6215265) (3.05750,0.6204807) (3.06000,0.6194429) (3.06250,0.6184130) (3.06500,0.6173913) (3.06750,0.6163776) (3.07000,0.6153722) (3.07250,0.6143751) (3.07500,0.6133863) (3.07750,0.6124059) (3.08000,0.6114340) (3.08250,0.6104706) (3.08500,0.6095159) (3.08750,0.6085699) (3.09000,0.6076326) (3.09250,0.6067043) (3.09500,0.6057848) (3.09750,0.6048744) (3.10000,0.6039731) (3.10250,0.6030809) (3.10500,0.6021981) (3.10750,0.6013246) (3.11000,0.6004605) (3.11250,0.5996060) (3.11500,0.5987611) (3.11750,0.5979259) (3.12000,0.5971005) (3.12250,0.5962850) (3.12500,0.5954796) (3.12750,0.5946842) (3.13000,0.5938991) (3.13250,0.5931243) (3.13500,0.5923599) (3.13750,0.5916060) (3.14000,0.5908628) (3.14250,0.5901303) (3.14500,0.5894087) (3.14750,0.5886980) (3.15000,0.5879985) (3.15250,0.5873102) (3.15500,0.5866332) (3.15750,0.5859677) (3.16000,0.5853138) (3.16250,0.5846716) (3.16500,0.5840413) (3.16750,0.5834230) (3.17000,0.5828169) (3.17250,0.5822231) (3.17500,0.5816416) (3.17750,0.5810728) (3.18000,0.5805167) (3.18250,0.5799735) (3.18500,0.5794434) (3.18750,0.5789265) (3.19000,0.5784230) (3.19250,0.5779330) (3.19500,0.5774568) (3.19750,0.5769944) (3.20000,0.5765462) (3.20250,0.5761123) (3.20500,0.5756928) (3.20750,0.5752880) (3.21000,0.5748980) (3.21250,0.5745232) (3.21500,0.5741636) (3.21750,0.5738195) (3.22000,0.5734912) (3.22250,0.5731788) (3.22500,0.5728826) (3.22750,0.5726028) (3.23000,0.5723397) (3.23250,0.5720934) (3.23500,0.5718644) (3.23750,0.5716528) (3.24000,0.5714589) (3.24250,0.5712830) (3.24500,0.5711253) (3.24750,0.5709862) (3.25000,0.5708660) (3.25250,0.5707649) (3.25500,0.5706833) (3.25750,0.5706215) (3.26000,0.5705799) (3.26250,0.5705587) (3.26500,0.5705584) (3.26750,0.5705792) (3.27000,0.5706217) (3.27250,0.5706861) (3.27500,0.5707729) (3.27750,0.5708824) (3.28000,0.5710151) (3.28250,0.5711714) (3.28500,0.5713517) (3.28750,0.5715566) (3.29000,0.5717864) (3.29250,0.5720418) (3.29500,0.5723231) (3.29750,0.5726308) (3.30000,0.5729656) (3.30250,0.5733280) (3.30500,0.5737185) (3.30750,0.5741377) (3.31000,0.5745862) (3.31250,0.5750647) (3.31500,0.5755738) (3.31750,0.5761142) (3.32000,0.5766865) (3.32250,0.5772915) (3.32500,0.5779299) (3.32750,0.5786024) (3.33000,0.5793099) (3.33250,0.5800532) (3.33500,0.5808332) (3.33750,0.5816506) (3.34000,0.5825065) (3.34250,0.5834017) (3.34500,0.5843373) (3.34750,0.5853142) (3.35000,0.5863336) (3.35250,0.5873964) (3.35500,0.5885040) (3.35750,0.5896574) (3.36000,0.5908578) (3.36250,0.5921066) (3.36500,0.5934051) (3.36750,0.5947547) (3.37000,0.5961568) (3.37250,0.5976130) (3.37500,0.5991247) (3.37750,0.6006937) (3.38000,0.6023217) (3.38250,0.6040104) (3.38500,0.6057618) (3.38750,0.6075777) (3.39000,0.6094602) (3.39250,0.6114115) (3.39500,0.6134338) (3.39750,0.6155294) (3.40000,0.6177009) (3.40250,0.6199507) (3.40500,0.6222816) (3.40750,0.6246965) (3.41000,0.6271984) (3.41250,0.6297903) (3.41500,0.6324757) (3.41750,0.6352579) (3.42000,0.6381408) (3.42250,0.6411282) (3.42500,0.6442242) (3.42750,0.6474332) (3.43000,0.6507598) (3.43250,0.6542088) (3.43500,0.6577854) (3.43750,0.6614951) (3.44000,0.6653438) (3.44250,0.6693377) (3.44500,0.6734833) (3.44750,0.6777877) (3.45000,0.6822585) (3.45250,0.6869037) (3.45500,0.6917319) (3.45750,0.6967524) (3.46000,0.7019750) (3.46250,0.7074105) (3.46500,0.7130702) (3.46750,0.7189667) (3.47000,0.7251133) (3.47250,0.7315245) (3.47500,0.7382160) (3.47750,0.7452050) (3.48000,0.7525101) (3.48250,0.7601517) (3.48500,0.7681520) (3.48750,0.7765357) (3.49000,0.7853295) (3.49250,0.7945634) (3.49500,0.8042702) (3.49750,0.8144867) (3.50000,0.8252538) (3.50250,0.8366172) (3.50500,0.8486285) (3.50750,0.8613459) (3.51000,0.8748355) (3.51250,0.8891728) (3.51500,0.9044445) (3.51750,0.9207509) (3.52000,0.9382085) (3.52250,0.9569543) (3.52500,0.9771503) (3.52750,0.9989903) (3.53000,1.0227086) (3.53250,1.0485920) (3.53500,1.0769965) (3.53750,1.1083722) (3.54000,1.1432985) (3.54250,1.1825397) (3.54500,1.2271321) (3.54750,1.2785306) (3.55000,1.3388698) (3.55250,1.4114654) (3.55500,1.5018859) (3.55750,1.6206081) (3.56000,1.7912934) (3.56250,2.0904534) (3.56500,nan) (3.56750,nan) (3.57000,nan) (3.57250,nan) (3.57500,nan) (3.57750,nan) (3.58000,nan) (3.58250,nan) (3.58500,nan) (3.58750,nan) (3.59000,nan) (3.59250,nan) (3.59500,nan) (3.59750,nan) (3.60000,nan) (3.60250,nan) (3.60500,nan) (3.60750,nan) (3.61000,nan) (3.61250,nan) (3.61500,nan) (3.61750,nan) (3.62000,nan) (3.62250,nan) (3.62500,nan) (3.62750,nan) (3.63000,nan) (3.63250,nan) (3.63500,nan) (3.63750,nan) (3.64000,nan) (3.64250,nan) (3.64500,nan) (3.64750,nan) (3.65000,nan) (3.65250,nan) (3.65500,nan) (3.65750,nan) (3.66000,nan) (3.66250,nan) (3.66500,nan) (3.66750,nan) (3.67000,nan) (3.67250,nan) (3.67500,nan) (3.67750,nan) (3.68000,nan) (3.68250,nan) (3.68500,nan) (3.68750,nan) (3.69000,nan) (3.69250,nan) (3.69500,nan) (3.69750,nan) (3.70000,nan) (3.70250,nan) (3.70500,nan) (3.70750,nan) (3.71000,nan) (3.71250,nan) (3.71500,nan) (3.71750,nan) (3.72000,nan) (3.72250,nan) (3.72500,nan) (3.72750,nan) (3.73000,nan) (3.73250,-1.7529820) (3.73500,-1.3914629) (3.73750,-1.2020037) (3.74000,-1.0741416) (3.74250,-0.9782594) (3.74500,-0.9019933) (3.74750,-0.8389885) (3.75000,-0.7855433) (3.75250,-0.7393116) (3.75500,-0.6987125) (3.75750,-0.6626303) (3.76000,-0.6302481) (3.76250,-0.6009502) (3.76500,-0.5742609) (3.76750,-0.5498050) (3.77000,-0.5272818) (3.77250,-0.5064462) (3.77500,-0.4870962) (3.77750,-0.4690631) (3.78000,-0.4522048) (3.78250,-0.4364003) (3.78500,-0.4215461) (3.78750,-0.4075525) (3.79000,-0.3943417) (3.79250,-0.3818456) (3.79500,-0.3700043) (3.79750,-0.3587649) (3.80000,-0.3480803) (3.80250,-0.3379086) (3.80500,-0.3282121) (3.80750,-0.3189572) (3.81000,-0.3101133) (3.81250,-0.3016530) (3.81500,-0.2935511) (3.81750,-0.2857849) (3.82000,-0.2783336) (3.82250,-0.2711782) (3.82500,-0.2643012) (3.82750,-0.2576867) (3.83000,-0.2513198) (3.83250,-0.2451869) (3.83500,-0.2392754) (3.83750,-0.2335736) (3.84000,-0.2280706) (3.84250,-0.2227565) (3.84500,-0.2176218) (3.84750,-0.2126578) (3.85000,-0.2078562) (3.85250,-0.2032095) (3.85500,-0.1987105) (3.85750,-0.1943526) (3.86000,-0.1901293) (3.86250,-0.1860348) (3.86500,-0.1820635) (3.86750,-0.1782102) (3.87000,-0.1744700) (3.87250,-0.1708382) (3.87500,-0.1673103) (3.87750,-0.1638823) (3.88000,-0.1605501) (3.88250,-0.1573101) (3.88500,-0.1541587) (3.88750,-0.1510925) (3.89000,-0.1481085) (3.89250,-0.1452034) (3.89500,-0.1423746) (3.89750,-0.1396191) (3.90000,-0.1369345) (3.90250,-0.1343182) (3.90500,-0.1317679) (3.90750,-0.1292813) (3.91000,-0.1268562) (3.91250,-0.1244906) (3.91500,-0.1221826) (3.91750,-0.1199301) (3.92000,-0.1177315) (3.92250,-0.1155850) (3.92500,-0.1134890) (3.92750,-0.1114418) (3.93000,-0.1094420) (3.93250,-0.1074881) (3.93500,-0.1055787) (3.93750,-0.1037124) (3.94000,-0.1018880) (3.94250,-0.1001043) (3.94500,-0.0983599) (3.94750,-0.0966539) (3.95000,-0.0949851) (3.95250,-0.0933524) (3.95500,-0.0917548) (3.95750,-0.0901913) (3.96000,-0.0886610) (3.96250,-0.0871630) (3.96500,-0.0856963) (3.96750,-0.0842602) (3.97000,-0.0828538) (3.97250,-0.0814763) (3.97500,-0.0801269) (3.97750,-0.0788050) (3.98000,-0.0775098) (3.98250,-0.0762406) (3.98500,-0.0749967) (3.98750,-0.0737776) (3.99000,-0.0725825) (3.99250,-0.0714109) (3.99500,-0.0702622) (3.99750,-0.0691358) (4.00000,-0.0680311)};\addlegendentry{R}
\addplot[thick,orange] coordinates {(-2.00000,nan) (-1.99750,nan) (-1.99500,nan) (-1.99250,nan) (-1.99000,nan) (-1.98750,nan) (-1.98500,nan) (-1.98250,nan) (-1.98000,nan) (-1.97750,nan) (-1.97500,nan) (-1.97250,nan) (-1.97000,nan) (-1.96750,nan) (-1.96500,nan) (-1.96250,nan) (-1.96000,nan) (-1.95750,nan) (-1.95500,nan) (-1.95250,nan) (-1.95000,nan) (-1.94750,nan) (-1.94500,nan) (-1.94250,nan) (-1.94000,nan) (-1.93750,nan) (-1.93500,nan) (-1.93250,nan) (-1.93000,nan) (-1.92750,nan) (-1.92500,nan) (-1.92250,nan) (-1.92000,nan) (-1.91750,nan) (-1.91500,nan) (-1.91250,nan) (-1.91000,nan) (-1.90750,nan) (-1.90500,nan) (-1.90250,nan) (-1.90000,nan) (-1.89750,nan) (-1.89500,nan) (-1.89250,nan) (-1.89000,nan) (-1.88750,nan) (-1.88500,nan) (-1.88250,nan) (-1.88000,nan) (-1.87750,nan) (-1.87500,nan) (-1.87250,nan) (-1.87000,nan) (-1.86750,nan) (-1.86500,nan) (-1.86250,nan) (-1.86000,nan) (-1.85750,nan) (-1.85500,nan) (-1.85250,nan) (-1.85000,nan) (-1.84750,nan) (-1.84500,nan) (-1.84250,nan) (-1.84000,nan) (-1.83750,nan) (-1.83500,nan) (-1.83250,nan) (-1.83000,nan) (-1.82750,nan) (-1.82500,nan) (-1.82250,nan) (-1.82000,nan) (-1.81750,nan) (-1.81500,nan) (-1.81250,nan) (-1.81000,nan) (-1.80750,nan) (-1.80500,nan) (-1.80250,nan) (-1.80000,nan) (-1.79750,nan) (-1.79500,nan) (-1.79250,nan) (-1.79000,nan) (-1.78750,nan) (-1.78500,nan) (-1.78250,nan) (-1.78000,nan) (-1.77750,nan) (-1.77500,nan) (-1.77250,nan) (-1.77000,nan) (-1.76750,nan) (-1.76500,nan) (-1.76250,nan) (-1.76000,nan) (-1.75750,nan) (-1.75500,nan) (-1.75250,nan) (-1.75000,nan) (-1.74750,nan) (-1.74500,nan) (-1.74250,nan) (-1.74000,nan) (-1.73750,nan) (-1.73500,nan) (-1.73250,nan) (-1.73000,nan) (-1.72750,nan) (-1.72500,nan) (-1.72250,nan) (-1.72000,nan) (-1.71750,nan) (-1.71500,nan) (-1.71250,nan) (-1.71000,nan) (-1.70750,nan) (-1.70500,nan) (-1.70250,nan) (-1.70000,nan) (-1.69750,nan) (-1.69500,nan) (-1.69250,nan) (-1.69000,nan) (-1.68750,nan) (-1.68500,nan) (-1.68250,nan) (-1.68000,nan) (-1.67750,nan) (-1.67500,nan) (-1.67250,nan) (-1.67000,nan) (-1.66750,nan) (-1.66500,nan) (-1.66250,nan) (-1.66000,nan) (-1.65750,nan) (-1.65500,nan) (-1.65250,nan) (-1.65000,nan) (-1.64750,nan) (-1.64500,nan) (-1.64250,nan) (-1.64000,nan) (-1.63750,nan) (-1.63500,nan) (-1.63250,nan) (-1.63000,nan) (-1.62750,nan) (-1.62500,nan) (-1.62250,nan) (-1.62000,nan) (-1.61750,nan) (-1.61500,nan) (-1.61250,nan) (-1.61000,nan) (-1.60750,nan) (-1.60500,nan) (-1.60250,nan) (-1.60000,nan) (-1.59750,nan) (-1.59500,nan) (-1.59250,nan) (-1.59000,nan) (-1.58750,nan) (-1.58500,nan) (-1.58250,nan) (-1.58000,nan) (-1.57750,nan) (-1.57500,nan) (-1.57250,nan) (-1.57000,nan) (-1.56750,nan) (-1.56500,nan) (-1.56250,nan) (-1.56000,nan) (-1.55750,nan) (-1.55500,nan) (-1.55250,nan) (-1.55000,nan) (-1.54750,nan) (-1.54500,nan) (-1.54250,nan) (-1.54000,nan) (-1.53750,nan) (-1.53500,nan) (-1.53250,nan) (-1.53000,nan) (-1.52750,nan) (-1.52500,nan) (-1.52250,nan) (-1.52000,nan) (-1.51750,nan) (-1.51500,nan) (-1.51250,nan) (-1.51000,nan) (-1.50750,nan) (-1.50500,nan) (-1.50250,nan) (-1.50000,nan) (-1.49750,nan) (-1.49500,nan) (-1.49250,nan) (-1.49000,nan) (-1.48750,nan) (-1.48500,nan) (-1.48250,nan) (-1.48000,nan) (-1.47750,nan) (-1.47500,nan) (-1.47250,nan) (-1.47000,nan) (-1.46750,nan) (-1.46500,nan) (-1.46250,nan) (-1.46000,nan) (-1.45750,nan) (-1.45500,nan) (-1.45250,nan) (-1.45000,nan) (-1.44750,nan) (-1.44500,nan) (-1.44250,nan) (-1.44000,nan) (-1.43750,nan) (-1.43500,nan) (-1.43250,nan) (-1.43000,nan) (-1.42750,nan) (-1.42500,nan) (-1.42250,nan) (-1.42000,nan) (-1.41750,nan) (-1.41500,nan) (-1.41250,nan) (-1.41000,nan) (-1.40750,nan) (-1.40500,nan) (-1.40250,nan) (-1.40000,nan) (-1.39750,nan) (-1.39500,nan) (-1.39250,nan) (-1.39000,nan) (-1.38750,nan) (-1.38500,nan) (-1.38250,nan) (-1.38000,nan) (-1.37750,nan) (-1.37500,nan) (-1.37250,nan) (-1.37000,nan) (-1.36750,nan) (-1.36500,nan) (-1.36250,nan) (-1.36000,nan) (-1.35750,nan) (-1.35500,nan) (-1.35250,nan) (-1.35000,nan) (-1.34750,nan) (-1.34500,nan) (-1.34250,nan) (-1.34000,nan) (-1.33750,nan) (-1.33500,nan) (-1.33250,nan) (-1.33000,nan) (-1.32750,nan) (-1.32500,nan) (-1.32250,nan) (-1.32000,nan) (-1.31750,nan) (-1.31500,nan) (-1.31250,nan) (-1.31000,nan) (-1.30750,nan) (-1.30500,nan) (-1.30250,nan) (-1.30000,nan) (-1.29750,nan) (-1.29500,nan) (-1.29250,nan) (-1.29000,nan) (-1.28750,nan) (-1.28500,nan) (-1.28250,nan) (-1.28000,nan) (-1.27750,nan) (-1.27500,nan) (-1.27250,nan) (-1.27000,nan) (-1.26750,nan) (-1.26500,nan) (-1.26250,nan) (-1.26000,nan) (-1.25750,nan) (-1.25500,nan) (-1.25250,nan) (-1.25000,nan) (-1.24750,nan) (-1.24500,nan) (-1.24250,nan) (-1.24000,nan) (-1.23750,nan) (-1.23500,nan) (-1.23250,nan) (-1.23000,nan) (-1.22750,nan) (-1.22500,nan) (-1.22250,nan) (-1.22000,nan) (-1.21750,nan) (-1.21500,nan) (-1.21250,nan) (-1.21000,nan) (-1.20750,nan) (-1.20500,nan) (-1.20250,nan) (-1.20000,nan) (-1.19750,nan) (-1.19500,nan) (-1.19250,nan) (-1.19000,nan) (-1.18750,nan) (-1.18500,nan) (-1.18250,nan) (-1.18000,nan) (-1.17750,nan) (-1.17500,nan) (-1.17250,nan) (-1.17000,nan) (-1.16750,nan) (-1.16500,nan) (-1.16250,nan) (-1.16000,nan) (-1.15750,nan) (-1.15500,nan) (-1.15250,nan) (-1.15000,nan) (-1.14750,nan) (-1.14500,nan) (-1.14250,nan) (-1.14000,nan) (-1.13750,nan) (-1.13500,nan) (-1.13250,nan) (-1.13000,nan) (-1.12750,nan) (-1.12500,nan) (-1.12250,nan) (-1.12000,nan) (-1.11750,nan) (-1.11500,nan) (-1.11250,nan) (-1.11000,nan) (-1.10750,nan) (-1.10500,nan) (-1.10250,nan) (-1.10000,nan) (-1.09750,nan) (-1.09500,nan) (-1.09250,nan) (-1.09000,nan) (-1.08750,nan) (-1.08500,nan) (-1.08250,nan) (-1.08000,nan) (-1.07750,nan) (-1.07500,nan) (-1.07250,nan) (-1.07000,nan) (-1.06750,nan) (-1.06500,nan) (-1.06250,nan) (-1.06000,nan) (-1.05750,nan) (-1.05500,nan) (-1.05250,nan) (-1.05000,nan) (-1.04750,nan) (-1.04500,nan) (-1.04250,nan) (-1.04000,nan) (-1.03750,nan) (-1.03500,nan) (-1.03250,nan) (-1.03000,nan) (-1.02750,nan) (-1.02500,nan) (-1.02250,nan) (-1.02000,nan) (-1.01750,nan) (-1.01500,nan) (-1.01250,nan) (-1.01000,nan) (-1.00750,nan) (-1.00500,nan) (-1.00250,nan) (-1.00000,nan) (-0.99750,nan) (-0.99500,nan) (-0.99250,nan) (-0.99000,nan) (-0.98750,nan) (-0.98500,nan) (-0.98250,nan) (-0.98000,nan) (-0.97750,nan) (-0.97500,nan) (-0.97250,nan) (-0.97000,nan) (-0.96750,nan) (-0.96500,nan) (-0.96250,nan) (-0.96000,nan) (-0.95750,nan) (-0.95500,nan) (-0.95250,nan) (-0.95000,nan) (-0.94750,nan) (-0.94500,nan) (-0.94250,nan) (-0.94000,nan) (-0.93750,nan) (-0.93500,nan) (-0.93250,nan) (-0.93000,nan) (-0.92750,nan) (-0.92500,nan) (-0.92250,nan) (-0.92000,nan) (-0.91750,nan) (-0.91500,nan) (-0.91250,nan) (-0.91000,nan) (-0.90750,nan) (-0.90500,nan) (-0.90250,nan) (-0.90000,nan) (-0.89750,nan) (-0.89500,nan) (-0.89250,nan) (-0.89000,nan) (-0.88750,nan) (-0.88500,nan) (-0.88250,nan) (-0.88000,nan) (-0.87750,nan) (-0.87500,nan) (-0.87250,nan) (-0.87000,nan) (-0.86750,nan) (-0.86500,nan) (-0.86250,nan) (-0.86000,nan) (-0.85750,nan) (-0.85500,nan) (-0.85250,nan) (-0.85000,nan) (-0.84750,nan) (-0.84500,nan) (-0.84250,nan) (-0.84000,nan) (-0.83750,nan) (-0.83500,nan) (-0.83250,nan) (-0.83000,nan) (-0.82750,nan) (-0.82500,nan) (-0.82250,nan) (-0.82000,nan) (-0.81750,nan) (-0.81500,nan) (-0.81250,nan) (-0.81000,nan) (-0.80750,nan) (-0.80500,nan) (-0.80250,nan) (-0.80000,nan) (-0.79750,nan) (-0.79500,nan) (-0.79250,nan) (-0.79000,nan) (-0.78750,nan) (-0.78500,nan) (-0.78250,nan) (-0.78000,nan) (-0.77750,nan) (-0.77500,nan) (-0.77250,nan) (-0.77000,nan) (-0.76750,nan) (-0.76500,nan) (-0.76250,nan) (-0.76000,nan) (-0.75750,nan) (-0.75500,nan) (-0.75250,nan) (-0.75000,nan) (-0.74750,nan) (-0.74500,nan) (-0.74250,nan) (-0.74000,nan) (-0.73750,nan) (-0.73500,nan) (-0.73250,nan) (-0.73000,nan) (-0.72750,nan) (-0.72500,nan) (-0.72250,nan) (-0.72000,nan) (-0.71750,nan) (-0.71500,nan) (-0.71250,nan) (-0.71000,nan) (-0.70750,nan) (-0.70500,nan) (-0.70250,nan) (-0.70000,nan) (-0.69750,nan) (-0.69500,nan) (-0.69250,nan) (-0.69000,nan) (-0.68750,nan) (-0.68500,nan) (-0.68250,nan) (-0.68000,nan) (-0.67750,nan) (-0.67500,nan) (-0.67250,nan) (-0.67000,nan) (-0.66750,nan) (-0.66500,nan) (-0.66250,nan) (-0.66000,nan) (-0.65750,nan) (-0.65500,nan) (-0.65250,nan) (-0.65000,nan) (-0.64750,nan) (-0.64500,nan) (-0.64250,nan) (-0.64000,nan) (-0.63750,nan) (-0.63500,nan) (-0.63250,nan) (-0.63000,nan) (-0.62750,nan) (-0.62500,nan) (-0.62250,nan) (-0.62000,nan) (-0.61750,nan) (-0.61500,nan) (-0.61250,nan) (-0.61000,nan) (-0.60750,nan) (-0.60500,nan) (-0.60250,nan) (-0.60000,nan) (-0.59750,nan) (-0.59500,nan) (-0.59250,nan) (-0.59000,nan) (-0.58750,nan) (-0.58500,nan) (-0.58250,nan) (-0.58000,nan) (-0.57750,nan) (-0.57500,nan) (-0.57250,nan) (-0.57000,nan) (-0.56750,nan) (-0.56500,nan) (-0.56250,nan) (-0.56000,nan) (-0.55750,nan) (-0.55500,nan) (-0.55250,nan) (-0.55000,nan) (-0.54750,nan) (-0.54500,nan) (-0.54250,nan) (-0.54000,nan) (-0.53750,nan) (-0.53500,nan) (-0.53250,nan) (-0.53000,nan) (-0.52750,nan) (-0.52500,nan) (-0.52250,nan) (-0.52000,nan) (-0.51750,nan) (-0.51500,nan) (-0.51250,nan) (-0.51000,nan) (-0.50750,nan) (-0.50500,nan) (-0.50250,nan) (-0.50000,nan) (-0.49750,nan) (-0.49500,nan) (-0.49250,nan) (-0.49000,nan) (-0.48750,nan) (-0.48500,nan) (-0.48250,nan) (-0.48000,nan) (-0.47750,nan) (-0.47500,nan) (-0.47250,nan) (-0.47000,nan) (-0.46750,nan) (-0.46500,nan) (-0.46250,nan) (-0.46000,nan) (-0.45750,nan) (-0.45500,nan) (-0.45250,nan) (-0.45000,nan) (-0.44750,nan) (-0.44500,nan) (-0.44250,nan) (-0.44000,nan) (-0.43750,nan) (-0.43500,nan) (-0.43250,nan) (-0.43000,nan) (-0.42750,nan) (-0.42500,nan) (-0.42250,nan) (-0.42000,nan) (-0.41750,nan) (-0.41500,nan) (-0.41250,nan) (-0.41000,nan) (-0.40750,nan) (-0.40500,nan) (-0.40250,nan) (-0.40000,nan) (-0.39750,nan) (-0.39500,nan) (-0.39250,nan) (-0.39000,nan) (-0.38750,nan) (-0.38500,nan) (-0.38250,nan) (-0.38000,nan) (-0.37750,nan) (-0.37500,nan) (-0.37250,nan) (-0.37000,nan) (-0.36750,nan) (-0.36500,nan) (-0.36250,nan) (-0.36000,nan) (-0.35750,nan) (-0.35500,nan) (-0.35250,nan) (-0.35000,nan) (-0.34750,nan) (-0.34500,nan) (-0.34250,nan) (-0.34000,nan) (-0.33750,nan) (-0.33500,nan) (-0.33250,nan) (-0.33000,nan) (-0.32750,nan) (-0.32500,nan) (-0.32250,nan) (-0.32000,nan) (-0.31750,nan) (-0.31500,nan) (-0.31250,nan) (-0.31000,nan) (-0.30750,nan) (-0.30500,nan) (-0.30250,nan) (-0.30000,nan) (-0.29750,nan) (-0.29500,nan) (-0.29250,nan) (-0.29000,nan) (-0.28750,nan) (-0.28500,nan) (-0.28250,nan) (-0.28000,nan) (-0.27750,nan) (-0.27500,nan) (-0.27250,nan) (-0.27000,nan) (-0.26750,nan) (-0.26500,nan) (-0.26250,nan) (-0.26000,nan) (-0.25750,nan) (-0.25500,nan) (-0.25250,nan) (-0.25000,nan) (-0.24750,nan) (-0.24500,nan) (-0.24250,nan) (-0.24000,nan) (-0.23750,nan) (-0.23500,nan) (-0.23250,nan) (-0.23000,nan) (-0.22750,nan) (-0.22500,nan) (-0.22250,nan) (-0.22000,nan) (-0.21750,nan) (-0.21500,nan) (-0.21250,nan) (-0.21000,nan) (-0.20750,nan) (-0.20500,nan) (-0.20250,nan) (-0.20000,nan) (-0.19750,nan) (-0.19500,nan) (-0.19250,nan) (-0.19000,nan) (-0.18750,nan) (-0.18500,nan) (-0.18250,nan) (-0.18000,nan) (-0.17750,nan) (-0.17500,nan) (-0.17250,nan) (-0.17000,nan) (-0.16750,nan) (-0.16500,nan) (-0.16250,nan) (-0.16000,nan) (-0.15750,nan) (-0.15500,nan) (-0.15250,nan) (-0.15000,nan) (-0.14750,nan) (-0.14500,nan) (-0.14250,nan) (-0.14000,nan) (-0.13750,nan) (-0.13500,nan) (-0.13250,nan) (-0.13000,nan) (-0.12750,nan) (-0.12500,nan) (-0.12250,nan) (-0.12000,nan) (-0.11750,nan) (-0.11500,nan) (-0.11250,nan) (-0.11000,nan) (-0.10750,nan) (-0.10500,nan) (-0.10250,nan) (-0.10000,nan) (-0.09750,nan) (-0.09500,nan) (-0.09250,nan) (-0.09000,nan) (-0.08750,nan) (-0.08500,nan) (-0.08250,nan) (-0.08000,nan) (-0.07750,nan) (-0.07500,nan) (-0.07250,nan) (-0.07000,nan) (-0.06750,nan) (-0.06500,nan) (-0.06250,nan) (-0.06000,nan) (-0.05750,nan) (-0.05500,nan) (-0.05250,nan) (-0.05000,nan) (-0.04750,nan) (-0.04500,nan) (-0.04250,nan) (-0.04000,nan) (-0.03750,nan) (-0.03500,nan) (-0.03250,nan) (-0.03000,nan) (-0.02750,nan) (-0.02500,nan) (-0.02250,nan) (-0.02000,nan) (-0.01750,nan) (-0.01500,nan) (-0.01250,nan) (-0.01000,nan) (-0.00750,nan) (-0.00500,nan) (-0.00250,nan) (0.00000,nan) (0.00250,nan) (0.00500,nan) (0.00750,nan) (0.01000,nan) (0.01250,nan) (0.01500,nan) (0.01750,nan) (0.02000,nan) (0.02250,nan) (0.02500,nan) (0.02750,nan) (0.03000,nan) (0.03250,nan) (0.03500,nan) (0.03750,nan) (0.04000,nan) (0.04250,nan) (0.04500,nan) (0.04750,nan) (0.05000,nan) (0.05250,nan) (0.05500,nan) (0.05750,nan) (0.06000,nan) (0.06250,nan) (0.06500,nan) (0.06750,nan) (0.07000,nan) (0.07250,nan) (0.07500,nan) (0.07750,nan) (0.08000,nan) (0.08250,nan) (0.08500,nan) (0.08750,nan) (0.09000,nan) (0.09250,nan) (0.09500,nan) (0.09750,nan) (0.10000,nan) (0.10250,nan) (0.10500,nan) (0.10750,nan) (0.11000,nan) (0.11250,nan) (0.11500,nan) (0.11750,nan) (0.12000,nan) (0.12250,nan) (0.12500,nan) (0.12750,nan) (0.13000,nan) (0.13250,nan) (0.13500,nan) (0.13750,nan) (0.14000,nan) (0.14250,nan) (0.14500,nan) (0.14750,nan) (0.15000,nan) (0.15250,nan) (0.15500,nan) (0.15750,nan) (0.16000,nan) (0.16250,nan) (0.16500,nan) (0.16750,nan) (0.17000,nan) (0.17250,nan) (0.17500,nan) (0.17750,nan) (0.18000,nan) (0.18250,nan) (0.18500,nan) (0.18750,nan) (0.19000,nan) (0.19250,nan) (0.19500,nan) (0.19750,nan) (0.20000,nan) (0.20250,nan) (0.20500,nan) (0.20750,nan) (0.21000,nan) (0.21250,nan) (0.21500,nan) (0.21750,nan) (0.22000,nan) (0.22250,nan) (0.22500,nan) (0.22750,nan) (0.23000,nan) (0.23250,nan) (0.23500,nan) (0.23750,nan) (0.24000,nan) (0.24250,nan) (0.24500,nan) (0.24750,nan) (0.25000,nan) (0.25250,nan) (0.25500,nan) (0.25750,nan) (0.26000,nan) (0.26250,nan) (0.26500,nan) (0.26750,nan) (0.27000,nan) (0.27250,nan) (0.27500,nan) (0.27750,nan) (0.28000,nan) (0.28250,nan) (0.28500,nan) (0.28750,nan) (0.29000,nan) (0.29250,nan) (0.29500,nan) (0.29750,nan) (0.30000,nan) (0.30250,nan) (0.30500,nan) (0.30750,nan) (0.31000,nan) (0.31250,nan) (0.31500,nan) (0.31750,nan) (0.32000,nan) (0.32250,nan) (0.32500,nan) (0.32750,nan) (0.33000,nan) (0.33250,nan) (0.33500,nan) (0.33750,nan) (0.34000,nan) (0.34250,nan) (0.34500,nan) (0.34750,nan) (0.35000,nan) (0.35250,nan) (0.35500,nan) (0.35750,nan) (0.36000,nan) (0.36250,nan) (0.36500,nan) (0.36750,nan) (0.37000,nan) (0.37250,nan) (0.37500,nan) (0.37750,nan) (0.38000,nan) (0.38250,nan) (0.38500,nan) (0.38750,nan) (0.39000,nan) (0.39250,nan) (0.39500,nan) (0.39750,nan) (0.40000,nan) (0.40250,nan) (0.40500,nan) (0.40750,nan) (0.41000,nan) (0.41250,nan) (0.41500,nan) (0.41750,nan) (0.42000,nan) (0.42250,nan) (0.42500,nan) (0.42750,nan) (0.43000,nan) (0.43250,nan) (0.43500,nan) (0.43750,nan) (0.44000,nan) (0.44250,nan) (0.44500,nan) (0.44750,nan) (0.45000,nan) (0.45250,nan) (0.45500,nan) (0.45750,nan) (0.46000,nan) (0.46250,nan) (0.46500,nan) (0.46750,nan) (0.47000,nan) (0.47250,nan) (0.47500,nan) (0.47750,nan) (0.48000,nan) (0.48250,nan) (0.48500,nan) (0.48750,nan) (0.49000,nan) (0.49250,nan) (0.49500,nan) (0.49750,nan) (0.50000,nan) (0.50250,nan) (0.50500,nan) (0.50750,nan) (0.51000,nan) (0.51250,nan) (0.51500,nan) (0.51750,nan) (0.52000,nan) (0.52250,nan) (0.52500,nan) (0.52750,nan) (0.53000,nan) (0.53250,nan) (0.53500,nan) (0.53750,nan) (0.54000,nan) (0.54250,nan) (0.54500,nan) (0.54750,nan) (0.55000,nan) (0.55250,nan) (0.55500,nan) (0.55750,nan) (0.56000,nan) (0.56250,nan) (0.56500,nan) (0.56750,nan) (0.57000,nan) (0.57250,nan) (0.57500,nan) (0.57750,nan) (0.58000,nan) (0.58250,nan) (0.58500,nan) (0.58750,nan) (0.59000,nan) (0.59250,nan) (0.59500,nan) (0.59750,nan) (0.60000,nan) (0.60250,nan) (0.60500,nan) (0.60750,nan) (0.61000,nan) (0.61250,nan) (0.61500,nan) (0.61750,nan) (0.62000,nan) (0.62250,nan) (0.62500,nan) (0.62750,nan) (0.63000,nan) (0.63250,nan) (0.63500,nan) (0.63750,nan) (0.64000,nan) (0.64250,nan) (0.64500,nan) (0.64750,nan) (0.65000,nan) (0.65250,nan) (0.65500,nan) (0.65750,nan) (0.66000,nan) (0.66250,nan) (0.66500,nan) (0.66750,nan) (0.67000,nan) (0.67250,nan) (0.67500,nan) (0.67750,nan) (0.68000,nan) (0.68250,nan) (0.68500,nan) (0.68750,nan) (0.69000,nan) (0.69250,nan) (0.69500,nan) (0.69750,nan) (0.70000,nan) (0.70250,nan) (0.70500,nan) (0.70750,nan) (0.71000,nan) (0.71250,nan) (0.71500,nan) (0.71750,nan) (0.72000,nan) (0.72250,nan) (0.72500,nan) (0.72750,nan) (0.73000,nan) (0.73250,nan) (0.73500,nan) (0.73750,nan) (0.74000,nan) (0.74250,nan) (0.74500,nan) (0.74750,nan) (0.75000,nan) (0.75250,nan) (0.75500,nan) (0.75750,nan) (0.76000,nan) (0.76250,nan) (0.76500,nan) (0.76750,nan) (0.77000,nan) (0.77250,nan) (0.77500,nan) (0.77750,nan) (0.78000,nan) (0.78250,nan) (0.78500,nan) (0.78750,nan) (0.79000,nan) (0.79250,nan) (0.79500,nan) (0.79750,nan) (0.80000,nan) (0.80250,nan) (0.80500,nan) (0.80750,nan) (0.81000,nan) (0.81250,nan) (0.81500,nan) (0.81750,nan) (0.82000,nan) (0.82250,nan) (0.82500,nan) (0.82750,nan) (0.83000,nan) (0.83250,nan) (0.83500,nan) (0.83750,nan) (0.84000,nan) (0.84250,nan) (0.84500,nan) (0.84750,nan) (0.85000,nan) (0.85250,nan) (0.85500,nan) (0.85750,nan) (0.86000,nan) (0.86250,nan) (0.86500,nan) (0.86750,nan) (0.87000,nan) (0.87250,nan) (0.87500,nan) (0.87750,nan) (0.88000,nan) (0.88250,nan) (0.88500,nan) (0.88750,nan) (0.89000,nan) (0.89250,nan) (0.89500,nan) (0.89750,nan) (0.90000,nan) (0.90250,nan) (0.90500,nan) (0.90750,nan) (0.91000,nan) (0.91250,nan) (0.91500,nan) (0.91750,nan) (0.92000,nan) (0.92250,nan) (0.92500,nan) (0.92750,nan) (0.93000,nan) (0.93250,nan) (0.93500,nan) (0.93750,nan) (0.94000,nan) (0.94250,nan) (0.94500,nan) (0.94750,nan) (0.95000,nan) (0.95250,nan) (0.95500,nan) (0.95750,nan) (0.96000,nan) (0.96250,nan) (0.96500,nan) (0.96750,nan) (0.97000,nan) (0.97250,nan) (0.97500,nan) (0.97750,nan) (0.98000,nan) (0.98250,nan) (0.98500,nan) (0.98750,nan) (0.99000,nan) (0.99250,nan) (0.99500,nan) (0.99750,nan) (1.00000,nan) (1.00250,nan) (1.00500,nan) (1.00750,nan) (1.01000,nan) (1.01250,nan) (1.01500,nan) (1.01750,nan) (1.02000,nan) (1.02250,nan) (1.02500,nan) (1.02750,nan) (1.03000,nan) (1.03250,nan) (1.03500,nan) (1.03750,nan) (1.04000,nan) (1.04250,nan) (1.04500,nan) (1.04750,nan) (1.05000,nan) (1.05250,nan) (1.05500,nan) (1.05750,nan) (1.06000,nan) (1.06250,nan) (1.06500,nan) (1.06750,nan) (1.07000,nan) (1.07250,nan) (1.07500,nan) (1.07750,nan) (1.08000,nan) (1.08250,nan) (1.08500,nan) (1.08750,nan) (1.09000,nan) (1.09250,nan) (1.09500,nan) (1.09750,nan) (1.10000,nan) (1.10250,nan) (1.10500,nan) (1.10750,nan) (1.11000,nan) (1.11250,nan) (1.11500,nan) (1.11750,nan) (1.12000,nan) (1.12250,nan) (1.12500,nan) (1.12750,nan) (1.13000,nan) (1.13250,nan) (1.13500,nan) (1.13750,nan) (1.14000,nan) (1.14250,nan) (1.14500,nan) (1.14750,nan) (1.15000,nan) (1.15250,nan) (1.15500,nan) (1.15750,nan) (1.16000,nan) (1.16250,nan) (1.16500,nan) (1.16750,nan) (1.17000,nan) (1.17250,nan) (1.17500,nan) (1.17750,nan) (1.18000,nan) (1.18250,nan) (1.18500,nan) (1.18750,nan) (1.19000,nan) (1.19250,nan) (1.19500,nan) (1.19750,nan) (1.20000,nan) (1.20250,nan) (1.20500,nan) (1.20750,nan) (1.21000,nan) (1.21250,nan) (1.21500,nan) (1.21750,nan) (1.22000,nan) (1.22250,nan) (1.22500,nan) (1.22750,nan) (1.23000,nan) (1.23250,nan) (1.23500,nan) (1.23750,nan) (1.24000,nan) (1.24250,nan) (1.24500,nan) (1.24750,nan) (1.25000,nan) (1.25250,nan) (1.25500,nan) (1.25750,nan) (1.26000,nan) (1.26250,nan) (1.26500,nan) (1.26750,nan) (1.27000,nan) (1.27250,nan) (1.27500,nan) (1.27750,nan) (1.28000,nan) (1.28250,nan) (1.28500,nan) (1.28750,nan) (1.29000,nan) (1.29250,nan) (1.29500,nan) (1.29750,nan) (1.30000,nan) (1.30250,nan) (1.30500,nan) (1.30750,nan) (1.31000,nan) (1.31250,nan) (1.31500,nan) (1.31750,nan) (1.32000,nan) (1.32250,nan) (1.32500,nan) (1.32750,nan) (1.33000,nan) (1.33250,nan) (1.33500,nan) (1.33750,nan) (1.34000,nan) (1.34250,nan) (1.34500,nan) (1.34750,nan) (1.35000,nan) (1.35250,nan) (1.35500,nan) (1.35750,nan) (1.36000,nan) (1.36250,nan) (1.36500,nan) (1.36750,nan) (1.37000,nan) (1.37250,nan) (1.37500,nan) (1.37750,nan) (1.38000,nan) (1.38250,nan) (1.38500,nan) (1.38750,nan) (1.39000,nan) (1.39250,nan) (1.39500,nan) (1.39750,nan) (1.40000,nan) (1.40250,nan) (1.40500,nan) (1.40750,nan) (1.41000,nan) (1.41250,nan) (1.41500,nan) (1.41750,nan) (1.42000,nan) (1.42250,nan) (1.42500,nan) (1.42750,nan) (1.43000,nan) (1.43250,nan) (1.43500,nan) (1.43750,nan) (1.44000,nan) (1.44250,nan) (1.44500,nan) (1.44750,nan) (1.45000,nan) (1.45250,nan) (1.45500,nan) (1.45750,nan) (1.46000,nan) (1.46250,nan) (1.46500,nan) (1.46750,nan) (1.47000,nan) (1.47250,nan) (1.47500,nan) (1.47750,nan) (1.48000,nan) (1.48250,nan) (1.48500,nan) (1.48750,nan) (1.49000,nan) (1.49250,nan) (1.49500,nan) (1.49750,nan) (1.50000,nan) (1.50250,nan) (1.50500,nan) (1.50750,nan) (1.51000,nan) (1.51250,nan) (1.51500,nan) (1.51750,nan) (1.52000,nan) (1.52250,nan) (1.52500,nan) (1.52750,nan) (1.53000,nan) (1.53250,nan) (1.53500,nan) (1.53750,nan) (1.54000,nan) (1.54250,nan) (1.54500,nan) (1.54750,nan) (1.55000,nan) (1.55250,nan) (1.55500,nan) (1.55750,nan) (1.56000,nan) (1.56250,nan) (1.56500,nan) (1.56750,nan) (1.57000,nan) (1.57250,nan) (1.57500,nan) (1.57750,nan) (1.58000,nan) (1.58250,nan) (1.58500,nan) (1.58750,nan) (1.59000,nan) (1.59250,nan) (1.59500,nan) (1.59750,nan) (1.60000,nan) (1.60250,nan) (1.60500,nan) (1.60750,nan) (1.61000,nan) (1.61250,nan) (1.61500,nan) (1.61750,nan) (1.62000,nan) (1.62250,nan) (1.62500,nan) (1.62750,nan) (1.63000,nan) (1.63250,nan) (1.63500,nan) (1.63750,nan) (1.64000,nan) (1.64250,nan) (1.64500,nan) (1.64750,nan) (1.65000,nan) (1.65250,nan) (1.65500,nan) (1.65750,nan) (1.66000,nan) (1.66250,nan) (1.66500,nan) (1.66750,nan) (1.67000,nan) (1.67250,nan) (1.67500,nan) (1.67750,nan) (1.68000,nan) (1.68250,nan) (1.68500,nan) (1.68750,nan) (1.69000,nan) (1.69250,nan) (1.69500,nan) (1.69750,nan) (1.70000,nan) (1.70250,nan) (1.70500,nan) (1.70750,nan) (1.71000,nan) (1.71250,nan) (1.71500,nan) (1.71750,nan) (1.72000,nan) (1.72250,nan) (1.72500,nan) (1.72750,nan) (1.73000,nan) (1.73250,nan) (1.73500,nan) (1.73750,nan) (1.74000,nan) (1.74250,nan) (1.74500,nan) (1.74750,nan) (1.75000,nan) (1.75250,nan) (1.75500,nan) (1.75750,nan) (1.76000,nan) (1.76250,nan) (1.76500,nan) (1.76750,nan) (1.77000,nan) (1.77250,nan) (1.77500,nan) (1.77750,nan) (1.78000,nan) (1.78250,nan) (1.78500,nan) (1.78750,nan) (1.79000,nan) (1.79250,nan) (1.79500,nan) (1.79750,nan) (1.80000,nan) (1.80250,nan) (1.80500,nan) (1.80750,nan) (1.81000,nan) (1.81250,nan) (1.81500,nan) (1.81750,nan) (1.82000,nan) (1.82250,nan) (1.82500,nan) (1.82750,nan) (1.83000,nan) (1.83250,nan) (1.83500,nan) (1.83750,nan) (1.84000,nan) (1.84250,nan) (1.84500,nan) (1.84750,nan) (1.85000,nan) (1.85250,nan) (1.85500,nan) (1.85750,nan) (1.86000,nan) (1.86250,nan) (1.86500,nan) (1.86750,nan) (1.87000,nan) (1.87250,nan) (1.87500,nan) (1.87750,nan) (1.88000,nan) (1.88250,nan) (1.88500,nan) (1.88750,nan) (1.89000,nan) (1.89250,nan) (1.89500,nan) (1.89750,nan) (1.90000,nan) (1.90250,nan) (1.90500,nan) (1.90750,nan) (1.91000,nan) (1.91250,nan) (1.91500,nan) (1.91750,nan) (1.92000,nan) (1.92250,nan) (1.92500,nan) (1.92750,nan) (1.93000,nan) (1.93250,nan) (1.93500,nan) (1.93750,nan) (1.94000,nan) (1.94250,nan) (1.94500,nan) (1.94750,nan) (1.95000,nan) (1.95250,nan) (1.95500,nan) (1.95750,nan) (1.96000,nan) (1.96250,nan) (1.96500,nan) (1.96750,nan) (1.97000,nan) (1.97250,nan) (1.97500,nan) (1.97750,nan) (1.98000,nan) (1.98250,nan) (1.98500,nan) (1.98750,nan) (1.99000,nan) (1.99250,nan) (1.99500,nan) (1.99750,nan) (2.00000,nan) (2.00250,nan) (2.00500,nan) (2.00750,nan) (2.01000,nan) (2.01250,nan) (2.01500,nan) (2.01750,nan) (2.02000,nan) (2.02250,nan) (2.02500,nan) (2.02750,nan) (2.03000,nan) (2.03250,nan) (2.03500,nan) (2.03750,nan) (2.04000,nan) (2.04250,nan) (2.04500,nan) (2.04750,nan) (2.05000,nan) (2.05250,nan) (2.05500,nan) (2.05750,nan) (2.06000,nan) (2.06250,nan) (2.06500,nan) (2.06750,nan) (2.07000,nan) (2.07250,nan) (2.07500,nan) (2.07750,nan) (2.08000,nan) (2.08250,nan) (2.08500,nan) (2.08750,nan) (2.09000,nan) (2.09250,nan) (2.09500,nan) (2.09750,nan) (2.10000,nan) (2.10250,nan) (2.10500,nan) (2.10750,nan) (2.11000,nan) (2.11250,nan) (2.11500,nan) (2.11750,nan) (2.12000,nan) (2.12250,nan) (2.12500,nan) (2.12750,nan) (2.13000,nan) (2.13250,nan) (2.13500,nan) (2.13750,nan) (2.14000,nan) (2.14250,nan) (2.14500,nan) (2.14750,nan) (2.15000,nan) (2.15250,nan) (2.15500,nan) (2.15750,nan) (2.16000,nan) (2.16250,nan) (2.16500,nan) (2.16750,nan) (2.17000,nan) (2.17250,nan) (2.17500,nan) (2.17750,nan) (2.18000,nan) (2.18250,nan) (2.18500,nan) (2.18750,nan) (2.19000,nan) (2.19250,nan) (2.19500,nan) (2.19750,nan) (2.20000,nan) (2.20250,nan) (2.20500,nan) (2.20750,nan) (2.21000,nan) (2.21250,nan) (2.21500,nan) (2.21750,nan) (2.22000,nan) (2.22250,nan) (2.22500,nan) (2.22750,nan) (2.23000,nan) (2.23250,nan) (2.23500,nan) (2.23750,nan) (2.24000,nan) (2.24250,nan) (2.24500,nan) (2.24750,nan) (2.25000,nan) (2.25250,nan) (2.25500,nan) (2.25750,nan) (2.26000,nan) (2.26250,nan) (2.26500,nan) (2.26750,nan) (2.27000,nan) (2.27250,nan) (2.27500,nan) (2.27750,nan) (2.28000,nan) (2.28250,nan) (2.28500,nan) (2.28750,nan) (2.29000,nan) (2.29250,nan) (2.29500,nan) (2.29750,nan) (2.30000,nan) (2.30250,nan) (2.30500,nan) (2.30750,nan) (2.31000,nan) (2.31250,nan) (2.31500,nan) (2.31750,nan) (2.32000,nan) (2.32250,nan) (2.32500,nan) (2.32750,nan) (2.33000,nan) (2.33250,nan) (2.33500,nan) (2.33750,nan) (2.34000,nan) (2.34250,nan) (2.34500,nan) (2.34750,nan) (2.35000,nan) (2.35250,nan) (2.35500,nan) (2.35750,nan) (2.36000,nan) (2.36250,nan) (2.36500,nan) (2.36750,nan) (2.37000,nan) (2.37250,nan) (2.37500,nan) (2.37750,nan) (2.38000,nan) (2.38250,nan) (2.38500,nan) (2.38750,nan) (2.39000,nan) (2.39250,nan) (2.39500,nan) (2.39750,nan) (2.40000,nan) (2.40250,nan) (2.40500,nan) (2.40750,nan) (2.41000,nan) (2.41250,nan) (2.41500,nan) (2.41750,nan) (2.42000,nan) (2.42250,nan) (2.42500,nan) (2.42750,nan) (2.43000,nan) (2.43250,nan) (2.43500,nan) (2.43750,nan) (2.44000,nan) (2.44250,nan) (2.44500,nan) (2.44750,nan) (2.45000,nan) (2.45250,nan) (2.45500,nan) (2.45750,nan) (2.46000,nan) (2.46250,nan) (2.46500,nan) (2.46750,nan) (2.47000,nan) (2.47250,nan) (2.47500,nan) (2.47750,nan) (2.48000,nan) (2.48250,nan) (2.48500,nan) (2.48750,nan) (2.49000,nan) (2.49250,nan) (2.49500,nan) (2.49750,nan) (2.50000,nan) (2.50250,nan) (2.50500,nan) (2.50750,nan) (2.51000,nan) (2.51250,nan) (2.51500,nan) (2.51750,nan) (2.52000,nan) (2.52250,nan) (2.52500,nan) (2.52750,nan) (2.53000,nan) (2.53250,nan) (2.53500,nan) (2.53750,nan) (2.54000,nan) (2.54250,nan) (2.54500,nan) (2.54750,nan) (2.55000,nan) (2.55250,nan) (2.55500,nan) (2.55750,nan) (2.56000,nan) (2.56250,nan) (2.56500,nan) (2.56750,nan) (2.57000,nan) (2.57250,nan) (2.57500,nan) (2.57750,nan) (2.58000,nan) (2.58250,nan) (2.58500,nan) (2.58750,nan) (2.59000,nan) (2.59250,nan) (2.59500,nan) (2.59750,nan) (2.60000,nan) (2.60250,nan) (2.60500,nan) (2.60750,nan) (2.61000,nan) (2.61250,nan) (2.61500,nan) (2.61750,nan) (2.62000,nan) (2.62250,nan) (2.62500,nan) (2.62750,nan) (2.63000,nan) (2.63250,nan) (2.63500,nan) (2.63750,nan) (2.64000,nan) (2.64250,nan) (2.64500,nan) (2.64750,nan) (2.65000,nan) (2.65250,nan) (2.65500,nan) (2.65750,nan) (2.66000,nan) (2.66250,nan) (2.66500,nan) (2.66750,nan) (2.67000,nan) (2.67250,nan) (2.67500,nan) (2.67750,nan) (2.68000,nan) (2.68250,nan) (2.68500,nan) (2.68750,nan) (2.69000,nan) (2.69250,nan) (2.69500,nan) (2.69750,nan) (2.70000,nan) (2.70250,nan) (2.70500,nan) (2.70750,nan) (2.71000,nan) (2.71250,nan) (2.71500,nan) (2.71750,nan) (2.72000,nan) (2.72250,nan) (2.72500,nan) (2.72750,nan) (2.73000,nan) (2.73250,nan) (2.73500,nan) (2.73750,nan) (2.74000,nan) (2.74250,nan) (2.74500,nan) (2.74750,nan) (2.75000,nan) (2.75250,nan) (2.75500,nan) (2.75750,nan) (2.76000,nan) (2.76250,nan) (2.76500,nan) (2.76750,nan) (2.77000,nan) (2.77250,nan) (2.77500,nan) (2.77750,nan) (2.78000,nan) (2.78250,nan) (2.78500,nan) (2.78750,nan) (2.79000,nan) (2.79250,nan) (2.79500,nan) (2.79750,nan) (2.80000,nan) (2.80250,nan) (2.80500,nan) (2.80750,nan) (2.81000,nan) (2.81250,nan) (2.81500,nan) (2.81750,nan) (2.82000,nan) (2.82250,nan) (2.82500,nan) (2.82750,nan) (2.83000,nan) (2.83250,nan) (2.83500,nan) (2.83750,nan) (2.84000,nan) (2.84250,nan) (2.84500,nan) (2.84750,nan) (2.85000,nan) (2.85250,nan) (2.85500,nan) (2.85750,nan) (2.86000,nan) (2.86250,nan) (2.86500,nan) (2.86750,nan) (2.87000,nan) (2.87250,nan) (2.87500,nan) (2.87750,nan) (2.88000,nan) (2.88250,nan) (2.88500,nan) (2.88750,nan) (2.89000,nan) (2.89250,nan) (2.89500,nan) (2.89750,nan) (2.90000,nan) (2.90250,nan) (2.90500,nan) (2.90750,nan) (2.91000,nan) (2.91250,nan) (2.91500,nan) (2.91750,nan) (2.92000,nan) (2.92250,nan) (2.92500,nan) (2.92750,nan) (2.93000,nan) (2.93250,nan) (2.93500,nan) (2.93750,nan) (2.94000,nan) (2.94250,nan) (2.94500,nan) (2.94750,nan) (2.95000,nan) (2.95250,nan) (2.95500,nan) (2.95750,nan) (2.96000,nan) (2.96250,nan) (2.96500,nan) (2.96750,nan) (2.97000,nan) (2.97250,nan) (2.97500,nan) (2.97750,nan) (2.98000,nan) (2.98250,nan) (2.98500,nan) (2.98750,nan) (2.99000,nan) (2.99250,nan) (2.99500,nan) (2.99750,nan) (3.00000,nan) (3.00250,nan) (3.00500,nan) (3.00750,nan) (3.01000,nan) (3.01250,nan) (3.01500,nan) (3.01750,nan) (3.02000,nan) (3.02250,nan) (3.02500,nan) (3.02750,nan) (3.03000,nan) (3.03250,nan) (3.03500,nan) (3.03750,nan) (3.04000,nan) (3.04250,nan) (3.04500,nan) (3.04750,nan) (3.05000,nan) (3.05250,nan) (3.05500,nan) (3.05750,nan) (3.06000,nan) (3.06250,nan) (3.06500,nan) (3.06750,nan) (3.07000,nan) (3.07250,nan) (3.07500,nan) (3.07750,nan) (3.08000,nan) (3.08250,nan) (3.08500,nan) (3.08750,nan) (3.09000,nan) (3.09250,nan) (3.09500,nan) (3.09750,nan) (3.10000,nan) (3.10250,nan) (3.10500,nan) (3.10750,nan) (3.11000,nan) (3.11250,nan) (3.11500,nan) (3.11750,nan) (3.12000,nan) (3.12250,nan) (3.12500,nan) (3.12750,nan) (3.13000,nan) (3.13250,nan) (3.13500,nan) (3.13750,nan) (3.14000,nan) (3.14250,nan) (3.14500,nan) (3.14750,nan) (3.15000,nan) (3.15250,nan) (3.15500,nan) (3.15750,nan) (3.16000,nan) (3.16250,nan) (3.16500,nan) (3.16750,nan) (3.17000,nan) (3.17250,nan) (3.17500,nan) (3.17750,nan) (3.18000,nan) (3.18250,nan) (3.18500,nan) (3.18750,nan) (3.19000,nan) (3.19250,nan) (3.19500,nan) (3.19750,nan) (3.20000,nan) (3.20250,nan) (3.20500,nan) (3.20750,nan) (3.21000,nan) (3.21250,nan) (3.21500,nan) (3.21750,nan) (3.22000,nan) (3.22250,nan) (3.22500,nan) (3.22750,nan) (3.23000,nan) (3.23250,nan) (3.23500,nan) (3.23750,-1.7954316) (3.24000,-1.6294579) (3.24250,-1.5103634) (3.24500,-1.4175567) (3.24750,-1.3415991) (3.25000,-1.2773697) (3.25250,-1.2217769) (3.25500,-1.1728105) (3.25750,-1.1290893) (3.26000,-1.0896236) (3.26250,-1.0536796) (3.26500,-1.0206985) (3.26750,-0.9902446) (3.27000,-0.9619718) (3.27250,-0.9356005) (3.27500,-0.9109018) (3.27750,-0.8876855) (3.28000,-0.8657923) (3.28250,-0.8450873) (3.28500,-0.8254549) (3.28750,-0.8067959) (3.29000,-0.7890242) (3.29250,-0.7720646) (3.29500,-0.7558509) (3.29750,-0.7403250) (3.30000,-0.7254350) (3.30250,-0.7111348) (3.30500,-0.6973832) (3.30750,-0.6841431) (3.31000,-0.6713811) (3.31250,-0.6590667) (3.31500,-0.6471725) (3.31750,-0.6356736) (3.32000,-0.6245469) (3.32250,-0.6137716) (3.32500,-0.6033284) (3.32750,-0.5931996) (3.33000,-0.5833689) (3.33250,-0.5738212) (3.33500,-0.5645425) (3.33750,-0.5555199) (3.34000,-0.5467413) (3.34250,-0.5381955) (3.34500,-0.5298719) (3.34750,-0.5217609) (3.35000,-0.5138533) (3.35250,-0.5061405) (3.35500,-0.4986144) (3.35750,-0.4912675) (3.36000,-0.4840926) (3.36250,-0.4770832) (3.36500,-0.4702327) (3.36750,-0.4635354) (3.37000,-0.4569855) (3.37250,-0.4505778) (3.37500,-0.4443071) (3.37750,-0.4381687) (3.38000,-0.4321580) (3.38250,-0.4262707) (3.38500,-0.4205027) (3.38750,-0.4148502) (3.39000,-0.4093094) (3.39250,-0.4038767) (3.39500,-0.3985488) (3.39750,-0.3933225) (3.40000,-0.3881946) (3.40250,-0.3831623) (3.40500,-0.3782228) (3.40750,-0.3733732) (3.41000,-0.3686112) (3.41250,-0.3639341) (3.41500,-0.3593397) (3.41750,-0.3548256) (3.42000,-0.3503898) (3.42250,-0.3460300) (3.42500,-0.3417442) (3.42750,-0.3375307) (3.43000,-0.3333874) (3.43250,-0.3293125) (3.43500,-0.3253045) (3.43750,-0.3213616) (3.44000,-0.3174822) (3.44250,-0.3136647) (3.44500,-0.3099077) (3.44750,-0.3062098) (3.45000,-0.3025695) (3.45250,-0.2989855) (3.45500,-0.2954566) (3.45750,-0.2919814) (3.46000,-0.2885588) (3.46250,-0.2851876) (3.46500,-0.2818666) (3.46750,-0.2785948) (3.47000,-0.2753711) (3.47250,-0.2721944) (3.47500,-0.2690639) (3.47750,-0.2659784) (3.48000,-0.2629372) (3.48250,-0.2599392) (3.48500,-0.2569835) (3.48750,-0.2540694) (3.49000,-0.2511960) (3.49250,-0.2483625) (3.49500,-0.2455681) (3.49750,-0.2428121) (3.50000,-0.2400936) (3.50250,-0.2374121) (3.50500,-0.2347667) (3.50750,-0.2321568) (3.51000,-0.2295818) (3.51250,-0.2270410) (3.51500,-0.2245337) (3.51750,-0.2220594) (3.52000,-0.2196175) (3.52250,-0.2172074) (3.52500,-0.2148284) (3.52750,-0.2124802) (3.53000,-0.2101620) (3.53250,-0.2078735) (3.53500,-0.2056141) (3.53750,-0.2033833) (3.54000,-0.2011806) (3.54250,-0.1990055) (3.54500,-0.1968576) (3.54750,-0.1947364) (3.55000,-0.1926415) (3.55250,-0.1905725) (3.55500,-0.1885289) (3.55750,-0.1865103) (3.56000,-0.1845164) (3.56250,-0.1825466) (3.56500,-0.1806007) (3.56750,-0.1786783) (3.57000,-0.1767789) (3.57250,-0.1749023) (3.57500,-0.1730481) (3.57750,-0.1712160) (3.58000,-0.1694055) (3.58250,-0.1676164) (3.58500,-0.1658484) (3.58750,-0.1641011) (3.59000,-0.1623742) (3.59250,-0.1606675) (3.59500,-0.1589806) (3.59750,-0.1573132) (3.60000,-0.1556651) (3.60250,-0.1540360) (3.60500,-0.1524255) (3.60750,-0.1508335) (3.61000,-0.1492597) (3.61250,-0.1477038) (3.61500,-0.1461656) (3.61750,-0.1446448) (3.62000,-0.1431411) (3.62250,-0.1416544) (3.62500,-0.1401844) (3.62750,-0.1387308) (3.63000,-0.1372935) (3.63250,-0.1358723) (3.63500,-0.1344668) (3.63750,-0.1330769) (3.64000,-0.1317024) (3.64250,-0.1303431) (3.64500,-0.1289988) (3.64750,-0.1276693) (3.65000,-0.1263543) (3.65250,-0.1250538) (3.65500,-0.1237674) (3.65750,-0.1224951) (3.66000,-0.1212367) (3.66250,-0.1199918) (3.66500,-0.1187605) (3.66750,-0.1175425) (3.67000,-0.1163377) (3.67250,-0.1151458) (3.67500,-0.1139667) (3.67750,-0.1128003) (3.68000,-0.1116465) (3.68250,-0.1105049) (3.68500,-0.1093755) (3.68750,-0.1082582) (3.69000,-0.1071528) (3.69250,-0.1060591) (3.69500,-0.1049770) (3.69750,-0.1039064) (3.70000,-0.1028471) (3.70250,-0.1017990) (3.70500,-0.1007620) (3.70750,-0.0997358) (3.71000,-0.0987205) (3.71250,-0.0977158) (3.71500,-0.0967216) (3.71750,-0.0957379) (3.72000,-0.0947644) (3.72250,-0.0938011) (3.72500,-0.0928479) (3.72750,-0.0919046) (3.73000,-0.0909711) (3.73250,-0.0900473) (3.73500,-0.0891330) (3.73750,-0.0882283) (3.74000,-0.0873329) (3.74250,-0.0864468) (3.74500,-0.0855698) (3.74750,-0.0847019) (3.75000,-0.0838429) (3.75250,-0.0829928) (3.75500,-0.0821514) (3.75750,-0.0813187) (3.76000,-0.0804945) (3.76250,-0.0796788) (3.76500,-0.0788714) (3.76750,-0.0780723) (3.77000,-0.0772814) (3.77250,-0.0764986) (3.77500,-0.0757237) (3.77750,-0.0749568) (3.78000,-0.0741976) (3.78250,-0.0734462) (3.78500,-0.0727025) (3.78750,-0.0719663) (3.79000,-0.0712376) (3.79250,-0.0705163) (3.79500,-0.0698023) (3.79750,-0.0690956) (3.80000,-0.0683960) (3.80250,-0.0677035) (3.80500,-0.0670180) (3.80750,-0.0663395) (3.81000,-0.0656678) (3.81250,-0.0650029) (3.81500,-0.0643447) (3.81750,-0.0636932) (3.82000,-0.0630482) (3.82250,-0.0624097) (3.82500,-0.0617777) (3.82750,-0.0611520) (3.83000,-0.0605326) (3.83250,-0.0599194) (3.83500,-0.0593124) (3.83750,-0.0587116) (3.84000,-0.0581167) (3.84250,-0.0575278) (3.84500,-0.0569449) (3.84750,-0.0563677) (3.85000,-0.0557964) (3.85250,-0.0552308) (3.85500,-0.0546709) (3.85750,-0.0541165) (3.86000,-0.0535678) (3.86250,-0.0530245) (3.86500,-0.0524867) (3.86750,-0.0519542) (3.87000,-0.0514271) (3.87250,-0.0509052) (3.87500,-0.0503886) (3.87750,-0.0498771) (3.88000,-0.0493707) (3.88250,-0.0488694) (3.88500,-0.0483731) (3.88750,-0.0478818) (3.89000,-0.0473953) (3.89250,-0.0469137) (3.89500,-0.0464370) (3.89750,-0.0459649) (3.90000,-0.0454976) (3.90250,-0.0450350) (3.90500,-0.0445769) (3.90750,-0.0441235) (3.91000,-0.0436745) (3.91250,-0.0432301) (3.91500,-0.0427900) (3.91750,-0.0423544) (3.92000,-0.0419231) (3.92250,-0.0414961) (3.92500,-0.0410733) (3.92750,-0.0406548) (3.93000,-0.0402404) (3.93250,-0.0398302) (3.93500,-0.0394240) (3.93750,-0.0390219) (3.94000,-0.0386238) (3.94250,-0.0382297) (3.94500,-0.0378395) (3.94750,-0.0374531) (3.95000,-0.0370707) (3.95250,-0.0366920) (3.95500,-0.0363171) (3.95750,-0.0359460) (3.96000,-0.0355785) (3.96250,-0.0352147) (3.96500,-0.0348546) (3.96750,-0.0344980) (3.97000,-0.0341450) (3.97250,-0.0337955) (3.97500,-0.0334495) (3.97750,-0.0331069) (3.98000,-0.0327677) (3.98250,-0.0324320) (3.98500,-0.0320996) (3.98750,-0.0317705) (3.99000,-0.0314446) (3.99250,-0.0311221) (3.99500,-0.0308027) (3.99750,-0.0304865) (4.00000,-0.0301735)};\addlegendentry{L}
\addplot[thick,dashed,gray] coordinates {(-2.00000,nan) (-1.99750,nan) (-1.99500,nan) (-1.99250,nan) (-1.99000,nan) (-1.98750,nan) (-1.98500,nan) (-1.98250,nan) (-1.98000,nan) (-1.97750,nan) (-1.97500,nan) (-1.97250,nan) (-1.97000,nan) (-1.96750,nan) (-1.96500,nan) (-1.96250,nan) (-1.96000,nan) (-1.95750,nan) (-1.95500,nan) (-1.95250,nan) (-1.95000,nan) (-1.94750,nan) (-1.94500,nan) (-1.94250,nan) (-1.94000,nan) (-1.93750,nan) (-1.93500,nan) (-1.93250,nan) (-1.93000,nan) (-1.92750,nan) (-1.92500,nan) (-1.92250,nan) (-1.92000,nan) (-1.91750,nan) (-1.91500,nan) (-1.91250,nan) (-1.91000,nan) (-1.90750,nan) (-1.90500,nan) (-1.90250,nan) (-1.90000,nan) (-1.89750,nan) (-1.89500,nan) (-1.89250,nan) (-1.89000,nan) (-1.88750,nan) (-1.88500,nan) (-1.88250,nan) (-1.88000,nan) (-1.87750,nan) (-1.87500,nan) (-1.87250,nan) (-1.87000,nan) (-1.86750,nan) (-1.86500,nan) (-1.86250,nan) (-1.86000,nan) (-1.85750,nan) (-1.85500,nan) (-1.85250,nan) (-1.85000,nan) (-1.84750,nan) (-1.84500,nan) (-1.84250,nan) (-1.84000,nan) (-1.83750,nan) (-1.83500,nan) (-1.83250,nan) (-1.83000,nan) (-1.82750,nan) (-1.82500,nan) (-1.82250,nan) (-1.82000,nan) (-1.81750,nan) (-1.81500,nan) (-1.81250,nan) (-1.81000,nan) (-1.80750,nan) (-1.80500,nan) (-1.80250,nan) (-1.80000,nan) (-1.79750,nan) (-1.79500,nan) (-1.79250,nan) (-1.79000,nan) (-1.78750,nan) (-1.78500,nan) (-1.78250,nan) (-1.78000,nan) (-1.77750,nan) (-1.77500,nan) (-1.77250,nan) (-1.77000,nan) (-1.76750,nan) (-1.76500,nan) (-1.76250,nan) (-1.76000,nan) (-1.75750,nan) (-1.75500,nan) (-1.75250,nan) (-1.75000,nan) (-1.74750,nan) (-1.74500,nan) (-1.74250,nan) (-1.74000,nan) (-1.73750,nan) (-1.73500,nan) (-1.73250,nan) (-1.73000,nan) (-1.72750,nan) (-1.72500,nan) (-1.72250,nan) (-1.72000,nan) (-1.71750,nan) (-1.71500,nan) (-1.71250,nan) (-1.71000,nan) (-1.70750,nan) (-1.70500,nan) (-1.70250,nan) (-1.70000,nan) (-1.69750,nan) (-1.69500,nan) (-1.69250,nan) (-1.69000,nan) (-1.68750,nan) (-1.68500,nan) (-1.68250,nan) (-1.68000,nan) (-1.67750,nan) (-1.67500,nan) (-1.67250,nan) (-1.67000,nan) (-1.66750,nan) (-1.66500,nan) (-1.66250,nan) (-1.66000,nan) (-1.65750,nan) (-1.65500,nan) (-1.65250,nan) (-1.65000,nan) (-1.64750,nan) (-1.64500,nan) (-1.64250,nan) (-1.64000,nan) (-1.63750,nan) (-1.63500,nan) (-1.63250,nan) (-1.63000,nan) (-1.62750,nan) (-1.62500,nan) (-1.62250,nan) (-1.62000,nan) (-1.61750,nan) (-1.61500,nan) (-1.61250,nan) (-1.61000,nan) (-1.60750,nan) (-1.60500,nan) (-1.60250,nan) (-1.60000,nan) (-1.59750,nan) (-1.59500,nan) (-1.59250,nan) (-1.59000,nan) (-1.58750,nan) (-1.58500,nan) (-1.58250,nan) (-1.58000,nan) (-1.57750,nan) (-1.57500,nan) (-1.57250,nan) (-1.57000,nan) (-1.56750,nan) (-1.56500,nan) (-1.56250,nan) (-1.56000,nan) (-1.55750,nan) (-1.55500,nan) (-1.55250,nan) (-1.55000,nan) (-1.54750,nan) (-1.54500,nan) (-1.54250,nan) (-1.54000,nan) (-1.53750,nan) (-1.53500,nan) (-1.53250,nan) (-1.53000,nan) (-1.52750,nan) (-1.52500,nan) (-1.52250,nan) (-1.52000,nan) (-1.51750,nan) (-1.51500,nan) (-1.51250,nan) (-1.51000,nan) (-1.50750,nan) (-1.50500,nan) (-1.50250,nan) (-1.50000,nan) (-1.49750,nan) (-1.49500,nan) (-1.49250,nan) (-1.49000,nan) (-1.48750,nan) (-1.48500,nan) (-1.48250,nan) (-1.48000,nan) (-1.47750,nan) (-1.47500,nan) (-1.47250,nan) (-1.47000,nan) (-1.46750,nan) (-1.46500,nan) (-1.46250,nan) (-1.46000,nan) (-1.45750,nan) (-1.45500,nan) (-1.45250,nan) (-1.45000,nan) (-1.44750,nan) (-1.44500,nan) (-1.44250,nan) (-1.44000,nan) (-1.43750,nan) (-1.43500,nan) (-1.43250,nan) (-1.43000,nan) (-1.42750,nan) (-1.42500,nan) (-1.42250,nan) (-1.42000,nan) (-1.41750,nan) (-1.41500,nan) (-1.41250,nan) (-1.41000,nan) (-1.40750,nan) (-1.40500,nan) (-1.40250,nan) (-1.40000,nan) (-1.39750,nan) (-1.39500,nan) (-1.39250,nan) (-1.39000,nan) (-1.38750,nan) (-1.38500,nan) (-1.38250,nan) (-1.38000,nan) (-1.37750,nan) (-1.37500,nan) (-1.37250,nan) (-1.37000,nan) (-1.36750,nan) (-1.36500,nan) (-1.36250,nan) (-1.36000,nan) (-1.35750,nan) (-1.35500,nan) (-1.35250,nan) (-1.35000,nan) (-1.34750,nan) (-1.34500,nan) (-1.34250,nan) (-1.34000,nan) (-1.33750,nan) (-1.33500,nan) (-1.33250,nan) (-1.33000,nan) (-1.32750,nan) (-1.32500,nan) (-1.32250,nan) (-1.32000,nan) (-1.31750,nan) (-1.31500,nan) (-1.31250,nan) (-1.31000,nan) (-1.30750,nan) (-1.30500,nan) (-1.30250,nan) (-1.30000,nan) (-1.29750,nan) (-1.29500,nan) (-1.29250,nan) (-1.29000,nan) (-1.28750,nan) (-1.28500,nan) (-1.28250,nan) (-1.28000,nan) (-1.27750,nan) (-1.27500,nan) (-1.27250,nan) (-1.27000,nan) (-1.26750,nan) (-1.26500,nan) (-1.26250,nan) (-1.26000,nan) (-1.25750,nan) (-1.25500,nan) (-1.25250,nan) (-1.25000,nan) (-1.24750,nan) (-1.24500,nan) (-1.24250,nan) (-1.24000,nan) (-1.23750,nan) (-1.23500,nan) (-1.23250,nan) (-1.23000,nan) (-1.22750,nan) (-1.22500,nan) (-1.22250,nan) (-1.22000,nan) (-1.21750,nan) (-1.21500,nan) (-1.21250,nan) (-1.21000,nan) (-1.20750,nan) (-1.20500,nan) (-1.20250,nan) (-1.20000,nan) (-1.19750,nan) (-1.19500,nan) (-1.19250,nan) (-1.19000,nan) (-1.18750,nan) (-1.18500,nan) (-1.18250,nan) (-1.18000,nan) (-1.17750,nan) (-1.17500,nan) (-1.17250,nan) (-1.17000,nan) (-1.16750,nan) (-1.16500,nan) (-1.16250,nan) (-1.16000,nan) (-1.15750,nan) (-1.15500,nan) (-1.15250,nan) (-1.15000,nan) (-1.14750,nan) (-1.14500,nan) (-1.14250,nan) (-1.14000,nan) (-1.13750,nan) (-1.13500,nan) (-1.13250,nan) (-1.13000,nan) (-1.12750,nan) (-1.12500,nan) (-1.12250,nan) (-1.12000,nan) (-1.11750,nan) (-1.11500,nan) (-1.11250,nan) (-1.11000,nan) (-1.10750,nan) (-1.10500,nan) (-1.10250,nan) (-1.10000,nan) (-1.09750,nan) (-1.09500,nan) (-1.09250,nan) (-1.09000,nan) (-1.08750,nan) (-1.08500,nan) (-1.08250,nan) (-1.08000,nan) (-1.07750,nan) (-1.07500,nan) (-1.07250,nan) (-1.07000,nan) (-1.06750,nan) (-1.06500,nan) (-1.06250,nan) (-1.06000,nan) (-1.05750,nan) (-1.05500,nan) (-1.05250,nan) (-1.05000,nan) (-1.04750,nan) (-1.04500,nan) (-1.04250,nan) (-1.04000,nan) (-1.03750,nan) (-1.03500,nan) (-1.03250,nan) (-1.03000,nan) (-1.02750,nan) (-1.02500,nan) (-1.02250,nan) (-1.02000,nan) (-1.01750,nan) (-1.01500,nan) (-1.01250,nan) (-1.01000,nan) (-1.00750,nan) (-1.00500,nan) (-1.00250,nan) (-1.00000,nan) (-0.99750,nan) (-0.99500,nan) (-0.99250,nan) (-0.99000,nan) (-0.98750,nan) (-0.98500,nan) (-0.98250,nan) (-0.98000,nan) (-0.97750,nan) (-0.97500,nan) (-0.97250,nan) (-0.97000,nan) (-0.96750,nan) (-0.96500,nan) (-0.96250,nan) (-0.96000,nan) (-0.95750,nan) (-0.95500,nan) (-0.95250,nan) (-0.95000,nan) (-0.94750,nan) (-0.94500,nan) (-0.94250,nan) (-0.94000,nan) (-0.93750,nan) (-0.93500,nan) (-0.93250,nan) (-0.93000,nan) (-0.92750,nan) (-0.92500,nan) (-0.92250,nan) (-0.92000,nan) (-0.91750,nan) (-0.91500,nan) (-0.91250,nan) (-0.91000,nan) (-0.90750,nan) (-0.90500,nan) (-0.90250,nan) (-0.90000,nan) (-0.89750,nan) (-0.89500,nan) (-0.89250,nan) (-0.89000,nan) (-0.88750,nan) (-0.88500,nan) (-0.88250,nan) (-0.88000,nan) (-0.87750,nan) (-0.87500,nan) (-0.87250,nan) (-0.87000,nan) (-0.86750,nan) (-0.86500,nan) (-0.86250,nan) (-0.86000,nan) (-0.85750,nan) (-0.85500,nan) (-0.85250,nan) (-0.85000,nan) (-0.84750,nan) (-0.84500,nan) (-0.84250,nan) (-0.84000,nan) (-0.83750,nan) (-0.83500,nan) (-0.83250,nan) (-0.83000,nan) (-0.82750,nan) (-0.82500,nan) (-0.82250,nan) (-0.82000,nan) (-0.81750,nan) (-0.81500,nan) (-0.81250,nan) (-0.81000,nan) (-0.80750,nan) (-0.80500,nan) (-0.80250,nan) (-0.80000,nan) (-0.79750,nan) (-0.79500,nan) (-0.79250,nan) (-0.79000,nan) (-0.78750,nan) (-0.78500,nan) (-0.78250,nan) (-0.78000,nan) (-0.77750,nan) (-0.77500,nan) (-0.77250,nan) (-0.77000,nan) (-0.76750,nan) (-0.76500,nan) (-0.76250,nan) (-0.76000,nan) (-0.75750,nan) (-0.75500,nan) (-0.75250,nan) (-0.75000,nan) (-0.74750,nan) (-0.74500,nan) (-0.74250,nan) (-0.74000,nan) (-0.73750,nan) (-0.73500,nan) (-0.73250,nan) (-0.73000,nan) (-0.72750,nan) (-0.72500,nan) (-0.72250,nan) (-0.72000,nan) (-0.71750,nan) (-0.71500,nan) (-0.71250,nan) (-0.71000,nan) (-0.70750,nan) (-0.70500,nan) (-0.70250,nan) (-0.70000,nan) (-0.69750,nan) (-0.69500,nan) (-0.69250,nan) (-0.69000,nan) (-0.68750,nan) (-0.68500,nan) (-0.68250,nan) (-0.68000,nan) (-0.67750,nan) (-0.67500,nan) (-0.67250,nan) (-0.67000,nan) (-0.66750,nan) (-0.66500,nan) (-0.66250,nan) (-0.66000,nan) (-0.65750,nan) (-0.65500,nan) (-0.65250,nan) (-0.65000,nan) (-0.64750,nan) (-0.64500,nan) (-0.64250,nan) (-0.64000,nan) (-0.63750,nan) (-0.63500,nan) (-0.63250,nan) (-0.63000,nan) (-0.62750,nan) (-0.62500,nan) (-0.62250,nan) (-0.62000,nan) (-0.61750,nan) (-0.61500,nan) (-0.61250,nan) (-0.61000,nan) (-0.60750,nan) (-0.60500,nan) (-0.60250,nan) (-0.60000,nan) (-0.59750,nan) (-0.59500,nan) (-0.59250,nan) (-0.59000,nan) (-0.58750,nan) (-0.58500,nan) (-0.58250,nan) (-0.58000,nan) (-0.57750,nan) (-0.57500,nan) (-0.57250,nan) (-0.57000,nan) (-0.56750,nan) (-0.56500,nan) (-0.56250,nan) (-0.56000,nan) (-0.55750,nan) (-0.55500,nan) (-0.55250,nan) (-0.55000,nan) (-0.54750,nan) (-0.54500,nan) (-0.54250,nan) (-0.54000,nan) (-0.53750,nan) (-0.53500,nan) (-0.53250,nan) (-0.53000,nan) (-0.52750,nan) (-0.52500,nan) (-0.52250,nan) (-0.52000,nan) (-0.51750,nan) (-0.51500,nan) (-0.51250,nan) (-0.51000,nan) (-0.50750,nan) (-0.50500,nan) (-0.50250,nan) (-0.50000,nan) (-0.49750,nan) (-0.49500,nan) (-0.49250,nan) (-0.49000,nan) (-0.48750,nan) (-0.48500,nan) (-0.48250,nan) (-0.48000,nan) (-0.47750,nan) (-0.47500,nan) (-0.47250,nan) (-0.47000,nan) (-0.46750,nan) (-0.46500,nan) (-0.46250,nan) (-0.46000,nan) (-0.45750,nan) (-0.45500,nan) (-0.45250,nan) (-0.45000,nan) (-0.44750,nan) (-0.44500,nan) (-0.44250,nan) (-0.44000,nan) (-0.43750,nan) (-0.43500,nan) (-0.43250,nan) (-0.43000,nan) (-0.42750,nan) (-0.42500,nan) (-0.42250,nan) (-0.42000,nan) (-0.41750,nan) (-0.41500,nan) (-0.41250,nan) (-0.41000,nan) (-0.40750,nan) (-0.40500,nan) (-0.40250,nan) (-0.40000,nan) (-0.39750,nan) (-0.39500,nan) (-0.39250,nan) (-0.39000,nan) (-0.38750,nan) (-0.38500,nan) (-0.38250,nan) (-0.38000,nan) (-0.37750,nan) (-0.37500,nan) (-0.37250,nan) (-0.37000,nan) (-0.36750,nan) (-0.36500,nan) (-0.36250,nan) (-0.36000,nan) (-0.35750,nan) (-0.35500,nan) (-0.35250,nan) (-0.35000,nan) (-0.34750,nan) (-0.34500,nan) (-0.34250,nan) (-0.34000,nan) (-0.33750,nan) (-0.33500,nan) (-0.33250,nan) (-0.33000,nan) (-0.32750,nan) (-0.32500,nan) (-0.32250,nan) (-0.32000,nan) (-0.31750,nan) (-0.31500,nan) (-0.31250,nan) (-0.31000,nan) (-0.30750,nan) (-0.30500,nan) (-0.30250,nan) (-0.30000,nan) (-0.29750,nan) (-0.29500,nan) (-0.29250,nan) (-0.29000,nan) (-0.28750,nan) (-0.28500,nan) (-0.28250,nan) (-0.28000,nan) (-0.27750,nan) (-0.27500,nan) (-0.27250,nan) (-0.27000,nan) (-0.26750,nan) (-0.26500,nan) (-0.26250,nan) (-0.26000,nan) (-0.25750,nan) (-0.25500,nan) (-0.25250,nan) (-0.25000,nan) (-0.24750,nan) (-0.24500,nan) (-0.24250,nan) (-0.24000,nan) (-0.23750,nan) (-0.23500,nan) (-0.23250,nan) (-0.23000,nan) (-0.22750,nan) (-0.22500,nan) (-0.22250,nan) (-0.22000,nan) (-0.21750,nan) (-0.21500,nan) (-0.21250,nan) (-0.21000,nan) (-0.20750,nan) (-0.20500,nan) (-0.20250,nan) (-0.20000,nan) (-0.19750,nan) (-0.19500,nan) (-0.19250,nan) (-0.19000,nan) (-0.18750,nan) (-0.18500,nan) (-0.18250,nan) (-0.18000,nan) (-0.17750,nan) (-0.17500,nan) (-0.17250,nan) (-0.17000,nan) (-0.16750,nan) (-0.16500,nan) (-0.16250,nan) (-0.16000,nan) (-0.15750,nan) (-0.15500,nan) (-0.15250,nan) (-0.15000,nan) (-0.14750,nan) (-0.14500,nan) (-0.14250,nan) (-0.14000,nan) (-0.13750,nan) (-0.13500,nan) (-0.13250,nan) (-0.13000,nan) (-0.12750,nan) (-0.12500,nan) (-0.12250,nan) (-0.12000,nan) (-0.11750,nan) (-0.11500,nan) (-0.11250,nan) (-0.11000,nan) (-0.10750,nan) (-0.10500,nan) (-0.10250,nan) (-0.10000,nan) (-0.09750,nan) (-0.09500,nan) (-0.09250,nan) (-0.09000,nan) (-0.08750,nan) (-0.08500,nan) (-0.08250,nan) (-0.08000,nan) (-0.07750,nan) (-0.07500,nan) (-0.07250,nan) (-0.07000,nan) (-0.06750,nan) (-0.06500,nan) (-0.06250,nan) (-0.06000,nan) (-0.05750,nan) (-0.05500,nan) (-0.05250,nan) (-0.05000,nan) (-0.04750,nan) (-0.04500,nan) (-0.04250,nan) (-0.04000,nan) (-0.03750,nan) (-0.03500,nan) (-0.03250,nan) (-0.03000,nan) (-0.02750,nan) (-0.02500,nan) (-0.02250,nan) (-0.02000,nan) (-0.01750,nan) (-0.01500,nan) (-0.01250,nan) (-0.01000,nan) (-0.00750,nan) (-0.00500,nan) (-0.00250,nan) (0.00000,nan) (0.00250,nan) (0.00500,nan) (0.00750,nan) (0.01000,nan) (0.01250,nan) (0.01500,nan) (0.01750,nan) (0.02000,nan) (0.02250,nan) (0.02500,nan) (0.02750,nan) (0.03000,nan) (0.03250,nan) (0.03500,nan) (0.03750,nan) (0.04000,nan) (0.04250,nan) (0.04500,nan) (0.04750,nan) (0.05000,nan) (0.05250,nan) (0.05500,nan) (0.05750,nan) (0.06000,nan) (0.06250,nan) (0.06500,nan) (0.06750,nan) (0.07000,nan) (0.07250,nan) (0.07500,nan) (0.07750,nan) (0.08000,nan) (0.08250,nan) (0.08500,nan) (0.08750,nan) (0.09000,nan) (0.09250,nan) (0.09500,nan) (0.09750,nan) (0.10000,nan) (0.10250,nan) (0.10500,nan) (0.10750,nan) (0.11000,nan) (0.11250,nan) (0.11500,nan) (0.11750,nan) (0.12000,nan) (0.12250,nan) (0.12500,nan) (0.12750,nan) (0.13000,nan) (0.13250,nan) (0.13500,nan) (0.13750,nan) (0.14000,nan) (0.14250,nan) (0.14500,nan) (0.14750,nan) (0.15000,nan) (0.15250,nan) (0.15500,nan) (0.15750,nan) (0.16000,nan) (0.16250,nan) (0.16500,nan) (0.16750,nan) (0.17000,nan) (0.17250,nan) (0.17500,nan) (0.17750,nan) (0.18000,nan) (0.18250,nan) (0.18500,nan) (0.18750,nan) (0.19000,nan) (0.19250,nan) (0.19500,nan) (0.19750,nan) (0.20000,nan) (0.20250,nan) (0.20500,nan) (0.20750,nan) (0.21000,nan) (0.21250,nan) (0.21500,nan) (0.21750,nan) (0.22000,nan) (0.22250,nan) (0.22500,nan) (0.22750,nan) (0.23000,nan) (0.23250,nan) (0.23500,nan) (0.23750,nan) (0.24000,nan) (0.24250,nan) (0.24500,nan) (0.24750,nan) (0.25000,nan) (0.25250,nan) (0.25500,nan) (0.25750,nan) (0.26000,nan) (0.26250,nan) (0.26500,nan) (0.26750,nan) (0.27000,nan) (0.27250,nan) (0.27500,nan) (0.27750,nan) (0.28000,nan) (0.28250,nan) (0.28500,nan) (0.28750,nan) (0.29000,nan) (0.29250,nan) (0.29500,nan) (0.29750,nan) (0.30000,nan) (0.30250,nan) (0.30500,nan) (0.30750,nan) (0.31000,nan) (0.31250,nan) (0.31500,nan) (0.31750,nan) (0.32000,nan) (0.32250,nan) (0.32500,nan) (0.32750,nan) (0.33000,nan) (0.33250,nan) (0.33500,nan) (0.33750,nan) (0.34000,nan) (0.34250,nan) (0.34500,nan) (0.34750,nan) (0.35000,nan) (0.35250,nan) (0.35500,nan) (0.35750,nan) (0.36000,nan) (0.36250,nan) (0.36500,nan) (0.36750,nan) (0.37000,nan) (0.37250,nan) (0.37500,nan) (0.37750,nan) (0.38000,nan) (0.38250,nan) (0.38500,nan) (0.38750,nan) (0.39000,nan) (0.39250,nan) (0.39500,nan) (0.39750,nan) (0.40000,nan) (0.40250,nan) (0.40500,nan) (0.40750,nan) (0.41000,nan) (0.41250,nan) (0.41500,nan) (0.41750,nan) (0.42000,nan) (0.42250,nan) (0.42500,nan) (0.42750,nan) (0.43000,nan) (0.43250,nan) (0.43500,nan) (0.43750,nan) (0.44000,nan) (0.44250,nan) (0.44500,nan) (0.44750,nan) (0.45000,nan) (0.45250,nan) (0.45500,nan) (0.45750,nan) (0.46000,nan) (0.46250,nan) (0.46500,nan) (0.46750,nan) (0.47000,nan) (0.47250,nan) (0.47500,nan) (0.47750,nan) (0.48000,nan) (0.48250,nan) (0.48500,nan) (0.48750,nan) (0.49000,nan) (0.49250,nan) (0.49500,nan) (0.49750,nan) (0.50000,nan) (0.50250,nan) (0.50500,nan) (0.50750,nan) (0.51000,nan) (0.51250,nan) (0.51500,nan) (0.51750,nan) (0.52000,nan) (0.52250,nan) (0.52500,nan) (0.52750,nan) (0.53000,nan) (0.53250,nan) (0.53500,nan) (0.53750,nan) (0.54000,nan) (0.54250,nan) (0.54500,nan) (0.54750,nan) (0.55000,nan) (0.55250,nan) (0.55500,nan) (0.55750,nan) (0.56000,nan) (0.56250,nan) (0.56500,nan) (0.56750,nan) (0.57000,nan) (0.57250,nan) (0.57500,nan) (0.57750,nan) (0.58000,nan) (0.58250,nan) (0.58500,nan) (0.58750,nan) (0.59000,nan) (0.59250,nan) (0.59500,nan) (0.59750,nan) (0.60000,nan) (0.60250,nan) (0.60500,nan) (0.60750,nan) (0.61000,nan) (0.61250,nan) (0.61500,nan) (0.61750,nan) (0.62000,nan) (0.62250,nan) (0.62500,nan) (0.62750,nan) (0.63000,nan) (0.63250,nan) (0.63500,nan) (0.63750,nan) (0.64000,nan) (0.64250,nan) (0.64500,nan) (0.64750,nan) (0.65000,nan) (0.65250,nan) (0.65500,nan) (0.65750,nan) (0.66000,nan) (0.66250,nan) (0.66500,nan) (0.66750,nan) (0.67000,nan) (0.67250,nan) (0.67500,nan) (0.67750,nan) (0.68000,nan) (0.68250,nan) (0.68500,nan) (0.68750,nan) (0.69000,nan) (0.69250,nan) (0.69500,nan) (0.69750,nan) (0.70000,nan) (0.70250,nan) (0.70500,nan) (0.70750,nan) (0.71000,nan) (0.71250,nan) (0.71500,nan) (0.71750,nan) (0.72000,nan) (0.72250,nan) (0.72500,nan) (0.72750,nan) (0.73000,nan) (0.73250,nan) (0.73500,nan) (0.73750,nan) (0.74000,nan) (0.74250,nan) (0.74500,nan) (0.74750,nan) (0.75000,nan) (0.75250,nan) (0.75500,nan) (0.75750,nan) (0.76000,nan) (0.76250,nan) (0.76500,nan) (0.76750,nan) (0.77000,nan) (0.77250,nan) (0.77500,nan) (0.77750,nan) (0.78000,nan) (0.78250,nan) (0.78500,nan) (0.78750,nan) (0.79000,nan) (0.79250,nan) (0.79500,nan) (0.79750,nan) (0.80000,nan) (0.80250,nan) (0.80500,nan) (0.80750,nan) (0.81000,nan) (0.81250,nan) (0.81500,nan) (0.81750,nan) (0.82000,nan) (0.82250,nan) (0.82500,nan) (0.82750,nan) (0.83000,nan) (0.83250,nan) (0.83500,nan) (0.83750,nan) (0.84000,nan) (0.84250,nan) (0.84500,nan) (0.84750,nan) (0.85000,nan) (0.85250,nan) (0.85500,nan) (0.85750,nan) (0.86000,nan) (0.86250,nan) (0.86500,nan) (0.86750,nan) (0.87000,nan) (0.87250,nan) (0.87500,nan) (0.87750,nan) (0.88000,nan) (0.88250,nan) (0.88500,nan) (0.88750,nan) (0.89000,nan) (0.89250,nan) (0.89500,nan) (0.89750,nan) (0.90000,nan) (0.90250,nan) (0.90500,nan) (0.90750,nan) (0.91000,nan) (0.91250,nan) (0.91500,nan) (0.91750,nan) (0.92000,nan) (0.92250,nan) (0.92500,nan) (0.92750,nan) (0.93000,nan) (0.93250,nan) (0.93500,nan) (0.93750,nan) (0.94000,nan) (0.94250,nan) (0.94500,nan) (0.94750,nan) (0.95000,nan) (0.95250,nan) (0.95500,nan) (0.95750,nan) (0.96000,nan) (0.96250,nan) (0.96500,nan) (0.96750,nan) (0.97000,nan) (0.97250,nan) (0.97500,nan) (0.97750,nan) (0.98000,nan) (0.98250,nan) (0.98500,nan) (0.98750,nan) (0.99000,nan) (0.99250,nan) (0.99500,nan) (0.99750,nan) (1.00000,nan) (1.00250,nan) (1.00500,nan) (1.00750,nan) (1.01000,nan) (1.01250,nan) (1.01500,nan) (1.01750,nan) (1.02000,nan) (1.02250,nan) (1.02500,nan) (1.02750,nan) (1.03000,nan) (1.03250,nan) (1.03500,nan) (1.03750,nan) (1.04000,nan) (1.04250,nan) (1.04500,nan) (1.04750,nan) (1.05000,nan) (1.05250,nan) (1.05500,nan) (1.05750,nan) (1.06000,nan) (1.06250,nan) (1.06500,nan) (1.06750,nan) (1.07000,nan) (1.07250,nan) (1.07500,nan) (1.07750,nan) (1.08000,nan) (1.08250,nan) (1.08500,nan) (1.08750,nan) (1.09000,nan) (1.09250,nan) (1.09500,nan) (1.09750,nan) (1.10000,nan) (1.10250,nan) (1.10500,nan) (1.10750,nan) (1.11000,nan) (1.11250,nan) (1.11500,nan) (1.11750,nan) (1.12000,nan) (1.12250,nan) (1.12500,nan) (1.12750,nan) (1.13000,nan) (1.13250,nan) (1.13500,nan) (1.13750,nan) (1.14000,nan) (1.14250,nan) (1.14500,nan) (1.14750,nan) (1.15000,nan) (1.15250,nan) (1.15500,nan) (1.15750,nan) (1.16000,nan) (1.16250,nan) (1.16500,nan) (1.16750,nan) (1.17000,nan) (1.17250,nan) (1.17500,nan) (1.17750,nan) (1.18000,nan) (1.18250,nan) (1.18500,nan) (1.18750,nan) (1.19000,nan) (1.19250,nan) (1.19500,nan) (1.19750,nan) (1.20000,nan) (1.20250,nan) (1.20500,nan) (1.20750,nan) (1.21000,nan) (1.21250,nan) (1.21500,nan) (1.21750,nan) (1.22000,nan) (1.22250,nan) (1.22500,nan) (1.22750,nan) (1.23000,nan) (1.23250,nan) (1.23500,nan) (1.23750,nan) (1.24000,nan) (1.24250,nan) (1.24500,nan) (1.24750,nan) (1.25000,nan) (1.25250,nan) (1.25500,nan) (1.25750,nan) (1.26000,nan) (1.26250,nan) (1.26500,nan) (1.26750,nan) (1.27000,nan) (1.27250,nan) (1.27500,nan) (1.27750,nan) (1.28000,nan) (1.28250,nan) (1.28500,nan) (1.28750,nan) (1.29000,nan) (1.29250,nan) (1.29500,nan) (1.29750,nan) (1.30000,nan) (1.30250,nan) (1.30500,nan) (1.30750,nan) (1.31000,nan) (1.31250,nan) (1.31500,nan) (1.31750,nan) (1.32000,nan) (1.32250,nan) (1.32500,nan) (1.32750,nan) (1.33000,nan) (1.33250,nan) (1.33500,nan) (1.33750,nan) (1.34000,nan) (1.34250,nan) (1.34500,nan) (1.34750,nan) (1.35000,nan) (1.35250,nan) (1.35500,nan) (1.35750,nan) (1.36000,nan) (1.36250,nan) (1.36500,nan) (1.36750,nan) (1.37000,nan) (1.37250,nan) (1.37500,nan) (1.37750,nan) (1.38000,nan) (1.38250,nan) (1.38500,nan) (1.38750,nan) (1.39000,nan) (1.39250,nan) (1.39500,nan) (1.39750,nan) (1.40000,nan) (1.40250,nan) (1.40500,nan) (1.40750,nan) (1.41000,nan) (1.41250,nan) (1.41500,nan) (1.41750,nan) (1.42000,nan) (1.42250,nan) (1.42500,nan) (1.42750,nan) (1.43000,nan) (1.43250,nan) (1.43500,nan) (1.43750,nan) (1.44000,nan) (1.44250,nan) (1.44500,nan) (1.44750,nan) (1.45000,nan) (1.45250,nan) (1.45500,nan) (1.45750,nan) (1.46000,nan) (1.46250,nan) (1.46500,nan) (1.46750,nan) (1.47000,nan) (1.47250,nan) (1.47500,nan) (1.47750,nan) (1.48000,nan) (1.48250,nan) (1.48500,nan) (1.48750,nan) (1.49000,nan) (1.49250,nan) (1.49500,nan) (1.49750,nan) (1.50000,nan) (1.50250,nan) (1.50500,nan) (1.50750,nan) (1.51000,nan) (1.51250,nan) (1.51500,nan) (1.51750,nan) (1.52000,nan) (1.52250,nan) (1.52500,nan) (1.52750,nan) (1.53000,nan) (1.53250,nan) (1.53500,nan) (1.53750,nan) (1.54000,nan) (1.54250,nan) (1.54500,nan) (1.54750,nan) (1.55000,nan) (1.55250,nan) (1.55500,nan) (1.55750,nan) (1.56000,nan) (1.56250,nan) (1.56500,nan) (1.56750,nan) (1.57000,nan) (1.57250,nan) (1.57500,nan) (1.57750,nan) (1.58000,nan) (1.58250,nan) (1.58500,nan) (1.58750,nan) (1.59000,nan) (1.59250,nan) (1.59500,nan) (1.59750,nan) (1.60000,nan) (1.60250,nan) (1.60500,nan) (1.60750,nan) (1.61000,nan) (1.61250,nan) (1.61500,nan) (1.61750,nan) (1.62000,nan) (1.62250,nan) (1.62500,nan) (1.62750,nan) (1.63000,nan) (1.63250,nan) (1.63500,nan) (1.63750,nan) (1.64000,nan) (1.64250,nan) (1.64500,nan) (1.64750,nan) (1.65000,nan) (1.65250,nan) (1.65500,nan) (1.65750,nan) (1.66000,nan) (1.66250,nan) (1.66500,nan) (1.66750,nan) (1.67000,nan) (1.67250,nan) (1.67500,nan) (1.67750,nan) (1.68000,nan) (1.68250,nan) (1.68500,nan) (1.68750,nan) (1.69000,nan) (1.69250,nan) (1.69500,nan) (1.69750,nan) (1.70000,nan) (1.70250,nan) (1.70500,nan) (1.70750,nan) (1.71000,nan) (1.71250,nan) (1.71500,nan) (1.71750,nan) (1.72000,nan) (1.72250,nan) (1.72500,nan) (1.72750,nan) (1.73000,nan) (1.73250,nan) (1.73500,nan) (1.73750,nan) (1.74000,nan) (1.74250,nan) (1.74500,nan) (1.74750,nan) (1.75000,nan) (1.75250,nan) (1.75500,nan) (1.75750,nan) (1.76000,nan) (1.76250,nan) (1.76500,nan) (1.76750,nan) (1.77000,nan) (1.77250,nan) (1.77500,nan) (1.77750,nan) (1.78000,nan) (1.78250,nan) (1.78500,nan) (1.78750,nan) (1.79000,nan) (1.79250,nan) (1.79500,nan) (1.79750,nan) (1.80000,nan) (1.80250,nan) (1.80500,nan) (1.80750,nan) (1.81000,nan) (1.81250,nan) (1.81500,nan) (1.81750,nan) (1.82000,nan) (1.82250,nan) (1.82500,nan) (1.82750,nan) (1.83000,nan) (1.83250,nan) (1.83500,nan) (1.83750,nan) (1.84000,nan) (1.84250,nan) (1.84500,nan) (1.84750,nan) (1.85000,nan) (1.85250,nan) (1.85500,nan) (1.85750,nan) (1.86000,nan) (1.86250,nan) (1.86500,nan) (1.86750,nan) (1.87000,nan) (1.87250,nan) (1.87500,nan) (1.87750,nan) (1.88000,nan) (1.88250,nan) (1.88500,nan) (1.88750,nan) (1.89000,nan) (1.89250,nan) (1.89500,nan) (1.89750,nan) (1.90000,nan) (1.90250,nan) (1.90500,nan) (1.90750,nan) (1.91000,nan) (1.91250,nan) (1.91500,nan) (1.91750,nan) (1.92000,nan) (1.92250,nan) (1.92500,nan) (1.92750,nan) (1.93000,nan) (1.93250,nan) (1.93500,nan) (1.93750,nan) (1.94000,nan) (1.94250,nan) (1.94500,nan) (1.94750,nan) (1.95000,nan) (1.95250,nan) (1.95500,nan) (1.95750,nan) (1.96000,nan) (1.96250,nan) (1.96500,nan) (1.96750,nan) (1.97000,nan) (1.97250,nan) (1.97500,nan) (1.97750,nan) (1.98000,nan) (1.98250,nan) (1.98500,nan) (1.98750,nan) (1.99000,nan) (1.99250,nan) (1.99500,nan) (1.99750,nan) (2.00000,nan) (2.00250,nan) (2.00500,nan) (2.00750,nan) (2.01000,nan) (2.01250,nan) (2.01500,nan) (2.01750,nan) (2.02000,nan) (2.02250,nan) (2.02500,nan) (2.02750,nan) (2.03000,nan) (2.03250,nan) (2.03500,nan) (2.03750,nan) (2.04000,nan) (2.04250,nan) (2.04500,nan) (2.04750,nan) (2.05000,nan) (2.05250,nan) (2.05500,nan) (2.05750,nan) (2.06000,nan) (2.06250,nan) (2.06500,nan) (2.06750,nan) (2.07000,nan) (2.07250,nan) (2.07500,nan) (2.07750,nan) (2.08000,nan) (2.08250,nan) (2.08500,nan) (2.08750,nan) (2.09000,nan) (2.09250,nan) (2.09500,nan) (2.09750,nan) (2.10000,nan) (2.10250,nan) (2.10500,nan) (2.10750,nan) (2.11000,nan) (2.11250,nan) (2.11500,nan) (2.11750,nan) (2.12000,nan) (2.12250,nan) (2.12500,nan) (2.12750,nan) (2.13000,nan) (2.13250,nan) (2.13500,nan) (2.13750,nan) (2.14000,nan) (2.14250,nan) (2.14500,nan) (2.14750,nan) (2.15000,nan) (2.15250,nan) (2.15500,nan) (2.15750,nan) (2.16000,nan) (2.16250,nan) (2.16500,nan) (2.16750,nan) (2.17000,nan) (2.17250,nan) (2.17500,nan) (2.17750,nan) (2.18000,nan) (2.18250,nan) (2.18500,nan) (2.18750,nan) (2.19000,nan) (2.19250,nan) (2.19500,nan) (2.19750,nan) (2.20000,nan) (2.20250,nan) (2.20500,nan) (2.20750,nan) (2.21000,nan) (2.21250,nan) (2.21500,nan) (2.21750,nan) (2.22000,nan) (2.22250,nan) (2.22500,nan) (2.22750,nan) (2.23000,nan) (2.23250,nan) (2.23500,nan) (2.23750,nan) (2.24000,nan) (2.24250,nan) (2.24500,nan) (2.24750,nan) (2.25000,nan) (2.25250,nan) (2.25500,nan) (2.25750,nan) (2.26000,nan) (2.26250,nan) (2.26500,nan) (2.26750,nan) (2.27000,nan) (2.27250,nan) (2.27500,nan) (2.27750,nan) (2.28000,nan) (2.28250,nan) (2.28500,nan) (2.28750,nan) (2.29000,nan) (2.29250,nan) (2.29500,nan) (2.29750,nan) (2.30000,nan) (2.30250,nan) (2.30500,nan) (2.30750,nan) (2.31000,nan) (2.31250,nan) (2.31500,nan) (2.31750,nan) (2.32000,nan) (2.32250,nan) (2.32500,nan) (2.32750,nan) (2.33000,nan) (2.33250,nan) (2.33500,nan) (2.33750,nan) (2.34000,nan) (2.34250,nan) (2.34500,nan) (2.34750,nan) (2.35000,nan) (2.35250,nan) (2.35500,nan) (2.35750,nan) (2.36000,nan) (2.36250,nan) (2.36500,nan) (2.36750,nan) (2.37000,nan) (2.37250,nan) (2.37500,nan) (2.37750,nan) (2.38000,nan) (2.38250,nan) (2.38500,nan) (2.38750,nan) (2.39000,nan) (2.39250,nan) (2.39500,nan) (2.39750,nan) (2.40000,nan) (2.40250,nan) (2.40500,nan) (2.40750,nan) (2.41000,nan) (2.41250,nan) (2.41500,nan) (2.41750,nan) (2.42000,nan) (2.42250,nan) (2.42500,nan) (2.42750,nan) (2.43000,nan) (2.43250,nan) (2.43500,nan) (2.43750,nan) (2.44000,nan) (2.44250,nan) (2.44500,nan) (2.44750,nan) (2.45000,nan) (2.45250,nan) (2.45500,nan) (2.45750,nan) (2.46000,nan) (2.46250,nan) (2.46500,nan) (2.46750,nan) (2.47000,nan) (2.47250,nan) (2.47500,nan) (2.47750,nan) (2.48000,nan) (2.48250,nan) (2.48500,nan) (2.48750,nan) (2.49000,nan) (2.49250,nan) (2.49500,nan) (2.49750,nan) (2.50000,nan) (2.50250,nan) (2.50500,nan) (2.50750,nan) (2.51000,nan) (2.51250,nan) (2.51500,nan) (2.51750,nan) (2.52000,nan) (2.52250,nan) (2.52500,nan) (2.52750,nan) (2.53000,nan) (2.53250,nan) (2.53500,nan) (2.53750,nan) (2.54000,nan) (2.54250,nan) (2.54500,nan) (2.54750,nan) (2.55000,nan) (2.55250,nan) (2.55500,nan) (2.55750,nan) (2.56000,nan) (2.56250,nan) (2.56500,nan) (2.56750,nan) (2.57000,nan) (2.57250,nan) (2.57500,nan) (2.57750,nan) (2.58000,nan) (2.58250,nan) (2.58500,nan) (2.58750,nan) (2.59000,nan) (2.59250,nan) (2.59500,nan) (2.59750,nan) (2.60000,nan) (2.60250,nan) (2.60500,nan) (2.60750,nan) (2.61000,nan) (2.61250,nan) (2.61500,nan) (2.61750,nan) (2.62000,nan) (2.62250,nan) (2.62500,nan) (2.62750,nan) (2.63000,nan) (2.63250,nan) (2.63500,nan) (2.63750,nan) (2.64000,nan) (2.64250,nan) (2.64500,nan) (2.64750,nan) (2.65000,nan) (2.65250,nan) (2.65500,nan) (2.65750,nan) (2.66000,nan) (2.66250,nan) (2.66500,nan) (2.66750,nan) (2.67000,nan) (2.67250,nan) (2.67500,nan) (2.67750,nan) (2.68000,nan) (2.68250,nan) (2.68500,nan) (2.68750,nan) (2.69000,nan) (2.69250,nan) (2.69500,nan) (2.69750,nan) (2.70000,nan) (2.70250,nan) (2.70500,nan) (2.70750,nan) (2.71000,nan) (2.71250,nan) (2.71500,nan) (2.71750,nan) (2.72000,nan) (2.72250,nan) (2.72500,nan) (2.72750,nan) (2.73000,nan) (2.73250,nan) (2.73500,nan) (2.73750,nan) (2.74000,nan) (2.74250,nan) (2.74500,nan) (2.74750,nan) (2.75000,nan) (2.75250,nan) (2.75500,nan) (2.75750,nan) (2.76000,nan) (2.76250,nan) (2.76500,nan) (2.76750,nan) (2.77000,nan) (2.77250,nan) (2.77500,nan) (2.77750,nan) (2.78000,nan) (2.78250,nan) (2.78500,nan) (2.78750,nan) (2.79000,nan) (2.79250,nan) (2.79500,nan) (2.79750,nan) (2.80000,nan) (2.80250,nan) (2.80500,nan) (2.80750,nan) (2.81000,nan) (2.81250,nan) (2.81500,nan) (2.81750,nan) (2.82000,nan) (2.82250,nan) (2.82500,nan) (2.82750,nan) (2.83000,nan) (2.83250,nan) (2.83500,nan) (2.83750,nan) (2.84000,nan) (2.84250,nan) (2.84500,nan) (2.84750,nan) (2.85000,nan) (2.85250,nan) (2.85500,nan) (2.85750,nan) (2.86000,nan) (2.86250,nan) (2.86500,nan) (2.86750,nan) (2.87000,nan) (2.87250,nan) (2.87500,nan) (2.87750,nan) (2.88000,nan) (2.88250,nan) (2.88500,nan) (2.88750,nan) (2.89000,nan) (2.89250,nan) (2.89500,nan) (2.89750,nan) (2.90000,nan) (2.90250,nan) (2.90500,nan) (2.90750,nan) (2.91000,nan) (2.91250,nan) (2.91500,nan) (2.91750,nan) (2.92000,nan) (2.92250,nan) (2.92500,nan) (2.92750,nan) (2.93000,nan) (2.93250,nan) (2.93500,nan) (2.93750,nan) (2.94000,nan) (2.94250,nan) (2.94500,nan) (2.94750,nan) (2.95000,nan) (2.95250,nan) (2.95500,nan) (2.95750,nan) (2.96000,nan) (2.96250,nan) (2.96500,nan) (2.96750,nan) (2.97000,nan) (2.97250,nan) (2.97500,nan) (2.97750,nan) (2.98000,nan) (2.98250,nan) (2.98500,nan) (2.98750,nan) (2.99000,nan) (2.99250,nan) (2.99500,nan) (2.99750,nan) (3.00000,nan) (3.00250,nan) (3.00500,nan) (3.00750,nan) (3.01000,nan) (3.01250,nan) (3.01500,nan) (3.01750,nan) (3.02000,nan) (3.02250,nan) (3.02500,nan) (3.02750,nan) (3.03000,nan) (3.03250,nan) (3.03500,nan) (3.03750,nan) (3.04000,nan) (3.04250,nan) (3.04500,nan) (3.04750,nan) (3.05000,nan) (3.05250,nan) (3.05500,nan) (3.05750,nan) (3.06000,nan) (3.06250,nan) (3.06500,nan) (3.06750,nan) (3.07000,nan) (3.07250,nan) (3.07500,nan) (3.07750,nan) (3.08000,nan) (3.08250,nan) (3.08500,nan) (3.08750,nan) (3.09000,nan) (3.09250,nan) (3.09500,nan) (3.09750,nan) (3.10000,nan) (3.10250,nan) (3.10500,nan) (3.10750,nan) (3.11000,nan) (3.11250,nan) (3.11500,nan) (3.11750,nan) (3.12000,nan) (3.12250,nan) (3.12500,nan) (3.12750,nan) (3.13000,nan) (3.13250,nan) (3.13500,nan) (3.13750,nan) (3.14000,nan) (3.14250,nan) (3.14500,nan) (3.14750,nan) (3.15000,nan) (3.15250,nan) (3.15500,nan) (3.15750,nan) (3.16000,nan) (3.16250,nan) (3.16500,nan) (3.16750,nan) (3.17000,nan) (3.17250,nan) (3.17500,nan) (3.17750,nan) (3.18000,nan) (3.18250,nan) (3.18500,nan) (3.18750,nan) (3.19000,nan) (3.19250,nan) (3.19500,nan) (3.19750,nan) (3.20000,nan) (3.20250,nan) (3.20500,nan) (3.20750,nan) (3.21000,nan) (3.21250,nan) (3.21500,nan) (3.21750,nan) (3.22000,nan) (3.22250,nan) (3.22500,nan) (3.22750,nan) (3.23000,nan) (3.23250,nan) (3.23500,nan) (3.23750,nan) (3.24000,nan) (3.24250,nan) (3.24500,nan) (3.24750,nan) (3.25000,nan) (3.25250,nan) (3.25500,nan) (3.25750,nan) (3.26000,nan) (3.26250,nan) (3.26500,nan) (3.26750,nan) (3.27000,nan) (3.27250,nan) (3.27500,nan) (3.27750,nan) (3.28000,nan) (3.28250,nan) (3.28500,nan) (3.28750,nan) (3.29000,nan) (3.29250,nan) (3.29500,nan) (3.29750,nan) (3.30000,nan) (3.30250,nan) (3.30500,nan) (3.30750,nan) (3.31000,nan) (3.31250,nan) (3.31500,nan) (3.31750,nan) (3.32000,nan) (3.32250,nan) (3.32500,nan) (3.32750,nan) (3.33000,nan) (3.33250,nan) (3.33500,nan) (3.33750,nan) (3.34000,nan) (3.34250,nan) (3.34500,nan) (3.34750,nan) (3.35000,nan) (3.35250,nan) (3.35500,nan) (3.35750,nan) (3.36000,nan) (3.36250,nan) (3.36500,nan) (3.36750,nan) (3.37000,nan) (3.37250,nan) (3.37500,nan) (3.37750,nan) (3.38000,nan) (3.38250,nan) (3.38500,nan) (3.38750,nan) (3.39000,nan) (3.39250,nan) (3.39500,nan) (3.39750,nan) (3.40000,nan) (3.40250,nan) (3.40500,nan) (3.40750,nan) (3.41000,nan) (3.41250,nan) (3.41500,nan) (3.41750,nan) (3.42000,nan) (3.42250,nan) (3.42500,nan) (3.42750,nan) (3.43000,nan) (3.43250,nan) (3.43500,nan) (3.43750,nan) (3.44000,nan) (3.44250,nan) (3.44500,nan) (3.44750,nan) (3.45000,nan) (3.45250,nan) (3.45500,nan) (3.45750,nan) (3.46000,nan) (3.46250,nan) (3.46500,nan) (3.46750,nan) (3.47000,nan) (3.47250,nan) (3.47500,nan) (3.47750,nan) (3.48000,nan) (3.48250,nan) (3.48500,-1.7746222) (3.48750,-1.5519554) (3.49000,-1.4069028) (3.49250,-1.2994578) (3.49500,-1.2142871) (3.49750,-1.1438644) (3.50000,-1.0839313) (3.50250,-1.0318430) (3.50500,-0.9858435) (3.50750,-0.9447071) (3.51000,-0.9075445) (3.51250,-0.8736897) (3.51500,-0.8426312) (3.51750,-0.8139674) (3.52000,-0.7873774) (3.52250,-0.7626007) (3.52500,-0.7394229) (3.52750,-0.7176652) (3.53000,-0.6971775) (3.53250,-0.6778319) (3.53500,-0.6595190) (3.53750,-0.6421444) (3.54000,-0.6256258) (3.54250,-0.6098915) (3.54500,-0.5948782) (3.54750,-0.5805300) (3.55000,-0.5667972) (3.55250,-0.5536355) (3.55500,-0.5410051) (3.55750,-0.5288702) (3.56000,-0.5171984) (3.56250,-0.5059607) (3.56500,-0.4951302) (3.56750,-0.4846828) (3.57000,-0.4745962) (3.57250,-0.4648502) (3.57500,-0.4554260) (3.57750,-0.4463065) (3.58000,-0.4374758) (3.58250,-0.4289192) (3.58500,-0.4206230) (3.58750,-0.4125747) (3.59000,-0.4047624) (3.59250,-0.3971753) (3.59500,-0.3898031) (3.59750,-0.3826363) (3.60000,-0.3756658) (3.60250,-0.3688834) (3.60500,-0.3622812) (3.60750,-0.3558517) (3.61000,-0.3495880) (3.61250,-0.3434836) (3.61500,-0.3375322) (3.61750,-0.3317281) (3.62000,-0.3260656) (3.62250,-0.3205396) (3.62500,-0.3151450) (3.62750,-0.3098773) (3.63000,-0.3047320) (3.63250,-0.2997047) (3.63500,-0.2947916) (3.63750,-0.2899887) (3.64000,-0.2852925) (3.64250,-0.2806994) (3.64500,-0.2762061) (3.64750,-0.2718095) (3.65000,-0.2675065) (3.65250,-0.2632943) (3.65500,-0.2591700) (3.65750,-0.2551311) (3.66000,-0.2511751) (3.66250,-0.2472994) (3.66500,-0.2435017) (3.66750,-0.2397799) (3.67000,-0.2361318) (3.67250,-0.2325554) (3.67500,-0.2290485) (3.67750,-0.2256095) (3.68000,-0.2222363) (3.68250,-0.2189273) (3.68500,-0.2156808) (3.68750,-0.2124952) (3.69000,-0.2093688) (3.69250,-0.2063001) (3.69500,-0.2032878) (3.69750,-0.2003303) (3.70000,-0.1974263) (3.70250,-0.1945746) (3.70500,-0.1917737) (3.70750,-0.1890226) (3.71000,-0.1863200) (3.71250,-0.1836648) (3.71500,-0.1810558) (3.71750,-0.1784920) (3.72000,-0.1759723) (3.72250,-0.1734957) (3.72500,-0.1710613) (3.72750,-0.1686682) (3.73000,-0.1663152) (3.73250,-0.1640017) (3.73500,-0.1617268) (3.73750,-0.1594895) (3.74000,-0.1572891) (3.74250,-0.1551248) (3.74500,-0.1529958) (3.74750,-0.1509014) (3.75000,-0.1488408) (3.75250,-0.1468134) (3.75500,-0.1448185) (3.75750,-0.1428554) (3.76000,-0.1409234) (3.76250,-0.1390219) (3.76500,-0.1371504) (3.76750,-0.1353082) (3.77000,-0.1334947) (3.77250,-0.1317094) (3.77500,-0.1299517) (3.77750,-0.1282211) (3.78000,-0.1265170) (3.78250,-0.1248390) (3.78500,-0.1231866) (3.78750,-0.1215592) (3.79000,-0.1199564) (3.79250,-0.1183778) (3.79500,-0.1168229) (3.79750,-0.1152912) (3.80000,-0.1137823) (3.80250,-0.1122959) (3.80500,-0.1108314) (3.80750,-0.1093886) (3.81000,-0.1079669) (3.81250,-0.1065661) (3.81500,-0.1051857) (3.81750,-0.1038255) (3.82000,-0.1024850) (3.82250,-0.1011639) (3.82500,-0.0998618) (3.82750,-0.0985785) (3.83000,-0.0973136) (3.83250,-0.0960668) (3.83500,-0.0948378) (3.83750,-0.0936262) (3.84000,-0.0924318) (3.84250,-0.0912543) (3.84500,-0.0900935) (3.84750,-0.0889489) (3.85000,-0.0878204) (3.85250,-0.0867077) (3.85500,-0.0856106) (3.85750,-0.0845287) (3.86000,-0.0834618) (3.86250,-0.0824097) (3.86500,-0.0813722) (3.86750,-0.0803490) (3.87000,-0.0793398) (3.87250,-0.0783445) (3.87500,-0.0773628) (3.87750,-0.0763946) (3.88000,-0.0754395) (3.88250,-0.0744974) (3.88500,-0.0735682) (3.88750,-0.0726515) (3.89000,-0.0717472) (3.89250,-0.0708551) (3.89500,-0.0699750) (3.89750,-0.0691067) (3.90000,-0.0682501) (3.90250,-0.0674049) (3.90500,-0.0665711) (3.90750,-0.0657483) (3.91000,-0.0649365) (3.91250,-0.0641355) (3.91500,-0.0633450) (3.91750,-0.0625651) (3.92000,-0.0617954) (3.92250,-0.0610359) (3.92500,-0.0602864) (3.92750,-0.0595468) (3.93000,-0.0588168) (3.93250,-0.0580964) (3.93500,-0.0573855) (3.93750,-0.0566838) (3.94000,-0.0559912) (3.94250,-0.0553077) (3.94500,-0.0546330) (3.94750,-0.0539671) (3.95000,-0.0533098) (3.95250,-0.0526610) (3.95500,-0.0520206) (3.95750,-0.0513885) (3.96000,-0.0507645) (3.96250,-0.0501485) (3.96500,-0.0495404) (3.96750,-0.0489402) (3.97000,-0.0483476) (3.97250,-0.0477625) (3.97500,-0.0471850) (3.97750,-0.0466148) (3.98000,-0.0460519) (3.98250,-0.0454961) (3.98500,-0.0449474) (3.98750,-0.0444057) (3.99000,-0.0438708) (3.99250,-0.0433427) (3.99500,-0.0428213) (3.99750,-0.0423064) (4.00000,-0.0417981)};\addlegendentry{O}
\addplot[thick,dotted,blue] coordinates {(-2.00000,nan) (-1.99750,nan) (-1.99500,nan) (-1.99250,nan) (-1.99000,nan) (-1.98750,nan) (-1.98500,nan) (-1.98250,nan) (-1.98000,nan) (-1.97750,nan) (-1.97500,nan) (-1.97250,nan) (-1.97000,nan) (-1.96750,nan) (-1.96500,nan) (-1.96250,nan) (-1.96000,nan) (-1.95750,nan) (-1.95500,nan) (-1.95250,nan) (-1.95000,nan) (-1.94750,nan) (-1.94500,nan) (-1.94250,nan) (-1.94000,nan) (-1.93750,nan) (-1.93500,nan) (-1.93250,nan) (-1.93000,nan) (-1.92750,nan) (-1.92500,nan) (-1.92250,nan) (-1.92000,nan) (-1.91750,nan) (-1.91500,nan) (-1.91250,nan) (-1.91000,nan) (-1.90750,nan) (-1.90500,nan) (-1.90250,nan) (-1.90000,nan) (-1.89750,nan) (-1.89500,nan) (-1.89250,nan) (-1.89000,nan) (-1.88750,nan) (-1.88500,nan) (-1.88250,nan) (-1.88000,nan) (-1.87750,nan) (-1.87500,nan) (-1.87250,nan) (-1.87000,nan) (-1.86750,nan) (-1.86500,nan) (-1.86250,nan) (-1.86000,nan) (-1.85750,nan) (-1.85500,nan) (-1.85250,nan) (-1.85000,nan) (-1.84750,nan) (-1.84500,nan) (-1.84250,nan) (-1.84000,nan) (-1.83750,nan) (-1.83500,nan) (-1.83250,nan) (-1.83000,nan) (-1.82750,nan) (-1.82500,nan) (-1.82250,nan) (-1.82000,nan) (-1.81750,nan) (-1.81500,nan) (-1.81250,nan) (-1.81000,nan) (-1.80750,nan) (-1.80500,nan) (-1.80250,nan) (-1.80000,nan) (-1.79750,nan) (-1.79500,nan) (-1.79250,nan) (-1.79000,nan) (-1.78750,nan) (-1.78500,nan) (-1.78250,nan) (-1.78000,nan) (-1.77750,nan) (-1.77500,nan) (-1.77250,nan) (-1.77000,nan) (-1.76750,nan) (-1.76500,nan) (-1.76250,nan) (-1.76000,nan) (-1.75750,nan) (-1.75500,nan) (-1.75250,nan) (-1.75000,nan) (-1.74750,nan) (-1.74500,nan) (-1.74250,nan) (-1.74000,nan) (-1.73750,nan) (-1.73500,nan) (-1.73250,nan) (-1.73000,nan) (-1.72750,nan) (-1.72500,nan) (-1.72250,nan) (-1.72000,nan) (-1.71750,nan) (-1.71500,nan) (-1.71250,nan) (-1.71000,nan) (-1.70750,nan) (-1.70500,nan) (-1.70250,nan) (-1.70000,nan) (-1.69750,nan) (-1.69500,nan) (-1.69250,nan) (-1.69000,nan) (-1.68750,nan) (-1.68500,nan) (-1.68250,nan) (-1.68000,nan) (-1.67750,nan) (-1.67500,nan) (-1.67250,nan) (-1.67000,nan) (-1.66750,nan) (-1.66500,nan) (-1.66250,nan) (-1.66000,nan) (-1.65750,nan) (-1.65500,nan) (-1.65250,nan) (-1.65000,nan) (-1.64750,nan) (-1.64500,nan) (-1.64250,nan) (-1.64000,nan) (-1.63750,nan) (-1.63500,nan) (-1.63250,nan) (-1.63000,nan) (-1.62750,nan) (-1.62500,nan) (-1.62250,nan) (-1.62000,nan) (-1.61750,nan) (-1.61500,nan) (-1.61250,nan) (-1.61000,nan) (-1.60750,nan) (-1.60500,nan) (-1.60250,nan) (-1.60000,nan) (-1.59750,nan) (-1.59500,nan) (-1.59250,nan) (-1.59000,nan) (-1.58750,nan) (-1.58500,nan) (-1.58250,nan) (-1.58000,nan) (-1.57750,nan) (-1.57500,nan) (-1.57250,nan) (-1.57000,nan) (-1.56750,nan) (-1.56500,nan) (-1.56250,nan) (-1.56000,nan) (-1.55750,nan) (-1.55500,nan) (-1.55250,nan) (-1.55000,nan) (-1.54750,nan) (-1.54500,nan) (-1.54250,nan) (-1.54000,nan) (-1.53750,nan) (-1.53500,nan) (-1.53250,nan) (-1.53000,nan) (-1.52750,nan) (-1.52500,nan) (-1.52250,nan) (-1.52000,nan) (-1.51750,nan) (-1.51500,nan) (-1.51250,nan) (-1.51000,nan) (-1.50750,nan) (-1.50500,nan) (-1.50250,nan) (-1.50000,nan) (-1.49750,nan) (-1.49500,nan) (-1.49250,nan) (-1.49000,nan) (-1.48750,nan) (-1.48500,nan) (-1.48250,nan) (-1.48000,nan) (-1.47750,nan) (-1.47500,nan) (-1.47250,nan) (-1.47000,nan) (-1.46750,nan) (-1.46500,nan) (-1.46250,nan) (-1.46000,nan) (-1.45750,nan) (-1.45500,nan) (-1.45250,nan) (-1.45000,nan) (-1.44750,nan) (-1.44500,nan) (-1.44250,nan) (-1.44000,nan) (-1.43750,nan) (-1.43500,nan) (-1.43250,nan) (-1.43000,nan) (-1.42750,nan) (-1.42500,nan) (-1.42250,nan) (-1.42000,nan) (-1.41750,nan) (-1.41500,nan) (-1.41250,nan) (-1.41000,nan) (-1.40750,nan) (-1.40500,nan) (-1.40250,nan) (-1.40000,nan) (-1.39750,nan) (-1.39500,nan) (-1.39250,nan) (-1.39000,nan) (-1.38750,nan) (-1.38500,nan) (-1.38250,nan) (-1.38000,nan) (-1.37750,nan) (-1.37500,nan) (-1.37250,nan) (-1.37000,nan) (-1.36750,nan) (-1.36500,nan) (-1.36250,nan) (-1.36000,nan) (-1.35750,nan) (-1.35500,nan) (-1.35250,nan) (-1.35000,nan) (-1.34750,nan) (-1.34500,nan) (-1.34250,nan) (-1.34000,nan) (-1.33750,nan) (-1.33500,nan) (-1.33250,nan) (-1.33000,nan) (-1.32750,nan) (-1.32500,nan) (-1.32250,nan) (-1.32000,nan) (-1.31750,nan) (-1.31500,nan) (-1.31250,nan) (-1.31000,nan) (-1.30750,nan) (-1.30500,nan) (-1.30250,nan) (-1.30000,nan) (-1.29750,nan) (-1.29500,nan) (-1.29250,nan) (-1.29000,nan) (-1.28750,nan) (-1.28500,nan) (-1.28250,nan) (-1.28000,nan) (-1.27750,nan) (-1.27500,nan) (-1.27250,nan) (-1.27000,nan) (-1.26750,nan) (-1.26500,nan) (-1.26250,nan) (-1.26000,nan) (-1.25750,nan) (-1.25500,nan) (-1.25250,nan) (-1.25000,nan) (-1.24750,nan) (-1.24500,nan) (-1.24250,nan) (-1.24000,nan) (-1.23750,nan) (-1.23500,nan) (-1.23250,nan) (-1.23000,nan) (-1.22750,nan) (-1.22500,nan) (-1.22250,nan) (-1.22000,nan) (-1.21750,nan) (-1.21500,nan) (-1.21250,nan) (-1.21000,nan) (-1.20750,nan) (-1.20500,nan) (-1.20250,nan) (-1.20000,nan) (-1.19750,nan) (-1.19500,nan) (-1.19250,nan) (-1.19000,nan) (-1.18750,nan) (-1.18500,nan) (-1.18250,nan) (-1.18000,nan) (-1.17750,nan) (-1.17500,nan) (-1.17250,nan) (-1.17000,nan) (-1.16750,nan) (-1.16500,nan) (-1.16250,nan) (-1.16000,nan) (-1.15750,nan) (-1.15500,nan) (-1.15250,nan) (-1.15000,nan) (-1.14750,nan) (-1.14500,nan) (-1.14250,nan) (-1.14000,nan) (-1.13750,nan) (-1.13500,nan) (-1.13250,nan) (-1.13000,nan) (-1.12750,nan) (-1.12500,nan) (-1.12250,nan) (-1.12000,nan) (-1.11750,nan) (-1.11500,nan) (-1.11250,nan) (-1.11000,nan) (-1.10750,nan) (-1.10500,nan) (-1.10250,nan) (-1.10000,nan) (-1.09750,nan) (-1.09500,nan) (-1.09250,nan) (-1.09000,nan) (-1.08750,nan) (-1.08500,nan) (-1.08250,nan) (-1.08000,nan) (-1.07750,nan) (-1.07500,nan) (-1.07250,nan) (-1.07000,nan) (-1.06750,nan) (-1.06500,nan) (-1.06250,nan) (-1.06000,nan) (-1.05750,nan) (-1.05500,nan) (-1.05250,nan) (-1.05000,nan) (-1.04750,nan) (-1.04500,nan) (-1.04250,nan) (-1.04000,nan) (-1.03750,nan) (-1.03500,nan) (-1.03250,nan) (-1.03000,nan) (-1.02750,nan) (-1.02500,nan) (-1.02250,nan) (-1.02000,nan) (-1.01750,nan) (-1.01500,nan) (-1.01250,nan) (-1.01000,nan) (-1.00750,nan) (-1.00500,nan) (-1.00250,nan) (-1.00000,nan) (-0.99750,nan) (-0.99500,nan) (-0.99250,nan) (-0.99000,nan) (-0.98750,nan) (-0.98500,nan) (-0.98250,nan) (-0.98000,nan) (-0.97750,nan) (-0.97500,nan) (-0.97250,nan) (-0.97000,nan) (-0.96750,nan) (-0.96500,nan) (-0.96250,nan) (-0.96000,nan) (-0.95750,nan) (-0.95500,nan) (-0.95250,nan) (-0.95000,nan) (-0.94750,nan) (-0.94500,nan) (-0.94250,nan) (-0.94000,nan) (-0.93750,nan) (-0.93500,nan) (-0.93250,nan) (-0.93000,nan) (-0.92750,nan) (-0.92500,nan) (-0.92250,nan) (-0.92000,nan) (-0.91750,nan) (-0.91500,nan) (-0.91250,nan) (-0.91000,nan) (-0.90750,nan) (-0.90500,nan) (-0.90250,nan) (-0.90000,nan) (-0.89750,nan) (-0.89500,nan) (-0.89250,nan) (-0.89000,nan) (-0.88750,nan) (-0.88500,nan) (-0.88250,nan) (-0.88000,nan) (-0.87750,nan) (-0.87500,nan) (-0.87250,nan) (-0.87000,nan) (-0.86750,nan) (-0.86500,nan) (-0.86250,nan) (-0.86000,nan) (-0.85750,nan) (-0.85500,nan) (-0.85250,nan) (-0.85000,nan) (-0.84750,nan) (-0.84500,nan) (-0.84250,nan) (-0.84000,nan) (-0.83750,nan) (-0.83500,nan) (-0.83250,nan) (-0.83000,nan) (-0.82750,nan) (-0.82500,nan) (-0.82250,nan) (-0.82000,nan) (-0.81750,nan) (-0.81500,nan) (-0.81250,nan) (-0.81000,nan) (-0.80750,nan) (-0.80500,nan) (-0.80250,nan) (-0.80000,nan) (-0.79750,nan) (-0.79500,nan) (-0.79250,nan) (-0.79000,nan) (-0.78750,nan) (-0.78500,nan) (-0.78250,nan) (-0.78000,nan) (-0.77750,nan) (-0.77500,nan) (-0.77250,nan) (-0.77000,nan) (-0.76750,nan) (-0.76500,nan) (-0.76250,nan) (-0.76000,nan) (-0.75750,nan) (-0.75500,nan) (-0.75250,nan) (-0.75000,nan) (-0.74750,nan) (-0.74500,nan) (-0.74250,nan) (-0.74000,nan) (-0.73750,nan) (-0.73500,nan) (-0.73250,nan) (-0.73000,nan) (-0.72750,nan) (-0.72500,nan) (-0.72250,nan) (-0.72000,nan) (-0.71750,nan) (-0.71500,nan) (-0.71250,nan) (-0.71000,nan) (-0.70750,nan) (-0.70500,nan) (-0.70250,nan) (-0.70000,nan) (-0.69750,nan) (-0.69500,nan) (-0.69250,nan) (-0.69000,nan) (-0.68750,nan) (-0.68500,nan) (-0.68250,nan) (-0.68000,nan) (-0.67750,nan) (-0.67500,nan) (-0.67250,nan) (-0.67000,nan) (-0.66750,nan) (-0.66500,nan) (-0.66250,nan) (-0.66000,nan) (-0.65750,nan) (-0.65500,nan) (-0.65250,nan) (-0.65000,nan) (-0.64750,nan) (-0.64500,nan) (-0.64250,nan) (-0.64000,nan) (-0.63750,nan) (-0.63500,nan) (-0.63250,nan) (-0.63000,nan) (-0.62750,nan) (-0.62500,nan) (-0.62250,nan) (-0.62000,nan) (-0.61750,nan) (-0.61500,nan) (-0.61250,nan) (-0.61000,nan) (-0.60750,nan) (-0.60500,nan) (-0.60250,nan) (-0.60000,nan) (-0.59750,nan) (-0.59500,nan) (-0.59250,nan) (-0.59000,nan) (-0.58750,nan) (-0.58500,nan) (-0.58250,nan) (-0.58000,nan) (-0.57750,nan) (-0.57500,nan) (-0.57250,nan) (-0.57000,nan) (-0.56750,nan) (-0.56500,nan) (-0.56250,nan) (-0.56000,nan) (-0.55750,nan) (-0.55500,nan) (-0.55250,nan) (-0.55000,nan) (-0.54750,nan) (-0.54500,nan) (-0.54250,nan) (-0.54000,nan) (-0.53750,nan) (-0.53500,nan) (-0.53250,nan) (-0.53000,nan) (-0.52750,nan) (-0.52500,nan) (-0.52250,nan) (-0.52000,nan) (-0.51750,nan) (-0.51500,nan) (-0.51250,nan) (-0.51000,nan) (-0.50750,nan) (-0.50500,nan) (-0.50250,nan) (-0.50000,nan) (-0.49750,nan) (-0.49500,nan) (-0.49250,nan) (-0.49000,nan) (-0.48750,nan) (-0.48500,nan) (-0.48250,nan) (-0.48000,nan) (-0.47750,nan) (-0.47500,nan) (-0.47250,nan) (-0.47000,nan) (-0.46750,nan) (-0.46500,nan) (-0.46250,nan) (-0.46000,nan) (-0.45750,nan) (-0.45500,nan) (-0.45250,nan) (-0.45000,nan) (-0.44750,nan) (-0.44500,nan) (-0.44250,nan) (-0.44000,nan) (-0.43750,nan) (-0.43500,nan) (-0.43250,nan) (-0.43000,nan) (-0.42750,nan) (-0.42500,nan) (-0.42250,nan) (-0.42000,nan) (-0.41750,nan) (-0.41500,nan) (-0.41250,nan) (-0.41000,nan) (-0.40750,nan) (-0.40500,nan) (-0.40250,nan) (-0.40000,nan) (-0.39750,nan) (-0.39500,nan) (-0.39250,nan) (-0.39000,nan) (-0.38750,nan) (-0.38500,nan) (-0.38250,nan) (-0.38000,nan) (-0.37750,nan) (-0.37500,nan) (-0.37250,nan) (-0.37000,nan) (-0.36750,nan) (-0.36500,nan) (-0.36250,nan) (-0.36000,nan) (-0.35750,nan) (-0.35500,nan) (-0.35250,nan) (-0.35000,nan) (-0.34750,nan) (-0.34500,nan) (-0.34250,nan) (-0.34000,nan) (-0.33750,nan) (-0.33500,nan) (-0.33250,nan) (-0.33000,nan) (-0.32750,nan) (-0.32500,nan) (-0.32250,nan) (-0.32000,nan) (-0.31750,nan) (-0.31500,nan) (-0.31250,nan) (-0.31000,nan) (-0.30750,nan) (-0.30500,nan) (-0.30250,nan) (-0.30000,nan) (-0.29750,nan) (-0.29500,nan) (-0.29250,nan) (-0.29000,nan) (-0.28750,nan) (-0.28500,nan) (-0.28250,nan) (-0.28000,nan) (-0.27750,nan) (-0.27500,nan) (-0.27250,nan) (-0.27000,nan) (-0.26750,nan) (-0.26500,nan) (-0.26250,nan) (-0.26000,nan) (-0.25750,nan) (-0.25500,nan) (-0.25250,nan) (-0.25000,nan) (-0.24750,nan) (-0.24500,nan) (-0.24250,nan) (-0.24000,nan) (-0.23750,nan) (-0.23500,nan) (-0.23250,nan) (-0.23000,nan) (-0.22750,nan) (-0.22500,nan) (-0.22250,nan) (-0.22000,nan) (-0.21750,nan) (-0.21500,nan) (-0.21250,nan) (-0.21000,nan) (-0.20750,nan) (-0.20500,nan) (-0.20250,nan) (-0.20000,nan) (-0.19750,nan) (-0.19500,nan) (-0.19250,nan) (-0.19000,nan) (-0.18750,nan) (-0.18500,nan) (-0.18250,nan) (-0.18000,nan) (-0.17750,nan) (-0.17500,nan) (-0.17250,nan) (-0.17000,nan) (-0.16750,nan) (-0.16500,nan) (-0.16250,nan) (-0.16000,nan) (-0.15750,nan) (-0.15500,nan) (-0.15250,nan) (-0.15000,nan) (-0.14750,nan) (-0.14500,nan) (-0.14250,nan) (-0.14000,nan) (-0.13750,nan) (-0.13500,nan) (-0.13250,nan) (-0.13000,nan) (-0.12750,nan) (-0.12500,nan) (-0.12250,nan) (-0.12000,nan) (-0.11750,nan) (-0.11500,nan) (-0.11250,nan) (-0.11000,nan) (-0.10750,nan) (-0.10500,nan) (-0.10250,nan) (-0.10000,nan) (-0.09750,nan) (-0.09500,nan) (-0.09250,nan) (-0.09000,nan) (-0.08750,nan) (-0.08500,nan) (-0.08250,nan) (-0.08000,nan) (-0.07750,nan) (-0.07500,nan) (-0.07250,nan) (-0.07000,nan) (-0.06750,nan) (-0.06500,nan) (-0.06250,nan) (-0.06000,nan) (-0.05750,nan) (-0.05500,nan) (-0.05250,nan) (-0.05000,nan) (-0.04750,nan) (-0.04500,nan) (-0.04250,nan) (-0.04000,nan) (-0.03750,nan) (-0.03500,nan) (-0.03250,nan) (-0.03000,nan) (-0.02750,nan) (-0.02500,nan) (-0.02250,nan) (-0.02000,nan) (-0.01750,nan) (-0.01500,nan) (-0.01250,nan) (-0.01000,nan) (-0.00750,nan) (-0.00500,nan) (-0.00250,nan) (0.00000,nan) (0.00250,nan) (0.00500,nan) (0.00750,nan) (0.01000,nan) (0.01250,nan) (0.01500,nan) (0.01750,nan) (0.02000,nan) (0.02250,nan) (0.02500,nan) (0.02750,nan) (0.03000,nan) (0.03250,nan) (0.03500,nan) (0.03750,nan) (0.04000,nan) (0.04250,nan) (0.04500,nan) (0.04750,nan) (0.05000,nan) (0.05250,nan) (0.05500,nan) (0.05750,nan) (0.06000,nan) (0.06250,nan) (0.06500,nan) (0.06750,nan) (0.07000,nan) (0.07250,nan) (0.07500,nan) (0.07750,nan) (0.08000,nan) (0.08250,nan) (0.08500,nan) (0.08750,nan) (0.09000,nan) (0.09250,nan) (0.09500,nan) (0.09750,nan) (0.10000,nan) (0.10250,nan) (0.10500,nan) (0.10750,nan) (0.11000,nan) (0.11250,nan) (0.11500,nan) (0.11750,nan) (0.12000,nan) (0.12250,nan) (0.12500,nan) (0.12750,nan) (0.13000,nan) (0.13250,nan) (0.13500,nan) (0.13750,nan) (0.14000,nan) (0.14250,nan) (0.14500,nan) (0.14750,nan) (0.15000,nan) (0.15250,nan) (0.15500,nan) (0.15750,nan) (0.16000,nan) (0.16250,nan) (0.16500,nan) (0.16750,nan) (0.17000,nan) (0.17250,nan) (0.17500,nan) (0.17750,nan) (0.18000,nan) (0.18250,nan) (0.18500,nan) (0.18750,nan) (0.19000,nan) (0.19250,nan) (0.19500,nan) (0.19750,nan) (0.20000,nan) (0.20250,nan) (0.20500,nan) (0.20750,nan) (0.21000,nan) (0.21250,nan) (0.21500,nan) (0.21750,nan) (0.22000,nan) (0.22250,nan) (0.22500,nan) (0.22750,nan) (0.23000,nan) (0.23250,nan) (0.23500,nan) (0.23750,nan) (0.24000,nan) (0.24250,nan) (0.24500,nan) (0.24750,nan) (0.25000,nan) (0.25250,nan) (0.25500,nan) (0.25750,nan) (0.26000,nan) (0.26250,nan) (0.26500,nan) (0.26750,nan) (0.27000,nan) (0.27250,nan) (0.27500,nan) (0.27750,nan) (0.28000,nan) (0.28250,nan) (0.28500,nan) (0.28750,nan) (0.29000,nan) (0.29250,nan) (0.29500,nan) (0.29750,nan) (0.30000,nan) (0.30250,nan) (0.30500,nan) (0.30750,nan) (0.31000,nan) (0.31250,nan) (0.31500,nan) (0.31750,nan) (0.32000,nan) (0.32250,nan) (0.32500,nan) (0.32750,nan) (0.33000,nan) (0.33250,nan) (0.33500,nan) (0.33750,nan) (0.34000,nan) (0.34250,nan) (0.34500,nan) (0.34750,nan) (0.35000,nan) (0.35250,nan) (0.35500,nan) (0.35750,nan) (0.36000,nan) (0.36250,nan) (0.36500,nan) (0.36750,nan) (0.37000,nan) (0.37250,nan) (0.37500,nan) (0.37750,nan) (0.38000,nan) (0.38250,nan) (0.38500,nan) (0.38750,nan) (0.39000,nan) (0.39250,nan) (0.39500,nan) (0.39750,nan) (0.40000,nan) (0.40250,nan) (0.40500,nan) (0.40750,nan) (0.41000,nan) (0.41250,nan) (0.41500,nan) (0.41750,nan) (0.42000,nan) (0.42250,nan) (0.42500,nan) (0.42750,nan) (0.43000,nan) (0.43250,nan) (0.43500,nan) (0.43750,nan) (0.44000,nan) (0.44250,nan) (0.44500,nan) (0.44750,nan) (0.45000,nan) (0.45250,nan) (0.45500,nan) (0.45750,nan) (0.46000,nan) (0.46250,nan) (0.46500,nan) (0.46750,nan) (0.47000,nan) (0.47250,nan) (0.47500,nan) (0.47750,nan) (0.48000,nan) (0.48250,nan) (0.48500,nan) (0.48750,nan) (0.49000,nan) (0.49250,nan) (0.49500,nan) (0.49750,nan) (0.50000,nan) (0.50250,nan) (0.50500,nan) (0.50750,nan) (0.51000,nan) (0.51250,nan) (0.51500,nan) (0.51750,nan) (0.52000,nan) (0.52250,nan) (0.52500,nan) (0.52750,nan) (0.53000,nan) (0.53250,nan) (0.53500,nan) (0.53750,nan) (0.54000,nan) (0.54250,nan) (0.54500,nan) (0.54750,nan) (0.55000,nan) (0.55250,nan) (0.55500,nan) (0.55750,nan) (0.56000,nan) (0.56250,nan) (0.56500,nan) (0.56750,nan) (0.57000,nan) (0.57250,nan) (0.57500,nan) (0.57750,nan) (0.58000,nan) (0.58250,nan) (0.58500,nan) (0.58750,nan) (0.59000,nan) (0.59250,nan) (0.59500,nan) (0.59750,nan) (0.60000,nan) (0.60250,nan) (0.60500,nan) (0.60750,nan) (0.61000,nan) (0.61250,nan) (0.61500,nan) (0.61750,nan) (0.62000,nan) (0.62250,nan) (0.62500,nan) (0.62750,nan) (0.63000,nan) (0.63250,nan) (0.63500,nan) (0.63750,nan) (0.64000,nan) (0.64250,nan) (0.64500,nan) (0.64750,nan) (0.65000,nan) (0.65250,nan) (0.65500,nan) (0.65750,nan) (0.66000,nan) (0.66250,nan) (0.66500,nan) (0.66750,nan) (0.67000,nan) (0.67250,nan) (0.67500,nan) (0.67750,nan) (0.68000,nan) (0.68250,nan) (0.68500,nan) (0.68750,nan) (0.69000,nan) (0.69250,nan) (0.69500,nan) (0.69750,nan) (0.70000,nan) (0.70250,nan) (0.70500,nan) (0.70750,nan) (0.71000,nan) (0.71250,nan) (0.71500,nan) (0.71750,nan) (0.72000,nan) (0.72250,nan) (0.72500,nan) (0.72750,nan) (0.73000,nan) (0.73250,nan) (0.73500,nan) (0.73750,nan) (0.74000,nan) (0.74250,nan) (0.74500,nan) (0.74750,nan) (0.75000,nan) (0.75250,nan) (0.75500,nan) (0.75750,nan) (0.76000,nan) (0.76250,nan) (0.76500,nan) (0.76750,nan) (0.77000,nan) (0.77250,nan) (0.77500,nan) (0.77750,nan) (0.78000,nan) (0.78250,nan) (0.78500,nan) (0.78750,nan) (0.79000,nan) (0.79250,nan) (0.79500,nan) (0.79750,nan) (0.80000,nan) (0.80250,nan) (0.80500,nan) (0.80750,nan) (0.81000,nan) (0.81250,nan) (0.81500,nan) (0.81750,nan) (0.82000,nan) (0.82250,nan) (0.82500,nan) (0.82750,nan) (0.83000,nan) (0.83250,nan) (0.83500,nan) (0.83750,nan) (0.84000,nan) (0.84250,nan) (0.84500,nan) (0.84750,nan) (0.85000,nan) (0.85250,nan) (0.85500,nan) (0.85750,nan) (0.86000,nan) (0.86250,nan) (0.86500,nan) (0.86750,nan) (0.87000,nan) (0.87250,nan) (0.87500,nan) (0.87750,nan) (0.88000,nan) (0.88250,nan) (0.88500,nan) (0.88750,nan) (0.89000,nan) (0.89250,nan) (0.89500,nan) (0.89750,nan) (0.90000,nan) (0.90250,nan) (0.90500,nan) (0.90750,nan) (0.91000,nan) (0.91250,nan) (0.91500,nan) (0.91750,nan) (0.92000,nan) (0.92250,nan) (0.92500,nan) (0.92750,nan) (0.93000,nan) (0.93250,nan) (0.93500,nan) (0.93750,nan) (0.94000,nan) (0.94250,nan) (0.94500,nan) (0.94750,nan) (0.95000,nan) (0.95250,nan) (0.95500,nan) (0.95750,nan) (0.96000,nan) (0.96250,nan) (0.96500,nan) (0.96750,nan) (0.97000,nan) (0.97250,nan) (0.97500,nan) (0.97750,nan) (0.98000,nan) (0.98250,nan) (0.98500,nan) (0.98750,nan) (0.99000,nan) (0.99250,nan) (0.99500,nan) (0.99750,nan) (1.00000,nan) (1.00250,nan) (1.00500,nan) (1.00750,nan) (1.01000,nan) (1.01250,nan) (1.01500,nan) (1.01750,nan) (1.02000,nan) (1.02250,nan) (1.02500,nan) (1.02750,nan) (1.03000,nan) (1.03250,nan) (1.03500,nan) (1.03750,nan) (1.04000,nan) (1.04250,nan) (1.04500,nan) (1.04750,nan) (1.05000,nan) (1.05250,nan) (1.05500,nan) (1.05750,nan) (1.06000,nan) (1.06250,nan) (1.06500,nan) (1.06750,nan) (1.07000,nan) (1.07250,nan) (1.07500,nan) (1.07750,nan) (1.08000,nan) (1.08250,nan) (1.08500,nan) (1.08750,nan) (1.09000,nan) (1.09250,nan) (1.09500,nan) (1.09750,nan) (1.10000,nan) (1.10250,nan) (1.10500,nan) (1.10750,nan) (1.11000,nan) (1.11250,nan) (1.11500,nan) (1.11750,nan) (1.12000,nan) (1.12250,nan) (1.12500,nan) (1.12750,nan) (1.13000,nan) (1.13250,nan) (1.13500,nan) (1.13750,nan) (1.14000,nan) (1.14250,nan) (1.14500,nan) (1.14750,nan) (1.15000,nan) (1.15250,nan) (1.15500,nan) (1.15750,nan) (1.16000,nan) (1.16250,nan) (1.16500,nan) (1.16750,nan) (1.17000,nan) (1.17250,nan) (1.17500,nan) (1.17750,nan) (1.18000,nan) (1.18250,nan) (1.18500,nan) (1.18750,nan) (1.19000,nan) (1.19250,nan) (1.19500,nan) (1.19750,nan) (1.20000,nan) (1.20250,nan) (1.20500,nan) (1.20750,nan) (1.21000,nan) (1.21250,nan) (1.21500,nan) (1.21750,nan) (1.22000,nan) (1.22250,nan) (1.22500,nan) (1.22750,nan) (1.23000,nan) (1.23250,nan) (1.23500,nan) (1.23750,nan) (1.24000,nan) (1.24250,nan) (1.24500,nan) (1.24750,nan) (1.25000,nan) (1.25250,nan) (1.25500,nan) (1.25750,nan) (1.26000,nan) (1.26250,nan) (1.26500,nan) (1.26750,nan) (1.27000,nan) (1.27250,nan) (1.27500,nan) (1.27750,nan) (1.28000,nan) (1.28250,nan) (1.28500,nan) (1.28750,nan) (1.29000,nan) (1.29250,nan) (1.29500,nan) (1.29750,nan) (1.30000,nan) (1.30250,nan) (1.30500,nan) (1.30750,nan) (1.31000,nan) (1.31250,nan) (1.31500,nan) (1.31750,nan) (1.32000,nan) (1.32250,nan) (1.32500,nan) (1.32750,nan) (1.33000,nan) (1.33250,nan) (1.33500,nan) (1.33750,nan) (1.34000,nan) (1.34250,nan) (1.34500,nan) (1.34750,nan) (1.35000,nan) (1.35250,nan) (1.35500,nan) (1.35750,nan) (1.36000,nan) (1.36250,nan) (1.36500,nan) (1.36750,nan) (1.37000,nan) (1.37250,nan) (1.37500,nan) (1.37750,nan) (1.38000,nan) (1.38250,nan) (1.38500,nan) (1.38750,nan) (1.39000,nan) (1.39250,nan) (1.39500,nan) (1.39750,nan) (1.40000,nan) (1.40250,nan) (1.40500,nan) (1.40750,nan) (1.41000,nan) (1.41250,nan) (1.41500,nan) (1.41750,nan) (1.42000,nan) (1.42250,nan) (1.42500,nan) (1.42750,nan) (1.43000,nan) (1.43250,nan) (1.43500,nan) (1.43750,nan) (1.44000,nan) (1.44250,nan) (1.44500,nan) (1.44750,nan) (1.45000,nan) (1.45250,nan) (1.45500,nan) (1.45750,nan) (1.46000,nan) (1.46250,nan) (1.46500,nan) (1.46750,nan) (1.47000,nan) (1.47250,nan) (1.47500,nan) (1.47750,nan) (1.48000,nan) (1.48250,nan) (1.48500,nan) (1.48750,nan) (1.49000,nan) (1.49250,nan) (1.49500,nan) (1.49750,nan) (1.50000,nan) (1.50250,nan) (1.50500,nan) (1.50750,nan) (1.51000,nan) (1.51250,nan) (1.51500,nan) (1.51750,nan) (1.52000,nan) (1.52250,nan) (1.52500,nan) (1.52750,nan) (1.53000,nan) (1.53250,nan) (1.53500,nan) (1.53750,nan) (1.54000,nan) (1.54250,nan) (1.54500,nan) (1.54750,nan) (1.55000,nan) (1.55250,nan) (1.55500,nan) (1.55750,nan) (1.56000,nan) (1.56250,nan) (1.56500,nan) (1.56750,nan) (1.57000,nan) (1.57250,nan) (1.57500,nan) (1.57750,nan) (1.58000,nan) (1.58250,nan) (1.58500,nan) (1.58750,nan) (1.59000,nan) (1.59250,nan) (1.59500,nan) (1.59750,nan) (1.60000,nan) (1.60250,nan) (1.60500,nan) (1.60750,nan) (1.61000,nan) (1.61250,nan) (1.61500,nan) (1.61750,nan) (1.62000,nan) (1.62250,nan) (1.62500,nan) (1.62750,nan) (1.63000,nan) (1.63250,nan) (1.63500,nan) (1.63750,nan) (1.64000,nan) (1.64250,nan) (1.64500,nan) (1.64750,nan) (1.65000,nan) (1.65250,nan) (1.65500,nan) (1.65750,nan) (1.66000,nan) (1.66250,nan) (1.66500,nan) (1.66750,nan) (1.67000,nan) (1.67250,nan) (1.67500,nan) (1.67750,nan) (1.68000,nan) (1.68250,nan) (1.68500,nan) (1.68750,nan) (1.69000,nan) (1.69250,nan) (1.69500,nan) (1.69750,nan) (1.70000,nan) (1.70250,nan) (1.70500,nan) (1.70750,nan) (1.71000,nan) (1.71250,nan) (1.71500,nan) (1.71750,nan) (1.72000,nan) (1.72250,nan) (1.72500,nan) (1.72750,nan) (1.73000,nan) (1.73250,nan) (1.73500,nan) (1.73750,nan) (1.74000,nan) (1.74250,nan) (1.74500,nan) (1.74750,nan) (1.75000,nan) (1.75250,nan) (1.75500,nan) (1.75750,nan) (1.76000,nan) (1.76250,nan) (1.76500,nan) (1.76750,nan) (1.77000,nan) (1.77250,nan) (1.77500,nan) (1.77750,nan) (1.78000,nan) (1.78250,nan) (1.78500,nan) (1.78750,nan) (1.79000,nan) (1.79250,nan) (1.79500,nan) (1.79750,nan) (1.80000,nan) (1.80250,nan) (1.80500,nan) (1.80750,nan) (1.81000,nan) (1.81250,nan) (1.81500,nan) (1.81750,nan) (1.82000,nan) (1.82250,nan) (1.82500,nan) (1.82750,nan) (1.83000,nan) (1.83250,nan) (1.83500,nan) (1.83750,nan) (1.84000,nan) (1.84250,nan) (1.84500,nan) (1.84750,nan) (1.85000,nan) (1.85250,nan) (1.85500,nan) (1.85750,nan) (1.86000,nan) (1.86250,nan) (1.86500,nan) (1.86750,nan) (1.87000,nan) (1.87250,nan) (1.87500,nan) (1.87750,nan) (1.88000,nan) (1.88250,nan) (1.88500,nan) (1.88750,nan) (1.89000,nan) (1.89250,nan) (1.89500,nan) (1.89750,nan) (1.90000,nan) (1.90250,nan) (1.90500,nan) (1.90750,nan) (1.91000,nan) (1.91250,nan) (1.91500,nan) (1.91750,nan) (1.92000,nan) (1.92250,nan) (1.92500,nan) (1.92750,nan) (1.93000,nan) (1.93250,nan) (1.93500,nan) (1.93750,nan) (1.94000,nan) (1.94250,nan) (1.94500,nan) (1.94750,nan) (1.95000,nan) (1.95250,nan) (1.95500,nan) (1.95750,nan) (1.96000,nan) (1.96250,nan) (1.96500,nan) (1.96750,nan) (1.97000,nan) (1.97250,nan) (1.97500,nan) (1.97750,nan) (1.98000,nan) (1.98250,nan) (1.98500,nan) (1.98750,nan) (1.99000,nan) (1.99250,nan) (1.99500,nan) (1.99750,nan) (2.00000,nan) (2.00250,nan) (2.00500,nan) (2.00750,nan) (2.01000,nan) (2.01250,nan) (2.01500,nan) (2.01750,nan) (2.02000,nan) (2.02250,nan) (2.02500,nan) (2.02750,nan) (2.03000,nan) (2.03250,nan) (2.03500,nan) (2.03750,nan) (2.04000,nan) (2.04250,nan) (2.04500,nan) (2.04750,nan) (2.05000,nan) (2.05250,nan) (2.05500,nan) (2.05750,nan) (2.06000,nan) (2.06250,nan) (2.06500,nan) (2.06750,nan) (2.07000,nan) (2.07250,nan) (2.07500,nan) (2.07750,nan) (2.08000,nan) (2.08250,nan) (2.08500,nan) (2.08750,nan) (2.09000,nan) (2.09250,nan) (2.09500,nan) (2.09750,nan) (2.10000,nan) (2.10250,nan) (2.10500,nan) (2.10750,nan) (2.11000,nan) (2.11250,nan) (2.11500,nan) (2.11750,nan) (2.12000,nan) (2.12250,nan) (2.12500,nan) (2.12750,nan) (2.13000,nan) (2.13250,nan) (2.13500,nan) (2.13750,nan) (2.14000,nan) (2.14250,nan) (2.14500,nan) (2.14750,nan) (2.15000,nan) (2.15250,nan) (2.15500,nan) (2.15750,nan) (2.16000,nan) (2.16250,nan) (2.16500,nan) (2.16750,nan) (2.17000,nan) (2.17250,nan) (2.17500,nan) (2.17750,nan) (2.18000,nan) (2.18250,nan) (2.18500,nan) (2.18750,nan) (2.19000,nan) (2.19250,nan) (2.19500,nan) (2.19750,nan) (2.20000,nan) (2.20250,nan) (2.20500,nan) (2.20750,nan) (2.21000,nan) (2.21250,nan) (2.21500,nan) (2.21750,nan) (2.22000,nan) (2.22250,nan) (2.22500,nan) (2.22750,nan) (2.23000,nan) (2.23250,nan) (2.23500,nan) (2.23750,nan) (2.24000,nan) (2.24250,nan) (2.24500,nan) (2.24750,nan) (2.25000,nan) (2.25250,nan) (2.25500,nan) (2.25750,nan) (2.26000,nan) (2.26250,nan) (2.26500,nan) (2.26750,nan) (2.27000,nan) (2.27250,nan) (2.27500,nan) (2.27750,nan) (2.28000,nan) (2.28250,nan) (2.28500,nan) (2.28750,nan) (2.29000,nan) (2.29250,nan) (2.29500,nan) (2.29750,nan) (2.30000,nan) (2.30250,nan) (2.30500,nan) (2.30750,nan) (2.31000,nan) (2.31250,nan) (2.31500,nan) (2.31750,nan) (2.32000,nan) (2.32250,nan) (2.32500,nan) (2.32750,nan) (2.33000,nan) (2.33250,nan) (2.33500,nan) (2.33750,nan) (2.34000,nan) (2.34250,nan) (2.34500,nan) (2.34750,nan) (2.35000,nan) (2.35250,nan) (2.35500,nan) (2.35750,nan) (2.36000,nan) (2.36250,nan) (2.36500,nan) (2.36750,nan) (2.37000,nan) (2.37250,nan) (2.37500,nan) (2.37750,nan) (2.38000,nan) (2.38250,nan) (2.38500,nan) (2.38750,nan) (2.39000,nan) (2.39250,nan) (2.39500,nan) (2.39750,nan) (2.40000,nan) (2.40250,nan) (2.40500,nan) (2.40750,nan) (2.41000,nan) (2.41250,nan) (2.41500,nan) (2.41750,nan) (2.42000,nan) (2.42250,nan) (2.42500,nan) (2.42750,nan) (2.43000,nan) (2.43250,nan) (2.43500,nan) (2.43750,nan) (2.44000,nan) (2.44250,nan) (2.44500,nan) (2.44750,nan) (2.45000,nan) (2.45250,nan) (2.45500,nan) (2.45750,nan) (2.46000,nan) (2.46250,nan) (2.46500,nan) (2.46750,nan) (2.47000,nan) (2.47250,nan) (2.47500,nan) (2.47750,nan) (2.48000,nan) (2.48250,nan) (2.48500,nan) (2.48750,nan) (2.49000,nan) (2.49250,nan) (2.49500,nan) (2.49750,nan) (2.50000,nan) (2.50250,nan) (2.50500,nan) (2.50750,nan) (2.51000,nan) (2.51250,nan) (2.51500,nan) (2.51750,nan) (2.52000,nan) (2.52250,nan) (2.52500,nan) (2.52750,nan) (2.53000,nan) (2.53250,nan) (2.53500,nan) (2.53750,nan) (2.54000,nan) (2.54250,nan) (2.54500,nan) (2.54750,nan) (2.55000,nan) (2.55250,nan) (2.55500,nan) (2.55750,nan) (2.56000,nan) (2.56250,nan) (2.56500,nan) (2.56750,nan) (2.57000,nan) (2.57250,nan) (2.57500,nan) (2.57750,nan) (2.58000,nan) (2.58250,nan) (2.58500,nan) (2.58750,nan) (2.59000,nan) (2.59250,nan) (2.59500,nan) (2.59750,nan) (2.60000,nan) (2.60250,nan) (2.60500,nan) (2.60750,nan) (2.61000,nan) (2.61250,nan) (2.61500,nan) (2.61750,nan) (2.62000,nan) (2.62250,nan) (2.62500,nan) (2.62750,nan) (2.63000,nan) (2.63250,nan) (2.63500,nan) (2.63750,nan) (2.64000,nan) (2.64250,nan) (2.64500,nan) (2.64750,nan) (2.65000,nan) (2.65250,nan) (2.65500,nan) (2.65750,nan) (2.66000,nan) (2.66250,nan) (2.66500,nan) (2.66750,nan) (2.67000,nan) (2.67250,nan) (2.67500,nan) (2.67750,nan) (2.68000,nan) (2.68250,nan) (2.68500,nan) (2.68750,nan) (2.69000,nan) (2.69250,nan) (2.69500,nan) (2.69750,nan) (2.70000,nan) (2.70250,nan) (2.70500,nan) (2.70750,nan) (2.71000,nan) (2.71250,nan) (2.71500,nan) (2.71750,nan) (2.72000,nan) (2.72250,nan) (2.72500,nan) (2.72750,nan) (2.73000,nan) (2.73250,nan) (2.73500,nan) (2.73750,nan) (2.74000,nan) (2.74250,nan) (2.74500,nan) (2.74750,nan) (2.75000,nan) (2.75250,nan) (2.75500,nan) (2.75750,nan) (2.76000,nan) (2.76250,nan) (2.76500,nan) (2.76750,nan) (2.77000,nan) (2.77250,nan) (2.77500,nan) (2.77750,nan) (2.78000,nan) (2.78250,nan) (2.78500,nan) (2.78750,nan) (2.79000,nan) (2.79250,nan) (2.79500,nan) (2.79750,nan) (2.80000,nan) (2.80250,nan) (2.80500,nan) (2.80750,nan) (2.81000,nan) (2.81250,nan) (2.81500,nan) (2.81750,nan) (2.82000,nan) (2.82250,nan) (2.82500,nan) (2.82750,nan) (2.83000,nan) (2.83250,nan) (2.83500,nan) (2.83750,nan) (2.84000,nan) (2.84250,nan) (2.84500,nan) (2.84750,nan) (2.85000,nan) (2.85250,nan) (2.85500,nan) (2.85750,nan) (2.86000,nan) (2.86250,nan) (2.86500,nan) (2.86750,nan) (2.87000,nan) (2.87250,nan) (2.87500,nan) (2.87750,nan) (2.88000,nan) (2.88250,nan) (2.88500,nan) (2.88750,nan) (2.89000,nan) (2.89250,nan) (2.89500,nan) (2.89750,nan) (2.90000,nan) (2.90250,nan) (2.90500,nan) (2.90750,nan) (2.91000,nan) (2.91250,nan) (2.91500,nan) (2.91750,nan) (2.92000,nan) (2.92250,nan) (2.92500,nan) (2.92750,nan) (2.93000,nan) (2.93250,nan) (2.93500,nan) (2.93750,nan) (2.94000,nan) (2.94250,nan) (2.94500,nan) (2.94750,nan) (2.95000,nan) (2.95250,nan) (2.95500,nan) (2.95750,nan) (2.96000,nan) (2.96250,nan) (2.96500,nan) (2.96750,nan) (2.97000,nan) (2.97250,nan) (2.97500,nan) (2.97750,nan) (2.98000,nan) (2.98250,nan) (2.98500,nan) (2.98750,nan) (2.99000,nan) (2.99250,nan) (2.99500,nan) (2.99750,nan) (3.00000,nan) (3.00250,nan) (3.00500,nan) (3.00750,nan) (3.01000,nan) (3.01250,nan) (3.01500,nan) (3.01750,nan) (3.02000,nan) (3.02250,nan) (3.02500,nan) (3.02750,nan) (3.03000,nan) (3.03250,nan) (3.03500,nan) (3.03750,nan) (3.04000,nan) (3.04250,nan) (3.04500,nan) (3.04750,nan) (3.05000,nan) (3.05250,nan) (3.05500,nan) (3.05750,nan) (3.06000,nan) (3.06250,nan) (3.06500,nan) (3.06750,nan) (3.07000,nan) (3.07250,nan) (3.07500,nan) (3.07750,nan) (3.08000,nan) (3.08250,nan) (3.08500,nan) (3.08750,nan) (3.09000,nan) (3.09250,nan) (3.09500,nan) (3.09750,nan) (3.10000,nan) (3.10250,nan) (3.10500,nan) (3.10750,nan) (3.11000,nan) (3.11250,nan) (3.11500,nan) (3.11750,nan) (3.12000,nan) (3.12250,nan) (3.12500,nan) (3.12750,nan) (3.13000,nan) (3.13250,nan) (3.13500,nan) (3.13750,nan) (3.14000,nan) (3.14250,nan) (3.14500,nan) (3.14750,nan) (3.15000,nan) (3.15250,nan) (3.15500,nan) (3.15750,nan) (3.16000,nan) (3.16250,nan) (3.16500,nan) (3.16750,nan) (3.17000,nan) (3.17250,nan) (3.17500,nan) (3.17750,nan) (3.18000,nan) (3.18250,nan) (3.18500,nan) (3.18750,nan) (3.19000,nan) (3.19250,nan) (3.19500,nan) (3.19750,nan) (3.20000,nan) (3.20250,nan) (3.20500,nan) (3.20750,nan) (3.21000,nan) (3.21250,nan) (3.21500,nan) (3.21750,nan) (3.22000,nan) (3.22250,nan) (3.22500,nan) (3.22750,nan) (3.23000,nan) (3.23250,nan) (3.23500,-1.7699428) (3.23750,-1.4962627) (3.24000,-1.3311538) (3.24250,-1.2129172) (3.24500,-1.1209616) (3.24750,-1.0458482) (3.25000,-0.9824560) (3.25250,-0.9276934) (3.25500,-0.8795501) (3.25750,-0.8366450) (3.26000,-0.7979881) (3.26250,-0.7628458) (3.26500,-0.7306591) (3.26750,-0.7009922) (3.27000,-0.6734991) (3.27250,-0.6479001) (3.27500,-0.6239661) (3.27750,-0.6015070) (3.28000,-0.5803634) (3.28250,-0.5604001) (3.28500,-0.5415019) (3.28750,-0.5235692) (3.29000,-0.5065159) (3.29250,-0.4902667) (3.29500,-0.4747554) (3.29750,-0.4599237) (3.30000,-0.4457198) (3.30250,-0.4320975) (3.30500,-0.4190155) (3.30750,-0.4064364) (3.31000,-0.3943269) (3.31250,-0.3826565) (3.31500,-0.3713977) (3.31750,-0.3605252) (3.32000,-0.3500162) (3.32250,-0.3398495) (3.32500,-0.3300060) (3.32750,-0.3204676) (3.33000,-0.3112181) (3.33250,-0.3022423) (3.33500,-0.2935259) (3.33750,-0.2850561) (3.34000,-0.2768205) (3.34250,-0.2688079) (3.34500,-0.2610076) (3.34750,-0.2534098) (3.35000,-0.2460052) (3.35250,-0.2387850) (3.35500,-0.2317411) (3.35750,-0.2248657) (3.36000,-0.2181517) (3.36250,-0.2115921) (3.36500,-0.2051805) (3.36750,-0.1989107) (3.37000,-0.1927770) (3.37250,-0.1867739) (3.37500,-0.1808961) (3.37750,-0.1751386) (3.38000,-0.1694968) (3.38250,-0.1639662) (3.38500,-0.1585423) (3.38750,-0.1532212) (3.39000,-0.1479989) (3.39250,-0.1428716) (3.39500,-0.1378357) (3.39750,-0.1328879) (3.40000,-0.1280247) (3.40250,-0.1232429) (3.40500,-0.1185396) (3.40750,-0.1139118) (3.41000,-0.1093565) (3.41250,-0.1048711) (3.41500,-0.1004529) (3.41750,-0.0960992) (3.42000,-0.0918077) (3.42250,-0.0875759) (3.42500,-0.0834014) (3.42750,-0.0792820) (3.43000,-0.0752154) (3.43250,-0.0711994) (3.43500,-0.0672320) (3.43750,-0.0633110) (3.44000,-0.0594345) (3.44250,-0.0556004) (3.44500,-0.0518068) (3.44750,-0.0480519) (3.45000,-0.0443336) (3.45250,-0.0406501) (3.45500,-0.0369997) (3.45750,-0.0333805) (3.46000,-0.0297906) (3.46250,-0.0262284) (3.46500,-0.0226920) (3.46750,-0.0191797) (3.47000,-0.0156898) (3.47250,-0.0122204) (3.47500,-0.0087698) (3.47750,-0.0053363) (3.48000,-0.0019180) (3.48250,0.0014867) (3.48500,0.0048797) (3.48750,0.0082627) (3.49000,0.0116377) (3.49250,0.0150065) (3.49500,0.0183709) (3.49750,0.0217328) (3.50000,0.0250943) (3.50250,0.0284573) (3.50500,0.0318239) (3.50750,0.0351961) (3.51000,0.0385761) (3.51250,0.0419661) (3.51500,0.0453684) (3.51750,0.0487854) (3.52000,0.0522194) (3.52250,0.0556731) (3.52500,0.0591490) (3.52750,0.0626498) (3.53000,0.0661786) (3.53250,0.0697381) (3.53500,0.0733315) (3.53750,0.0769622) (3.54000,0.0806335) (3.54250,0.0843490) (3.54500,0.0881126) (3.54750,0.0919283) (3.55000,0.0958002) (3.55250,0.0997331) (3.55500,0.1037315) (3.55750,0.1078007) (3.56000,0.1119460) (3.56250,0.1161733) (3.56500,0.1204887) (3.56750,0.1248989) (3.57000,0.1294109) (3.57250,0.1340326) (3.57500,0.1387721) (3.57750,0.1436384) (3.58000,0.1486412) (3.58250,0.1537909) (3.58500,0.1590990) (3.58750,0.1645779) (3.59000,0.1702412) (3.59250,0.1761038) (3.59500,0.1821821) (3.59750,0.1884940) (3.60000,0.1950595) (3.60250,0.2019006) (3.60500,0.2090418) (3.60750,0.2165104) (3.61000,0.2243370) (3.61250,0.2325560) (3.61500,0.2412061) (3.61750,0.2503313) (3.62000,0.2599815) (3.62250,0.2702139) (3.62500,0.2810943) (3.62750,0.2926988) (3.63000,0.3051162) (3.63250,0.3184505) (3.63500,0.3328246) (3.63750,0.3483853) (3.64000,0.3653091) (3.64250,0.3838111) (3.64500,0.4041563) (3.64750,0.4266760) (3.65000,0.4517912) (3.65250,0.4800471) (3.65500,0.5121661) (3.65750,0.5491316) (3.66000,0.5923285) (3.66250,0.6437913) (3.66500,0.7066795) (3.66750,0.7862830) (3.67000,0.8924821) (3.67250,1.0472454) (3.67500,1.3189154) (3.67750,2.6035798) (3.68000,nan) (3.68250,nan) (3.68500,nan) (3.68750,nan) (3.69000,nan) (3.69250,nan) (3.69500,nan) (3.69750,nan) (3.70000,nan) (3.70250,nan) (3.70500,nan) (3.70750,nan) (3.71000,nan) (3.71250,nan) (3.71500,nan) (3.71750,nan) (3.72000,nan) (3.72250,nan) (3.72500,nan) (3.72750,nan) (3.73000,nan) (3.73250,-1.4612883) (3.73500,-1.1115604) (3.73750,-0.9331704) (3.74000,-0.8157202) (3.74250,-0.7296496) (3.74500,-0.6626449) (3.74750,-0.6083959) (3.75000,-0.5632407) (3.75250,-0.5248686) (3.75500,-0.4917308) (3.75750,-0.4627400) (3.76000,-0.4371056) (3.76250,-0.4142350) (3.76500,-0.3936741) (3.76750,-0.3750672) (3.77000,-0.3581312) (3.77250,-0.3426374) (3.77500,-0.3283986) (3.77750,-0.3152599) (3.78000,-0.3030919) (3.78250,-0.2917852) (3.78500,-0.2812470) (3.78750,-0.2713978) (3.79000,-0.2621689) (3.79250,-0.2535008) (3.79500,-0.2453418) (3.79750,-0.2376462) (3.80000,-0.2303740) (3.80250,-0.2234898) (3.80500,-0.2169621) (3.80750,-0.2107627) (3.81000,-0.2048667) (3.81250,-0.1992514) (3.81500,-0.1938965) (3.81750,-0.1887838) (3.82000,-0.1838965) (3.82250,-0.1792196) (3.82500,-0.1747393) (3.82750,-0.1704432) (3.83000,-0.1663197) (3.83250,-0.1623584) (3.83500,-0.1585496) (3.83750,-0.1548844) (3.84000,-0.1513547) (3.84250,-0.1479530) (3.84500,-0.1446721) (3.84750,-0.1415057) (3.85000,-0.1384477) (3.85250,-0.1354926) (3.85500,-0.1326351) (3.85750,-0.1298705) (3.86000,-0.1271942) (3.86250,-0.1246020) (3.86500,-0.1220899) (3.86750,-0.1196542) (3.87000,-0.1172916) (3.87250,-0.1149986) (3.87500,-0.1127723) (3.87750,-0.1106099) (3.88000,-0.1085085) (3.88250,-0.1064656) (3.88500,-0.1044789) (3.88750,-0.1025460) (3.89000,-0.1006648) (3.89250,-0.0988333) (3.89500,-0.0970495) (3.89750,-0.0953117) (3.90000,-0.0936180) (3.90250,-0.0919669) (3.90500,-0.0903569) (3.90750,-0.0887863) (3.91000,-0.0872538) (3.91250,-0.0857581) (3.91500,-0.0842980) (3.91750,-0.0828721) (3.92000,-0.0814793) (3.92250,-0.0801185) (3.92500,-0.0787888) (3.92750,-0.0774889) (3.93000,-0.0762181) (3.93250,-0.0749753) (3.93500,-0.0737598) (3.93750,-0.0725705) (3.94000,-0.0714069) (3.94250,-0.0702680) (3.94500,-0.0691531) (3.94750,-0.0680615) (3.95000,-0.0669925) (3.95250,-0.0659456) (3.95500,-0.0649199) (3.95750,-0.0639150) (3.96000,-0.0629303) (3.96250,-0.0619651) (3.96500,-0.0610190) (3.96750,-0.0600914) (3.97000,-0.0591819) (3.97250,-0.0582899) (3.97500,-0.0574150) (3.97750,-0.0565568) (3.98000,-0.0557147) (3.98250,-0.0548884) (3.98500,-0.0540776) (3.98750,-0.0532817) (3.99000,-0.0525005) (3.99250,-0.0517335) (3.99500,-0.0509804) (3.99750,-0.0502409) (4.00000,-0.0495147)};\addlegendentry{X}
\end{axis}\end{tikzpicture}\end{center}
\begin{center}\begin{tikzpicture}\begin{axis}[width=.82\textwidth,height=5cm,xlabel={$\log_{10}(\omega/\Omega_{ci})$},ylabel={$\log_{10}n_{\rm ref}^2$},xmin=-2,xmax=4,ymin=-2,ymax=3,unbounded coords=jump,title={$n_e=5\times10^{20}\,\mathrm m^{-3}$},legend style={font=\small,at={(1.02,1)},anchor=north west}]
\addplot[thick,structurecolor] coordinates {(-2.00000,nan) (-1.99750,nan) (-1.99500,nan) (-1.99250,nan) (-1.99000,nan) (-1.98750,nan) (-1.98500,nan) (-1.98250,nan) (-1.98000,nan) (-1.97750,nan) (-1.97500,nan) (-1.97250,nan) (-1.97000,nan) (-1.96750,nan) (-1.96500,nan) (-1.96250,nan) (-1.96000,nan) (-1.95750,nan) (-1.95500,nan) (-1.95250,nan) (-1.95000,nan) (-1.94750,nan) (-1.94500,nan) (-1.94250,nan) (-1.94000,nan) (-1.93750,nan) (-1.93500,nan) (-1.93250,nan) (-1.93000,nan) (-1.92750,nan) (-1.92500,nan) (-1.92250,nan) (-1.92000,nan) (-1.91750,nan) (-1.91500,nan) (-1.91250,nan) (-1.91000,nan) (-1.90750,nan) (-1.90500,nan) (-1.90250,nan) (-1.90000,nan) (-1.89750,nan) (-1.89500,nan) (-1.89250,nan) (-1.89000,nan) (-1.88750,nan) (-1.88500,nan) (-1.88250,nan) (-1.88000,nan) (-1.87750,nan) (-1.87500,nan) (-1.87250,nan) (-1.87000,nan) (-1.86750,nan) (-1.86500,nan) (-1.86250,nan) (-1.86000,nan) (-1.85750,nan) (-1.85500,nan) (-1.85250,nan) (-1.85000,nan) (-1.84750,nan) (-1.84500,nan) (-1.84250,nan) (-1.84000,nan) (-1.83750,nan) (-1.83500,nan) (-1.83250,nan) (-1.83000,nan) (-1.82750,nan) (-1.82500,nan) (-1.82250,nan) (-1.82000,nan) (-1.81750,nan) (-1.81500,nan) (-1.81250,nan) (-1.81000,nan) (-1.80750,nan) (-1.80500,nan) (-1.80250,nan) (-1.80000,nan) (-1.79750,nan) (-1.79500,nan) (-1.79250,nan) (-1.79000,nan) (-1.78750,nan) (-1.78500,nan) (-1.78250,nan) (-1.78000,nan) (-1.77750,nan) (-1.77500,nan) (-1.77250,nan) (-1.77000,nan) (-1.76750,nan) (-1.76500,nan) (-1.76250,nan) (-1.76000,nan) (-1.75750,nan) (-1.75500,nan) (-1.75250,nan) (-1.75000,nan) (-1.74750,nan) (-1.74500,nan) (-1.74250,nan) (-1.74000,nan) (-1.73750,nan) (-1.73500,nan) (-1.73250,nan) (-1.73000,nan) (-1.72750,nan) (-1.72500,nan) (-1.72250,nan) (-1.72000,nan) (-1.71750,nan) (-1.71500,nan) (-1.71250,nan) (-1.71000,nan) (-1.70750,nan) (-1.70500,nan) (-1.70250,nan) (-1.70000,nan) (-1.69750,nan) (-1.69500,nan) (-1.69250,nan) (-1.69000,nan) (-1.68750,nan) (-1.68500,nan) (-1.68250,nan) (-1.68000,nan) (-1.67750,nan) (-1.67500,nan) (-1.67250,nan) (-1.67000,nan) (-1.66750,nan) (-1.66500,nan) (-1.66250,nan) (-1.66000,nan) (-1.65750,nan) (-1.65500,nan) (-1.65250,nan) (-1.65000,nan) (-1.64750,nan) (-1.64500,nan) (-1.64250,nan) (-1.64000,nan) (-1.63750,nan) (-1.63500,nan) (-1.63250,nan) (-1.63000,nan) (-1.62750,nan) (-1.62500,nan) (-1.62250,nan) (-1.62000,nan) (-1.61750,nan) (-1.61500,nan) (-1.61250,nan) (-1.61000,nan) (-1.60750,nan) (-1.60500,nan) (-1.60250,nan) (-1.60000,nan) (-1.59750,nan) (-1.59500,nan) (-1.59250,nan) (-1.59000,nan) (-1.58750,nan) (-1.58500,nan) (-1.58250,nan) (-1.58000,nan) (-1.57750,nan) (-1.57500,nan) (-1.57250,nan) (-1.57000,nan) (-1.56750,nan) (-1.56500,nan) (-1.56250,nan) (-1.56000,nan) (-1.55750,nan) (-1.55500,nan) (-1.55250,nan) (-1.55000,nan) (-1.54750,nan) (-1.54500,nan) (-1.54250,nan) (-1.54000,nan) (-1.53750,nan) (-1.53500,nan) (-1.53250,nan) (-1.53000,nan) (-1.52750,nan) (-1.52500,nan) (-1.52250,nan) (-1.52000,nan) (-1.51750,nan) (-1.51500,nan) (-1.51250,nan) (-1.51000,nan) (-1.50750,nan) (-1.50500,nan) (-1.50250,nan) (-1.50000,nan) (-1.49750,nan) (-1.49500,nan) (-1.49250,nan) (-1.49000,nan) (-1.48750,nan) (-1.48500,nan) (-1.48250,nan) (-1.48000,nan) (-1.47750,nan) (-1.47500,nan) (-1.47250,nan) (-1.47000,nan) (-1.46750,nan) (-1.46500,nan) (-1.46250,nan) (-1.46000,nan) (-1.45750,nan) (-1.45500,nan) (-1.45250,nan) (-1.45000,nan) (-1.44750,nan) (-1.44500,nan) (-1.44250,nan) (-1.44000,nan) (-1.43750,nan) (-1.43500,nan) (-1.43250,nan) (-1.43000,nan) (-1.42750,nan) (-1.42500,nan) (-1.42250,nan) (-1.42000,nan) (-1.41750,nan) (-1.41500,nan) (-1.41250,nan) (-1.41000,nan) (-1.40750,nan) (-1.40500,nan) (-1.40250,nan) (-1.40000,nan) (-1.39750,nan) (-1.39500,nan) (-1.39250,nan) (-1.39000,nan) (-1.38750,nan) (-1.38500,nan) (-1.38250,nan) (-1.38000,nan) (-1.37750,nan) (-1.37500,nan) (-1.37250,nan) (-1.37000,nan) (-1.36750,nan) (-1.36500,nan) (-1.36250,nan) (-1.36000,nan) (-1.35750,nan) (-1.35500,nan) (-1.35250,nan) (-1.35000,nan) (-1.34750,nan) (-1.34500,nan) (-1.34250,nan) (-1.34000,nan) (-1.33750,nan) (-1.33500,nan) (-1.33250,nan) (-1.33000,nan) (-1.32750,nan) (-1.32500,nan) (-1.32250,nan) (-1.32000,nan) (-1.31750,nan) (-1.31500,nan) (-1.31250,nan) (-1.31000,nan) (-1.30750,nan) (-1.30500,nan) (-1.30250,nan) (-1.30000,nan) (-1.29750,nan) (-1.29500,nan) (-1.29250,nan) (-1.29000,nan) (-1.28750,nan) (-1.28500,nan) (-1.28250,nan) (-1.28000,nan) (-1.27750,nan) (-1.27500,nan) (-1.27250,nan) (-1.27000,nan) (-1.26750,nan) (-1.26500,nan) (-1.26250,nan) (-1.26000,nan) (-1.25750,nan) (-1.25500,nan) (-1.25250,nan) (-1.25000,nan) (-1.24750,nan) (-1.24500,nan) (-1.24250,nan) (-1.24000,nan) (-1.23750,nan) (-1.23500,nan) (-1.23250,nan) (-1.23000,nan) (-1.22750,nan) (-1.22500,nan) (-1.22250,nan) (-1.22000,nan) (-1.21750,nan) (-1.21500,nan) (-1.21250,nan) (-1.21000,nan) (-1.20750,nan) (-1.20500,nan) (-1.20250,nan) (-1.20000,nan) (-1.19750,nan) (-1.19500,nan) (-1.19250,nan) (-1.19000,nan) (-1.18750,nan) (-1.18500,nan) (-1.18250,nan) (-1.18000,nan) (-1.17750,nan) (-1.17500,nan) (-1.17250,nan) (-1.17000,nan) (-1.16750,nan) (-1.16500,nan) (-1.16250,nan) (-1.16000,nan) (-1.15750,nan) (-1.15500,nan) (-1.15250,nan) (-1.15000,nan) (-1.14750,nan) (-1.14500,nan) (-1.14250,nan) (-1.14000,nan) (-1.13750,nan) (-1.13500,nan) (-1.13250,nan) (-1.13000,nan) (-1.12750,nan) (-1.12500,nan) (-1.12250,nan) (-1.12000,nan) (-1.11750,nan) (-1.11500,nan) (-1.11250,nan) (-1.11000,nan) (-1.10750,nan) (-1.10500,nan) (-1.10250,nan) (-1.10000,nan) (-1.09750,nan) (-1.09500,nan) (-1.09250,nan) (-1.09000,nan) (-1.08750,nan) (-1.08500,nan) (-1.08250,nan) (-1.08000,nan) (-1.07750,nan) (-1.07500,nan) (-1.07250,nan) (-1.07000,nan) (-1.06750,nan) (-1.06500,nan) (-1.06250,nan) (-1.06000,nan) (-1.05750,nan) (-1.05500,nan) (-1.05250,nan) (-1.05000,nan) (-1.04750,nan) (-1.04500,nan) (-1.04250,nan) (-1.04000,nan) (-1.03750,nan) (-1.03500,nan) (-1.03250,nan) (-1.03000,nan) (-1.02750,nan) (-1.02500,nan) (-1.02250,nan) (-1.02000,nan) (-1.01750,nan) (-1.01500,nan) (-1.01250,nan) (-1.01000,nan) (-1.00750,nan) (-1.00500,nan) (-1.00250,nan) (-1.00000,nan) (-0.99750,nan) (-0.99500,nan) (-0.99250,nan) (-0.99000,nan) (-0.98750,nan) (-0.98500,nan) (-0.98250,nan) (-0.98000,nan) (-0.97750,nan) (-0.97500,nan) (-0.97250,nan) (-0.97000,nan) (-0.96750,nan) (-0.96500,nan) (-0.96250,nan) (-0.96000,nan) (-0.95750,nan) (-0.95500,nan) (-0.95250,nan) (-0.95000,nan) (-0.94750,nan) (-0.94500,nan) (-0.94250,nan) (-0.94000,nan) (-0.93750,nan) (-0.93500,nan) (-0.93250,nan) (-0.93000,nan) (-0.92750,nan) (-0.92500,nan) (-0.92250,nan) (-0.92000,nan) (-0.91750,nan) (-0.91500,nan) (-0.91250,nan) (-0.91000,nan) (-0.90750,nan) (-0.90500,nan) (-0.90250,nan) (-0.90000,nan) (-0.89750,nan) (-0.89500,nan) (-0.89250,nan) (-0.89000,nan) (-0.88750,nan) (-0.88500,nan) (-0.88250,nan) (-0.88000,nan) (-0.87750,nan) (-0.87500,nan) (-0.87250,nan) (-0.87000,nan) (-0.86750,nan) (-0.86500,nan) (-0.86250,nan) (-0.86000,nan) (-0.85750,nan) (-0.85500,nan) (-0.85250,nan) (-0.85000,nan) (-0.84750,nan) (-0.84500,nan) (-0.84250,nan) (-0.84000,nan) (-0.83750,nan) (-0.83500,nan) (-0.83250,nan) (-0.83000,nan) (-0.82750,nan) (-0.82500,nan) (-0.82250,nan) (-0.82000,nan) (-0.81750,nan) (-0.81500,nan) (-0.81250,nan) (-0.81000,nan) (-0.80750,nan) (-0.80500,nan) (-0.80250,nan) (-0.80000,nan) (-0.79750,nan) (-0.79500,nan) (-0.79250,nan) (-0.79000,nan) (-0.78750,nan) (-0.78500,nan) (-0.78250,nan) (-0.78000,nan) (-0.77750,nan) (-0.77500,nan) (-0.77250,nan) (-0.77000,nan) (-0.76750,nan) (-0.76500,nan) (-0.76250,nan) (-0.76000,nan) (-0.75750,nan) (-0.75500,nan) (-0.75250,nan) (-0.75000,nan) (-0.74750,nan) (-0.74500,nan) (-0.74250,nan) (-0.74000,nan) (-0.73750,nan) (-0.73500,nan) (-0.73250,nan) (-0.73000,nan) (-0.72750,nan) (-0.72500,nan) (-0.72250,nan) (-0.72000,nan) (-0.71750,nan) (-0.71500,nan) (-0.71250,nan) (-0.71000,nan) (-0.70750,nan) (-0.70500,nan) (-0.70250,nan) (-0.70000,nan) (-0.69750,nan) (-0.69500,nan) (-0.69250,nan) (-0.69000,nan) (-0.68750,nan) (-0.68500,nan) (-0.68250,nan) (-0.68000,nan) (-0.67750,nan) (-0.67500,nan) (-0.67250,nan) (-0.67000,nan) (-0.66750,nan) (-0.66500,nan) (-0.66250,nan) (-0.66000,nan) (-0.65750,nan) (-0.65500,nan) (-0.65250,nan) (-0.65000,nan) (-0.64750,nan) (-0.64500,nan) (-0.64250,nan) (-0.64000,nan) (-0.63750,nan) (-0.63500,nan) (-0.63250,nan) (-0.63000,nan) (-0.62750,nan) (-0.62500,nan) (-0.62250,nan) (-0.62000,nan) (-0.61750,nan) (-0.61500,nan) (-0.61250,nan) (-0.61000,nan) (-0.60750,nan) (-0.60500,nan) (-0.60250,nan) (-0.60000,nan) (-0.59750,nan) (-0.59500,nan) (-0.59250,nan) (-0.59000,nan) (-0.58750,nan) (-0.58500,nan) (-0.58250,nan) (-0.58000,nan) (-0.57750,nan) (-0.57500,nan) (-0.57250,nan) (-0.57000,nan) (-0.56750,nan) (-0.56500,nan) (-0.56250,nan) (-0.56000,nan) (-0.55750,nan) (-0.55500,nan) (-0.55250,nan) (-0.55000,nan) (-0.54750,nan) (-0.54500,nan) (-0.54250,nan) (-0.54000,nan) (-0.53750,nan) (-0.53500,nan) (-0.53250,nan) (-0.53000,nan) (-0.52750,nan) (-0.52500,nan) (-0.52250,nan) (-0.52000,nan) (-0.51750,nan) (-0.51500,nan) (-0.51250,nan) (-0.51000,nan) (-0.50750,nan) (-0.50500,nan) (-0.50250,nan) (-0.50000,nan) (-0.49750,nan) (-0.49500,nan) (-0.49250,nan) (-0.49000,nan) (-0.48750,nan) (-0.48500,nan) (-0.48250,nan) (-0.48000,nan) (-0.47750,nan) (-0.47500,nan) (-0.47250,nan) (-0.47000,nan) (-0.46750,nan) (-0.46500,nan) (-0.46250,nan) (-0.46000,nan) (-0.45750,nan) (-0.45500,nan) (-0.45250,nan) (-0.45000,nan) (-0.44750,nan) (-0.44500,nan) (-0.44250,nan) (-0.44000,nan) (-0.43750,nan) (-0.43500,nan) (-0.43250,nan) (-0.43000,nan) (-0.42750,nan) (-0.42500,nan) (-0.42250,nan) (-0.42000,nan) (-0.41750,nan) (-0.41500,nan) (-0.41250,nan) (-0.41000,nan) (-0.40750,nan) (-0.40500,nan) (-0.40250,nan) (-0.40000,nan) (-0.39750,nan) (-0.39500,nan) (-0.39250,nan) (-0.39000,nan) (-0.38750,nan) (-0.38500,nan) (-0.38250,nan) (-0.38000,nan) (-0.37750,nan) (-0.37500,nan) (-0.37250,nan) (-0.37000,nan) (-0.36750,nan) (-0.36500,nan) (-0.36250,nan) (-0.36000,nan) (-0.35750,nan) (-0.35500,nan) (-0.35250,nan) (-0.35000,nan) (-0.34750,nan) (-0.34500,nan) (-0.34250,nan) (-0.34000,nan) (-0.33750,nan) (-0.33500,nan) (-0.33250,nan) (-0.33000,nan) (-0.32750,nan) (-0.32500,nan) (-0.32250,nan) (-0.32000,nan) (-0.31750,nan) (-0.31500,nan) (-0.31250,nan) (-0.31000,nan) (-0.30750,nan) (-0.30500,nan) (-0.30250,nan) (-0.30000,nan) (-0.29750,nan) (-0.29500,nan) (-0.29250,nan) (-0.29000,nan) (-0.28750,nan) (-0.28500,nan) (-0.28250,nan) (-0.28000,nan) (-0.27750,nan) (-0.27500,nan) (-0.27250,nan) (-0.27000,nan) (-0.26750,nan) (-0.26500,nan) (-0.26250,nan) (-0.26000,nan) (-0.25750,nan) (-0.25500,nan) (-0.25250,nan) (-0.25000,nan) (-0.24750,nan) (-0.24500,nan) (-0.24250,nan) (-0.24000,nan) (-0.23750,nan) (-0.23500,nan) (-0.23250,nan) (-0.23000,nan) (-0.22750,nan) (-0.22500,nan) (-0.22250,nan) (-0.22000,nan) (-0.21750,nan) (-0.21500,nan) (-0.21250,nan) (-0.21000,nan) (-0.20750,nan) (-0.20500,nan) (-0.20250,nan) (-0.20000,nan) (-0.19750,nan) (-0.19500,nan) (-0.19250,nan) (-0.19000,nan) (-0.18750,nan) (-0.18500,nan) (-0.18250,nan) (-0.18000,nan) (-0.17750,nan) (-0.17500,nan) (-0.17250,nan) (-0.17000,nan) (-0.16750,nan) (-0.16500,nan) (-0.16250,nan) (-0.16000,nan) (-0.15750,nan) (-0.15500,nan) (-0.15250,nan) (-0.15000,nan) (-0.14750,nan) (-0.14500,nan) (-0.14250,nan) (-0.14000,nan) (-0.13750,nan) (-0.13500,nan) (-0.13250,nan) (-0.13000,nan) (-0.12750,nan) (-0.12500,nan) (-0.12250,nan) (-0.12000,nan) (-0.11750,nan) (-0.11500,nan) (-0.11250,nan) (-0.11000,nan) (-0.10750,nan) (-0.10500,nan) (-0.10250,nan) (-0.10000,nan) (-0.09750,nan) (-0.09500,nan) (-0.09250,nan) (-0.09000,nan) (-0.08750,nan) (-0.08500,nan) (-0.08250,nan) (-0.08000,nan) (-0.07750,nan) (-0.07500,nan) (-0.07250,nan) (-0.07000,nan) (-0.06750,nan) (-0.06500,nan) (-0.06250,nan) (-0.06000,nan) (-0.05750,nan) (-0.05500,nan) (-0.05250,nan) (-0.05000,nan) (-0.04750,nan) (-0.04500,nan) (-0.04250,nan) (-0.04000,nan) (-0.03750,nan) (-0.03500,nan) (-0.03250,nan) (-0.03000,nan) (-0.02750,nan) (-0.02500,nan) (-0.02250,nan) (-0.02000,nan) (-0.01750,nan) (-0.01500,nan) (-0.01250,nan) (-0.01000,nan) (-0.00750,nan) (-0.00500,nan) (-0.00250,nan) (0.00000,nan) (0.00250,nan) (0.00500,nan) (0.00750,nan) (0.01000,nan) (0.01250,nan) (0.01500,nan) (0.01750,nan) (0.02000,nan) (0.02250,nan) (0.02500,nan) (0.02750,nan) (0.03000,nan) (0.03250,nan) (0.03500,nan) (0.03750,nan) (0.04000,nan) (0.04250,nan) (0.04500,nan) (0.04750,nan) (0.05000,nan) (0.05250,nan) (0.05500,nan) (0.05750,nan) (0.06000,nan) (0.06250,nan) (0.06500,nan) (0.06750,nan) (0.07000,nan) (0.07250,nan) (0.07500,nan) (0.07750,nan) (0.08000,nan) (0.08250,nan) (0.08500,nan) (0.08750,nan) (0.09000,nan) (0.09250,nan) (0.09500,nan) (0.09750,nan) (0.10000,nan) (0.10250,nan) (0.10500,nan) (0.10750,nan) (0.11000,nan) (0.11250,nan) (0.11500,nan) (0.11750,nan) (0.12000,nan) (0.12250,nan) (0.12500,nan) (0.12750,nan) (0.13000,nan) (0.13250,nan) (0.13500,nan) (0.13750,nan) (0.14000,nan) (0.14250,nan) (0.14500,nan) (0.14750,nan) (0.15000,nan) (0.15250,nan) (0.15500,nan) (0.15750,nan) (0.16000,nan) (0.16250,nan) (0.16500,nan) (0.16750,nan) (0.17000,nan) (0.17250,nan) (0.17500,nan) (0.17750,nan) (0.18000,nan) (0.18250,nan) (0.18500,nan) (0.18750,nan) (0.19000,nan) (0.19250,nan) (0.19500,nan) (0.19750,nan) (0.20000,nan) (0.20250,nan) (0.20500,nan) (0.20750,nan) (0.21000,nan) (0.21250,nan) (0.21500,nan) (0.21750,nan) (0.22000,nan) (0.22250,nan) (0.22500,nan) (0.22750,nan) (0.23000,nan) (0.23250,nan) (0.23500,nan) (0.23750,nan) (0.24000,nan) (0.24250,nan) (0.24500,nan) (0.24750,nan) (0.25000,nan) (0.25250,nan) (0.25500,nan) (0.25750,nan) (0.26000,nan) (0.26250,nan) (0.26500,nan) (0.26750,nan) (0.27000,nan) (0.27250,nan) (0.27500,nan) (0.27750,nan) (0.28000,nan) (0.28250,nan) (0.28500,nan) (0.28750,nan) (0.29000,nan) (0.29250,nan) (0.29500,nan) (0.29750,nan) (0.30000,nan) (0.30250,nan) (0.30500,nan) (0.30750,nan) (0.31000,nan) (0.31250,nan) (0.31500,nan) (0.31750,nan) (0.32000,nan) (0.32250,nan) (0.32500,nan) (0.32750,nan) (0.33000,nan) (0.33250,nan) (0.33500,nan) (0.33750,nan) (0.34000,nan) (0.34250,nan) (0.34500,nan) (0.34750,nan) (0.35000,nan) (0.35250,nan) (0.35500,nan) (0.35750,nan) (0.36000,nan) (0.36250,nan) (0.36500,nan) (0.36750,nan) (0.37000,nan) (0.37250,nan) (0.37500,nan) (0.37750,nan) (0.38000,nan) (0.38250,nan) (0.38500,nan) (0.38750,nan) (0.39000,nan) (0.39250,nan) (0.39500,nan) (0.39750,nan) (0.40000,nan) (0.40250,nan) (0.40500,nan) (0.40750,nan) (0.41000,nan) (0.41250,nan) (0.41500,nan) (0.41750,nan) (0.42000,nan) (0.42250,nan) (0.42500,nan) (0.42750,nan) (0.43000,nan) (0.43250,nan) (0.43500,nan) (0.43750,nan) (0.44000,nan) (0.44250,nan) (0.44500,nan) (0.44750,nan) (0.45000,nan) (0.45250,nan) (0.45500,nan) (0.45750,nan) (0.46000,nan) (0.46250,nan) (0.46500,nan) (0.46750,nan) (0.47000,nan) (0.47250,nan) (0.47500,nan) (0.47750,nan) (0.48000,nan) (0.48250,nan) (0.48500,nan) (0.48750,nan) (0.49000,nan) (0.49250,nan) (0.49500,nan) (0.49750,nan) (0.50000,nan) (0.50250,nan) (0.50500,nan) (0.50750,nan) (0.51000,nan) (0.51250,nan) (0.51500,nan) (0.51750,nan) (0.52000,nan) (0.52250,nan) (0.52500,nan) (0.52750,nan) (0.53000,nan) (0.53250,nan) (0.53500,nan) (0.53750,nan) (0.54000,nan) (0.54250,nan) (0.54500,nan) (0.54750,nan) (0.55000,nan) (0.55250,nan) (0.55500,nan) (0.55750,nan) (0.56000,nan) (0.56250,nan) (0.56500,nan) (0.56750,nan) (0.57000,nan) (0.57250,nan) (0.57500,nan) (0.57750,nan) (0.58000,nan) (0.58250,nan) (0.58500,nan) (0.58750,nan) (0.59000,nan) (0.59250,nan) (0.59500,nan) (0.59750,nan) (0.60000,nan) (0.60250,nan) (0.60500,nan) (0.60750,nan) (0.61000,nan) (0.61250,nan) (0.61500,nan) (0.61750,nan) (0.62000,nan) (0.62250,nan) (0.62500,nan) (0.62750,nan) (0.63000,nan) (0.63250,nan) (0.63500,nan) (0.63750,nan) (0.64000,nan) (0.64250,nan) (0.64500,nan) (0.64750,nan) (0.65000,nan) (0.65250,nan) (0.65500,nan) (0.65750,nan) (0.66000,nan) (0.66250,nan) (0.66500,nan) (0.66750,nan) (0.67000,nan) (0.67250,nan) (0.67500,nan) (0.67750,nan) (0.68000,nan) (0.68250,nan) (0.68500,nan) (0.68750,nan) (0.69000,nan) (0.69250,nan) (0.69500,nan) (0.69750,nan) (0.70000,nan) (0.70250,nan) (0.70500,nan) (0.70750,nan) (0.71000,nan) (0.71250,nan) (0.71500,nan) (0.71750,nan) (0.72000,nan) (0.72250,2.9988552) (0.72500,2.9967582) (0.72750,2.9946592) (0.73000,2.9925584) (0.73250,2.9904557) (0.73500,2.9883512) (0.73750,2.9862448) (0.74000,2.9841365) (0.74250,2.9820264) (0.74500,2.9799145) (0.74750,2.9778007) (0.75000,2.9756852) (0.75250,2.9735678) (0.75500,2.9714486) (0.75750,2.9693276) (0.76000,2.9672049) (0.76250,2.9650803) (0.76500,2.9629540) (0.76750,2.9608259) (0.77000,2.9586961) (0.77250,2.9565645) (0.77500,2.9544312) (0.77750,2.9522961) (0.78000,2.9501593) (0.78250,2.9480208) (0.78500,2.9458806) (0.78750,2.9437387) (0.79000,2.9415951) (0.79250,2.9394497) (0.79500,2.9373027) (0.79750,2.9351541) (0.80000,2.9330037) (0.80250,2.9308517) (0.80500,2.9286981) (0.80750,2.9265428) (0.81000,2.9243859) (0.81250,2.9222273) (0.81500,2.9200671) (0.81750,2.9179053) (0.82000,2.9157419) (0.82250,2.9135769) (0.82500,2.9114103) (0.82750,2.9092421) (0.83000,2.9070723) (0.83250,2.9049009) (0.83500,2.9027280) (0.83750,2.9005536) (0.84000,2.8983775) (0.84250,2.8961999) (0.84500,2.8940208) (0.84750,2.8918402) (0.85000,2.8896580) (0.85250,2.8874743) (0.85500,2.8852891) (0.85750,2.8831025) (0.86000,2.8809143) (0.86250,2.8787246) (0.86500,2.8765334) (0.86750,2.8743408) (0.87000,2.8721467) (0.87250,2.8699511) (0.87500,2.8677541) (0.87750,2.8655556) (0.88000,2.8633557) (0.88250,2.8611544) (0.88500,2.8589516) (0.88750,2.8567474) (0.89000,2.8545418) (0.89250,2.8523348) (0.89500,2.8501264) (0.89750,2.8479166) (0.90000,2.8457054) (0.90250,2.8434928) (0.90500,2.8412789) (0.90750,2.8390636) (0.91000,2.8368469) (0.91250,2.8346289) (0.91500,2.8324095) (0.91750,2.8301888) (0.92000,2.8279667) (0.92250,2.8257434) (0.92500,2.8235187) (0.92750,2.8212927) (0.93000,2.8190654) (0.93250,2.8168367) (0.93500,2.8146068) (0.93750,2.8123756) (0.94000,2.8101432) (0.94250,2.8079094) (0.94500,2.8056744) (0.94750,2.8034381) (0.95000,2.8012006) (0.95250,2.7989618) (0.95500,2.7967217) (0.95750,2.7944805) (0.96000,2.7922380) (0.96250,2.7899942) (0.96500,2.7877493) (0.96750,2.7855031) (0.97000,2.7832557) (0.97250,2.7810072) (0.97500,2.7787574) (0.97750,2.7765065) (0.98000,2.7742543) (0.98250,2.7720010) (0.98500,2.7697466) (0.98750,2.7674909) (0.99000,2.7652341) (0.99250,2.7629762) (0.99500,2.7607171) (0.99750,2.7584569) (1.00000,2.7561955) (1.00250,2.7539330) (1.00500,2.7516694) (1.00750,2.7494047) (1.01000,2.7471388) (1.01250,2.7448719) (1.01500,2.7426038) (1.01750,2.7403347) (1.02000,2.7380645) (1.02250,2.7357932) (1.02500,2.7335208) (1.02750,2.7312474) (1.03000,2.7289729) (1.03250,2.7266973) (1.03500,2.7244207) (1.03750,2.7221431) (1.04000,2.7198644) (1.04250,2.7175846) (1.04500,2.7153039) (1.04750,2.7130221) (1.05000,2.7107393) (1.05250,2.7084555) (1.05500,2.7061707) (1.05750,2.7038849) (1.06000,2.7015980) (1.06250,2.6993102) (1.06500,2.6970215) (1.06750,2.6947317) (1.07000,2.6924410) (1.07250,2.6901493) (1.07500,2.6878566) (1.07750,2.6855630) (1.08000,2.6832684) (1.08250,2.6809729) (1.08500,2.6786764) (1.08750,2.6763790) (1.09000,2.6740807) (1.09250,2.6717815) (1.09500,2.6694813) (1.09750,2.6671802) (1.10000,2.6648782) (1.10250,2.6625753) (1.10500,2.6602716) (1.10750,2.6579669) (1.11000,2.6556613) (1.11250,2.6533549) (1.11500,2.6510475) (1.11750,2.6487393) (1.12000,2.6464303) (1.12250,2.6441203) (1.12500,2.6418096) (1.12750,2.6394979) (1.13000,2.6371854) (1.13250,2.6348721) (1.13500,2.6325580) (1.13750,2.6302430) (1.14000,2.6279272) (1.14250,2.6256105) (1.14500,2.6232931) (1.14750,2.6209748) (1.15000,2.6186558) (1.15250,2.6163359) (1.15500,2.6140152) (1.15750,2.6116938) (1.16000,2.6093715) (1.16250,2.6070485) (1.16500,2.6047247) (1.16750,2.6024002) (1.17000,2.6000748) (1.17250,2.5977487) (1.17500,2.5954219) (1.17750,2.5930943) (1.18000,2.5907659) (1.18250,2.5884368) (1.18500,2.5861070) (1.18750,2.5837764) (1.19000,2.5814451) (1.19250,2.5791131) (1.19500,2.5767804) (1.19750,2.5744469) (1.20000,2.5721127) (1.20250,2.5697779) (1.20500,2.5674423) (1.20750,2.5651060) (1.21000,2.5627691) (1.21250,2.5604314) (1.21500,2.5580931) (1.21750,2.5557541) (1.22000,2.5534144) (1.22250,2.5510741) (1.22500,2.5487331) (1.22750,2.5463914) (1.23000,2.5440491) (1.23250,2.5417061) (1.23500,2.5393625) (1.23750,2.5370182) (1.24000,2.5346733) (1.24250,2.5323278) (1.24500,2.5299817) (1.24750,2.5276349) (1.25000,2.5252875) (1.25250,2.5229395) (1.25500,2.5205908) (1.25750,2.5182416) (1.26000,2.5158918) (1.26250,2.5135413) (1.26500,2.5111903) (1.26750,2.5088387) (1.27000,2.5064865) (1.27250,2.5041337) (1.27500,2.5017803) (1.27750,2.4994264) (1.28000,2.4970719) (1.28250,2.4947169) (1.28500,2.4923612) (1.28750,2.4900051) (1.29000,2.4876483) (1.29250,2.4852911) (1.29500,2.4829332) (1.29750,2.4805749) (1.30000,2.4782160) (1.30250,2.4758566) (1.30500,2.4734966) (1.30750,2.4711362) (1.31000,2.4687752) (1.31250,2.4664137) (1.31500,2.4640517) (1.31750,2.4616892) (1.32000,2.4593261) (1.32250,2.4569626) (1.32500,2.4545986) (1.32750,2.4522341) (1.33000,2.4498691) (1.33250,2.4475037) (1.33500,2.4451377) (1.33750,2.4427713) (1.34000,2.4404044) (1.34250,2.4380371) (1.34500,2.4356692) (1.34750,2.4333010) (1.35000,2.4309322) (1.35250,2.4285631) (1.35500,2.4261934) (1.35750,2.4238234) (1.36000,2.4214529) (1.36250,2.4190819) (1.36500,2.4167105) (1.36750,2.4143387) (1.37000,2.4119665) (1.37250,2.4095939) (1.37500,2.4072208) (1.37750,2.4048473) (1.38000,2.4024735) (1.38250,2.4000992) (1.38500,2.3977245) (1.38750,2.3953494) (1.39000,2.3929740) (1.39250,2.3905981) (1.39500,2.3882218) (1.39750,2.3858452) (1.40000,2.3834682) (1.40250,2.3810908) (1.40500,2.3787131) (1.40750,2.3763350) (1.41000,2.3739565) (1.41250,2.3715776) (1.41500,2.3691984) (1.41750,2.3668189) (1.42000,2.3644390) (1.42250,2.3620588) (1.42500,2.3596782) (1.42750,2.3572973) (1.43000,2.3549160) (1.43250,2.3525344) (1.43500,2.3501525) (1.43750,2.3477703) (1.44000,2.3453877) (1.44250,2.3430048) (1.44500,2.3406217) (1.44750,2.3382382) (1.45000,2.3358544) (1.45250,2.3334703) (1.45500,2.3310859) (1.45750,2.3287012) (1.46000,2.3263162) (1.46250,2.3239309) (1.46500,2.3215454) (1.46750,2.3191595) (1.47000,2.3167734) (1.47250,2.3143870) (1.47500,2.3120004) (1.47750,2.3096135) (1.48000,2.3072263) (1.48250,2.3048388) (1.48500,2.3024511) (1.48750,2.3000632) (1.49000,2.2976750) (1.49250,2.2952865) (1.49500,2.2928978) (1.49750,2.2905089) (1.50000,2.2881197) (1.50250,2.2857303) (1.50500,2.2833407) (1.50750,2.2809509) (1.51000,2.2785608) (1.51250,2.2761705) (1.51500,2.2737800) (1.51750,2.2713893) (1.52000,2.2689984) (1.52250,2.2666072) (1.52500,2.2642159) (1.52750,2.2618244) (1.53000,2.2594327) (1.53250,2.2570407) (1.53500,2.2546486) (1.53750,2.2522564) (1.54000,2.2498639) (1.54250,2.2474712) (1.54500,2.2450784) (1.54750,2.2426854) (1.55000,2.2402923) (1.55250,2.2378990) (1.55500,2.2355055) (1.55750,2.2331118) (1.56000,2.2307180) (1.56250,2.2283241) (1.56500,2.2259300) (1.56750,2.2235358) (1.57000,2.2211414) (1.57250,2.2187469) (1.57500,2.2163522) (1.57750,2.2139575) (1.58000,2.2115625) (1.58250,2.2091675) (1.58500,2.2067724) (1.58750,2.2043771) (1.59000,2.2019817) (1.59250,2.1995862) (1.59500,2.1971906) (1.59750,2.1947949) (1.60000,2.1923991) (1.60250,2.1900032) (1.60500,2.1876072) (1.60750,2.1852111) (1.61000,2.1828149) (1.61250,2.1804186) (1.61500,2.1780223) (1.61750,2.1756259) (1.62000,2.1732294) (1.62250,2.1708328) (1.62500,2.1684361) (1.62750,2.1660394) (1.63000,2.1636427) (1.63250,2.1612458) (1.63500,2.1588490) (1.63750,2.1564520) (1.64000,2.1540550) (1.64250,2.1516580) (1.64500,2.1492609) (1.64750,2.1468638) (1.65000,2.1444667) (1.65250,2.1420695) (1.65500,2.1396723) (1.65750,2.1372751) (1.66000,2.1348778) (1.66250,2.1324805) (1.66500,2.1300833) (1.66750,2.1276859) (1.67000,2.1252886) (1.67250,2.1228913) (1.67500,2.1204940) (1.67750,2.1180967) (1.68000,2.1156994) (1.68250,2.1133020) (1.68500,2.1109047) (1.68750,2.1085075) (1.69000,2.1061102) (1.69250,2.1037129) (1.69500,2.1013157) (1.69750,2.0989185) (1.70000,2.0965213) (1.70250,2.0941242) (1.70500,2.0917271) (1.70750,2.0893301) (1.71000,2.0869330) (1.71250,2.0845361) (1.71500,2.0821391) (1.71750,2.0797423) (1.72000,2.0773455) (1.72250,2.0749487) (1.72500,2.0725520) (1.72750,2.0701554) (1.73000,2.0677589) (1.73250,2.0653624) (1.73500,2.0629660) (1.73750,2.0605696) (1.74000,2.0581734) (1.74250,2.0557773) (1.74500,2.0533812) (1.74750,2.0509852) (1.75000,2.0485893) (1.75250,2.0461936) (1.75500,2.0437979) (1.75750,2.0414023) (1.76000,2.0390069) (1.76250,2.0366115) (1.76500,2.0342163) (1.76750,2.0318212) (1.77000,2.0294262) (1.77250,2.0270313) (1.77500,2.0246366) (1.77750,2.0222420) (1.78000,2.0198476) (1.78250,2.0174533) (1.78500,2.0150591) (1.78750,2.0126651) (1.79000,2.0102712) (1.79250,2.0078775) (1.79500,2.0054839) (1.79750,2.0030905) (1.80000,2.0006972) (1.80250,1.9983042) (1.80500,1.9959113) (1.80750,1.9935185) (1.81000,1.9911260) (1.81250,1.9887336) (1.81500,1.9863414) (1.81750,1.9839494) (1.82000,1.9815576) (1.82250,1.9791660) (1.82500,1.9767746) (1.82750,1.9743834) (1.83000,1.9719924) (1.83250,1.9696016) (1.83500,1.9672111) (1.83750,1.9648207) (1.84000,1.9624306) (1.84250,1.9600406) (1.84500,1.9576509) (1.84750,1.9552615) (1.85000,1.9528723) (1.85250,1.9504833) (1.85500,1.9480945) (1.85750,1.9457060) (1.86000,1.9433178) (1.86250,1.9409298) (1.86500,1.9385420) (1.86750,1.9361546) (1.87000,1.9337673) (1.87250,1.9313804) (1.87500,1.9289937) (1.87750,1.9266073) (1.88000,1.9242212) (1.88250,1.9218353) (1.88500,1.9194497) (1.88750,1.9170645) (1.89000,1.9146795) (1.89250,1.9122948) (1.89500,1.9099104) (1.89750,1.9075263) (1.90000,1.9051426) (1.90250,1.9027591) (1.90500,1.9003760) (1.90750,1.8979931) (1.91000,1.8956106) (1.91250,1.8932284) (1.91500,1.8908466) (1.91750,1.8884651) (1.92000,1.8860839) (1.92250,1.8837031) (1.92500,1.8813226) (1.92750,1.8789424) (1.93000,1.8765627) (1.93250,1.8741832) (1.93500,1.8718042) (1.93750,1.8694255) (1.94000,1.8670471) (1.94250,1.8646692) (1.94500,1.8622916) (1.94750,1.8599144) (1.95000,1.8575376) (1.95250,1.8551612) (1.95500,1.8527851) (1.95750,1.8504095) (1.96000,1.8480343) (1.96250,1.8456594) (1.96500,1.8432850) (1.96750,1.8409110) (1.97000,1.8385374) (1.97250,1.8361642) (1.97500,1.8337915) (1.97750,1.8314192) (1.98000,1.8290473) (1.98250,1.8266759) (1.98500,1.8243049) (1.98750,1.8219343) (1.99000,1.8195642) (1.99250,1.8171946) (1.99500,1.8148254) (1.99750,1.8124567) (2.00000,1.8100884) (2.00250,1.8077206) (2.00500,1.8053533) (2.00750,1.8029865) (2.01000,1.8006202) (2.01250,1.7982543) (2.01500,1.7958890) (2.01750,1.7935241) (2.02000,1.7911597) (2.02250,1.7887959) (2.02500,1.7864326) (2.02750,1.7840697) (2.03000,1.7817074) (2.03250,1.7793456) (2.03500,1.7769844) (2.03750,1.7746237) (2.04000,1.7722635) (2.04250,1.7699039) (2.04500,1.7675448) (2.04750,1.7651862) (2.05000,1.7628282) (2.05250,1.7604708) (2.05500,1.7581140) (2.05750,1.7557577) (2.06000,1.7534020) (2.06250,1.7510468) (2.06500,1.7486923) (2.06750,1.7463383) (2.07000,1.7439850) (2.07250,1.7416322) (2.07500,1.7392800) (2.07750,1.7369285) (2.08000,1.7345775) (2.08250,1.7322272) (2.08500,1.7298775) (2.08750,1.7275284) (2.09000,1.7251800) (2.09250,1.7228322) (2.09500,1.7204850) (2.09750,1.7181385) (2.10000,1.7157926) (2.10250,1.7134474) (2.10500,1.7111029) (2.10750,1.7087590) (2.11000,1.7064158) (2.11250,1.7040733) (2.11500,1.7017314) (2.11750,1.6993902) (2.12000,1.6970498) (2.12250,1.6947100) (2.12500,1.6923709) (2.12750,1.6900326) (2.13000,1.6876949) (2.13250,1.6853580) (2.13500,1.6830218) (2.13750,1.6806863) (2.14000,1.6783516) (2.14250,1.6760176) (2.14500,1.6736843) (2.14750,1.6713518) (2.15000,1.6690201) (2.15250,1.6666891) (2.15500,1.6643589) (2.15750,1.6620294) (2.16000,1.6597008) (2.16250,1.6573729) (2.16500,1.6550458) (2.16750,1.6527195) (2.17000,1.6503940) (2.17250,1.6480693) (2.17500,1.6457454) (2.17750,1.6434224) (2.18000,1.6411001) (2.18250,1.6387787) (2.18500,1.6364581) (2.18750,1.6341384) (2.19000,1.6318195) (2.19250,1.6295015) (2.19500,1.6271843) (2.19750,1.6248680) (2.20000,1.6225525) (2.20250,1.6202380) (2.20500,1.6179243) (2.20750,1.6156115) (2.21000,1.6132996) (2.21250,1.6109886) (2.21500,1.6086785) (2.21750,1.6063693) (2.22000,1.6040610) (2.22250,1.6017537) (2.22500,1.5994473) (2.22750,1.5971418) (2.23000,1.5948373) (2.23250,1.5925337) (2.23500,1.5902311) (2.23750,1.5879294) (2.24000,1.5856287) (2.24250,1.5833290) (2.24500,1.5810303) (2.24750,1.5787326) (2.25000,1.5764358) (2.25250,1.5741401) (2.25500,1.5718453) (2.25750,1.5695516) (2.26000,1.5672589) (2.26250,1.5649672) (2.26500,1.5626766) (2.26750,1.5603870) (2.27000,1.5580985) (2.27250,1.5558110) (2.27500,1.5535246) (2.27750,1.5512392) (2.28000,1.5489549) (2.28250,1.5466717) (2.28500,1.5443896) (2.28750,1.5421086) (2.29000,1.5398287) (2.29250,1.5375500) (2.29500,1.5352723) (2.29750,1.5329957) (2.30000,1.5307203) (2.30250,1.5284461) (2.30500,1.5261730) (2.30750,1.5239010) (2.31000,1.5216302) (2.31250,1.5193606) (2.31500,1.5170921) (2.31750,1.5148248) (2.32000,1.5125588) (2.32250,1.5102939) (2.32500,1.5080302) (2.32750,1.5057678) (2.33000,1.5035066) (2.33250,1.5012466) (2.33500,1.4989878) (2.33750,1.4967303) (2.34000,1.4944741) (2.34250,1.4922191) (2.34500,1.4899654) (2.34750,1.4877129) (2.35000,1.4854618) (2.35250,1.4832119) (2.35500,1.4809634) (2.35750,1.4787161) (2.36000,1.4764702) (2.36250,1.4742256) (2.36500,1.4719823) (2.36750,1.4697404) (2.37000,1.4674999) (2.37250,1.4652607) (2.37500,1.4630228) (2.37750,1.4607864) (2.38000,1.4585513) (2.38250,1.4563177) (2.38500,1.4540854) (2.38750,1.4518545) (2.39000,1.4496251) (2.39250,1.4473971) (2.39500,1.4451705) (2.39750,1.4429454) (2.40000,1.4407218) (2.40250,1.4384996) (2.40500,1.4362789) (2.40750,1.4340597) (2.41000,1.4318419) (2.41250,1.4296257) (2.41500,1.4274110) (2.41750,1.4251978) (2.42000,1.4229861) (2.42250,1.4207760) (2.42500,1.4185674) (2.42750,1.4163604) (2.43000,1.4141549) (2.43250,1.4119510) (2.43500,1.4097488) (2.43750,1.4075481) (2.44000,1.4053490) (2.44250,1.4031515) (2.44500,1.4009557) (2.44750,1.3987615) (2.45000,1.3965690) (2.45250,1.3943781) (2.45500,1.3921888) (2.45750,1.3900013) (2.46000,1.3878154) (2.46250,1.3856313) (2.46500,1.3834488) (2.46750,1.3812681) (2.47000,1.3790891) (2.47250,1.3769118) (2.47500,1.3747363) (2.47750,1.3725626) (2.48000,1.3703906) (2.48250,1.3682204) (2.48500,1.3660520) (2.48750,1.3638854) (2.49000,1.3617206) (2.49250,1.3595577) (2.49500,1.3573966) (2.49750,1.3552373) (2.50000,1.3530799) (2.50250,1.3509244) (2.50500,1.3487707) (2.50750,1.3466189) (2.51000,1.3444691) (2.51250,1.3423212) (2.51500,1.3401752) (2.51750,1.3380311) (2.52000,1.3358890) (2.52250,1.3337488) (2.52500,1.3316107) (2.52750,1.3294745) (2.53000,1.3273403) (2.53250,1.3252082) (2.53500,1.3230780) (2.53750,1.3209499) (2.54000,1.3188239) (2.54250,1.3166999) (2.54500,1.3145780) (2.54750,1.3124581) (2.55000,1.3103404) (2.55250,1.3082248) (2.55500,1.3061113) (2.55750,1.3039999) (2.56000,1.3018907) (2.56250,1.2997837) (2.56500,1.2976789) (2.56750,1.2955762) (2.57000,1.2934757) (2.57250,1.2913775) (2.57500,1.2892815) (2.57750,1.2871877) (2.58000,1.2850962) (2.58250,1.2830070) (2.58500,1.2809200) (2.58750,1.2788354) (2.59000,1.2767530) (2.59250,1.2746730) (2.59500,1.2725954) (2.59750,1.2705201) (2.60000,1.2684471) (2.60250,1.2663766) (2.60500,1.2643084) (2.60750,1.2622427) (2.61000,1.2601794) (2.61250,1.2581185) (2.61500,1.2560601) (2.61750,1.2540042) (2.62000,1.2519508) (2.62250,1.2498998) (2.62500,1.2478514) (2.62750,1.2458056) (2.63000,1.2437622) (2.63250,1.2417215) (2.63500,1.2396833) (2.63750,1.2376477) (2.64000,1.2356148) (2.64250,1.2335844) (2.64500,1.2315568) (2.64750,1.2295317) (2.65000,1.2275094) (2.65250,1.2254898) (2.65500,1.2234728) (2.65750,1.2214586) (2.66000,1.2194472) (2.66250,1.2174385) (2.66500,1.2154326) (2.66750,1.2134295) (2.67000,1.2114292) (2.67250,1.2094317) (2.67500,1.2074371) (2.67750,1.2054454) (2.68000,1.2034565) (2.68250,1.2014706) (2.68500,1.1994875) (2.68750,1.1975074) (2.69000,1.1955303) (2.69250,1.1935562) (2.69500,1.1915850) (2.69750,1.1896169) (2.70000,1.1876518) (2.70250,1.1856897) (2.70500,1.1837307) (2.70750,1.1817748) (2.71000,1.1798220) (2.71250,1.1778724) (2.71500,1.1759259) (2.71750,1.1739825) (2.72000,1.1720424) (2.72250,1.1701055) (2.72500,1.1681718) (2.72750,1.1662413) (2.73000,1.1643141) (2.73250,1.1623902) (2.73500,1.1604697) (2.73750,1.1585524) (2.74000,1.1566385) (2.74250,1.1547280) (2.74500,1.1528209) (2.74750,1.1509172) (2.75000,1.1490170) (2.75250,1.1471202) (2.75500,1.1452269) (2.75750,1.1433371) (2.76000,1.1414509) (2.76250,1.1395682) (2.76500,1.1376891) (2.76750,1.1358135) (2.77000,1.1339417) (2.77250,1.1320734) (2.77500,1.1302088) (2.77750,1.1283480) (2.78000,1.1264908) (2.78250,1.1246374) (2.78500,1.1227878) (2.78750,1.1209419) (2.79000,1.1190999) (2.79250,1.1172617) (2.79500,1.1154274) (2.79750,1.1135970) (2.80000,1.1117705) (2.80250,1.1099479) (2.80500,1.1081294) (2.80750,1.1063148) (2.81000,1.1045043) (2.81250,1.1026978) (2.81500,1.1008954) (2.81750,1.0990971) (2.82000,1.0973029) (2.82250,1.0955129) (2.82500,1.0937271) (2.82750,1.0919455) (2.83000,1.0901682) (2.83250,1.0883951) (2.83500,1.0866264) (2.83750,1.0848620) (2.84000,1.0831019) (2.84250,1.0813463) (2.84500,1.0795950) (2.84750,1.0778483) (2.85000,1.0761060) (2.85250,1.0743683) (2.85500,1.0726351) (2.85750,1.0709065) (2.86000,1.0691825) (2.86250,1.0674631) (2.86500,1.0657484) (2.86750,1.0640385) (2.87000,1.0623333) (2.87250,1.0606328) (2.87500,1.0589372) (2.87750,1.0572465) (2.88000,1.0555606) (2.88250,1.0538796) (2.88500,1.0522036) (2.88750,1.0505326) (2.89000,1.0488666) (2.89250,1.0472057) (2.89500,1.0455499) (2.89750,1.0438992) (2.90000,1.0422537) (2.90250,1.0406133) (2.90500,1.0389783) (2.90750,1.0373485) (2.91000,1.0357241) (2.91250,1.0341050) (2.91500,1.0324914) (2.91750,1.0308832) (2.92000,1.0292804) (2.92250,1.0276832) (2.92500,1.0260916) (2.92750,1.0245056) (2.93000,1.0229253) (2.93250,1.0213506) (2.93500,1.0197817) (2.93750,1.0182186) (2.94000,1.0166613) (2.94250,1.0151099) (2.94500,1.0135645) (2.94750,1.0120250) (2.95000,1.0104915) (2.95250,1.0089640) (2.95500,1.0074427) (2.95750,1.0059276) (2.96000,1.0044186) (2.96250,1.0029160) (2.96500,1.0014196) (2.96750,0.9999296) (2.97000,0.9984460) (2.97250,0.9969688) (2.97500,0.9954982) (2.97750,0.9940342) (2.98000,0.9925767) (2.98250,0.9911260) (2.98500,0.9896820) (2.98750,0.9882447) (2.99000,0.9868144) (2.99250,0.9853909) (2.99500,0.9839744) (2.99750,0.9825649) (3.00000,0.9811624) (3.00250,0.9797672) (3.00500,0.9783791) (3.00750,0.9769983) (3.01000,0.9756248) (3.01250,0.9742587) (3.01500,0.9729001) (3.01750,0.9715489) (3.02000,0.9702054) (3.02250,0.9688695) (3.02500,0.9675414) (3.02750,0.9662210) (3.03000,0.9649085) (3.03250,0.9636039) (3.03500,0.9623073) (3.03750,0.9610188) (3.04000,0.9597385) (3.04250,0.9584664) (3.04500,0.9572025) (3.04750,0.9559471) (3.05000,0.9547001) (3.05250,0.9534617) (3.05500,0.9522319) (3.05750,0.9510107) (3.06000,0.9497984) (3.06250,0.9485949) (3.06500,0.9474004) (3.06750,0.9462149) (3.07000,0.9450386) (3.07250,0.9438715) (3.07500,0.9427137) (3.07750,0.9415653) (3.08000,0.9404263) (3.08250,0.9392970) (3.08500,0.9381774) (3.08750,0.9370676) (3.09000,0.9359676) (3.09250,0.9348776) (3.09500,0.9337978) (3.09750,0.9327281) (3.10000,0.9316687) (3.10250,0.9306198) (3.10500,0.9295814) (3.10750,0.9285536) (3.11000,0.9275366) (3.11250,0.9265304) (3.11500,0.9255352) (3.11750,0.9245512) (3.12000,0.9235783) (3.12250,0.9226169) (3.12500,0.9216669) (3.12750,0.9207285) (3.13000,0.9198018) (3.13250,0.9188871) (3.13500,0.9179843) (3.13750,0.9170937) (3.14000,0.9162154) (3.14250,0.9153495) (3.14500,0.9144962) (3.14750,0.9136557) (3.15000,0.9128280) (3.15250,0.9120134) (3.15500,0.9112119) (3.15750,0.9104238) (3.16000,0.9096493) (3.16250,0.9088884) (3.16500,0.9081414) (3.16750,0.9074084) (3.17000,0.9066897) (3.17250,0.9059853) (3.17500,0.9052955) (3.17750,0.9046205) (3.18000,0.9039604) (3.18250,0.9033155) (3.18500,0.9026860) (3.18750,0.9020720) (3.19000,0.9014738) (3.19250,0.9008915) (3.19500,0.9003255) (3.19750,0.8997759) (3.20000,0.8992429) (3.20250,0.8987268) (3.20500,0.8982278) (3.20750,0.8977462) (3.21000,0.8972822) (3.21250,0.8968361) (3.21500,0.8964081) (3.21750,0.8959985) (3.22000,0.8956075) (3.22250,0.8952355) (3.22500,0.8948827) (3.22750,0.8945495) (3.23000,0.8942360) (3.23250,0.8939427) (3.23500,0.8936697) (3.23750,0.8934176) (3.24000,0.8931865) (3.24250,0.8929768) (3.24500,0.8927889) (3.24750,0.8926231) (3.25000,0.8924797) (3.25250,0.8923592) (3.25500,0.8922619) (3.25750,0.8921883) (3.26000,0.8921386) (3.26250,0.8921134) (3.26500,0.8921130) (3.26750,0.8921379) (3.27000,0.8921885) (3.27250,0.8922653) (3.27500,0.8923687) (3.27750,0.8924993) (3.28000,0.8926575) (3.28250,0.8928438) (3.28500,0.8930588) (3.28750,0.8933029) (3.29000,0.8935768) (3.29250,0.8938811) (3.29500,0.8942162) (3.29750,0.8945829) (3.30000,0.8949816) (3.30250,0.8954132) (3.30500,0.8958782) (3.30750,0.8963772) (3.31000,0.8969112) (3.31250,0.8974806) (3.31500,0.8980863) (3.31750,0.8987291) (3.32000,0.8994097) (3.32250,0.9001290) (3.32500,0.9008878) (3.32750,0.9016870) (3.33000,0.9025275) (3.33250,0.9034102) (3.33500,0.9043361) (3.33750,0.9053062) (3.34000,0.9063215) (3.34250,0.9073831) (3.34500,0.9084922) (3.34750,0.9096498) (3.35000,0.9108572) (3.35250,0.9121155) (3.35500,0.9134261) (3.35750,0.9147903) (3.36000,0.9162095) (3.36250,0.9176851) (3.36500,0.9192187) (3.36750,0.9208116) (3.37000,0.9224656) (3.37250,0.9241824) (3.37500,0.9259636) (3.37750,0.9278111) (3.38000,0.9297268) (3.38250,0.9317126) (3.38500,0.9337707) (3.38750,0.9359031) (3.39000,0.9381121) (3.39250,0.9404000) (3.39500,0.9427693) (3.39750,0.9452226) (3.40000,0.9477624) (3.40250,0.9503917) (3.40500,0.9531134) (3.40750,0.9559305) (3.41000,0.9588463) (3.41250,0.9618642) (3.41500,0.9649877) (3.41750,0.9682207) (3.42000,0.9715671) (3.42250,0.9750310) (3.42500,0.9786169) (3.42750,0.9823293) (3.43000,0.9861733) (3.43250,0.9901539) (3.43500,0.9942768) (3.43750,0.9985476) (3.44000,1.0029726) (3.44250,1.0075583) (3.44500,1.0123117) (3.44750,1.0172402) (3.45000,1.0223517) (3.45250,1.0276547) (3.45500,1.0331582) (3.45750,1.0388718) (3.46000,1.0448059) (3.46250,1.0509715) (3.46500,1.0573807) (3.46750,1.0640463) (3.47000,1.0709821) (3.47250,1.0782032) (3.47500,1.0857258) (3.47750,1.0935677) (3.48000,1.1017479) (3.48250,1.1102874) (3.48500,1.1192092) (3.48750,1.1285383) (3.49000,1.1383024) (3.49250,1.1485318) (3.49500,1.1592602) (3.49750,1.1705251) (3.50000,1.1823680) (3.50250,1.1948355) (3.50500,1.2079799) (3.50750,1.2218603) (3.51000,1.2365437) (3.51250,1.2521063) (3.51500,1.2686359) (3.51750,1.2862336) (3.52000,1.3050170) (3.52250,1.3251241) (3.52500,1.3467180) (3.52750,1.3699937) (3.53000,1.3951864) (3.53250,1.4225843) (3.53500,1.4525448) (3.53750,1.4855191) (3.54000,1.5220882) (3.54250,1.5630176) (3.54500,1.6093452) (3.54750,1.6625276) (3.55000,1.7247009) (3.55250,1.7991825) (3.55500,1.8915429) (3.55750,2.0122607) (3.56000,2.1849994) (3.56250,2.4862724) (3.56500,nan) (3.56750,nan) (3.57000,nan) (3.57250,nan) (3.57500,nan) (3.57750,nan) (3.58000,nan) (3.58250,nan) (3.58500,nan) (3.58750,nan) (3.59000,nan) (3.59250,nan) (3.59500,nan) (3.59750,nan) (3.60000,nan) (3.60250,nan) (3.60500,nan) (3.60750,nan) (3.61000,nan) (3.61250,nan) (3.61500,nan) (3.61750,nan) (3.62000,nan) (3.62250,nan) (3.62500,nan) (3.62750,nan) (3.63000,nan) (3.63250,nan) (3.63500,nan) (3.63750,nan) (3.64000,nan) (3.64250,nan) (3.64500,nan) (3.64750,nan) (3.65000,nan) (3.65250,nan) (3.65500,nan) (3.65750,nan) (3.66000,nan) (3.66250,nan) (3.66500,nan) (3.66750,nan) (3.67000,nan) (3.67250,nan) (3.67500,nan) (3.67750,nan) (3.68000,nan) (3.68250,nan) (3.68500,nan) (3.68750,nan) (3.69000,nan) (3.69250,nan) (3.69500,nan) (3.69750,nan) (3.70000,nan) (3.70250,nan) (3.70500,nan) (3.70750,nan) (3.71000,nan) (3.71250,nan) (3.71500,nan) (3.71750,nan) (3.72000,nan) (3.72250,nan) (3.72500,nan) (3.72750,nan) (3.73000,nan) (3.73250,nan) (3.73500,nan) (3.73750,nan) (3.74000,nan) (3.74250,nan) (3.74500,nan) (3.74750,nan) (3.75000,nan) (3.75250,nan) (3.75500,nan) (3.75750,nan) (3.76000,nan) (3.76250,nan) (3.76500,nan) (3.76750,nan) (3.77000,nan) (3.77250,nan) (3.77500,nan) (3.77750,nan) (3.78000,nan) (3.78250,nan) (3.78500,nan) (3.78750,nan) (3.79000,nan) (3.79250,nan) (3.79500,nan) (3.79750,nan) (3.80000,nan) (3.80250,nan) (3.80500,nan) (3.80750,nan) (3.81000,nan) (3.81250,nan) (3.81500,nan) (3.81750,nan) (3.82000,nan) (3.82250,nan) (3.82500,nan) (3.82750,nan) (3.83000,nan) (3.83250,nan) (3.83500,nan) (3.83750,nan) (3.84000,nan) (3.84250,nan) (3.84500,-1.8335524) (3.84750,-1.4937278) (3.85000,-1.3087819) (3.85250,-1.1819125) (3.85500,-1.0857430) (3.85750,-1.0085999) (3.86000,-0.9444108) (3.86250,-0.8896106) (3.86500,-0.8419267) (3.86750,-0.7998212) (3.87000,-0.7622048) (3.87250,-0.7282778) (3.87500,-0.6974357) (3.87750,-0.6692104) (3.88000,-0.6432323) (3.88250,-0.6192041) (3.88500,-0.5968835) (3.88750,-0.5760701) (3.89000,-0.5565964) (3.89250,-0.5383212) (3.89500,-0.5211240) (3.89750,-0.5049013) (3.90000,-0.4895637) (3.90250,-0.4750333) (3.90500,-0.4612418) (3.90750,-0.4481291) (3.91000,-0.4356418) (3.91250,-0.4237327) (3.91500,-0.4123595) (3.91750,-0.4014841) (3.92000,-0.3910723) (3.92250,-0.3810934) (3.92500,-0.3715192) (3.92750,-0.3623245) (3.93000,-0.3534858) (3.93250,-0.3449821) (3.93500,-0.3367939) (3.93750,-0.3289033) (3.94000,-0.3212939) (3.94250,-0.3139505) (3.94500,-0.3068590) (3.94750,-0.3000065) (3.95000,-0.2933809) (3.95250,-0.2869710) (3.95500,-0.2807663) (3.95750,-0.2747571) (3.96000,-0.2689343) (3.96250,-0.2632894) (3.96500,-0.2578143) (3.96750,-0.2525017) (3.97000,-0.2473445) (3.97250,-0.2423361) (3.97500,-0.2374702) (3.97750,-0.2327411) (3.98000,-0.2281431) (3.98250,-0.2236711) (3.98500,-0.2193201) (3.98750,-0.2150855) (3.99000,-0.2109628) (3.99250,-0.2069478) (3.99500,-0.2030366) (3.99750,-0.1992254) (4.00000,-0.1955106)};\addlegendentry{R}
\addplot[thick,orange] coordinates {(-2.00000,nan) (-1.99750,nan) (-1.99500,nan) (-1.99250,nan) (-1.99000,nan) (-1.98750,nan) (-1.98500,nan) (-1.98250,nan) (-1.98000,nan) (-1.97750,nan) (-1.97500,nan) (-1.97250,nan) (-1.97000,nan) (-1.96750,nan) (-1.96500,nan) (-1.96250,nan) (-1.96000,nan) (-1.95750,nan) (-1.95500,nan) (-1.95250,nan) (-1.95000,nan) (-1.94750,nan) (-1.94500,nan) (-1.94250,nan) (-1.94000,nan) (-1.93750,nan) (-1.93500,nan) (-1.93250,nan) (-1.93000,nan) (-1.92750,nan) (-1.92500,nan) (-1.92250,nan) (-1.92000,nan) (-1.91750,nan) (-1.91500,nan) (-1.91250,nan) (-1.91000,nan) (-1.90750,nan) (-1.90500,nan) (-1.90250,nan) (-1.90000,nan) (-1.89750,nan) (-1.89500,nan) (-1.89250,nan) (-1.89000,nan) (-1.88750,nan) (-1.88500,nan) (-1.88250,nan) (-1.88000,nan) (-1.87750,nan) (-1.87500,nan) (-1.87250,nan) (-1.87000,nan) (-1.86750,nan) (-1.86500,nan) (-1.86250,nan) (-1.86000,nan) (-1.85750,nan) (-1.85500,nan) (-1.85250,nan) (-1.85000,nan) (-1.84750,nan) (-1.84500,nan) (-1.84250,nan) (-1.84000,nan) (-1.83750,nan) (-1.83500,nan) (-1.83250,nan) (-1.83000,nan) (-1.82750,nan) (-1.82500,nan) (-1.82250,nan) (-1.82000,nan) (-1.81750,nan) (-1.81500,nan) (-1.81250,nan) (-1.81000,nan) (-1.80750,nan) (-1.80500,nan) (-1.80250,nan) (-1.80000,nan) (-1.79750,nan) (-1.79500,nan) (-1.79250,nan) (-1.79000,nan) (-1.78750,nan) (-1.78500,nan) (-1.78250,nan) (-1.78000,nan) (-1.77750,nan) (-1.77500,nan) (-1.77250,nan) (-1.77000,nan) (-1.76750,nan) (-1.76500,nan) (-1.76250,nan) (-1.76000,nan) (-1.75750,nan) (-1.75500,nan) (-1.75250,nan) (-1.75000,nan) (-1.74750,nan) (-1.74500,nan) (-1.74250,nan) (-1.74000,nan) (-1.73750,nan) (-1.73500,nan) (-1.73250,nan) (-1.73000,nan) (-1.72750,nan) (-1.72500,nan) (-1.72250,nan) (-1.72000,nan) (-1.71750,nan) (-1.71500,nan) (-1.71250,nan) (-1.71000,nan) (-1.70750,nan) (-1.70500,nan) (-1.70250,nan) (-1.70000,nan) (-1.69750,nan) (-1.69500,nan) (-1.69250,nan) (-1.69000,nan) (-1.68750,nan) (-1.68500,nan) (-1.68250,nan) (-1.68000,nan) (-1.67750,nan) (-1.67500,nan) (-1.67250,nan) (-1.67000,nan) (-1.66750,nan) (-1.66500,nan) (-1.66250,nan) (-1.66000,nan) (-1.65750,nan) (-1.65500,nan) (-1.65250,nan) (-1.65000,nan) (-1.64750,nan) (-1.64500,nan) (-1.64250,nan) (-1.64000,nan) (-1.63750,nan) (-1.63500,nan) (-1.63250,nan) (-1.63000,nan) (-1.62750,nan) (-1.62500,nan) (-1.62250,nan) (-1.62000,nan) (-1.61750,nan) (-1.61500,nan) (-1.61250,nan) (-1.61000,nan) (-1.60750,nan) (-1.60500,nan) (-1.60250,nan) (-1.60000,nan) (-1.59750,nan) (-1.59500,nan) (-1.59250,nan) (-1.59000,nan) (-1.58750,nan) (-1.58500,nan) (-1.58250,nan) (-1.58000,nan) (-1.57750,nan) (-1.57500,nan) (-1.57250,nan) (-1.57000,nan) (-1.56750,nan) (-1.56500,nan) (-1.56250,nan) (-1.56000,nan) (-1.55750,nan) (-1.55500,nan) (-1.55250,nan) (-1.55000,nan) (-1.54750,nan) (-1.54500,nan) (-1.54250,nan) (-1.54000,nan) (-1.53750,nan) (-1.53500,nan) (-1.53250,nan) (-1.53000,nan) (-1.52750,nan) (-1.52500,nan) (-1.52250,nan) (-1.52000,nan) (-1.51750,nan) (-1.51500,nan) (-1.51250,nan) (-1.51000,nan) (-1.50750,nan) (-1.50500,nan) (-1.50250,nan) (-1.50000,nan) (-1.49750,nan) (-1.49500,nan) (-1.49250,nan) (-1.49000,nan) (-1.48750,nan) (-1.48500,nan) (-1.48250,nan) (-1.48000,nan) (-1.47750,nan) (-1.47500,nan) (-1.47250,nan) (-1.47000,nan) (-1.46750,nan) (-1.46500,nan) (-1.46250,nan) (-1.46000,nan) (-1.45750,nan) (-1.45500,nan) (-1.45250,nan) (-1.45000,nan) (-1.44750,nan) (-1.44500,nan) (-1.44250,nan) (-1.44000,nan) (-1.43750,nan) (-1.43500,nan) (-1.43250,nan) (-1.43000,nan) (-1.42750,nan) (-1.42500,nan) (-1.42250,nan) (-1.42000,nan) (-1.41750,nan) (-1.41500,nan) (-1.41250,nan) (-1.41000,nan) (-1.40750,nan) (-1.40500,nan) (-1.40250,nan) (-1.40000,nan) (-1.39750,nan) (-1.39500,nan) (-1.39250,nan) (-1.39000,nan) (-1.38750,nan) (-1.38500,nan) (-1.38250,nan) (-1.38000,nan) (-1.37750,nan) (-1.37500,nan) (-1.37250,nan) (-1.37000,nan) (-1.36750,nan) (-1.36500,nan) (-1.36250,nan) (-1.36000,nan) (-1.35750,nan) (-1.35500,nan) (-1.35250,nan) (-1.35000,nan) (-1.34750,nan) (-1.34500,nan) (-1.34250,nan) (-1.34000,nan) (-1.33750,nan) (-1.33500,nan) (-1.33250,nan) (-1.33000,nan) (-1.32750,nan) (-1.32500,nan) (-1.32250,nan) (-1.32000,nan) (-1.31750,nan) (-1.31500,nan) (-1.31250,nan) (-1.31000,nan) (-1.30750,nan) (-1.30500,nan) (-1.30250,nan) (-1.30000,nan) (-1.29750,nan) (-1.29500,nan) (-1.29250,nan) (-1.29000,nan) (-1.28750,nan) (-1.28500,nan) (-1.28250,nan) (-1.28000,nan) (-1.27750,nan) (-1.27500,nan) (-1.27250,nan) (-1.27000,nan) (-1.26750,nan) (-1.26500,nan) (-1.26250,nan) (-1.26000,nan) (-1.25750,nan) (-1.25500,nan) (-1.25250,nan) (-1.25000,nan) (-1.24750,nan) (-1.24500,nan) (-1.24250,nan) (-1.24000,nan) (-1.23750,nan) (-1.23500,nan) (-1.23250,nan) (-1.23000,nan) (-1.22750,nan) (-1.22500,nan) (-1.22250,nan) (-1.22000,nan) (-1.21750,nan) (-1.21500,nan) (-1.21250,nan) (-1.21000,nan) (-1.20750,nan) (-1.20500,nan) (-1.20250,nan) (-1.20000,nan) (-1.19750,nan) (-1.19500,nan) (-1.19250,nan) (-1.19000,nan) (-1.18750,nan) (-1.18500,nan) (-1.18250,nan) (-1.18000,nan) (-1.17750,nan) (-1.17500,nan) (-1.17250,nan) (-1.17000,nan) (-1.16750,nan) (-1.16500,nan) (-1.16250,nan) (-1.16000,nan) (-1.15750,nan) (-1.15500,nan) (-1.15250,nan) (-1.15000,nan) (-1.14750,nan) (-1.14500,nan) (-1.14250,nan) (-1.14000,nan) (-1.13750,nan) (-1.13500,nan) (-1.13250,nan) (-1.13000,nan) (-1.12750,nan) (-1.12500,nan) (-1.12250,nan) (-1.12000,nan) (-1.11750,nan) (-1.11500,nan) (-1.11250,nan) (-1.11000,nan) (-1.10750,nan) (-1.10500,nan) (-1.10250,nan) (-1.10000,nan) (-1.09750,nan) (-1.09500,nan) (-1.09250,nan) (-1.09000,nan) (-1.08750,nan) (-1.08500,nan) (-1.08250,nan) (-1.08000,nan) (-1.07750,nan) (-1.07500,nan) (-1.07250,nan) (-1.07000,nan) (-1.06750,nan) (-1.06500,nan) (-1.06250,nan) (-1.06000,nan) (-1.05750,nan) (-1.05500,nan) (-1.05250,nan) (-1.05000,nan) (-1.04750,nan) (-1.04500,nan) (-1.04250,nan) (-1.04000,nan) (-1.03750,nan) (-1.03500,nan) (-1.03250,nan) (-1.03000,nan) (-1.02750,nan) (-1.02500,nan) (-1.02250,nan) (-1.02000,nan) (-1.01750,nan) (-1.01500,nan) (-1.01250,nan) (-1.01000,nan) (-1.00750,nan) (-1.00500,nan) (-1.00250,nan) (-1.00000,nan) (-0.99750,nan) (-0.99500,nan) (-0.99250,nan) (-0.99000,nan) (-0.98750,nan) (-0.98500,nan) (-0.98250,nan) (-0.98000,nan) (-0.97750,nan) (-0.97500,nan) (-0.97250,nan) (-0.97000,nan) (-0.96750,nan) (-0.96500,nan) (-0.96250,nan) (-0.96000,nan) (-0.95750,nan) (-0.95500,nan) (-0.95250,nan) (-0.95000,nan) (-0.94750,nan) (-0.94500,nan) (-0.94250,nan) (-0.94000,nan) (-0.93750,nan) (-0.93500,nan) (-0.93250,nan) (-0.93000,nan) (-0.92750,nan) (-0.92500,nan) (-0.92250,nan) (-0.92000,nan) (-0.91750,nan) (-0.91500,nan) (-0.91250,nan) (-0.91000,nan) (-0.90750,nan) (-0.90500,nan) (-0.90250,nan) (-0.90000,nan) (-0.89750,nan) (-0.89500,nan) (-0.89250,nan) (-0.89000,nan) (-0.88750,nan) (-0.88500,nan) (-0.88250,nan) (-0.88000,nan) (-0.87750,nan) (-0.87500,nan) (-0.87250,nan) (-0.87000,nan) (-0.86750,nan) (-0.86500,nan) (-0.86250,nan) (-0.86000,nan) (-0.85750,nan) (-0.85500,nan) (-0.85250,nan) (-0.85000,nan) (-0.84750,nan) (-0.84500,nan) (-0.84250,nan) (-0.84000,nan) (-0.83750,nan) (-0.83500,nan) (-0.83250,nan) (-0.83000,nan) (-0.82750,nan) (-0.82500,nan) (-0.82250,nan) (-0.82000,nan) (-0.81750,nan) (-0.81500,nan) (-0.81250,nan) (-0.81000,nan) (-0.80750,nan) (-0.80500,nan) (-0.80250,nan) (-0.80000,nan) (-0.79750,nan) (-0.79500,nan) (-0.79250,nan) (-0.79000,nan) (-0.78750,nan) (-0.78500,nan) (-0.78250,nan) (-0.78000,nan) (-0.77750,nan) (-0.77500,nan) (-0.77250,nan) (-0.77000,nan) (-0.76750,nan) (-0.76500,nan) (-0.76250,nan) (-0.76000,nan) (-0.75750,nan) (-0.75500,nan) (-0.75250,nan) (-0.75000,nan) (-0.74750,nan) (-0.74500,nan) (-0.74250,nan) (-0.74000,nan) (-0.73750,nan) (-0.73500,nan) (-0.73250,nan) (-0.73000,nan) (-0.72750,nan) (-0.72500,nan) (-0.72250,nan) (-0.72000,nan) (-0.71750,nan) (-0.71500,nan) (-0.71250,nan) (-0.71000,nan) (-0.70750,nan) (-0.70500,nan) (-0.70250,nan) (-0.70000,nan) (-0.69750,nan) (-0.69500,nan) (-0.69250,nan) (-0.69000,nan) (-0.68750,nan) (-0.68500,nan) (-0.68250,nan) (-0.68000,nan) (-0.67750,nan) (-0.67500,nan) (-0.67250,nan) (-0.67000,nan) (-0.66750,nan) (-0.66500,nan) (-0.66250,nan) (-0.66000,nan) (-0.65750,nan) (-0.65500,nan) (-0.65250,nan) (-0.65000,nan) (-0.64750,nan) (-0.64500,nan) (-0.64250,nan) (-0.64000,nan) (-0.63750,nan) (-0.63500,nan) (-0.63250,nan) (-0.63000,nan) (-0.62750,nan) (-0.62500,nan) (-0.62250,nan) (-0.62000,nan) (-0.61750,nan) (-0.61500,nan) (-0.61250,nan) (-0.61000,nan) (-0.60750,nan) (-0.60500,nan) (-0.60250,nan) (-0.60000,nan) (-0.59750,nan) (-0.59500,nan) (-0.59250,nan) (-0.59000,nan) (-0.58750,nan) (-0.58500,nan) (-0.58250,nan) (-0.58000,nan) (-0.57750,nan) (-0.57500,nan) (-0.57250,nan) (-0.57000,nan) (-0.56750,nan) (-0.56500,nan) (-0.56250,nan) (-0.56000,nan) (-0.55750,nan) (-0.55500,nan) (-0.55250,nan) (-0.55000,nan) (-0.54750,nan) (-0.54500,nan) (-0.54250,nan) (-0.54000,nan) (-0.53750,nan) (-0.53500,nan) (-0.53250,nan) (-0.53000,nan) (-0.52750,nan) (-0.52500,nan) (-0.52250,nan) (-0.52000,nan) (-0.51750,nan) (-0.51500,nan) (-0.51250,nan) (-0.51000,nan) (-0.50750,nan) (-0.50500,nan) (-0.50250,nan) (-0.50000,nan) (-0.49750,nan) (-0.49500,nan) (-0.49250,nan) (-0.49000,nan) (-0.48750,nan) (-0.48500,nan) (-0.48250,nan) (-0.48000,nan) (-0.47750,nan) (-0.47500,nan) (-0.47250,nan) (-0.47000,nan) (-0.46750,nan) (-0.46500,nan) (-0.46250,nan) (-0.46000,nan) (-0.45750,nan) (-0.45500,nan) (-0.45250,nan) (-0.45000,nan) (-0.44750,nan) (-0.44500,nan) (-0.44250,nan) (-0.44000,nan) (-0.43750,nan) (-0.43500,nan) (-0.43250,nan) (-0.43000,nan) (-0.42750,nan) (-0.42500,nan) (-0.42250,nan) (-0.42000,nan) (-0.41750,nan) (-0.41500,nan) (-0.41250,nan) (-0.41000,nan) (-0.40750,nan) (-0.40500,nan) (-0.40250,nan) (-0.40000,nan) (-0.39750,nan) (-0.39500,nan) (-0.39250,nan) (-0.39000,nan) (-0.38750,nan) (-0.38500,nan) (-0.38250,nan) (-0.38000,nan) (-0.37750,nan) (-0.37500,nan) (-0.37250,nan) (-0.37000,nan) (-0.36750,nan) (-0.36500,nan) (-0.36250,nan) (-0.36000,nan) (-0.35750,nan) (-0.35500,nan) (-0.35250,nan) (-0.35000,nan) (-0.34750,nan) (-0.34500,nan) (-0.34250,nan) (-0.34000,nan) (-0.33750,nan) (-0.33500,nan) (-0.33250,nan) (-0.33000,nan) (-0.32750,nan) (-0.32500,nan) (-0.32250,nan) (-0.32000,nan) (-0.31750,nan) (-0.31500,nan) (-0.31250,nan) (-0.31000,nan) (-0.30750,nan) (-0.30500,nan) (-0.30250,nan) (-0.30000,nan) (-0.29750,nan) (-0.29500,nan) (-0.29250,nan) (-0.29000,nan) (-0.28750,nan) (-0.28500,nan) (-0.28250,nan) (-0.28000,nan) (-0.27750,nan) (-0.27500,nan) (-0.27250,nan) (-0.27000,nan) (-0.26750,nan) (-0.26500,nan) (-0.26250,nan) (-0.26000,nan) (-0.25750,nan) (-0.25500,nan) (-0.25250,nan) (-0.25000,nan) (-0.24750,nan) (-0.24500,nan) (-0.24250,nan) (-0.24000,nan) (-0.23750,nan) (-0.23500,nan) (-0.23250,nan) (-0.23000,nan) (-0.22750,nan) (-0.22500,nan) (-0.22250,nan) (-0.22000,nan) (-0.21750,nan) (-0.21500,nan) (-0.21250,nan) (-0.21000,nan) (-0.20750,nan) (-0.20500,nan) (-0.20250,nan) (-0.20000,nan) (-0.19750,nan) (-0.19500,nan) (-0.19250,nan) (-0.19000,nan) (-0.18750,nan) (-0.18500,nan) (-0.18250,nan) (-0.18000,nan) (-0.17750,nan) (-0.17500,nan) (-0.17250,nan) (-0.17000,nan) (-0.16750,nan) (-0.16500,nan) (-0.16250,nan) (-0.16000,nan) (-0.15750,nan) (-0.15500,nan) (-0.15250,nan) (-0.15000,nan) (-0.14750,nan) (-0.14500,nan) (-0.14250,nan) (-0.14000,nan) (-0.13750,nan) (-0.13500,nan) (-0.13250,nan) (-0.13000,nan) (-0.12750,nan) (-0.12500,nan) (-0.12250,nan) (-0.12000,nan) (-0.11750,nan) (-0.11500,nan) (-0.11250,nan) (-0.11000,nan) (-0.10750,nan) (-0.10500,nan) (-0.10250,nan) (-0.10000,nan) (-0.09750,nan) (-0.09500,nan) (-0.09250,nan) (-0.09000,nan) (-0.08750,nan) (-0.08500,nan) (-0.08250,nan) (-0.08000,nan) (-0.07750,nan) (-0.07500,nan) (-0.07250,nan) (-0.07000,nan) (-0.06750,nan) (-0.06500,nan) (-0.06250,nan) (-0.06000,nan) (-0.05750,nan) (-0.05500,nan) (-0.05250,nan) (-0.05000,nan) (-0.04750,nan) (-0.04500,nan) (-0.04250,nan) (-0.04000,nan) (-0.03750,nan) (-0.03500,nan) (-0.03250,nan) (-0.03000,nan) (-0.02750,nan) (-0.02500,nan) (-0.02250,nan) (-0.02000,nan) (-0.01750,nan) (-0.01500,nan) (-0.01250,nan) (-0.01000,nan) (-0.00750,nan) (-0.00500,nan) (-0.00250,nan) (0.00000,nan) (0.00250,nan) (0.00500,nan) (0.00750,nan) (0.01000,nan) (0.01250,nan) (0.01500,nan) (0.01750,nan) (0.02000,nan) (0.02250,nan) (0.02500,nan) (0.02750,nan) (0.03000,nan) (0.03250,nan) (0.03500,nan) (0.03750,nan) (0.04000,nan) (0.04250,nan) (0.04500,nan) (0.04750,nan) (0.05000,nan) (0.05250,nan) (0.05500,nan) (0.05750,nan) (0.06000,nan) (0.06250,nan) (0.06500,nan) (0.06750,nan) (0.07000,nan) (0.07250,nan) (0.07500,nan) (0.07750,nan) (0.08000,nan) (0.08250,nan) (0.08500,nan) (0.08750,nan) (0.09000,nan) (0.09250,nan) (0.09500,nan) (0.09750,nan) (0.10000,nan) (0.10250,nan) (0.10500,nan) (0.10750,nan) (0.11000,nan) (0.11250,nan) (0.11500,nan) (0.11750,nan) (0.12000,nan) (0.12250,nan) (0.12500,nan) (0.12750,nan) (0.13000,nan) (0.13250,nan) (0.13500,nan) (0.13750,nan) (0.14000,nan) (0.14250,nan) (0.14500,nan) (0.14750,nan) (0.15000,nan) (0.15250,nan) (0.15500,nan) (0.15750,nan) (0.16000,nan) (0.16250,nan) (0.16500,nan) (0.16750,nan) (0.17000,nan) (0.17250,nan) (0.17500,nan) (0.17750,nan) (0.18000,nan) (0.18250,nan) (0.18500,nan) (0.18750,nan) (0.19000,nan) (0.19250,nan) (0.19500,nan) (0.19750,nan) (0.20000,nan) (0.20250,nan) (0.20500,nan) (0.20750,nan) (0.21000,nan) (0.21250,nan) (0.21500,nan) (0.21750,nan) (0.22000,nan) (0.22250,nan) (0.22500,nan) (0.22750,nan) (0.23000,nan) (0.23250,nan) (0.23500,nan) (0.23750,nan) (0.24000,nan) (0.24250,nan) (0.24500,nan) (0.24750,nan) (0.25000,nan) (0.25250,nan) (0.25500,nan) (0.25750,nan) (0.26000,nan) (0.26250,nan) (0.26500,nan) (0.26750,nan) (0.27000,nan) (0.27250,nan) (0.27500,nan) (0.27750,nan) (0.28000,nan) (0.28250,nan) (0.28500,nan) (0.28750,nan) (0.29000,nan) (0.29250,nan) (0.29500,nan) (0.29750,nan) (0.30000,nan) (0.30250,nan) (0.30500,nan) (0.30750,nan) (0.31000,nan) (0.31250,nan) (0.31500,nan) (0.31750,nan) (0.32000,nan) (0.32250,nan) (0.32500,nan) (0.32750,nan) (0.33000,nan) (0.33250,nan) (0.33500,nan) (0.33750,nan) (0.34000,nan) (0.34250,nan) (0.34500,nan) (0.34750,nan) (0.35000,nan) (0.35250,nan) (0.35500,nan) (0.35750,nan) (0.36000,nan) (0.36250,nan) (0.36500,nan) (0.36750,nan) (0.37000,nan) (0.37250,nan) (0.37500,nan) (0.37750,nan) (0.38000,nan) (0.38250,nan) (0.38500,nan) (0.38750,nan) (0.39000,nan) (0.39250,nan) (0.39500,nan) (0.39750,nan) (0.40000,nan) (0.40250,nan) (0.40500,nan) (0.40750,nan) (0.41000,nan) (0.41250,nan) (0.41500,nan) (0.41750,nan) (0.42000,nan) (0.42250,nan) (0.42500,nan) (0.42750,nan) (0.43000,nan) (0.43250,nan) (0.43500,nan) (0.43750,nan) (0.44000,nan) (0.44250,nan) (0.44500,nan) (0.44750,nan) (0.45000,nan) (0.45250,nan) (0.45500,nan) (0.45750,nan) (0.46000,nan) (0.46250,nan) (0.46500,nan) (0.46750,nan) (0.47000,nan) (0.47250,nan) (0.47500,nan) (0.47750,nan) (0.48000,nan) (0.48250,nan) (0.48500,nan) (0.48750,nan) (0.49000,nan) (0.49250,nan) (0.49500,nan) (0.49750,nan) (0.50000,nan) (0.50250,nan) (0.50500,nan) (0.50750,nan) (0.51000,nan) (0.51250,nan) (0.51500,nan) (0.51750,nan) (0.52000,nan) (0.52250,nan) (0.52500,nan) (0.52750,nan) (0.53000,nan) (0.53250,nan) (0.53500,nan) (0.53750,nan) (0.54000,nan) (0.54250,nan) (0.54500,nan) (0.54750,nan) (0.55000,nan) (0.55250,nan) (0.55500,nan) (0.55750,nan) (0.56000,nan) (0.56250,nan) (0.56500,nan) (0.56750,nan) (0.57000,nan) (0.57250,nan) (0.57500,nan) (0.57750,nan) (0.58000,nan) (0.58250,nan) (0.58500,nan) (0.58750,nan) (0.59000,nan) (0.59250,nan) (0.59500,nan) (0.59750,nan) (0.60000,nan) (0.60250,nan) (0.60500,nan) (0.60750,nan) (0.61000,nan) (0.61250,nan) (0.61500,nan) (0.61750,nan) (0.62000,nan) (0.62250,nan) (0.62500,nan) (0.62750,nan) (0.63000,nan) (0.63250,nan) (0.63500,nan) (0.63750,nan) (0.64000,nan) (0.64250,nan) (0.64500,nan) (0.64750,nan) (0.65000,nan) (0.65250,nan) (0.65500,nan) (0.65750,nan) (0.66000,nan) (0.66250,nan) (0.66500,nan) (0.66750,nan) (0.67000,nan) (0.67250,nan) (0.67500,nan) (0.67750,nan) (0.68000,nan) (0.68250,nan) (0.68500,nan) (0.68750,nan) (0.69000,nan) (0.69250,nan) (0.69500,nan) (0.69750,nan) (0.70000,nan) (0.70250,nan) (0.70500,nan) (0.70750,nan) (0.71000,nan) (0.71250,nan) (0.71500,nan) (0.71750,nan) (0.72000,nan) (0.72250,nan) (0.72500,nan) (0.72750,nan) (0.73000,nan) (0.73250,nan) (0.73500,nan) (0.73750,nan) (0.74000,nan) (0.74250,nan) (0.74500,nan) (0.74750,nan) (0.75000,nan) (0.75250,nan) (0.75500,nan) (0.75750,nan) (0.76000,nan) (0.76250,nan) (0.76500,nan) (0.76750,nan) (0.77000,nan) (0.77250,nan) (0.77500,nan) (0.77750,nan) (0.78000,nan) (0.78250,nan) (0.78500,nan) (0.78750,nan) (0.79000,nan) (0.79250,nan) (0.79500,nan) (0.79750,nan) (0.80000,nan) (0.80250,nan) (0.80500,nan) (0.80750,nan) (0.81000,nan) (0.81250,nan) (0.81500,nan) (0.81750,nan) (0.82000,nan) (0.82250,nan) (0.82500,nan) (0.82750,nan) (0.83000,nan) (0.83250,nan) (0.83500,nan) (0.83750,nan) (0.84000,nan) (0.84250,nan) (0.84500,nan) (0.84750,nan) (0.85000,nan) (0.85250,nan) (0.85500,nan) (0.85750,nan) (0.86000,nan) (0.86250,nan) (0.86500,nan) (0.86750,nan) (0.87000,nan) (0.87250,nan) (0.87500,nan) (0.87750,nan) (0.88000,nan) (0.88250,nan) (0.88500,nan) (0.88750,nan) (0.89000,nan) (0.89250,nan) (0.89500,nan) (0.89750,nan) (0.90000,nan) (0.90250,nan) (0.90500,nan) (0.90750,nan) (0.91000,nan) (0.91250,nan) (0.91500,nan) (0.91750,nan) (0.92000,nan) (0.92250,nan) (0.92500,nan) (0.92750,nan) (0.93000,nan) (0.93250,nan) (0.93500,nan) (0.93750,nan) (0.94000,nan) (0.94250,nan) (0.94500,nan) (0.94750,nan) (0.95000,nan) (0.95250,nan) (0.95500,nan) (0.95750,nan) (0.96000,nan) (0.96250,nan) (0.96500,nan) (0.96750,nan) (0.97000,nan) (0.97250,nan) (0.97500,nan) (0.97750,nan) (0.98000,nan) (0.98250,nan) (0.98500,nan) (0.98750,nan) (0.99000,nan) (0.99250,nan) (0.99500,nan) (0.99750,nan) (1.00000,nan) (1.00250,nan) (1.00500,nan) (1.00750,nan) (1.01000,nan) (1.01250,nan) (1.01500,nan) (1.01750,nan) (1.02000,nan) (1.02250,nan) (1.02500,nan) (1.02750,nan) (1.03000,nan) (1.03250,nan) (1.03500,nan) (1.03750,nan) (1.04000,nan) (1.04250,nan) (1.04500,nan) (1.04750,nan) (1.05000,nan) (1.05250,nan) (1.05500,nan) (1.05750,nan) (1.06000,nan) (1.06250,nan) (1.06500,nan) (1.06750,nan) (1.07000,nan) (1.07250,nan) (1.07500,nan) (1.07750,nan) (1.08000,nan) (1.08250,nan) (1.08500,nan) (1.08750,nan) (1.09000,nan) (1.09250,nan) (1.09500,nan) (1.09750,nan) (1.10000,nan) (1.10250,nan) (1.10500,nan) (1.10750,nan) (1.11000,nan) (1.11250,nan) (1.11500,nan) (1.11750,nan) (1.12000,nan) (1.12250,nan) (1.12500,nan) (1.12750,nan) (1.13000,nan) (1.13250,nan) (1.13500,nan) (1.13750,nan) (1.14000,nan) (1.14250,nan) (1.14500,nan) (1.14750,nan) (1.15000,nan) (1.15250,nan) (1.15500,nan) (1.15750,nan) (1.16000,nan) (1.16250,nan) (1.16500,nan) (1.16750,nan) (1.17000,nan) (1.17250,nan) (1.17500,nan) (1.17750,nan) (1.18000,nan) (1.18250,nan) (1.18500,nan) (1.18750,nan) (1.19000,nan) (1.19250,nan) (1.19500,nan) (1.19750,nan) (1.20000,nan) (1.20250,nan) (1.20500,nan) (1.20750,nan) (1.21000,nan) (1.21250,nan) (1.21500,nan) (1.21750,nan) (1.22000,nan) (1.22250,nan) (1.22500,nan) (1.22750,nan) (1.23000,nan) (1.23250,nan) (1.23500,nan) (1.23750,nan) (1.24000,nan) (1.24250,nan) (1.24500,nan) (1.24750,nan) (1.25000,nan) (1.25250,nan) (1.25500,nan) (1.25750,nan) (1.26000,nan) (1.26250,nan) (1.26500,nan) (1.26750,nan) (1.27000,nan) (1.27250,nan) (1.27500,nan) (1.27750,nan) (1.28000,nan) (1.28250,nan) (1.28500,nan) (1.28750,nan) (1.29000,nan) (1.29250,nan) (1.29500,nan) (1.29750,nan) (1.30000,nan) (1.30250,nan) (1.30500,nan) (1.30750,nan) (1.31000,nan) (1.31250,nan) (1.31500,nan) (1.31750,nan) (1.32000,nan) (1.32250,nan) (1.32500,nan) (1.32750,nan) (1.33000,nan) (1.33250,nan) (1.33500,nan) (1.33750,nan) (1.34000,nan) (1.34250,nan) (1.34500,nan) (1.34750,nan) (1.35000,nan) (1.35250,nan) (1.35500,nan) (1.35750,nan) (1.36000,nan) (1.36250,nan) (1.36500,nan) (1.36750,nan) (1.37000,nan) (1.37250,nan) (1.37500,nan) (1.37750,nan) (1.38000,nan) (1.38250,nan) (1.38500,nan) (1.38750,nan) (1.39000,nan) (1.39250,nan) (1.39500,nan) (1.39750,nan) (1.40000,nan) (1.40250,nan) (1.40500,nan) (1.40750,nan) (1.41000,nan) (1.41250,nan) (1.41500,nan) (1.41750,nan) (1.42000,nan) (1.42250,nan) (1.42500,nan) (1.42750,nan) (1.43000,nan) (1.43250,nan) (1.43500,nan) (1.43750,nan) (1.44000,nan) (1.44250,nan) (1.44500,nan) (1.44750,nan) (1.45000,nan) (1.45250,nan) (1.45500,nan) (1.45750,nan) (1.46000,nan) (1.46250,nan) (1.46500,nan) (1.46750,nan) (1.47000,nan) (1.47250,nan) (1.47500,nan) (1.47750,nan) (1.48000,nan) (1.48250,nan) (1.48500,nan) (1.48750,nan) (1.49000,nan) (1.49250,nan) (1.49500,nan) (1.49750,nan) (1.50000,nan) (1.50250,nan) (1.50500,nan) (1.50750,nan) (1.51000,nan) (1.51250,nan) (1.51500,nan) (1.51750,nan) (1.52000,nan) (1.52250,nan) (1.52500,nan) (1.52750,nan) (1.53000,nan) (1.53250,nan) (1.53500,nan) (1.53750,nan) (1.54000,nan) (1.54250,nan) (1.54500,nan) (1.54750,nan) (1.55000,nan) (1.55250,nan) (1.55500,nan) (1.55750,nan) (1.56000,nan) (1.56250,nan) (1.56500,nan) (1.56750,nan) (1.57000,nan) (1.57250,nan) (1.57500,nan) (1.57750,nan) (1.58000,nan) (1.58250,nan) (1.58500,nan) (1.58750,nan) (1.59000,nan) (1.59250,nan) (1.59500,nan) (1.59750,nan) (1.60000,nan) (1.60250,nan) (1.60500,nan) (1.60750,nan) (1.61000,nan) (1.61250,nan) (1.61500,nan) (1.61750,nan) (1.62000,nan) (1.62250,nan) (1.62500,nan) (1.62750,nan) (1.63000,nan) (1.63250,nan) (1.63500,nan) (1.63750,nan) (1.64000,nan) (1.64250,nan) (1.64500,nan) (1.64750,nan) (1.65000,nan) (1.65250,nan) (1.65500,nan) (1.65750,nan) (1.66000,nan) (1.66250,nan) (1.66500,nan) (1.66750,nan) (1.67000,nan) (1.67250,nan) (1.67500,nan) (1.67750,nan) (1.68000,nan) (1.68250,nan) (1.68500,nan) (1.68750,nan) (1.69000,nan) (1.69250,nan) (1.69500,nan) (1.69750,nan) (1.70000,nan) (1.70250,nan) (1.70500,nan) (1.70750,nan) (1.71000,nan) (1.71250,nan) (1.71500,nan) (1.71750,nan) (1.72000,nan) (1.72250,nan) (1.72500,nan) (1.72750,nan) (1.73000,nan) (1.73250,nan) (1.73500,nan) (1.73750,nan) (1.74000,nan) (1.74250,nan) (1.74500,nan) (1.74750,nan) (1.75000,nan) (1.75250,nan) (1.75500,nan) (1.75750,nan) (1.76000,nan) (1.76250,nan) (1.76500,nan) (1.76750,nan) (1.77000,nan) (1.77250,nan) (1.77500,nan) (1.77750,nan) (1.78000,nan) (1.78250,nan) (1.78500,nan) (1.78750,nan) (1.79000,nan) (1.79250,nan) (1.79500,nan) (1.79750,nan) (1.80000,nan) (1.80250,nan) (1.80500,nan) (1.80750,nan) (1.81000,nan) (1.81250,nan) (1.81500,nan) (1.81750,nan) (1.82000,nan) (1.82250,nan) (1.82500,nan) (1.82750,nan) (1.83000,nan) (1.83250,nan) (1.83500,nan) (1.83750,nan) (1.84000,nan) (1.84250,nan) (1.84500,nan) (1.84750,nan) (1.85000,nan) (1.85250,nan) (1.85500,nan) (1.85750,nan) (1.86000,nan) (1.86250,nan) (1.86500,nan) (1.86750,nan) (1.87000,nan) (1.87250,nan) (1.87500,nan) (1.87750,nan) (1.88000,nan) (1.88250,nan) (1.88500,nan) (1.88750,nan) (1.89000,nan) (1.89250,nan) (1.89500,nan) (1.89750,nan) (1.90000,nan) (1.90250,nan) (1.90500,nan) (1.90750,nan) (1.91000,nan) (1.91250,nan) (1.91500,nan) (1.91750,nan) (1.92000,nan) (1.92250,nan) (1.92500,nan) (1.92750,nan) (1.93000,nan) (1.93250,nan) (1.93500,nan) (1.93750,nan) (1.94000,nan) (1.94250,nan) (1.94500,nan) (1.94750,nan) (1.95000,nan) (1.95250,nan) (1.95500,nan) (1.95750,nan) (1.96000,nan) (1.96250,nan) (1.96500,nan) (1.96750,nan) (1.97000,nan) (1.97250,nan) (1.97500,nan) (1.97750,nan) (1.98000,nan) (1.98250,nan) (1.98500,nan) (1.98750,nan) (1.99000,nan) (1.99250,nan) (1.99500,nan) (1.99750,nan) (2.00000,nan) (2.00250,nan) (2.00500,nan) (2.00750,nan) (2.01000,nan) (2.01250,nan) (2.01500,nan) (2.01750,nan) (2.02000,nan) (2.02250,nan) (2.02500,nan) (2.02750,nan) (2.03000,nan) (2.03250,nan) (2.03500,nan) (2.03750,nan) (2.04000,nan) (2.04250,nan) (2.04500,nan) (2.04750,nan) (2.05000,nan) (2.05250,nan) (2.05500,nan) (2.05750,nan) (2.06000,nan) (2.06250,nan) (2.06500,nan) (2.06750,nan) (2.07000,nan) (2.07250,nan) (2.07500,nan) (2.07750,nan) (2.08000,nan) (2.08250,nan) (2.08500,nan) (2.08750,nan) (2.09000,nan) (2.09250,nan) (2.09500,nan) (2.09750,nan) (2.10000,nan) (2.10250,nan) (2.10500,nan) (2.10750,nan) (2.11000,nan) (2.11250,nan) (2.11500,nan) (2.11750,nan) (2.12000,nan) (2.12250,nan) (2.12500,nan) (2.12750,nan) (2.13000,nan) (2.13250,nan) (2.13500,nan) (2.13750,nan) (2.14000,nan) (2.14250,nan) (2.14500,nan) (2.14750,nan) (2.15000,nan) (2.15250,nan) (2.15500,nan) (2.15750,nan) (2.16000,nan) (2.16250,nan) (2.16500,nan) (2.16750,nan) (2.17000,nan) (2.17250,nan) (2.17500,nan) (2.17750,nan) (2.18000,nan) (2.18250,nan) (2.18500,nan) (2.18750,nan) (2.19000,nan) (2.19250,nan) (2.19500,nan) (2.19750,nan) (2.20000,nan) (2.20250,nan) (2.20500,nan) (2.20750,nan) (2.21000,nan) (2.21250,nan) (2.21500,nan) (2.21750,nan) (2.22000,nan) (2.22250,nan) (2.22500,nan) (2.22750,nan) (2.23000,nan) (2.23250,nan) (2.23500,nan) (2.23750,nan) (2.24000,nan) (2.24250,nan) (2.24500,nan) (2.24750,nan) (2.25000,nan) (2.25250,nan) (2.25500,nan) (2.25750,nan) (2.26000,nan) (2.26250,nan) (2.26500,nan) (2.26750,nan) (2.27000,nan) (2.27250,nan) (2.27500,nan) (2.27750,nan) (2.28000,nan) (2.28250,nan) (2.28500,nan) (2.28750,nan) (2.29000,nan) (2.29250,nan) (2.29500,nan) (2.29750,nan) (2.30000,nan) (2.30250,nan) (2.30500,nan) (2.30750,nan) (2.31000,nan) (2.31250,nan) (2.31500,nan) (2.31750,nan) (2.32000,nan) (2.32250,nan) (2.32500,nan) (2.32750,nan) (2.33000,nan) (2.33250,nan) (2.33500,nan) (2.33750,nan) (2.34000,nan) (2.34250,nan) (2.34500,nan) (2.34750,nan) (2.35000,nan) (2.35250,nan) (2.35500,nan) (2.35750,nan) (2.36000,nan) (2.36250,nan) (2.36500,nan) (2.36750,nan) (2.37000,nan) (2.37250,nan) (2.37500,nan) (2.37750,nan) (2.38000,nan) (2.38250,nan) (2.38500,nan) (2.38750,nan) (2.39000,nan) (2.39250,nan) (2.39500,nan) (2.39750,nan) (2.40000,nan) (2.40250,nan) (2.40500,nan) (2.40750,nan) (2.41000,nan) (2.41250,nan) (2.41500,nan) (2.41750,nan) (2.42000,nan) (2.42250,nan) (2.42500,nan) (2.42750,nan) (2.43000,nan) (2.43250,nan) (2.43500,nan) (2.43750,nan) (2.44000,nan) (2.44250,nan) (2.44500,nan) (2.44750,nan) (2.45000,nan) (2.45250,nan) (2.45500,nan) (2.45750,nan) (2.46000,nan) (2.46250,nan) (2.46500,nan) (2.46750,nan) (2.47000,nan) (2.47250,nan) (2.47500,nan) (2.47750,nan) (2.48000,nan) (2.48250,nan) (2.48500,nan) (2.48750,nan) (2.49000,nan) (2.49250,nan) (2.49500,nan) (2.49750,nan) (2.50000,nan) (2.50250,nan) (2.50500,nan) (2.50750,nan) (2.51000,nan) (2.51250,nan) (2.51500,nan) (2.51750,nan) (2.52000,nan) (2.52250,nan) (2.52500,nan) (2.52750,nan) (2.53000,nan) (2.53250,nan) (2.53500,nan) (2.53750,nan) (2.54000,nan) (2.54250,nan) (2.54500,nan) (2.54750,nan) (2.55000,nan) (2.55250,nan) (2.55500,nan) (2.55750,nan) (2.56000,nan) (2.56250,nan) (2.56500,nan) (2.56750,nan) (2.57000,nan) (2.57250,nan) (2.57500,nan) (2.57750,nan) (2.58000,nan) (2.58250,nan) (2.58500,nan) (2.58750,nan) (2.59000,nan) (2.59250,nan) (2.59500,nan) (2.59750,nan) (2.60000,nan) (2.60250,nan) (2.60500,nan) (2.60750,nan) (2.61000,nan) (2.61250,nan) (2.61500,nan) (2.61750,nan) (2.62000,nan) (2.62250,nan) (2.62500,nan) (2.62750,nan) (2.63000,nan) (2.63250,nan) (2.63500,nan) (2.63750,nan) (2.64000,nan) (2.64250,nan) (2.64500,nan) (2.64750,nan) (2.65000,nan) (2.65250,nan) (2.65500,nan) (2.65750,nan) (2.66000,nan) (2.66250,nan) (2.66500,nan) (2.66750,nan) (2.67000,nan) (2.67250,nan) (2.67500,nan) (2.67750,nan) (2.68000,nan) (2.68250,nan) (2.68500,nan) (2.68750,nan) (2.69000,nan) (2.69250,nan) (2.69500,nan) (2.69750,nan) (2.70000,nan) (2.70250,nan) (2.70500,nan) (2.70750,nan) (2.71000,nan) (2.71250,nan) (2.71500,nan) (2.71750,nan) (2.72000,nan) (2.72250,nan) (2.72500,nan) (2.72750,nan) (2.73000,nan) (2.73250,nan) (2.73500,nan) (2.73750,nan) (2.74000,nan) (2.74250,nan) (2.74500,nan) (2.74750,nan) (2.75000,nan) (2.75250,nan) (2.75500,nan) (2.75750,nan) (2.76000,nan) (2.76250,nan) (2.76500,nan) (2.76750,nan) (2.77000,nan) (2.77250,nan) (2.77500,nan) (2.77750,nan) (2.78000,nan) (2.78250,nan) (2.78500,nan) (2.78750,nan) (2.79000,nan) (2.79250,nan) (2.79500,nan) (2.79750,nan) (2.80000,nan) (2.80250,nan) (2.80500,nan) (2.80750,nan) (2.81000,nan) (2.81250,nan) (2.81500,nan) (2.81750,nan) (2.82000,nan) (2.82250,nan) (2.82500,nan) (2.82750,nan) (2.83000,nan) (2.83250,nan) (2.83500,nan) (2.83750,nan) (2.84000,nan) (2.84250,nan) (2.84500,nan) (2.84750,nan) (2.85000,nan) (2.85250,nan) (2.85500,nan) (2.85750,nan) (2.86000,nan) (2.86250,nan) (2.86500,nan) (2.86750,nan) (2.87000,nan) (2.87250,nan) (2.87500,nan) (2.87750,nan) (2.88000,nan) (2.88250,nan) (2.88500,nan) (2.88750,nan) (2.89000,nan) (2.89250,nan) (2.89500,nan) (2.89750,nan) (2.90000,nan) (2.90250,nan) (2.90500,nan) (2.90750,nan) (2.91000,nan) (2.91250,nan) (2.91500,nan) (2.91750,nan) (2.92000,nan) (2.92250,nan) (2.92500,nan) (2.92750,nan) (2.93000,nan) (2.93250,nan) (2.93500,nan) (2.93750,nan) (2.94000,nan) (2.94250,nan) (2.94500,nan) (2.94750,nan) (2.95000,nan) (2.95250,nan) (2.95500,nan) (2.95750,nan) (2.96000,nan) (2.96250,nan) (2.96500,nan) (2.96750,nan) (2.97000,nan) (2.97250,nan) (2.97500,nan) (2.97750,nan) (2.98000,nan) (2.98250,nan) (2.98500,nan) (2.98750,nan) (2.99000,nan) (2.99250,nan) (2.99500,nan) (2.99750,nan) (3.00000,nan) (3.00250,nan) (3.00500,nan) (3.00750,nan) (3.01000,nan) (3.01250,nan) (3.01500,nan) (3.01750,nan) (3.02000,nan) (3.02250,nan) (3.02500,nan) (3.02750,nan) (3.03000,nan) (3.03250,nan) (3.03500,nan) (3.03750,nan) (3.04000,nan) (3.04250,nan) (3.04500,nan) (3.04750,nan) (3.05000,nan) (3.05250,nan) (3.05500,nan) (3.05750,nan) (3.06000,nan) (3.06250,nan) (3.06500,nan) (3.06750,nan) (3.07000,nan) (3.07250,nan) (3.07500,nan) (3.07750,nan) (3.08000,nan) (3.08250,nan) (3.08500,nan) (3.08750,nan) (3.09000,nan) (3.09250,nan) (3.09500,nan) (3.09750,nan) (3.10000,nan) (3.10250,nan) (3.10500,nan) (3.10750,nan) (3.11000,nan) (3.11250,nan) (3.11500,nan) (3.11750,nan) (3.12000,nan) (3.12250,nan) (3.12500,nan) (3.12750,nan) (3.13000,nan) (3.13250,nan) (3.13500,nan) (3.13750,nan) (3.14000,nan) (3.14250,nan) (3.14500,nan) (3.14750,nan) (3.15000,nan) (3.15250,nan) (3.15500,nan) (3.15750,nan) (3.16000,nan) (3.16250,nan) (3.16500,nan) (3.16750,nan) (3.17000,nan) (3.17250,nan) (3.17500,nan) (3.17750,nan) (3.18000,nan) (3.18250,nan) (3.18500,nan) (3.18750,nan) (3.19000,nan) (3.19250,nan) (3.19500,nan) (3.19750,nan) (3.20000,nan) (3.20250,nan) (3.20500,nan) (3.20750,nan) (3.21000,nan) (3.21250,nan) (3.21500,nan) (3.21750,nan) (3.22000,nan) (3.22250,nan) (3.22500,nan) (3.22750,nan) (3.23000,nan) (3.23250,nan) (3.23500,nan) (3.23750,nan) (3.24000,nan) (3.24250,nan) (3.24500,nan) (3.24750,nan) (3.25000,nan) (3.25250,nan) (3.25500,nan) (3.25750,nan) (3.26000,nan) (3.26250,nan) (3.26500,nan) (3.26750,nan) (3.27000,nan) (3.27250,nan) (3.27500,nan) (3.27750,nan) (3.28000,nan) (3.28250,nan) (3.28500,nan) (3.28750,nan) (3.29000,nan) (3.29250,nan) (3.29500,nan) (3.29750,nan) (3.30000,nan) (3.30250,nan) (3.30500,nan) (3.30750,nan) (3.31000,nan) (3.31250,nan) (3.31500,nan) (3.31750,nan) (3.32000,nan) (3.32250,nan) (3.32500,nan) (3.32750,nan) (3.33000,nan) (3.33250,nan) (3.33500,nan) (3.33750,nan) (3.34000,nan) (3.34250,nan) (3.34500,nan) (3.34750,nan) (3.35000,nan) (3.35250,nan) (3.35500,nan) (3.35750,nan) (3.36000,nan) (3.36250,nan) (3.36500,nan) (3.36750,nan) (3.37000,nan) (3.37250,nan) (3.37500,nan) (3.37750,nan) (3.38000,nan) (3.38250,nan) (3.38500,nan) (3.38750,nan) (3.39000,nan) (3.39250,nan) (3.39500,nan) (3.39750,nan) (3.40000,nan) (3.40250,nan) (3.40500,nan) (3.40750,nan) (3.41000,nan) (3.41250,nan) (3.41500,nan) (3.41750,nan) (3.42000,nan) (3.42250,nan) (3.42500,nan) (3.42750,nan) (3.43000,nan) (3.43250,nan) (3.43500,nan) (3.43750,nan) (3.44000,nan) (3.44250,nan) (3.44500,nan) (3.44750,nan) (3.45000,nan) (3.45250,nan) (3.45500,nan) (3.45750,nan) (3.46000,nan) (3.46250,nan) (3.46500,nan) (3.46750,nan) (3.47000,nan) (3.47250,nan) (3.47500,nan) (3.47750,nan) (3.48000,nan) (3.48250,nan) (3.48500,nan) (3.48750,nan) (3.49000,nan) (3.49250,nan) (3.49500,nan) (3.49750,nan) (3.50000,nan) (3.50250,nan) (3.50500,nan) (3.50750,nan) (3.51000,nan) (3.51250,nan) (3.51500,nan) (3.51750,nan) (3.52000,nan) (3.52250,-1.7927223) (3.52500,-1.6118214) (3.52750,-1.4853267) (3.53000,-1.3881445) (3.53250,-1.3093278) (3.53500,-1.2431075) (3.53750,-1.1860672) (3.54000,-1.1360157) (3.54250,-1.0914634) (3.54500,-1.0513514) (3.54750,-1.0148998) (3.55000,-0.9815177) (3.55250,-0.9507466) (3.55500,-0.9222237) (3.55750,-0.8956568) (3.56000,-0.8708073) (3.56250,-0.8474777) (3.56500,-0.8255025) (3.56750,-0.8047422) (3.57000,-0.7850775) (3.57250,-0.7664060) (3.57500,-0.7486389) (3.57750,-0.7316988) (3.58000,-0.7155179) (3.58250,-0.7000363) (3.58500,-0.6852010) (3.58750,-0.6709648) (3.59000,-0.6572854) (3.59250,-0.6441248) (3.59500,-0.6314489) (3.59750,-0.6192267) (3.60000,-0.6074303) (3.60250,-0.5960339) (3.60500,-0.5850144) (3.60750,-0.5743503) (3.61000,-0.5640221) (3.61250,-0.5540118) (3.61500,-0.5443027) (3.61750,-0.5348796) (3.62000,-0.5257281) (3.62250,-0.5168352) (3.62500,-0.5081885) (3.62750,-0.4997768) (3.63000,-0.4915893) (3.63250,-0.4836161) (3.63500,-0.4758480) (3.63750,-0.4682762) (3.64000,-0.4608926) (3.64250,-0.4536895) (3.64500,-0.4466597) (3.64750,-0.4397965) (3.65000,-0.4330934) (3.65250,-0.4265443) (3.65500,-0.4201436) (3.65750,-0.4138859) (3.66000,-0.4077660) (3.66250,-0.4017792) (3.66500,-0.3959208) (3.66750,-0.3901864) (3.67000,-0.3845719) (3.67250,-0.3790733) (3.67500,-0.3736870) (3.67750,-0.3684093) (3.68000,-0.3632367) (3.68250,-0.3581661) (3.68500,-0.3531944) (3.68750,-0.3483185) (3.69000,-0.3435356) (3.69250,-0.3388430) (3.69500,-0.3342381) (3.69750,-0.3297185) (3.70000,-0.3252816) (3.70250,-0.3209252) (3.70500,-0.3166472) (3.70750,-0.3124453) (3.71000,-0.3083176) (3.71250,-0.3042622) (3.71500,-0.3002770) (3.71750,-0.2963604) (3.72000,-0.2925106) (3.72250,-0.2887259) (3.72500,-0.2850046) (3.72750,-0.2813453) (3.73000,-0.2777465) (3.73250,-0.2742066) (3.73500,-0.2707243) (3.73750,-0.2672982) (3.74000,-0.2639271) (3.74250,-0.2606096) (3.74500,-0.2573446) (3.74750,-0.2541308) (3.75000,-0.2509671) (3.75250,-0.2478524) (3.75500,-0.2447857) (3.75750,-0.2417658) (3.76000,-0.2387919) (3.76250,-0.2358628) (3.76500,-0.2329778) (3.76750,-0.2301358) (3.77000,-0.2273359) (3.77250,-0.2245773) (3.77500,-0.2218592) (3.77750,-0.2191808) (3.78000,-0.2165412) (3.78250,-0.2139397) (3.78500,-0.2113755) (3.78750,-0.2088480) (3.79000,-0.2063563) (3.79250,-0.2038999) (3.79500,-0.2014780) (3.79750,-0.1990900) (3.80000,-0.1967352) (3.80250,-0.1944131) (3.80500,-0.1921231) (3.80750,-0.1898645) (3.81000,-0.1876368) (3.81250,-0.1854394) (3.81500,-0.1832719) (3.81750,-0.1811336) (3.82000,-0.1790240) (3.82250,-0.1769427) (3.82500,-0.1748891) (3.82750,-0.1728628) (3.83000,-0.1708633) (3.83250,-0.1688901) (3.83500,-0.1669429) (3.83750,-0.1650211) (3.84000,-0.1631244) (3.84250,-0.1612523) (3.84500,-0.1594044) (3.84750,-0.1575803) (3.85000,-0.1557797) (3.85250,-0.1540021) (3.85500,-0.1522472) (3.85750,-0.1505146) (3.86000,-0.1488040) (3.86250,-0.1471150) (3.86500,-0.1454473) (3.86750,-0.1438006) (3.87000,-0.1421745) (3.87250,-0.1405686) (3.87500,-0.1389828) (3.87750,-0.1374166) (3.88000,-0.1358699) (3.88250,-0.1343422) (3.88500,-0.1328333) (3.88750,-0.1313430) (3.89000,-0.1298709) (3.89250,-0.1284167) (3.89500,-0.1269803) (3.89750,-0.1255614) (3.90000,-0.1241596) (3.90250,-0.1227748) (3.90500,-0.1214067) (3.90750,-0.1200550) (3.91000,-0.1187196) (3.91250,-0.1174002) (3.91500,-0.1160965) (3.91750,-0.1148084) (3.92000,-0.1135356) (3.92250,-0.1122779) (3.92500,-0.1110351) (3.92750,-0.1098070) (3.93000,-0.1085934) (3.93250,-0.1073940) (3.93500,-0.1062088) (3.93750,-0.1050375) (3.94000,-0.1038798) (3.94250,-0.1027357) (3.94500,-0.1016050) (3.94750,-0.1004874) (3.95000,-0.0993828) (3.95250,-0.0982910) (3.95500,-0.0972118) (3.95750,-0.0961451) (3.96000,-0.0950908) (3.96250,-0.0940485) (3.96500,-0.0930183) (3.96750,-0.0919999) (3.97000,-0.0909932) (3.97250,-0.0899980) (3.97500,-0.0890141) (3.97750,-0.0880415) (3.98000,-0.0870800) (3.98250,-0.0861294) (3.98500,-0.0851897) (3.98750,-0.0842606) (3.99000,-0.0833420) (3.99250,-0.0824338) (3.99500,-0.0815359) (3.99750,-0.0806481) (4.00000,-0.0797704)};\addlegendentry{L}
\addplot[thick,dashed,gray] coordinates {(-2.00000,nan) (-1.99750,nan) (-1.99500,nan) (-1.99250,nan) (-1.99000,nan) (-1.98750,nan) (-1.98500,nan) (-1.98250,nan) (-1.98000,nan) (-1.97750,nan) (-1.97500,nan) (-1.97250,nan) (-1.97000,nan) (-1.96750,nan) (-1.96500,nan) (-1.96250,nan) (-1.96000,nan) (-1.95750,nan) (-1.95500,nan) (-1.95250,nan) (-1.95000,nan) (-1.94750,nan) (-1.94500,nan) (-1.94250,nan) (-1.94000,nan) (-1.93750,nan) (-1.93500,nan) (-1.93250,nan) (-1.93000,nan) (-1.92750,nan) (-1.92500,nan) (-1.92250,nan) (-1.92000,nan) (-1.91750,nan) (-1.91500,nan) (-1.91250,nan) (-1.91000,nan) (-1.90750,nan) (-1.90500,nan) (-1.90250,nan) (-1.90000,nan) (-1.89750,nan) (-1.89500,nan) (-1.89250,nan) (-1.89000,nan) (-1.88750,nan) (-1.88500,nan) (-1.88250,nan) (-1.88000,nan) (-1.87750,nan) (-1.87500,nan) (-1.87250,nan) (-1.87000,nan) (-1.86750,nan) (-1.86500,nan) (-1.86250,nan) (-1.86000,nan) (-1.85750,nan) (-1.85500,nan) (-1.85250,nan) (-1.85000,nan) (-1.84750,nan) (-1.84500,nan) (-1.84250,nan) (-1.84000,nan) (-1.83750,nan) (-1.83500,nan) (-1.83250,nan) (-1.83000,nan) (-1.82750,nan) (-1.82500,nan) (-1.82250,nan) (-1.82000,nan) (-1.81750,nan) (-1.81500,nan) (-1.81250,nan) (-1.81000,nan) (-1.80750,nan) (-1.80500,nan) (-1.80250,nan) (-1.80000,nan) (-1.79750,nan) (-1.79500,nan) (-1.79250,nan) (-1.79000,nan) (-1.78750,nan) (-1.78500,nan) (-1.78250,nan) (-1.78000,nan) (-1.77750,nan) (-1.77500,nan) (-1.77250,nan) (-1.77000,nan) (-1.76750,nan) (-1.76500,nan) (-1.76250,nan) (-1.76000,nan) (-1.75750,nan) (-1.75500,nan) (-1.75250,nan) (-1.75000,nan) (-1.74750,nan) (-1.74500,nan) (-1.74250,nan) (-1.74000,nan) (-1.73750,nan) (-1.73500,nan) (-1.73250,nan) (-1.73000,nan) (-1.72750,nan) (-1.72500,nan) (-1.72250,nan) (-1.72000,nan) (-1.71750,nan) (-1.71500,nan) (-1.71250,nan) (-1.71000,nan) (-1.70750,nan) (-1.70500,nan) (-1.70250,nan) (-1.70000,nan) (-1.69750,nan) (-1.69500,nan) (-1.69250,nan) (-1.69000,nan) (-1.68750,nan) (-1.68500,nan) (-1.68250,nan) (-1.68000,nan) (-1.67750,nan) (-1.67500,nan) (-1.67250,nan) (-1.67000,nan) (-1.66750,nan) (-1.66500,nan) (-1.66250,nan) (-1.66000,nan) (-1.65750,nan) (-1.65500,nan) (-1.65250,nan) (-1.65000,nan) (-1.64750,nan) (-1.64500,nan) (-1.64250,nan) (-1.64000,nan) (-1.63750,nan) (-1.63500,nan) (-1.63250,nan) (-1.63000,nan) (-1.62750,nan) (-1.62500,nan) (-1.62250,nan) (-1.62000,nan) (-1.61750,nan) (-1.61500,nan) (-1.61250,nan) (-1.61000,nan) (-1.60750,nan) (-1.60500,nan) (-1.60250,nan) (-1.60000,nan) (-1.59750,nan) (-1.59500,nan) (-1.59250,nan) (-1.59000,nan) (-1.58750,nan) (-1.58500,nan) (-1.58250,nan) (-1.58000,nan) (-1.57750,nan) (-1.57500,nan) (-1.57250,nan) (-1.57000,nan) (-1.56750,nan) (-1.56500,nan) (-1.56250,nan) (-1.56000,nan) (-1.55750,nan) (-1.55500,nan) (-1.55250,nan) (-1.55000,nan) (-1.54750,nan) (-1.54500,nan) (-1.54250,nan) (-1.54000,nan) (-1.53750,nan) (-1.53500,nan) (-1.53250,nan) (-1.53000,nan) (-1.52750,nan) (-1.52500,nan) (-1.52250,nan) (-1.52000,nan) (-1.51750,nan) (-1.51500,nan) (-1.51250,nan) (-1.51000,nan) (-1.50750,nan) (-1.50500,nan) (-1.50250,nan) (-1.50000,nan) (-1.49750,nan) (-1.49500,nan) (-1.49250,nan) (-1.49000,nan) (-1.48750,nan) (-1.48500,nan) (-1.48250,nan) (-1.48000,nan) (-1.47750,nan) (-1.47500,nan) (-1.47250,nan) (-1.47000,nan) (-1.46750,nan) (-1.46500,nan) (-1.46250,nan) (-1.46000,nan) (-1.45750,nan) (-1.45500,nan) (-1.45250,nan) (-1.45000,nan) (-1.44750,nan) (-1.44500,nan) (-1.44250,nan) (-1.44000,nan) (-1.43750,nan) (-1.43500,nan) (-1.43250,nan) (-1.43000,nan) (-1.42750,nan) (-1.42500,nan) (-1.42250,nan) (-1.42000,nan) (-1.41750,nan) (-1.41500,nan) (-1.41250,nan) (-1.41000,nan) (-1.40750,nan) (-1.40500,nan) (-1.40250,nan) (-1.40000,nan) (-1.39750,nan) (-1.39500,nan) (-1.39250,nan) (-1.39000,nan) (-1.38750,nan) (-1.38500,nan) (-1.38250,nan) (-1.38000,nan) (-1.37750,nan) (-1.37500,nan) (-1.37250,nan) (-1.37000,nan) (-1.36750,nan) (-1.36500,nan) (-1.36250,nan) (-1.36000,nan) (-1.35750,nan) (-1.35500,nan) (-1.35250,nan) (-1.35000,nan) (-1.34750,nan) (-1.34500,nan) (-1.34250,nan) (-1.34000,nan) (-1.33750,nan) (-1.33500,nan) (-1.33250,nan) (-1.33000,nan) (-1.32750,nan) (-1.32500,nan) (-1.32250,nan) (-1.32000,nan) (-1.31750,nan) (-1.31500,nan) (-1.31250,nan) (-1.31000,nan) (-1.30750,nan) (-1.30500,nan) (-1.30250,nan) (-1.30000,nan) (-1.29750,nan) (-1.29500,nan) (-1.29250,nan) (-1.29000,nan) (-1.28750,nan) (-1.28500,nan) (-1.28250,nan) (-1.28000,nan) (-1.27750,nan) (-1.27500,nan) (-1.27250,nan) (-1.27000,nan) (-1.26750,nan) (-1.26500,nan) (-1.26250,nan) (-1.26000,nan) (-1.25750,nan) (-1.25500,nan) (-1.25250,nan) (-1.25000,nan) (-1.24750,nan) (-1.24500,nan) (-1.24250,nan) (-1.24000,nan) (-1.23750,nan) (-1.23500,nan) (-1.23250,nan) (-1.23000,nan) (-1.22750,nan) (-1.22500,nan) (-1.22250,nan) (-1.22000,nan) (-1.21750,nan) (-1.21500,nan) (-1.21250,nan) (-1.21000,nan) (-1.20750,nan) (-1.20500,nan) (-1.20250,nan) (-1.20000,nan) (-1.19750,nan) (-1.19500,nan) (-1.19250,nan) (-1.19000,nan) (-1.18750,nan) (-1.18500,nan) (-1.18250,nan) (-1.18000,nan) (-1.17750,nan) (-1.17500,nan) (-1.17250,nan) (-1.17000,nan) (-1.16750,nan) (-1.16500,nan) (-1.16250,nan) (-1.16000,nan) (-1.15750,nan) (-1.15500,nan) (-1.15250,nan) (-1.15000,nan) (-1.14750,nan) (-1.14500,nan) (-1.14250,nan) (-1.14000,nan) (-1.13750,nan) (-1.13500,nan) (-1.13250,nan) (-1.13000,nan) (-1.12750,nan) (-1.12500,nan) (-1.12250,nan) (-1.12000,nan) (-1.11750,nan) (-1.11500,nan) (-1.11250,nan) (-1.11000,nan) (-1.10750,nan) (-1.10500,nan) (-1.10250,nan) (-1.10000,nan) (-1.09750,nan) (-1.09500,nan) (-1.09250,nan) (-1.09000,nan) (-1.08750,nan) (-1.08500,nan) (-1.08250,nan) (-1.08000,nan) (-1.07750,nan) (-1.07500,nan) (-1.07250,nan) (-1.07000,nan) (-1.06750,nan) (-1.06500,nan) (-1.06250,nan) (-1.06000,nan) (-1.05750,nan) (-1.05500,nan) (-1.05250,nan) (-1.05000,nan) (-1.04750,nan) (-1.04500,nan) (-1.04250,nan) (-1.04000,nan) (-1.03750,nan) (-1.03500,nan) (-1.03250,nan) (-1.03000,nan) (-1.02750,nan) (-1.02500,nan) (-1.02250,nan) (-1.02000,nan) (-1.01750,nan) (-1.01500,nan) (-1.01250,nan) (-1.01000,nan) (-1.00750,nan) (-1.00500,nan) (-1.00250,nan) (-1.00000,nan) (-0.99750,nan) (-0.99500,nan) (-0.99250,nan) (-0.99000,nan) (-0.98750,nan) (-0.98500,nan) (-0.98250,nan) (-0.98000,nan) (-0.97750,nan) (-0.97500,nan) (-0.97250,nan) (-0.97000,nan) (-0.96750,nan) (-0.96500,nan) (-0.96250,nan) (-0.96000,nan) (-0.95750,nan) (-0.95500,nan) (-0.95250,nan) (-0.95000,nan) (-0.94750,nan) (-0.94500,nan) (-0.94250,nan) (-0.94000,nan) (-0.93750,nan) (-0.93500,nan) (-0.93250,nan) (-0.93000,nan) (-0.92750,nan) (-0.92500,nan) (-0.92250,nan) (-0.92000,nan) (-0.91750,nan) (-0.91500,nan) (-0.91250,nan) (-0.91000,nan) (-0.90750,nan) (-0.90500,nan) (-0.90250,nan) (-0.90000,nan) (-0.89750,nan) (-0.89500,nan) (-0.89250,nan) (-0.89000,nan) (-0.88750,nan) (-0.88500,nan) (-0.88250,nan) (-0.88000,nan) (-0.87750,nan) (-0.87500,nan) (-0.87250,nan) (-0.87000,nan) (-0.86750,nan) (-0.86500,nan) (-0.86250,nan) (-0.86000,nan) (-0.85750,nan) (-0.85500,nan) (-0.85250,nan) (-0.85000,nan) (-0.84750,nan) (-0.84500,nan) (-0.84250,nan) (-0.84000,nan) (-0.83750,nan) (-0.83500,nan) (-0.83250,nan) (-0.83000,nan) (-0.82750,nan) (-0.82500,nan) (-0.82250,nan) (-0.82000,nan) (-0.81750,nan) (-0.81500,nan) (-0.81250,nan) (-0.81000,nan) (-0.80750,nan) (-0.80500,nan) (-0.80250,nan) (-0.80000,nan) (-0.79750,nan) (-0.79500,nan) (-0.79250,nan) (-0.79000,nan) (-0.78750,nan) (-0.78500,nan) (-0.78250,nan) (-0.78000,nan) (-0.77750,nan) (-0.77500,nan) (-0.77250,nan) (-0.77000,nan) (-0.76750,nan) (-0.76500,nan) (-0.76250,nan) (-0.76000,nan) (-0.75750,nan) (-0.75500,nan) (-0.75250,nan) (-0.75000,nan) (-0.74750,nan) (-0.74500,nan) (-0.74250,nan) (-0.74000,nan) (-0.73750,nan) (-0.73500,nan) (-0.73250,nan) (-0.73000,nan) (-0.72750,nan) (-0.72500,nan) (-0.72250,nan) (-0.72000,nan) (-0.71750,nan) (-0.71500,nan) (-0.71250,nan) (-0.71000,nan) (-0.70750,nan) (-0.70500,nan) (-0.70250,nan) (-0.70000,nan) (-0.69750,nan) (-0.69500,nan) (-0.69250,nan) (-0.69000,nan) (-0.68750,nan) (-0.68500,nan) (-0.68250,nan) (-0.68000,nan) (-0.67750,nan) (-0.67500,nan) (-0.67250,nan) (-0.67000,nan) (-0.66750,nan) (-0.66500,nan) (-0.66250,nan) (-0.66000,nan) (-0.65750,nan) (-0.65500,nan) (-0.65250,nan) (-0.65000,nan) (-0.64750,nan) (-0.64500,nan) (-0.64250,nan) (-0.64000,nan) (-0.63750,nan) (-0.63500,nan) (-0.63250,nan) (-0.63000,nan) (-0.62750,nan) (-0.62500,nan) (-0.62250,nan) (-0.62000,nan) (-0.61750,nan) (-0.61500,nan) (-0.61250,nan) (-0.61000,nan) (-0.60750,nan) (-0.60500,nan) (-0.60250,nan) (-0.60000,nan) (-0.59750,nan) (-0.59500,nan) (-0.59250,nan) (-0.59000,nan) (-0.58750,nan) (-0.58500,nan) (-0.58250,nan) (-0.58000,nan) (-0.57750,nan) (-0.57500,nan) (-0.57250,nan) (-0.57000,nan) (-0.56750,nan) (-0.56500,nan) (-0.56250,nan) (-0.56000,nan) (-0.55750,nan) (-0.55500,nan) (-0.55250,nan) (-0.55000,nan) (-0.54750,nan) (-0.54500,nan) (-0.54250,nan) (-0.54000,nan) (-0.53750,nan) (-0.53500,nan) (-0.53250,nan) (-0.53000,nan) (-0.52750,nan) (-0.52500,nan) (-0.52250,nan) (-0.52000,nan) (-0.51750,nan) (-0.51500,nan) (-0.51250,nan) (-0.51000,nan) (-0.50750,nan) (-0.50500,nan) (-0.50250,nan) (-0.50000,nan) (-0.49750,nan) (-0.49500,nan) (-0.49250,nan) (-0.49000,nan) (-0.48750,nan) (-0.48500,nan) (-0.48250,nan) (-0.48000,nan) (-0.47750,nan) (-0.47500,nan) (-0.47250,nan) (-0.47000,nan) (-0.46750,nan) (-0.46500,nan) (-0.46250,nan) (-0.46000,nan) (-0.45750,nan) (-0.45500,nan) (-0.45250,nan) (-0.45000,nan) (-0.44750,nan) (-0.44500,nan) (-0.44250,nan) (-0.44000,nan) (-0.43750,nan) (-0.43500,nan) (-0.43250,nan) (-0.43000,nan) (-0.42750,nan) (-0.42500,nan) (-0.42250,nan) (-0.42000,nan) (-0.41750,nan) (-0.41500,nan) (-0.41250,nan) (-0.41000,nan) (-0.40750,nan) (-0.40500,nan) (-0.40250,nan) (-0.40000,nan) (-0.39750,nan) (-0.39500,nan) (-0.39250,nan) (-0.39000,nan) (-0.38750,nan) (-0.38500,nan) (-0.38250,nan) (-0.38000,nan) (-0.37750,nan) (-0.37500,nan) (-0.37250,nan) (-0.37000,nan) (-0.36750,nan) (-0.36500,nan) (-0.36250,nan) (-0.36000,nan) (-0.35750,nan) (-0.35500,nan) (-0.35250,nan) (-0.35000,nan) (-0.34750,nan) (-0.34500,nan) (-0.34250,nan) (-0.34000,nan) (-0.33750,nan) (-0.33500,nan) (-0.33250,nan) (-0.33000,nan) (-0.32750,nan) (-0.32500,nan) (-0.32250,nan) (-0.32000,nan) (-0.31750,nan) (-0.31500,nan) (-0.31250,nan) (-0.31000,nan) (-0.30750,nan) (-0.30500,nan) (-0.30250,nan) (-0.30000,nan) (-0.29750,nan) (-0.29500,nan) (-0.29250,nan) (-0.29000,nan) (-0.28750,nan) (-0.28500,nan) (-0.28250,nan) (-0.28000,nan) (-0.27750,nan) (-0.27500,nan) (-0.27250,nan) (-0.27000,nan) (-0.26750,nan) (-0.26500,nan) (-0.26250,nan) (-0.26000,nan) (-0.25750,nan) (-0.25500,nan) (-0.25250,nan) (-0.25000,nan) (-0.24750,nan) (-0.24500,nan) (-0.24250,nan) (-0.24000,nan) (-0.23750,nan) (-0.23500,nan) (-0.23250,nan) (-0.23000,nan) (-0.22750,nan) (-0.22500,nan) (-0.22250,nan) (-0.22000,nan) (-0.21750,nan) (-0.21500,nan) (-0.21250,nan) (-0.21000,nan) (-0.20750,nan) (-0.20500,nan) (-0.20250,nan) (-0.20000,nan) (-0.19750,nan) (-0.19500,nan) (-0.19250,nan) (-0.19000,nan) (-0.18750,nan) (-0.18500,nan) (-0.18250,nan) (-0.18000,nan) (-0.17750,nan) (-0.17500,nan) (-0.17250,nan) (-0.17000,nan) (-0.16750,nan) (-0.16500,nan) (-0.16250,nan) (-0.16000,nan) (-0.15750,nan) (-0.15500,nan) (-0.15250,nan) (-0.15000,nan) (-0.14750,nan) (-0.14500,nan) (-0.14250,nan) (-0.14000,nan) (-0.13750,nan) (-0.13500,nan) (-0.13250,nan) (-0.13000,nan) (-0.12750,nan) (-0.12500,nan) (-0.12250,nan) (-0.12000,nan) (-0.11750,nan) (-0.11500,nan) (-0.11250,nan) (-0.11000,nan) (-0.10750,nan) (-0.10500,nan) (-0.10250,nan) (-0.10000,nan) (-0.09750,nan) (-0.09500,nan) (-0.09250,nan) (-0.09000,nan) (-0.08750,nan) (-0.08500,nan) (-0.08250,nan) (-0.08000,nan) (-0.07750,nan) (-0.07500,nan) (-0.07250,nan) (-0.07000,nan) (-0.06750,nan) (-0.06500,nan) (-0.06250,nan) (-0.06000,nan) (-0.05750,nan) (-0.05500,nan) (-0.05250,nan) (-0.05000,nan) (-0.04750,nan) (-0.04500,nan) (-0.04250,nan) (-0.04000,nan) (-0.03750,nan) (-0.03500,nan) (-0.03250,nan) (-0.03000,nan) (-0.02750,nan) (-0.02500,nan) (-0.02250,nan) (-0.02000,nan) (-0.01750,nan) (-0.01500,nan) (-0.01250,nan) (-0.01000,nan) (-0.00750,nan) (-0.00500,nan) (-0.00250,nan) (0.00000,nan) (0.00250,nan) (0.00500,nan) (0.00750,nan) (0.01000,nan) (0.01250,nan) (0.01500,nan) (0.01750,nan) (0.02000,nan) (0.02250,nan) (0.02500,nan) (0.02750,nan) (0.03000,nan) (0.03250,nan) (0.03500,nan) (0.03750,nan) (0.04000,nan) (0.04250,nan) (0.04500,nan) (0.04750,nan) (0.05000,nan) (0.05250,nan) (0.05500,nan) (0.05750,nan) (0.06000,nan) (0.06250,nan) (0.06500,nan) (0.06750,nan) (0.07000,nan) (0.07250,nan) (0.07500,nan) (0.07750,nan) (0.08000,nan) (0.08250,nan) (0.08500,nan) (0.08750,nan) (0.09000,nan) (0.09250,nan) (0.09500,nan) (0.09750,nan) (0.10000,nan) (0.10250,nan) (0.10500,nan) (0.10750,nan) (0.11000,nan) (0.11250,nan) (0.11500,nan) (0.11750,nan) (0.12000,nan) (0.12250,nan) (0.12500,nan) (0.12750,nan) (0.13000,nan) (0.13250,nan) (0.13500,nan) (0.13750,nan) (0.14000,nan) (0.14250,nan) (0.14500,nan) (0.14750,nan) (0.15000,nan) (0.15250,nan) (0.15500,nan) (0.15750,nan) (0.16000,nan) (0.16250,nan) (0.16500,nan) (0.16750,nan) (0.17000,nan) (0.17250,nan) (0.17500,nan) (0.17750,nan) (0.18000,nan) (0.18250,nan) (0.18500,nan) (0.18750,nan) (0.19000,nan) (0.19250,nan) (0.19500,nan) (0.19750,nan) (0.20000,nan) (0.20250,nan) (0.20500,nan) (0.20750,nan) (0.21000,nan) (0.21250,nan) (0.21500,nan) (0.21750,nan) (0.22000,nan) (0.22250,nan) (0.22500,nan) (0.22750,nan) (0.23000,nan) (0.23250,nan) (0.23500,nan) (0.23750,nan) (0.24000,nan) (0.24250,nan) (0.24500,nan) (0.24750,nan) (0.25000,nan) (0.25250,nan) (0.25500,nan) (0.25750,nan) (0.26000,nan) (0.26250,nan) (0.26500,nan) (0.26750,nan) (0.27000,nan) (0.27250,nan) (0.27500,nan) (0.27750,nan) (0.28000,nan) (0.28250,nan) (0.28500,nan) (0.28750,nan) (0.29000,nan) (0.29250,nan) (0.29500,nan) (0.29750,nan) (0.30000,nan) (0.30250,nan) (0.30500,nan) (0.30750,nan) (0.31000,nan) (0.31250,nan) (0.31500,nan) (0.31750,nan) (0.32000,nan) (0.32250,nan) (0.32500,nan) (0.32750,nan) (0.33000,nan) (0.33250,nan) (0.33500,nan) (0.33750,nan) (0.34000,nan) (0.34250,nan) (0.34500,nan) (0.34750,nan) (0.35000,nan) (0.35250,nan) (0.35500,nan) (0.35750,nan) (0.36000,nan) (0.36250,nan) (0.36500,nan) (0.36750,nan) (0.37000,nan) (0.37250,nan) (0.37500,nan) (0.37750,nan) (0.38000,nan) (0.38250,nan) (0.38500,nan) (0.38750,nan) (0.39000,nan) (0.39250,nan) (0.39500,nan) (0.39750,nan) (0.40000,nan) (0.40250,nan) (0.40500,nan) (0.40750,nan) (0.41000,nan) (0.41250,nan) (0.41500,nan) (0.41750,nan) (0.42000,nan) (0.42250,nan) (0.42500,nan) (0.42750,nan) (0.43000,nan) (0.43250,nan) (0.43500,nan) (0.43750,nan) (0.44000,nan) (0.44250,nan) (0.44500,nan) (0.44750,nan) (0.45000,nan) (0.45250,nan) (0.45500,nan) (0.45750,nan) (0.46000,nan) (0.46250,nan) (0.46500,nan) (0.46750,nan) (0.47000,nan) (0.47250,nan) (0.47500,nan) (0.47750,nan) (0.48000,nan) (0.48250,nan) (0.48500,nan) (0.48750,nan) (0.49000,nan) (0.49250,nan) (0.49500,nan) (0.49750,nan) (0.50000,nan) (0.50250,nan) (0.50500,nan) (0.50750,nan) (0.51000,nan) (0.51250,nan) (0.51500,nan) (0.51750,nan) (0.52000,nan) (0.52250,nan) (0.52500,nan) (0.52750,nan) (0.53000,nan) (0.53250,nan) (0.53500,nan) (0.53750,nan) (0.54000,nan) (0.54250,nan) (0.54500,nan) (0.54750,nan) (0.55000,nan) (0.55250,nan) (0.55500,nan) (0.55750,nan) (0.56000,nan) (0.56250,nan) (0.56500,nan) (0.56750,nan) (0.57000,nan) (0.57250,nan) (0.57500,nan) (0.57750,nan) (0.58000,nan) (0.58250,nan) (0.58500,nan) (0.58750,nan) (0.59000,nan) (0.59250,nan) (0.59500,nan) (0.59750,nan) (0.60000,nan) (0.60250,nan) (0.60500,nan) (0.60750,nan) (0.61000,nan) (0.61250,nan) (0.61500,nan) (0.61750,nan) (0.62000,nan) (0.62250,nan) (0.62500,nan) (0.62750,nan) (0.63000,nan) (0.63250,nan) (0.63500,nan) (0.63750,nan) (0.64000,nan) (0.64250,nan) (0.64500,nan) (0.64750,nan) (0.65000,nan) (0.65250,nan) (0.65500,nan) (0.65750,nan) (0.66000,nan) (0.66250,nan) (0.66500,nan) (0.66750,nan) (0.67000,nan) (0.67250,nan) (0.67500,nan) (0.67750,nan) (0.68000,nan) (0.68250,nan) (0.68500,nan) (0.68750,nan) (0.69000,nan) (0.69250,nan) (0.69500,nan) (0.69750,nan) (0.70000,nan) (0.70250,nan) (0.70500,nan) (0.70750,nan) (0.71000,nan) (0.71250,nan) (0.71500,nan) (0.71750,nan) (0.72000,nan) (0.72250,nan) (0.72500,nan) (0.72750,nan) (0.73000,nan) (0.73250,nan) (0.73500,nan) (0.73750,nan) (0.74000,nan) (0.74250,nan) (0.74500,nan) (0.74750,nan) (0.75000,nan) (0.75250,nan) (0.75500,nan) (0.75750,nan) (0.76000,nan) (0.76250,nan) (0.76500,nan) (0.76750,nan) (0.77000,nan) (0.77250,nan) (0.77500,nan) (0.77750,nan) (0.78000,nan) (0.78250,nan) (0.78500,nan) (0.78750,nan) (0.79000,nan) (0.79250,nan) (0.79500,nan) (0.79750,nan) (0.80000,nan) (0.80250,nan) (0.80500,nan) (0.80750,nan) (0.81000,nan) (0.81250,nan) (0.81500,nan) (0.81750,nan) (0.82000,nan) (0.82250,nan) (0.82500,nan) (0.82750,nan) (0.83000,nan) (0.83250,nan) (0.83500,nan) (0.83750,nan) (0.84000,nan) (0.84250,nan) (0.84500,nan) (0.84750,nan) (0.85000,nan) (0.85250,nan) (0.85500,nan) (0.85750,nan) (0.86000,nan) (0.86250,nan) (0.86500,nan) (0.86750,nan) (0.87000,nan) (0.87250,nan) (0.87500,nan) (0.87750,nan) (0.88000,nan) (0.88250,nan) (0.88500,nan) (0.88750,nan) (0.89000,nan) (0.89250,nan) (0.89500,nan) (0.89750,nan) (0.90000,nan) (0.90250,nan) (0.90500,nan) (0.90750,nan) (0.91000,nan) (0.91250,nan) (0.91500,nan) (0.91750,nan) (0.92000,nan) (0.92250,nan) (0.92500,nan) (0.92750,nan) (0.93000,nan) (0.93250,nan) (0.93500,nan) (0.93750,nan) (0.94000,nan) (0.94250,nan) (0.94500,nan) (0.94750,nan) (0.95000,nan) (0.95250,nan) (0.95500,nan) (0.95750,nan) (0.96000,nan) (0.96250,nan) (0.96500,nan) (0.96750,nan) (0.97000,nan) (0.97250,nan) (0.97500,nan) (0.97750,nan) (0.98000,nan) (0.98250,nan) (0.98500,nan) (0.98750,nan) (0.99000,nan) (0.99250,nan) (0.99500,nan) (0.99750,nan) (1.00000,nan) (1.00250,nan) (1.00500,nan) (1.00750,nan) (1.01000,nan) (1.01250,nan) (1.01500,nan) (1.01750,nan) (1.02000,nan) (1.02250,nan) (1.02500,nan) (1.02750,nan) (1.03000,nan) (1.03250,nan) (1.03500,nan) (1.03750,nan) (1.04000,nan) (1.04250,nan) (1.04500,nan) (1.04750,nan) (1.05000,nan) (1.05250,nan) (1.05500,nan) (1.05750,nan) (1.06000,nan) (1.06250,nan) (1.06500,nan) (1.06750,nan) (1.07000,nan) (1.07250,nan) (1.07500,nan) (1.07750,nan) (1.08000,nan) (1.08250,nan) (1.08500,nan) (1.08750,nan) (1.09000,nan) (1.09250,nan) (1.09500,nan) (1.09750,nan) (1.10000,nan) (1.10250,nan) (1.10500,nan) (1.10750,nan) (1.11000,nan) (1.11250,nan) (1.11500,nan) (1.11750,nan) (1.12000,nan) (1.12250,nan) (1.12500,nan) (1.12750,nan) (1.13000,nan) (1.13250,nan) (1.13500,nan) (1.13750,nan) (1.14000,nan) (1.14250,nan) (1.14500,nan) (1.14750,nan) (1.15000,nan) (1.15250,nan) (1.15500,nan) (1.15750,nan) (1.16000,nan) (1.16250,nan) (1.16500,nan) (1.16750,nan) (1.17000,nan) (1.17250,nan) (1.17500,nan) (1.17750,nan) (1.18000,nan) (1.18250,nan) (1.18500,nan) (1.18750,nan) (1.19000,nan) (1.19250,nan) (1.19500,nan) (1.19750,nan) (1.20000,nan) (1.20250,nan) (1.20500,nan) (1.20750,nan) (1.21000,nan) (1.21250,nan) (1.21500,nan) (1.21750,nan) (1.22000,nan) (1.22250,nan) (1.22500,nan) (1.22750,nan) (1.23000,nan) (1.23250,nan) (1.23500,nan) (1.23750,nan) (1.24000,nan) (1.24250,nan) (1.24500,nan) (1.24750,nan) (1.25000,nan) (1.25250,nan) (1.25500,nan) (1.25750,nan) (1.26000,nan) (1.26250,nan) (1.26500,nan) (1.26750,nan) (1.27000,nan) (1.27250,nan) (1.27500,nan) (1.27750,nan) (1.28000,nan) (1.28250,nan) (1.28500,nan) (1.28750,nan) (1.29000,nan) (1.29250,nan) (1.29500,nan) (1.29750,nan) (1.30000,nan) (1.30250,nan) (1.30500,nan) (1.30750,nan) (1.31000,nan) (1.31250,nan) (1.31500,nan) (1.31750,nan) (1.32000,nan) (1.32250,nan) (1.32500,nan) (1.32750,nan) (1.33000,nan) (1.33250,nan) (1.33500,nan) (1.33750,nan) (1.34000,nan) (1.34250,nan) (1.34500,nan) (1.34750,nan) (1.35000,nan) (1.35250,nan) (1.35500,nan) (1.35750,nan) (1.36000,nan) (1.36250,nan) (1.36500,nan) (1.36750,nan) (1.37000,nan) (1.37250,nan) (1.37500,nan) (1.37750,nan) (1.38000,nan) (1.38250,nan) (1.38500,nan) (1.38750,nan) (1.39000,nan) (1.39250,nan) (1.39500,nan) (1.39750,nan) (1.40000,nan) (1.40250,nan) (1.40500,nan) (1.40750,nan) (1.41000,nan) (1.41250,nan) (1.41500,nan) (1.41750,nan) (1.42000,nan) (1.42250,nan) (1.42500,nan) (1.42750,nan) (1.43000,nan) (1.43250,nan) (1.43500,nan) (1.43750,nan) (1.44000,nan) (1.44250,nan) (1.44500,nan) (1.44750,nan) (1.45000,nan) (1.45250,nan) (1.45500,nan) (1.45750,nan) (1.46000,nan) (1.46250,nan) (1.46500,nan) (1.46750,nan) (1.47000,nan) (1.47250,nan) (1.47500,nan) (1.47750,nan) (1.48000,nan) (1.48250,nan) (1.48500,nan) (1.48750,nan) (1.49000,nan) (1.49250,nan) (1.49500,nan) (1.49750,nan) (1.50000,nan) (1.50250,nan) (1.50500,nan) (1.50750,nan) (1.51000,nan) (1.51250,nan) (1.51500,nan) (1.51750,nan) (1.52000,nan) (1.52250,nan) (1.52500,nan) (1.52750,nan) (1.53000,nan) (1.53250,nan) (1.53500,nan) (1.53750,nan) (1.54000,nan) (1.54250,nan) (1.54500,nan) (1.54750,nan) (1.55000,nan) (1.55250,nan) (1.55500,nan) (1.55750,nan) (1.56000,nan) (1.56250,nan) (1.56500,nan) (1.56750,nan) (1.57000,nan) (1.57250,nan) (1.57500,nan) (1.57750,nan) (1.58000,nan) (1.58250,nan) (1.58500,nan) (1.58750,nan) (1.59000,nan) (1.59250,nan) (1.59500,nan) (1.59750,nan) (1.60000,nan) (1.60250,nan) (1.60500,nan) (1.60750,nan) (1.61000,nan) (1.61250,nan) (1.61500,nan) (1.61750,nan) (1.62000,nan) (1.62250,nan) (1.62500,nan) (1.62750,nan) (1.63000,nan) (1.63250,nan) (1.63500,nan) (1.63750,nan) (1.64000,nan) (1.64250,nan) (1.64500,nan) (1.64750,nan) (1.65000,nan) (1.65250,nan) (1.65500,nan) (1.65750,nan) (1.66000,nan) (1.66250,nan) (1.66500,nan) (1.66750,nan) (1.67000,nan) (1.67250,nan) (1.67500,nan) (1.67750,nan) (1.68000,nan) (1.68250,nan) (1.68500,nan) (1.68750,nan) (1.69000,nan) (1.69250,nan) (1.69500,nan) (1.69750,nan) (1.70000,nan) (1.70250,nan) (1.70500,nan) (1.70750,nan) (1.71000,nan) (1.71250,nan) (1.71500,nan) (1.71750,nan) (1.72000,nan) (1.72250,nan) (1.72500,nan) (1.72750,nan) (1.73000,nan) (1.73250,nan) (1.73500,nan) (1.73750,nan) (1.74000,nan) (1.74250,nan) (1.74500,nan) (1.74750,nan) (1.75000,nan) (1.75250,nan) (1.75500,nan) (1.75750,nan) (1.76000,nan) (1.76250,nan) (1.76500,nan) (1.76750,nan) (1.77000,nan) (1.77250,nan) (1.77500,nan) (1.77750,nan) (1.78000,nan) (1.78250,nan) (1.78500,nan) (1.78750,nan) (1.79000,nan) (1.79250,nan) (1.79500,nan) (1.79750,nan) (1.80000,nan) (1.80250,nan) (1.80500,nan) (1.80750,nan) (1.81000,nan) (1.81250,nan) (1.81500,nan) (1.81750,nan) (1.82000,nan) (1.82250,nan) (1.82500,nan) (1.82750,nan) (1.83000,nan) (1.83250,nan) (1.83500,nan) (1.83750,nan) (1.84000,nan) (1.84250,nan) (1.84500,nan) (1.84750,nan) (1.85000,nan) (1.85250,nan) (1.85500,nan) (1.85750,nan) (1.86000,nan) (1.86250,nan) (1.86500,nan) (1.86750,nan) (1.87000,nan) (1.87250,nan) (1.87500,nan) (1.87750,nan) (1.88000,nan) (1.88250,nan) (1.88500,nan) (1.88750,nan) (1.89000,nan) (1.89250,nan) (1.89500,nan) (1.89750,nan) (1.90000,nan) (1.90250,nan) (1.90500,nan) (1.90750,nan) (1.91000,nan) (1.91250,nan) (1.91500,nan) (1.91750,nan) (1.92000,nan) (1.92250,nan) (1.92500,nan) (1.92750,nan) (1.93000,nan) (1.93250,nan) (1.93500,nan) (1.93750,nan) (1.94000,nan) (1.94250,nan) (1.94500,nan) (1.94750,nan) (1.95000,nan) (1.95250,nan) (1.95500,nan) (1.95750,nan) (1.96000,nan) (1.96250,nan) (1.96500,nan) (1.96750,nan) (1.97000,nan) (1.97250,nan) (1.97500,nan) (1.97750,nan) (1.98000,nan) (1.98250,nan) (1.98500,nan) (1.98750,nan) (1.99000,nan) (1.99250,nan) (1.99500,nan) (1.99750,nan) (2.00000,nan) (2.00250,nan) (2.00500,nan) (2.00750,nan) (2.01000,nan) (2.01250,nan) (2.01500,nan) (2.01750,nan) (2.02000,nan) (2.02250,nan) (2.02500,nan) (2.02750,nan) (2.03000,nan) (2.03250,nan) (2.03500,nan) (2.03750,nan) (2.04000,nan) (2.04250,nan) (2.04500,nan) (2.04750,nan) (2.05000,nan) (2.05250,nan) (2.05500,nan) (2.05750,nan) (2.06000,nan) (2.06250,nan) (2.06500,nan) (2.06750,nan) (2.07000,nan) (2.07250,nan) (2.07500,nan) (2.07750,nan) (2.08000,nan) (2.08250,nan) (2.08500,nan) (2.08750,nan) (2.09000,nan) (2.09250,nan) (2.09500,nan) (2.09750,nan) (2.10000,nan) (2.10250,nan) (2.10500,nan) (2.10750,nan) (2.11000,nan) (2.11250,nan) (2.11500,nan) (2.11750,nan) (2.12000,nan) (2.12250,nan) (2.12500,nan) (2.12750,nan) (2.13000,nan) (2.13250,nan) (2.13500,nan) (2.13750,nan) (2.14000,nan) (2.14250,nan) (2.14500,nan) (2.14750,nan) (2.15000,nan) (2.15250,nan) (2.15500,nan) (2.15750,nan) (2.16000,nan) (2.16250,nan) (2.16500,nan) (2.16750,nan) (2.17000,nan) (2.17250,nan) (2.17500,nan) (2.17750,nan) (2.18000,nan) (2.18250,nan) (2.18500,nan) (2.18750,nan) (2.19000,nan) (2.19250,nan) (2.19500,nan) (2.19750,nan) (2.20000,nan) (2.20250,nan) (2.20500,nan) (2.20750,nan) (2.21000,nan) (2.21250,nan) (2.21500,nan) (2.21750,nan) (2.22000,nan) (2.22250,nan) (2.22500,nan) (2.22750,nan) (2.23000,nan) (2.23250,nan) (2.23500,nan) (2.23750,nan) (2.24000,nan) (2.24250,nan) (2.24500,nan) (2.24750,nan) (2.25000,nan) (2.25250,nan) (2.25500,nan) (2.25750,nan) (2.26000,nan) (2.26250,nan) (2.26500,nan) (2.26750,nan) (2.27000,nan) (2.27250,nan) (2.27500,nan) (2.27750,nan) (2.28000,nan) (2.28250,nan) (2.28500,nan) (2.28750,nan) (2.29000,nan) (2.29250,nan) (2.29500,nan) (2.29750,nan) (2.30000,nan) (2.30250,nan) (2.30500,nan) (2.30750,nan) (2.31000,nan) (2.31250,nan) (2.31500,nan) (2.31750,nan) (2.32000,nan) (2.32250,nan) (2.32500,nan) (2.32750,nan) (2.33000,nan) (2.33250,nan) (2.33500,nan) (2.33750,nan) (2.34000,nan) (2.34250,nan) (2.34500,nan) (2.34750,nan) (2.35000,nan) (2.35250,nan) (2.35500,nan) (2.35750,nan) (2.36000,nan) (2.36250,nan) (2.36500,nan) (2.36750,nan) (2.37000,nan) (2.37250,nan) (2.37500,nan) (2.37750,nan) (2.38000,nan) (2.38250,nan) (2.38500,nan) (2.38750,nan) (2.39000,nan) (2.39250,nan) (2.39500,nan) (2.39750,nan) (2.40000,nan) (2.40250,nan) (2.40500,nan) (2.40750,nan) (2.41000,nan) (2.41250,nan) (2.41500,nan) (2.41750,nan) (2.42000,nan) (2.42250,nan) (2.42500,nan) (2.42750,nan) (2.43000,nan) (2.43250,nan) (2.43500,nan) (2.43750,nan) (2.44000,nan) (2.44250,nan) (2.44500,nan) (2.44750,nan) (2.45000,nan) (2.45250,nan) (2.45500,nan) (2.45750,nan) (2.46000,nan) (2.46250,nan) (2.46500,nan) (2.46750,nan) (2.47000,nan) (2.47250,nan) (2.47500,nan) (2.47750,nan) (2.48000,nan) (2.48250,nan) (2.48500,nan) (2.48750,nan) (2.49000,nan) (2.49250,nan) (2.49500,nan) (2.49750,nan) (2.50000,nan) (2.50250,nan) (2.50500,nan) (2.50750,nan) (2.51000,nan) (2.51250,nan) (2.51500,nan) (2.51750,nan) (2.52000,nan) (2.52250,nan) (2.52500,nan) (2.52750,nan) (2.53000,nan) (2.53250,nan) (2.53500,nan) (2.53750,nan) (2.54000,nan) (2.54250,nan) (2.54500,nan) (2.54750,nan) (2.55000,nan) (2.55250,nan) (2.55500,nan) (2.55750,nan) (2.56000,nan) (2.56250,nan) (2.56500,nan) (2.56750,nan) (2.57000,nan) (2.57250,nan) (2.57500,nan) (2.57750,nan) (2.58000,nan) (2.58250,nan) (2.58500,nan) (2.58750,nan) (2.59000,nan) (2.59250,nan) (2.59500,nan) (2.59750,nan) (2.60000,nan) (2.60250,nan) (2.60500,nan) (2.60750,nan) (2.61000,nan) (2.61250,nan) (2.61500,nan) (2.61750,nan) (2.62000,nan) (2.62250,nan) (2.62500,nan) (2.62750,nan) (2.63000,nan) (2.63250,nan) (2.63500,nan) (2.63750,nan) (2.64000,nan) (2.64250,nan) (2.64500,nan) (2.64750,nan) (2.65000,nan) (2.65250,nan) (2.65500,nan) (2.65750,nan) (2.66000,nan) (2.66250,nan) (2.66500,nan) (2.66750,nan) (2.67000,nan) (2.67250,nan) (2.67500,nan) (2.67750,nan) (2.68000,nan) (2.68250,nan) (2.68500,nan) (2.68750,nan) (2.69000,nan) (2.69250,nan) (2.69500,nan) (2.69750,nan) (2.70000,nan) (2.70250,nan) (2.70500,nan) (2.70750,nan) (2.71000,nan) (2.71250,nan) (2.71500,nan) (2.71750,nan) (2.72000,nan) (2.72250,nan) (2.72500,nan) (2.72750,nan) (2.73000,nan) (2.73250,nan) (2.73500,nan) (2.73750,nan) (2.74000,nan) (2.74250,nan) (2.74500,nan) (2.74750,nan) (2.75000,nan) (2.75250,nan) (2.75500,nan) (2.75750,nan) (2.76000,nan) (2.76250,nan) (2.76500,nan) (2.76750,nan) (2.77000,nan) (2.77250,nan) (2.77500,nan) (2.77750,nan) (2.78000,nan) (2.78250,nan) (2.78500,nan) (2.78750,nan) (2.79000,nan) (2.79250,nan) (2.79500,nan) (2.79750,nan) (2.80000,nan) (2.80250,nan) (2.80500,nan) (2.80750,nan) (2.81000,nan) (2.81250,nan) (2.81500,nan) (2.81750,nan) (2.82000,nan) (2.82250,nan) (2.82500,nan) (2.82750,nan) (2.83000,nan) (2.83250,nan) (2.83500,nan) (2.83750,nan) (2.84000,nan) (2.84250,nan) (2.84500,nan) (2.84750,nan) (2.85000,nan) (2.85250,nan) (2.85500,nan) (2.85750,nan) (2.86000,nan) (2.86250,nan) (2.86500,nan) (2.86750,nan) (2.87000,nan) (2.87250,nan) (2.87500,nan) (2.87750,nan) (2.88000,nan) (2.88250,nan) (2.88500,nan) (2.88750,nan) (2.89000,nan) (2.89250,nan) (2.89500,nan) (2.89750,nan) (2.90000,nan) (2.90250,nan) (2.90500,nan) (2.90750,nan) (2.91000,nan) (2.91250,nan) (2.91500,nan) (2.91750,nan) (2.92000,nan) (2.92250,nan) (2.92500,nan) (2.92750,nan) (2.93000,nan) (2.93250,nan) (2.93500,nan) (2.93750,nan) (2.94000,nan) (2.94250,nan) (2.94500,nan) (2.94750,nan) (2.95000,nan) (2.95250,nan) (2.95500,nan) (2.95750,nan) (2.96000,nan) (2.96250,nan) (2.96500,nan) (2.96750,nan) (2.97000,nan) (2.97250,nan) (2.97500,nan) (2.97750,nan) (2.98000,nan) (2.98250,nan) (2.98500,nan) (2.98750,nan) (2.99000,nan) (2.99250,nan) (2.99500,nan) (2.99750,nan) (3.00000,nan) (3.00250,nan) (3.00500,nan) (3.00750,nan) (3.01000,nan) (3.01250,nan) (3.01500,nan) (3.01750,nan) (3.02000,nan) (3.02250,nan) (3.02500,nan) (3.02750,nan) (3.03000,nan) (3.03250,nan) (3.03500,nan) (3.03750,nan) (3.04000,nan) (3.04250,nan) (3.04500,nan) (3.04750,nan) (3.05000,nan) (3.05250,nan) (3.05500,nan) (3.05750,nan) (3.06000,nan) (3.06250,nan) (3.06500,nan) (3.06750,nan) (3.07000,nan) (3.07250,nan) (3.07500,nan) (3.07750,nan) (3.08000,nan) (3.08250,nan) (3.08500,nan) (3.08750,nan) (3.09000,nan) (3.09250,nan) (3.09500,nan) (3.09750,nan) (3.10000,nan) (3.10250,nan) (3.10500,nan) (3.10750,nan) (3.11000,nan) (3.11250,nan) (3.11500,nan) (3.11750,nan) (3.12000,nan) (3.12250,nan) (3.12500,nan) (3.12750,nan) (3.13000,nan) (3.13250,nan) (3.13500,nan) (3.13750,nan) (3.14000,nan) (3.14250,nan) (3.14500,nan) (3.14750,nan) (3.15000,nan) (3.15250,nan) (3.15500,nan) (3.15750,nan) (3.16000,nan) (3.16250,nan) (3.16500,nan) (3.16750,nan) (3.17000,nan) (3.17250,nan) (3.17500,nan) (3.17750,nan) (3.18000,nan) (3.18250,nan) (3.18500,nan) (3.18750,nan) (3.19000,nan) (3.19250,nan) (3.19500,nan) (3.19750,nan) (3.20000,nan) (3.20250,nan) (3.20500,nan) (3.20750,nan) (3.21000,nan) (3.21250,nan) (3.21500,nan) (3.21750,nan) (3.22000,nan) (3.22250,nan) (3.22500,nan) (3.22750,nan) (3.23000,nan) (3.23250,nan) (3.23500,nan) (3.23750,nan) (3.24000,nan) (3.24250,nan) (3.24500,nan) (3.24750,nan) (3.25000,nan) (3.25250,nan) (3.25500,nan) (3.25750,nan) (3.26000,nan) (3.26250,nan) (3.26500,nan) (3.26750,nan) (3.27000,nan) (3.27250,nan) (3.27500,nan) (3.27750,nan) (3.28000,nan) (3.28250,nan) (3.28500,nan) (3.28750,nan) (3.29000,nan) (3.29250,nan) (3.29500,nan) (3.29750,nan) (3.30000,nan) (3.30250,nan) (3.30500,nan) (3.30750,nan) (3.31000,nan) (3.31250,nan) (3.31500,nan) (3.31750,nan) (3.32000,nan) (3.32250,nan) (3.32500,nan) (3.32750,nan) (3.33000,nan) (3.33250,nan) (3.33500,nan) (3.33750,nan) (3.34000,nan) (3.34250,nan) (3.34500,nan) (3.34750,nan) (3.35000,nan) (3.35250,nan) (3.35500,nan) (3.35750,nan) (3.36000,nan) (3.36250,nan) (3.36500,nan) (3.36750,nan) (3.37000,nan) (3.37250,nan) (3.37500,nan) (3.37750,nan) (3.38000,nan) (3.38250,nan) (3.38500,nan) (3.38750,nan) (3.39000,nan) (3.39250,nan) (3.39500,nan) (3.39750,nan) (3.40000,nan) (3.40250,nan) (3.40500,nan) (3.40750,nan) (3.41000,nan) (3.41250,nan) (3.41500,nan) (3.41750,nan) (3.42000,nan) (3.42250,nan) (3.42500,nan) (3.42750,nan) (3.43000,nan) (3.43250,nan) (3.43500,nan) (3.43750,nan) (3.44000,nan) (3.44250,nan) (3.44500,nan) (3.44750,nan) (3.45000,nan) (3.45250,nan) (3.45500,nan) (3.45750,nan) (3.46000,nan) (3.46250,nan) (3.46500,nan) (3.46750,nan) (3.47000,nan) (3.47250,nan) (3.47500,nan) (3.47750,nan) (3.48000,nan) (3.48250,nan) (3.48500,nan) (3.48750,nan) (3.49000,nan) (3.49250,nan) (3.49500,nan) (3.49750,nan) (3.50000,nan) (3.50250,nan) (3.50500,nan) (3.50750,nan) (3.51000,nan) (3.51250,nan) (3.51500,nan) (3.51750,nan) (3.52000,nan) (3.52250,nan) (3.52500,nan) (3.52750,nan) (3.53000,nan) (3.53250,nan) (3.53500,nan) (3.53750,nan) (3.54000,nan) (3.54250,nan) (3.54500,nan) (3.54750,nan) (3.55000,nan) (3.55250,nan) (3.55500,nan) (3.55750,nan) (3.56000,nan) (3.56250,nan) (3.56500,nan) (3.56750,nan) (3.57000,nan) (3.57250,nan) (3.57500,nan) (3.57750,nan) (3.58000,nan) (3.58250,nan) (3.58500,nan) (3.58750,nan) (3.59000,nan) (3.59250,nan) (3.59500,nan) (3.59750,nan) (3.60000,nan) (3.60250,nan) (3.60500,nan) (3.60750,nan) (3.61000,nan) (3.61250,nan) (3.61500,nan) (3.61750,nan) (3.62000,nan) (3.62250,nan) (3.62500,nan) (3.62750,nan) (3.63000,nan) (3.63250,nan) (3.63500,nan) (3.63750,nan) (3.64000,nan) (3.64250,nan) (3.64500,nan) (3.64750,nan) (3.65000,nan) (3.65250,nan) (3.65500,nan) (3.65750,nan) (3.66000,nan) (3.66250,nan) (3.66500,nan) (3.66750,nan) (3.67000,nan) (3.67250,nan) (3.67500,nan) (3.67750,nan) (3.68000,nan) (3.68250,-1.9946393) (3.68500,-1.6684671) (3.68750,-1.4860300) (3.69000,-1.3592244) (3.69250,-1.2622037) (3.69500,-1.1837787) (3.69750,-1.1180784) (3.70000,-1.0616361) (3.70250,-1.0122330) (3.70500,-0.9683631) (3.70750,-0.9289570) (3.71000,-0.8932278) (3.71250,-0.8605799) (3.71500,-0.8305515) (3.71750,-0.8027772) (3.72000,-0.7769630) (3.72250,-0.7528687) (3.72500,-0.7302959) (3.72750,-0.7090782) (3.73000,-0.6890754) (3.73250,-0.6701676) (3.73500,-0.6522520) (3.73750,-0.6352394) (3.74000,-0.6190521) (3.74250,-0.6036221) (3.74500,-0.5888892) (3.74750,-0.5748002) (3.75000,-0.5613078) (3.75250,-0.5483695) (3.75500,-0.5359474) (3.75750,-0.5240071) (3.76000,-0.5125175) (3.76250,-0.5014506) (3.76500,-0.4907808) (3.76750,-0.4804846) (3.77000,-0.4705407) (3.77250,-0.4609295) (3.77500,-0.4516329) (3.77750,-0.4426342) (3.78000,-0.4339181) (3.78250,-0.4254703) (3.78500,-0.4172778) (3.78750,-0.4093281) (3.79000,-0.4016099) (3.79250,-0.3941125) (3.79500,-0.3868260) (3.79750,-0.3797412) (3.80000,-0.3728492) (3.80250,-0.3661419) (3.80500,-0.3596117) (3.80750,-0.3532514) (3.81000,-0.3470541) (3.81250,-0.3410135) (3.81500,-0.3351234) (3.81750,-0.3293783) (3.82000,-0.3237726) (3.82250,-0.3183014) (3.82500,-0.3129597) (3.82750,-0.3077429) (3.83000,-0.3026467) (3.83250,-0.2976669) (3.83500,-0.2927997) (3.83750,-0.2880412) (3.84000,-0.2833878) (3.84250,-0.2788362) (3.84500,-0.2743832) (3.84750,-0.2700255) (3.85000,-0.2657602) (3.85250,-0.2615846) (3.85500,-0.2574958) (3.85750,-0.2534913) (3.86000,-0.2495687) (3.86250,-0.2457254) (3.86500,-0.2419593) (3.86750,-0.2382681) (3.87000,-0.2346498) (3.87250,-0.2311022) (3.87500,-0.2276235) (3.87750,-0.2242119) (3.88000,-0.2208654) (3.88250,-0.2175823) (3.88500,-0.2143611) (3.88750,-0.2112000) (3.89000,-0.2080976) (3.89250,-0.2050523) (3.89500,-0.2020627) (3.89750,-0.1991275) (3.90000,-0.1962452) (3.90250,-0.1934145) (3.90500,-0.1906343) (3.90750,-0.1879033) (3.91000,-0.1852204) (3.91250,-0.1825843) (3.91500,-0.1799941) (3.91750,-0.1774486) (3.92000,-0.1749468) (3.92250,-0.1724877) (3.92500,-0.1700704) (3.92750,-0.1676939) (3.93000,-0.1653574) (3.93250,-0.1630598) (3.93500,-0.1608005) (3.93750,-0.1585785) (3.94000,-0.1563931) (3.94250,-0.1542434) (3.94500,-0.1521287) (3.94750,-0.1500484) (3.95000,-0.1480016) (3.95250,-0.1459876) (3.95500,-0.1440059) (3.95750,-0.1420556) (3.96000,-0.1401363) (3.96250,-0.1382473) (3.96500,-0.1363879) (3.96750,-0.1345576) (3.97000,-0.1327557) (3.97250,-0.1309819) (3.97500,-0.1292354) (3.97750,-0.1275158) (3.98000,-0.1258225) (3.98250,-0.1241551) (3.98500,-0.1225131) (3.98750,-0.1208959) (3.99000,-0.1193031) (3.99250,-0.1177343) (3.99500,-0.1161890) (3.99750,-0.1146668) (4.00000,-0.1131672)};\addlegendentry{O}
\addplot[thick,dotted,blue] coordinates {(-2.00000,nan) (-1.99750,nan) (-1.99500,nan) (-1.99250,nan) (-1.99000,nan) (-1.98750,nan) (-1.98500,nan) (-1.98250,nan) (-1.98000,nan) (-1.97750,nan) (-1.97500,nan) (-1.97250,nan) (-1.97000,nan) (-1.96750,nan) (-1.96500,nan) (-1.96250,nan) (-1.96000,nan) (-1.95750,nan) (-1.95500,nan) (-1.95250,nan) (-1.95000,nan) (-1.94750,nan) (-1.94500,nan) (-1.94250,nan) (-1.94000,nan) (-1.93750,nan) (-1.93500,nan) (-1.93250,nan) (-1.93000,nan) (-1.92750,nan) (-1.92500,nan) (-1.92250,nan) (-1.92000,nan) (-1.91750,nan) (-1.91500,nan) (-1.91250,nan) (-1.91000,nan) (-1.90750,nan) (-1.90500,nan) (-1.90250,nan) (-1.90000,nan) (-1.89750,nan) (-1.89500,nan) (-1.89250,nan) (-1.89000,nan) (-1.88750,nan) (-1.88500,nan) (-1.88250,nan) (-1.88000,nan) (-1.87750,nan) (-1.87500,nan) (-1.87250,nan) (-1.87000,nan) (-1.86750,nan) (-1.86500,nan) (-1.86250,nan) (-1.86000,nan) (-1.85750,nan) (-1.85500,nan) (-1.85250,nan) (-1.85000,nan) (-1.84750,nan) (-1.84500,nan) (-1.84250,nan) (-1.84000,nan) (-1.83750,nan) (-1.83500,nan) (-1.83250,nan) (-1.83000,nan) (-1.82750,nan) (-1.82500,nan) (-1.82250,nan) (-1.82000,nan) (-1.81750,nan) (-1.81500,nan) (-1.81250,nan) (-1.81000,nan) (-1.80750,nan) (-1.80500,nan) (-1.80250,nan) (-1.80000,nan) (-1.79750,nan) (-1.79500,nan) (-1.79250,nan) (-1.79000,nan) (-1.78750,nan) (-1.78500,nan) (-1.78250,nan) (-1.78000,nan) (-1.77750,nan) (-1.77500,nan) (-1.77250,nan) (-1.77000,nan) (-1.76750,nan) (-1.76500,nan) (-1.76250,nan) (-1.76000,nan) (-1.75750,nan) (-1.75500,nan) (-1.75250,nan) (-1.75000,nan) (-1.74750,nan) (-1.74500,nan) (-1.74250,nan) (-1.74000,nan) (-1.73750,nan) (-1.73500,nan) (-1.73250,nan) (-1.73000,nan) (-1.72750,nan) (-1.72500,nan) (-1.72250,nan) (-1.72000,nan) (-1.71750,nan) (-1.71500,nan) (-1.71250,nan) (-1.71000,nan) (-1.70750,nan) (-1.70500,nan) (-1.70250,nan) (-1.70000,nan) (-1.69750,nan) (-1.69500,nan) (-1.69250,nan) (-1.69000,nan) (-1.68750,nan) (-1.68500,nan) (-1.68250,nan) (-1.68000,nan) (-1.67750,nan) (-1.67500,nan) (-1.67250,nan) (-1.67000,nan) (-1.66750,nan) (-1.66500,nan) (-1.66250,nan) (-1.66000,nan) (-1.65750,nan) (-1.65500,nan) (-1.65250,nan) (-1.65000,nan) (-1.64750,nan) (-1.64500,nan) (-1.64250,nan) (-1.64000,nan) (-1.63750,nan) (-1.63500,nan) (-1.63250,nan) (-1.63000,nan) (-1.62750,nan) (-1.62500,nan) (-1.62250,nan) (-1.62000,nan) (-1.61750,nan) (-1.61500,nan) (-1.61250,nan) (-1.61000,nan) (-1.60750,nan) (-1.60500,nan) (-1.60250,nan) (-1.60000,nan) (-1.59750,nan) (-1.59500,nan) (-1.59250,nan) (-1.59000,nan) (-1.58750,nan) (-1.58500,nan) (-1.58250,nan) (-1.58000,nan) (-1.57750,nan) (-1.57500,nan) (-1.57250,nan) (-1.57000,nan) (-1.56750,nan) (-1.56500,nan) (-1.56250,nan) (-1.56000,nan) (-1.55750,nan) (-1.55500,nan) (-1.55250,nan) (-1.55000,nan) (-1.54750,nan) (-1.54500,nan) (-1.54250,nan) (-1.54000,nan) (-1.53750,nan) (-1.53500,nan) (-1.53250,nan) (-1.53000,nan) (-1.52750,nan) (-1.52500,nan) (-1.52250,nan) (-1.52000,nan) (-1.51750,nan) (-1.51500,nan) (-1.51250,nan) (-1.51000,nan) (-1.50750,nan) (-1.50500,nan) (-1.50250,nan) (-1.50000,nan) (-1.49750,nan) (-1.49500,nan) (-1.49250,nan) (-1.49000,nan) (-1.48750,nan) (-1.48500,nan) (-1.48250,nan) (-1.48000,nan) (-1.47750,nan) (-1.47500,nan) (-1.47250,nan) (-1.47000,nan) (-1.46750,nan) (-1.46500,nan) (-1.46250,nan) (-1.46000,nan) (-1.45750,nan) (-1.45500,nan) (-1.45250,nan) (-1.45000,nan) (-1.44750,nan) (-1.44500,nan) (-1.44250,nan) (-1.44000,nan) (-1.43750,nan) (-1.43500,nan) (-1.43250,nan) (-1.43000,nan) (-1.42750,nan) (-1.42500,nan) (-1.42250,nan) (-1.42000,nan) (-1.41750,nan) (-1.41500,nan) (-1.41250,nan) (-1.41000,nan) (-1.40750,nan) (-1.40500,nan) (-1.40250,nan) (-1.40000,nan) (-1.39750,nan) (-1.39500,nan) (-1.39250,nan) (-1.39000,nan) (-1.38750,nan) (-1.38500,nan) (-1.38250,nan) (-1.38000,nan) (-1.37750,nan) (-1.37500,nan) (-1.37250,nan) (-1.37000,nan) (-1.36750,nan) (-1.36500,nan) (-1.36250,nan) (-1.36000,nan) (-1.35750,nan) (-1.35500,nan) (-1.35250,nan) (-1.35000,nan) (-1.34750,nan) (-1.34500,nan) (-1.34250,nan) (-1.34000,nan) (-1.33750,nan) (-1.33500,nan) (-1.33250,nan) (-1.33000,nan) (-1.32750,nan) (-1.32500,nan) (-1.32250,nan) (-1.32000,nan) (-1.31750,nan) (-1.31500,nan) (-1.31250,nan) (-1.31000,nan) (-1.30750,nan) (-1.30500,nan) (-1.30250,nan) (-1.30000,nan) (-1.29750,nan) (-1.29500,nan) (-1.29250,nan) (-1.29000,nan) (-1.28750,nan) (-1.28500,nan) (-1.28250,nan) (-1.28000,nan) (-1.27750,nan) (-1.27500,nan) (-1.27250,nan) (-1.27000,nan) (-1.26750,nan) (-1.26500,nan) (-1.26250,nan) (-1.26000,nan) (-1.25750,nan) (-1.25500,nan) (-1.25250,nan) (-1.25000,nan) (-1.24750,nan) (-1.24500,nan) (-1.24250,nan) (-1.24000,nan) (-1.23750,nan) (-1.23500,nan) (-1.23250,nan) (-1.23000,nan) (-1.22750,nan) (-1.22500,nan) (-1.22250,nan) (-1.22000,nan) (-1.21750,nan) (-1.21500,nan) (-1.21250,nan) (-1.21000,nan) (-1.20750,nan) (-1.20500,nan) (-1.20250,nan) (-1.20000,nan) (-1.19750,nan) (-1.19500,nan) (-1.19250,nan) (-1.19000,nan) (-1.18750,nan) (-1.18500,nan) (-1.18250,nan) (-1.18000,nan) (-1.17750,nan) (-1.17500,nan) (-1.17250,nan) (-1.17000,nan) (-1.16750,nan) (-1.16500,nan) (-1.16250,nan) (-1.16000,nan) (-1.15750,nan) (-1.15500,nan) (-1.15250,nan) (-1.15000,nan) (-1.14750,nan) (-1.14500,nan) (-1.14250,nan) (-1.14000,nan) (-1.13750,nan) (-1.13500,nan) (-1.13250,nan) (-1.13000,nan) (-1.12750,nan) (-1.12500,nan) (-1.12250,nan) (-1.12000,nan) (-1.11750,nan) (-1.11500,nan) (-1.11250,nan) (-1.11000,nan) (-1.10750,nan) (-1.10500,nan) (-1.10250,nan) (-1.10000,nan) (-1.09750,nan) (-1.09500,nan) (-1.09250,nan) (-1.09000,nan) (-1.08750,nan) (-1.08500,nan) (-1.08250,nan) (-1.08000,nan) (-1.07750,nan) (-1.07500,nan) (-1.07250,nan) (-1.07000,nan) (-1.06750,nan) (-1.06500,nan) (-1.06250,nan) (-1.06000,nan) (-1.05750,nan) (-1.05500,nan) (-1.05250,nan) (-1.05000,nan) (-1.04750,nan) (-1.04500,nan) (-1.04250,nan) (-1.04000,nan) (-1.03750,nan) (-1.03500,nan) (-1.03250,nan) (-1.03000,nan) (-1.02750,nan) (-1.02500,nan) (-1.02250,nan) (-1.02000,nan) (-1.01750,nan) (-1.01500,nan) (-1.01250,nan) (-1.01000,nan) (-1.00750,nan) (-1.00500,nan) (-1.00250,nan) (-1.00000,nan) (-0.99750,nan) (-0.99500,nan) (-0.99250,nan) (-0.99000,nan) (-0.98750,nan) (-0.98500,nan) (-0.98250,nan) (-0.98000,nan) (-0.97750,nan) (-0.97500,nan) (-0.97250,nan) (-0.97000,nan) (-0.96750,nan) (-0.96500,nan) (-0.96250,nan) (-0.96000,nan) (-0.95750,nan) (-0.95500,nan) (-0.95250,nan) (-0.95000,nan) (-0.94750,nan) (-0.94500,nan) (-0.94250,nan) (-0.94000,nan) (-0.93750,nan) (-0.93500,nan) (-0.93250,nan) (-0.93000,nan) (-0.92750,nan) (-0.92500,nan) (-0.92250,nan) (-0.92000,nan) (-0.91750,nan) (-0.91500,nan) (-0.91250,nan) (-0.91000,nan) (-0.90750,nan) (-0.90500,nan) (-0.90250,nan) (-0.90000,nan) (-0.89750,nan) (-0.89500,nan) (-0.89250,nan) (-0.89000,nan) (-0.88750,nan) (-0.88500,nan) (-0.88250,nan) (-0.88000,nan) (-0.87750,nan) (-0.87500,nan) (-0.87250,nan) (-0.87000,nan) (-0.86750,nan) (-0.86500,nan) (-0.86250,nan) (-0.86000,nan) (-0.85750,nan) (-0.85500,nan) (-0.85250,nan) (-0.85000,nan) (-0.84750,nan) (-0.84500,nan) (-0.84250,nan) (-0.84000,nan) (-0.83750,nan) (-0.83500,nan) (-0.83250,nan) (-0.83000,nan) (-0.82750,nan) (-0.82500,nan) (-0.82250,nan) (-0.82000,nan) (-0.81750,nan) (-0.81500,nan) (-0.81250,nan) (-0.81000,nan) (-0.80750,nan) (-0.80500,nan) (-0.80250,nan) (-0.80000,nan) (-0.79750,nan) (-0.79500,nan) (-0.79250,nan) (-0.79000,nan) (-0.78750,nan) (-0.78500,nan) (-0.78250,nan) (-0.78000,nan) (-0.77750,nan) (-0.77500,nan) (-0.77250,nan) (-0.77000,nan) (-0.76750,nan) (-0.76500,nan) (-0.76250,nan) (-0.76000,nan) (-0.75750,nan) (-0.75500,nan) (-0.75250,nan) (-0.75000,nan) (-0.74750,nan) (-0.74500,nan) (-0.74250,nan) (-0.74000,nan) (-0.73750,nan) (-0.73500,nan) (-0.73250,nan) (-0.73000,nan) (-0.72750,nan) (-0.72500,nan) (-0.72250,nan) (-0.72000,nan) (-0.71750,nan) (-0.71500,nan) (-0.71250,nan) (-0.71000,nan) (-0.70750,nan) (-0.70500,nan) (-0.70250,nan) (-0.70000,nan) (-0.69750,nan) (-0.69500,nan) (-0.69250,nan) (-0.69000,nan) (-0.68750,nan) (-0.68500,nan) (-0.68250,nan) (-0.68000,nan) (-0.67750,nan) (-0.67500,nan) (-0.67250,nan) (-0.67000,nan) (-0.66750,nan) (-0.66500,nan) (-0.66250,nan) (-0.66000,nan) (-0.65750,nan) (-0.65500,nan) (-0.65250,nan) (-0.65000,nan) (-0.64750,nan) (-0.64500,nan) (-0.64250,nan) (-0.64000,nan) (-0.63750,nan) (-0.63500,nan) (-0.63250,nan) (-0.63000,nan) (-0.62750,nan) (-0.62500,nan) (-0.62250,nan) (-0.62000,nan) (-0.61750,nan) (-0.61500,nan) (-0.61250,nan) (-0.61000,nan) (-0.60750,nan) (-0.60500,nan) (-0.60250,nan) (-0.60000,nan) (-0.59750,nan) (-0.59500,nan) (-0.59250,nan) (-0.59000,nan) (-0.58750,nan) (-0.58500,nan) (-0.58250,nan) (-0.58000,nan) (-0.57750,nan) (-0.57500,nan) (-0.57250,nan) (-0.57000,nan) (-0.56750,nan) (-0.56500,nan) (-0.56250,nan) (-0.56000,nan) (-0.55750,nan) (-0.55500,nan) (-0.55250,nan) (-0.55000,nan) (-0.54750,nan) (-0.54500,nan) (-0.54250,nan) (-0.54000,nan) (-0.53750,nan) (-0.53500,nan) (-0.53250,nan) (-0.53000,nan) (-0.52750,nan) (-0.52500,nan) (-0.52250,nan) (-0.52000,nan) (-0.51750,nan) (-0.51500,nan) (-0.51250,nan) (-0.51000,nan) (-0.50750,nan) (-0.50500,nan) (-0.50250,nan) (-0.50000,nan) (-0.49750,nan) (-0.49500,nan) (-0.49250,nan) (-0.49000,nan) (-0.48750,nan) (-0.48500,nan) (-0.48250,nan) (-0.48000,nan) (-0.47750,nan) (-0.47500,nan) (-0.47250,nan) (-0.47000,nan) (-0.46750,nan) (-0.46500,nan) (-0.46250,nan) (-0.46000,nan) (-0.45750,nan) (-0.45500,nan) (-0.45250,nan) (-0.45000,nan) (-0.44750,nan) (-0.44500,nan) (-0.44250,nan) (-0.44000,nan) (-0.43750,nan) (-0.43500,nan) (-0.43250,nan) (-0.43000,nan) (-0.42750,nan) (-0.42500,nan) (-0.42250,nan) (-0.42000,nan) (-0.41750,nan) (-0.41500,nan) (-0.41250,nan) (-0.41000,nan) (-0.40750,nan) (-0.40500,nan) (-0.40250,nan) (-0.40000,nan) (-0.39750,nan) (-0.39500,nan) (-0.39250,nan) (-0.39000,nan) (-0.38750,nan) (-0.38500,nan) (-0.38250,nan) (-0.38000,nan) (-0.37750,nan) (-0.37500,nan) (-0.37250,nan) (-0.37000,nan) (-0.36750,nan) (-0.36500,nan) (-0.36250,nan) (-0.36000,nan) (-0.35750,nan) (-0.35500,nan) (-0.35250,nan) (-0.35000,nan) (-0.34750,nan) (-0.34500,nan) (-0.34250,nan) (-0.34000,nan) (-0.33750,nan) (-0.33500,nan) (-0.33250,nan) (-0.33000,nan) (-0.32750,nan) (-0.32500,nan) (-0.32250,nan) (-0.32000,nan) (-0.31750,nan) (-0.31500,nan) (-0.31250,nan) (-0.31000,nan) (-0.30750,nan) (-0.30500,nan) (-0.30250,nan) (-0.30000,nan) (-0.29750,nan) (-0.29500,nan) (-0.29250,nan) (-0.29000,nan) (-0.28750,nan) (-0.28500,nan) (-0.28250,nan) (-0.28000,nan) (-0.27750,nan) (-0.27500,nan) (-0.27250,nan) (-0.27000,nan) (-0.26750,nan) (-0.26500,nan) (-0.26250,nan) (-0.26000,nan) (-0.25750,nan) (-0.25500,nan) (-0.25250,nan) (-0.25000,nan) (-0.24750,nan) (-0.24500,nan) (-0.24250,nan) (-0.24000,nan) (-0.23750,nan) (-0.23500,nan) (-0.23250,nan) (-0.23000,nan) (-0.22750,nan) (-0.22500,nan) (-0.22250,nan) (-0.22000,nan) (-0.21750,nan) (-0.21500,nan) (-0.21250,nan) (-0.21000,nan) (-0.20750,nan) (-0.20500,nan) (-0.20250,nan) (-0.20000,nan) (-0.19750,nan) (-0.19500,nan) (-0.19250,nan) (-0.19000,nan) (-0.18750,nan) (-0.18500,nan) (-0.18250,nan) (-0.18000,nan) (-0.17750,nan) (-0.17500,nan) (-0.17250,nan) (-0.17000,nan) (-0.16750,nan) (-0.16500,nan) (-0.16250,nan) (-0.16000,nan) (-0.15750,nan) (-0.15500,nan) (-0.15250,nan) (-0.15000,nan) (-0.14750,nan) (-0.14500,nan) (-0.14250,nan) (-0.14000,nan) (-0.13750,nan) (-0.13500,nan) (-0.13250,nan) (-0.13000,nan) (-0.12750,nan) (-0.12500,nan) (-0.12250,nan) (-0.12000,nan) (-0.11750,nan) (-0.11500,nan) (-0.11250,nan) (-0.11000,nan) (-0.10750,nan) (-0.10500,nan) (-0.10250,nan) (-0.10000,nan) (-0.09750,nan) (-0.09500,nan) (-0.09250,nan) (-0.09000,nan) (-0.08750,nan) (-0.08500,nan) (-0.08250,nan) (-0.08000,nan) (-0.07750,nan) (-0.07500,nan) (-0.07250,nan) (-0.07000,nan) (-0.06750,nan) (-0.06500,nan) (-0.06250,nan) (-0.06000,nan) (-0.05750,nan) (-0.05500,nan) (-0.05250,nan) (-0.05000,nan) (-0.04750,nan) (-0.04500,nan) (-0.04250,nan) (-0.04000,nan) (-0.03750,nan) (-0.03500,nan) (-0.03250,nan) (-0.03000,nan) (-0.02750,nan) (-0.02500,nan) (-0.02250,nan) (-0.02000,nan) (-0.01750,nan) (-0.01500,nan) (-0.01250,nan) (-0.01000,nan) (-0.00750,nan) (-0.00500,nan) (-0.00250,nan) (0.00000,nan) (0.00250,nan) (0.00500,nan) (0.00750,nan) (0.01000,nan) (0.01250,nan) (0.01500,nan) (0.01750,nan) (0.02000,nan) (0.02250,nan) (0.02500,nan) (0.02750,nan) (0.03000,nan) (0.03250,nan) (0.03500,nan) (0.03750,nan) (0.04000,nan) (0.04250,nan) (0.04500,nan) (0.04750,nan) (0.05000,nan) (0.05250,nan) (0.05500,nan) (0.05750,nan) (0.06000,nan) (0.06250,nan) (0.06500,nan) (0.06750,nan) (0.07000,nan) (0.07250,nan) (0.07500,nan) (0.07750,nan) (0.08000,nan) (0.08250,nan) (0.08500,nan) (0.08750,nan) (0.09000,nan) (0.09250,nan) (0.09500,nan) (0.09750,nan) (0.10000,nan) (0.10250,nan) (0.10500,nan) (0.10750,nan) (0.11000,nan) (0.11250,nan) (0.11500,nan) (0.11750,nan) (0.12000,nan) (0.12250,nan) (0.12500,nan) (0.12750,nan) (0.13000,nan) (0.13250,nan) (0.13500,nan) (0.13750,nan) (0.14000,nan) (0.14250,nan) (0.14500,nan) (0.14750,nan) (0.15000,nan) (0.15250,nan) (0.15500,nan) (0.15750,nan) (0.16000,nan) (0.16250,nan) (0.16500,nan) (0.16750,nan) (0.17000,nan) (0.17250,nan) (0.17500,nan) (0.17750,nan) (0.18000,nan) (0.18250,nan) (0.18500,nan) (0.18750,nan) (0.19000,nan) (0.19250,nan) (0.19500,nan) (0.19750,nan) (0.20000,nan) (0.20250,nan) (0.20500,nan) (0.20750,nan) (0.21000,nan) (0.21250,nan) (0.21500,nan) (0.21750,nan) (0.22000,nan) (0.22250,nan) (0.22500,nan) (0.22750,nan) (0.23000,nan) (0.23250,nan) (0.23500,nan) (0.23750,nan) (0.24000,nan) (0.24250,nan) (0.24500,nan) (0.24750,nan) (0.25000,nan) (0.25250,nan) (0.25500,nan) (0.25750,nan) (0.26000,nan) (0.26250,nan) (0.26500,nan) (0.26750,nan) (0.27000,nan) (0.27250,nan) (0.27500,nan) (0.27750,nan) (0.28000,nan) (0.28250,nan) (0.28500,nan) (0.28750,nan) (0.29000,nan) (0.29250,nan) (0.29500,nan) (0.29750,nan) (0.30000,nan) (0.30250,nan) (0.30500,nan) (0.30750,nan) (0.31000,nan) (0.31250,nan) (0.31500,nan) (0.31750,nan) (0.32000,nan) (0.32250,nan) (0.32500,nan) (0.32750,nan) (0.33000,nan) (0.33250,nan) (0.33500,nan) (0.33750,nan) (0.34000,nan) (0.34250,nan) (0.34500,nan) (0.34750,nan) (0.35000,nan) (0.35250,nan) (0.35500,nan) (0.35750,nan) (0.36000,nan) (0.36250,nan) (0.36500,nan) (0.36750,nan) (0.37000,nan) (0.37250,nan) (0.37500,nan) (0.37750,nan) (0.38000,nan) (0.38250,nan) (0.38500,nan) (0.38750,nan) (0.39000,nan) (0.39250,nan) (0.39500,nan) (0.39750,nan) (0.40000,nan) (0.40250,nan) (0.40500,nan) (0.40750,nan) (0.41000,nan) (0.41250,nan) (0.41500,nan) (0.41750,nan) (0.42000,nan) (0.42250,nan) (0.42500,nan) (0.42750,nan) (0.43000,nan) (0.43250,nan) (0.43500,nan) (0.43750,nan) (0.44000,nan) (0.44250,nan) (0.44500,nan) (0.44750,nan) (0.45000,nan) (0.45250,nan) (0.45500,nan) (0.45750,nan) (0.46000,nan) (0.46250,nan) (0.46500,nan) (0.46750,nan) (0.47000,nan) (0.47250,nan) (0.47500,nan) (0.47750,nan) (0.48000,nan) (0.48250,nan) (0.48500,nan) (0.48750,nan) (0.49000,nan) (0.49250,nan) (0.49500,nan) (0.49750,nan) (0.50000,nan) (0.50250,nan) (0.50500,nan) (0.50750,nan) (0.51000,nan) (0.51250,nan) (0.51500,nan) (0.51750,nan) (0.52000,nan) (0.52250,nan) (0.52500,nan) (0.52750,nan) (0.53000,nan) (0.53250,nan) (0.53500,nan) (0.53750,nan) (0.54000,nan) (0.54250,nan) (0.54500,nan) (0.54750,nan) (0.55000,nan) (0.55250,nan) (0.55500,nan) (0.55750,nan) (0.56000,nan) (0.56250,nan) (0.56500,nan) (0.56750,nan) (0.57000,nan) (0.57250,nan) (0.57500,nan) (0.57750,nan) (0.58000,nan) (0.58250,nan) (0.58500,nan) (0.58750,nan) (0.59000,nan) (0.59250,nan) (0.59500,nan) (0.59750,nan) (0.60000,nan) (0.60250,nan) (0.60500,nan) (0.60750,nan) (0.61000,nan) (0.61250,nan) (0.61500,nan) (0.61750,nan) (0.62000,nan) (0.62250,nan) (0.62500,nan) (0.62750,nan) (0.63000,nan) (0.63250,nan) (0.63500,nan) (0.63750,nan) (0.64000,nan) (0.64250,nan) (0.64500,nan) (0.64750,nan) (0.65000,nan) (0.65250,nan) (0.65500,nan) (0.65750,nan) (0.66000,nan) (0.66250,nan) (0.66500,nan) (0.66750,nan) (0.67000,nan) (0.67250,nan) (0.67500,nan) (0.67750,nan) (0.68000,nan) (0.68250,nan) (0.68500,nan) (0.68750,nan) (0.69000,nan) (0.69250,nan) (0.69500,nan) (0.69750,nan) (0.70000,nan) (0.70250,nan) (0.70500,nan) (0.70750,nan) (0.71000,nan) (0.71250,nan) (0.71500,nan) (0.71750,nan) (0.72000,nan) (0.72250,nan) (0.72500,nan) (0.72750,nan) (0.73000,nan) (0.73250,nan) (0.73500,nan) (0.73750,nan) (0.74000,nan) (0.74250,nan) (0.74500,nan) (0.74750,nan) (0.75000,nan) (0.75250,nan) (0.75500,nan) (0.75750,nan) (0.76000,nan) (0.76250,nan) (0.76500,nan) (0.76750,nan) (0.77000,nan) (0.77250,nan) (0.77500,nan) (0.77750,nan) (0.78000,nan) (0.78250,nan) (0.78500,nan) (0.78750,nan) (0.79000,nan) (0.79250,nan) (0.79500,nan) (0.79750,nan) (0.80000,nan) (0.80250,nan) (0.80500,nan) (0.80750,nan) (0.81000,nan) (0.81250,nan) (0.81500,nan) (0.81750,nan) (0.82000,nan) (0.82250,nan) (0.82500,nan) (0.82750,nan) (0.83000,nan) (0.83250,nan) (0.83500,nan) (0.83750,nan) (0.84000,nan) (0.84250,nan) (0.84500,nan) (0.84750,nan) (0.85000,nan) (0.85250,nan) (0.85500,nan) (0.85750,nan) (0.86000,nan) (0.86250,nan) (0.86500,nan) (0.86750,nan) (0.87000,nan) (0.87250,nan) (0.87500,nan) (0.87750,nan) (0.88000,nan) (0.88250,nan) (0.88500,nan) (0.88750,nan) (0.89000,nan) (0.89250,nan) (0.89500,nan) (0.89750,nan) (0.90000,nan) (0.90250,nan) (0.90500,nan) (0.90750,nan) (0.91000,nan) (0.91250,nan) (0.91500,nan) (0.91750,nan) (0.92000,nan) (0.92250,nan) (0.92500,nan) (0.92750,nan) (0.93000,nan) (0.93250,nan) (0.93500,nan) (0.93750,nan) (0.94000,nan) (0.94250,nan) (0.94500,nan) (0.94750,nan) (0.95000,nan) (0.95250,nan) (0.95500,nan) (0.95750,nan) (0.96000,nan) (0.96250,nan) (0.96500,nan) (0.96750,nan) (0.97000,nan) (0.97250,nan) (0.97500,nan) (0.97750,nan) (0.98000,nan) (0.98250,nan) (0.98500,nan) (0.98750,nan) (0.99000,nan) (0.99250,nan) (0.99500,nan) (0.99750,nan) (1.00000,nan) (1.00250,nan) (1.00500,nan) (1.00750,nan) (1.01000,nan) (1.01250,nan) (1.01500,nan) (1.01750,nan) (1.02000,nan) (1.02250,nan) (1.02500,nan) (1.02750,nan) (1.03000,nan) (1.03250,nan) (1.03500,nan) (1.03750,nan) (1.04000,nan) (1.04250,nan) (1.04500,nan) (1.04750,nan) (1.05000,nan) (1.05250,nan) (1.05500,nan) (1.05750,nan) (1.06000,nan) (1.06250,nan) (1.06500,nan) (1.06750,nan) (1.07000,nan) (1.07250,nan) (1.07500,nan) (1.07750,nan) (1.08000,nan) (1.08250,nan) (1.08500,nan) (1.08750,nan) (1.09000,nan) (1.09250,nan) (1.09500,nan) (1.09750,nan) (1.10000,nan) (1.10250,nan) (1.10500,nan) (1.10750,nan) (1.11000,nan) (1.11250,nan) (1.11500,nan) (1.11750,nan) (1.12000,nan) (1.12250,nan) (1.12500,nan) (1.12750,nan) (1.13000,nan) (1.13250,nan) (1.13500,nan) (1.13750,nan) (1.14000,nan) (1.14250,nan) (1.14500,nan) (1.14750,nan) (1.15000,nan) (1.15250,nan) (1.15500,nan) (1.15750,nan) (1.16000,nan) (1.16250,nan) (1.16500,nan) (1.16750,nan) (1.17000,nan) (1.17250,nan) (1.17500,nan) (1.17750,nan) (1.18000,nan) (1.18250,nan) (1.18500,nan) (1.18750,nan) (1.19000,nan) (1.19250,nan) (1.19500,nan) (1.19750,nan) (1.20000,nan) (1.20250,nan) (1.20500,nan) (1.20750,nan) (1.21000,nan) (1.21250,nan) (1.21500,nan) (1.21750,nan) (1.22000,nan) (1.22250,nan) (1.22500,nan) (1.22750,nan) (1.23000,nan) (1.23250,nan) (1.23500,nan) (1.23750,nan) (1.24000,nan) (1.24250,nan) (1.24500,nan) (1.24750,nan) (1.25000,nan) (1.25250,nan) (1.25500,nan) (1.25750,nan) (1.26000,nan) (1.26250,nan) (1.26500,nan) (1.26750,nan) (1.27000,nan) (1.27250,nan) (1.27500,nan) (1.27750,nan) (1.28000,nan) (1.28250,nan) (1.28500,nan) (1.28750,nan) (1.29000,nan) (1.29250,nan) (1.29500,nan) (1.29750,nan) (1.30000,nan) (1.30250,nan) (1.30500,nan) (1.30750,nan) (1.31000,nan) (1.31250,nan) (1.31500,nan) (1.31750,nan) (1.32000,nan) (1.32250,nan) (1.32500,nan) (1.32750,nan) (1.33000,nan) (1.33250,nan) (1.33500,nan) (1.33750,nan) (1.34000,nan) (1.34250,nan) (1.34500,nan) (1.34750,nan) (1.35000,nan) (1.35250,nan) (1.35500,nan) (1.35750,nan) (1.36000,nan) (1.36250,nan) (1.36500,nan) (1.36750,nan) (1.37000,nan) (1.37250,nan) (1.37500,nan) (1.37750,nan) (1.38000,nan) (1.38250,nan) (1.38500,nan) (1.38750,nan) (1.39000,nan) (1.39250,nan) (1.39500,nan) (1.39750,nan) (1.40000,nan) (1.40250,nan) (1.40500,nan) (1.40750,nan) (1.41000,nan) (1.41250,nan) (1.41500,nan) (1.41750,nan) (1.42000,nan) (1.42250,nan) (1.42500,nan) (1.42750,nan) (1.43000,nan) (1.43250,nan) (1.43500,nan) (1.43750,nan) (1.44000,nan) (1.44250,nan) (1.44500,nan) (1.44750,nan) (1.45000,nan) (1.45250,nan) (1.45500,nan) (1.45750,nan) (1.46000,nan) (1.46250,nan) (1.46500,nan) (1.46750,nan) (1.47000,nan) (1.47250,nan) (1.47500,nan) (1.47750,nan) (1.48000,nan) (1.48250,nan) (1.48500,nan) (1.48750,nan) (1.49000,nan) (1.49250,nan) (1.49500,nan) (1.49750,nan) (1.50000,nan) (1.50250,nan) (1.50500,nan) (1.50750,nan) (1.51000,nan) (1.51250,nan) (1.51500,nan) (1.51750,nan) (1.52000,nan) (1.52250,nan) (1.52500,nan) (1.52750,nan) (1.53000,nan) (1.53250,nan) (1.53500,nan) (1.53750,nan) (1.54000,nan) (1.54250,nan) (1.54500,nan) (1.54750,nan) (1.55000,nan) (1.55250,nan) (1.55500,nan) (1.55750,nan) (1.56000,nan) (1.56250,nan) (1.56500,nan) (1.56750,nan) (1.57000,nan) (1.57250,nan) (1.57500,nan) (1.57750,nan) (1.58000,nan) (1.58250,nan) (1.58500,nan) (1.58750,nan) (1.59000,nan) (1.59250,nan) (1.59500,nan) (1.59750,nan) (1.60000,nan) (1.60250,nan) (1.60500,nan) (1.60750,nan) (1.61000,nan) (1.61250,nan) (1.61500,nan) (1.61750,nan) (1.62000,nan) (1.62250,nan) (1.62500,nan) (1.62750,nan) (1.63000,nan) (1.63250,nan) (1.63500,nan) (1.63750,nan) (1.64000,nan) (1.64250,nan) (1.64500,nan) (1.64750,nan) (1.65000,nan) (1.65250,nan) (1.65500,nan) (1.65750,nan) (1.66000,nan) (1.66250,nan) (1.66500,nan) (1.66750,nan) (1.67000,nan) (1.67250,nan) (1.67500,nan) (1.67750,nan) (1.68000,nan) (1.68250,nan) (1.68500,nan) (1.68750,nan) (1.69000,nan) (1.69250,nan) (1.69500,nan) (1.69750,nan) (1.70000,nan) (1.70250,nan) (1.70500,nan) (1.70750,nan) (1.71000,nan) (1.71250,nan) (1.71500,nan) (1.71750,nan) (1.72000,nan) (1.72250,nan) (1.72500,nan) (1.72750,nan) (1.73000,nan) (1.73250,nan) (1.73500,nan) (1.73750,nan) (1.74000,nan) (1.74250,nan) (1.74500,nan) (1.74750,nan) (1.75000,nan) (1.75250,nan) (1.75500,nan) (1.75750,nan) (1.76000,nan) (1.76250,nan) (1.76500,nan) (1.76750,nan) (1.77000,nan) (1.77250,nan) (1.77500,nan) (1.77750,nan) (1.78000,nan) (1.78250,nan) (1.78500,nan) (1.78750,nan) (1.79000,nan) (1.79250,nan) (1.79500,nan) (1.79750,nan) (1.80000,nan) (1.80250,nan) (1.80500,nan) (1.80750,nan) (1.81000,nan) (1.81250,nan) (1.81500,nan) (1.81750,nan) (1.82000,nan) (1.82250,nan) (1.82500,nan) (1.82750,nan) (1.83000,nan) (1.83250,nan) (1.83500,nan) (1.83750,nan) (1.84000,nan) (1.84250,nan) (1.84500,nan) (1.84750,nan) (1.85000,nan) (1.85250,nan) (1.85500,nan) (1.85750,nan) (1.86000,nan) (1.86250,nan) (1.86500,nan) (1.86750,nan) (1.87000,nan) (1.87250,nan) (1.87500,nan) (1.87750,nan) (1.88000,nan) (1.88250,nan) (1.88500,nan) (1.88750,nan) (1.89000,nan) (1.89250,nan) (1.89500,nan) (1.89750,nan) (1.90000,nan) (1.90250,nan) (1.90500,nan) (1.90750,nan) (1.91000,nan) (1.91250,nan) (1.91500,nan) (1.91750,nan) (1.92000,nan) (1.92250,nan) (1.92500,nan) (1.92750,nan) (1.93000,nan) (1.93250,nan) (1.93500,nan) (1.93750,nan) (1.94000,nan) (1.94250,nan) (1.94500,nan) (1.94750,nan) (1.95000,nan) (1.95250,nan) (1.95500,nan) (1.95750,nan) (1.96000,nan) (1.96250,nan) (1.96500,nan) (1.96750,nan) (1.97000,nan) (1.97250,nan) (1.97500,nan) (1.97750,nan) (1.98000,nan) (1.98250,nan) (1.98500,nan) (1.98750,nan) (1.99000,nan) (1.99250,nan) (1.99500,nan) (1.99750,nan) (2.00000,nan) (2.00250,nan) (2.00500,nan) (2.00750,nan) (2.01000,nan) (2.01250,nan) (2.01500,nan) (2.01750,nan) (2.02000,nan) (2.02250,nan) (2.02500,nan) (2.02750,nan) (2.03000,nan) (2.03250,nan) (2.03500,nan) (2.03750,nan) (2.04000,nan) (2.04250,nan) (2.04500,nan) (2.04750,nan) (2.05000,nan) (2.05250,nan) (2.05500,nan) (2.05750,nan) (2.06000,nan) (2.06250,nan) (2.06500,nan) (2.06750,nan) (2.07000,nan) (2.07250,nan) (2.07500,nan) (2.07750,nan) (2.08000,nan) (2.08250,nan) (2.08500,nan) (2.08750,nan) (2.09000,nan) (2.09250,nan) (2.09500,nan) (2.09750,nan) (2.10000,nan) (2.10250,nan) (2.10500,nan) (2.10750,nan) (2.11000,nan) (2.11250,nan) (2.11500,nan) (2.11750,nan) (2.12000,nan) (2.12250,nan) (2.12500,nan) (2.12750,nan) (2.13000,nan) (2.13250,nan) (2.13500,nan) (2.13750,nan) (2.14000,nan) (2.14250,nan) (2.14500,nan) (2.14750,nan) (2.15000,nan) (2.15250,nan) (2.15500,nan) (2.15750,nan) (2.16000,nan) (2.16250,nan) (2.16500,nan) (2.16750,nan) (2.17000,nan) (2.17250,nan) (2.17500,nan) (2.17750,nan) (2.18000,nan) (2.18250,nan) (2.18500,nan) (2.18750,nan) (2.19000,nan) (2.19250,nan) (2.19500,nan) (2.19750,nan) (2.20000,nan) (2.20250,nan) (2.20500,nan) (2.20750,nan) (2.21000,nan) (2.21250,nan) (2.21500,nan) (2.21750,nan) (2.22000,nan) (2.22250,nan) (2.22500,nan) (2.22750,nan) (2.23000,nan) (2.23250,nan) (2.23500,nan) (2.23750,nan) (2.24000,nan) (2.24250,nan) (2.24500,nan) (2.24750,nan) (2.25000,nan) (2.25250,nan) (2.25500,nan) (2.25750,nan) (2.26000,nan) (2.26250,nan) (2.26500,nan) (2.26750,nan) (2.27000,nan) (2.27250,nan) (2.27500,nan) (2.27750,nan) (2.28000,nan) (2.28250,nan) (2.28500,nan) (2.28750,nan) (2.29000,nan) (2.29250,nan) (2.29500,nan) (2.29750,nan) (2.30000,nan) (2.30250,nan) (2.30500,nan) (2.30750,nan) (2.31000,nan) (2.31250,nan) (2.31500,nan) (2.31750,nan) (2.32000,nan) (2.32250,nan) (2.32500,nan) (2.32750,nan) (2.33000,nan) (2.33250,nan) (2.33500,nan) (2.33750,nan) (2.34000,nan) (2.34250,nan) (2.34500,nan) (2.34750,nan) (2.35000,nan) (2.35250,nan) (2.35500,nan) (2.35750,nan) (2.36000,nan) (2.36250,nan) (2.36500,nan) (2.36750,nan) (2.37000,nan) (2.37250,nan) (2.37500,nan) (2.37750,nan) (2.38000,nan) (2.38250,nan) (2.38500,nan) (2.38750,nan) (2.39000,nan) (2.39250,nan) (2.39500,nan) (2.39750,nan) (2.40000,nan) (2.40250,nan) (2.40500,nan) (2.40750,nan) (2.41000,nan) (2.41250,nan) (2.41500,nan) (2.41750,nan) (2.42000,nan) (2.42250,nan) (2.42500,nan) (2.42750,nan) (2.43000,nan) (2.43250,nan) (2.43500,nan) (2.43750,nan) (2.44000,nan) (2.44250,nan) (2.44500,nan) (2.44750,nan) (2.45000,nan) (2.45250,nan) (2.45500,nan) (2.45750,nan) (2.46000,nan) (2.46250,nan) (2.46500,nan) (2.46750,nan) (2.47000,nan) (2.47250,nan) (2.47500,nan) (2.47750,nan) (2.48000,nan) (2.48250,nan) (2.48500,nan) (2.48750,nan) (2.49000,nan) (2.49250,nan) (2.49500,nan) (2.49750,nan) (2.50000,nan) (2.50250,nan) (2.50500,nan) (2.50750,nan) (2.51000,nan) (2.51250,nan) (2.51500,nan) (2.51750,nan) (2.52000,nan) (2.52250,nan) (2.52500,nan) (2.52750,nan) (2.53000,nan) (2.53250,nan) (2.53500,nan) (2.53750,nan) (2.54000,nan) (2.54250,nan) (2.54500,nan) (2.54750,nan) (2.55000,nan) (2.55250,nan) (2.55500,nan) (2.55750,nan) (2.56000,nan) (2.56250,nan) (2.56500,nan) (2.56750,nan) (2.57000,nan) (2.57250,nan) (2.57500,nan) (2.57750,nan) (2.58000,nan) (2.58250,nan) (2.58500,nan) (2.58750,nan) (2.59000,nan) (2.59250,nan) (2.59500,nan) (2.59750,nan) (2.60000,nan) (2.60250,nan) (2.60500,nan) (2.60750,nan) (2.61000,nan) (2.61250,nan) (2.61500,nan) (2.61750,nan) (2.62000,nan) (2.62250,nan) (2.62500,nan) (2.62750,nan) (2.63000,nan) (2.63250,nan) (2.63500,nan) (2.63750,nan) (2.64000,nan) (2.64250,nan) (2.64500,nan) (2.64750,nan) (2.65000,nan) (2.65250,nan) (2.65500,nan) (2.65750,nan) (2.66000,nan) (2.66250,nan) (2.66500,nan) (2.66750,nan) (2.67000,nan) (2.67250,nan) (2.67500,nan) (2.67750,nan) (2.68000,nan) (2.68250,nan) (2.68500,nan) (2.68750,nan) (2.69000,nan) (2.69250,nan) (2.69500,nan) (2.69750,nan) (2.70000,nan) (2.70250,nan) (2.70500,nan) (2.70750,nan) (2.71000,nan) (2.71250,nan) (2.71500,nan) (2.71750,nan) (2.72000,nan) (2.72250,nan) (2.72500,nan) (2.72750,nan) (2.73000,nan) (2.73250,nan) (2.73500,nan) (2.73750,nan) (2.74000,nan) (2.74250,nan) (2.74500,nan) (2.74750,nan) (2.75000,nan) (2.75250,nan) (2.75500,nan) (2.75750,nan) (2.76000,nan) (2.76250,nan) (2.76500,nan) (2.76750,nan) (2.77000,nan) (2.77250,nan) (2.77500,nan) (2.77750,nan) (2.78000,nan) (2.78250,nan) (2.78500,nan) (2.78750,nan) (2.79000,nan) (2.79250,nan) (2.79500,nan) (2.79750,nan) (2.80000,nan) (2.80250,nan) (2.80500,nan) (2.80750,nan) (2.81000,nan) (2.81250,nan) (2.81500,nan) (2.81750,nan) (2.82000,nan) (2.82250,nan) (2.82500,nan) (2.82750,nan) (2.83000,nan) (2.83250,nan) (2.83500,nan) (2.83750,nan) (2.84000,nan) (2.84250,nan) (2.84500,nan) (2.84750,nan) (2.85000,nan) (2.85250,nan) (2.85500,nan) (2.85750,nan) (2.86000,nan) (2.86250,nan) (2.86500,nan) (2.86750,nan) (2.87000,nan) (2.87250,nan) (2.87500,nan) (2.87750,nan) (2.88000,nan) (2.88250,nan) (2.88500,nan) (2.88750,nan) (2.89000,nan) (2.89250,nan) (2.89500,nan) (2.89750,nan) (2.90000,nan) (2.90250,nan) (2.90500,nan) (2.90750,nan) (2.91000,nan) (2.91250,nan) (2.91500,nan) (2.91750,nan) (2.92000,nan) (2.92250,nan) (2.92500,nan) (2.92750,nan) (2.93000,nan) (2.93250,nan) (2.93500,nan) (2.93750,nan) (2.94000,nan) (2.94250,nan) (2.94500,nan) (2.94750,nan) (2.95000,nan) (2.95250,nan) (2.95500,nan) (2.95750,nan) (2.96000,nan) (2.96250,nan) (2.96500,nan) (2.96750,nan) (2.97000,nan) (2.97250,nan) (2.97500,nan) (2.97750,nan) (2.98000,nan) (2.98250,nan) (2.98500,nan) (2.98750,nan) (2.99000,nan) (2.99250,nan) (2.99500,nan) (2.99750,nan) (3.00000,nan) (3.00250,nan) (3.00500,nan) (3.00750,nan) (3.01000,nan) (3.01250,nan) (3.01500,nan) (3.01750,nan) (3.02000,nan) (3.02250,nan) (3.02500,nan) (3.02750,nan) (3.03000,nan) (3.03250,nan) (3.03500,nan) (3.03750,nan) (3.04000,nan) (3.04250,nan) (3.04500,nan) (3.04750,nan) (3.05000,nan) (3.05250,nan) (3.05500,nan) (3.05750,nan) (3.06000,nan) (3.06250,nan) (3.06500,nan) (3.06750,nan) (3.07000,nan) (3.07250,nan) (3.07500,nan) (3.07750,nan) (3.08000,nan) (3.08250,nan) (3.08500,nan) (3.08750,nan) (3.09000,nan) (3.09250,nan) (3.09500,nan) (3.09750,nan) (3.10000,nan) (3.10250,nan) (3.10500,nan) (3.10750,nan) (3.11000,nan) (3.11250,nan) (3.11500,nan) (3.11750,nan) (3.12000,nan) (3.12250,nan) (3.12500,nan) (3.12750,nan) (3.13000,nan) (3.13250,nan) (3.13500,nan) (3.13750,nan) (3.14000,nan) (3.14250,nan) (3.14500,nan) (3.14750,nan) (3.15000,nan) (3.15250,nan) (3.15500,nan) (3.15750,nan) (3.16000,nan) (3.16250,nan) (3.16500,nan) (3.16750,nan) (3.17000,nan) (3.17250,nan) (3.17500,nan) (3.17750,nan) (3.18000,nan) (3.18250,nan) (3.18500,nan) (3.18750,nan) (3.19000,nan) (3.19250,nan) (3.19500,nan) (3.19750,nan) (3.20000,nan) (3.20250,nan) (3.20500,nan) (3.20750,nan) (3.21000,nan) (3.21250,nan) (3.21500,nan) (3.21750,nan) (3.22000,nan) (3.22250,nan) (3.22500,nan) (3.22750,nan) (3.23000,nan) (3.23250,nan) (3.23500,nan) (3.23750,nan) (3.24000,nan) (3.24250,nan) (3.24500,nan) (3.24750,nan) (3.25000,nan) (3.25250,nan) (3.25500,nan) (3.25750,nan) (3.26000,nan) (3.26250,nan) (3.26500,nan) (3.26750,nan) (3.27000,nan) (3.27250,nan) (3.27500,nan) (3.27750,nan) (3.28000,nan) (3.28250,nan) (3.28500,nan) (3.28750,nan) (3.29000,nan) (3.29250,nan) (3.29500,nan) (3.29750,nan) (3.30000,nan) (3.30250,nan) (3.30500,nan) (3.30750,nan) (3.31000,nan) (3.31250,nan) (3.31500,nan) (3.31750,nan) (3.32000,nan) (3.32250,nan) (3.32500,nan) (3.32750,nan) (3.33000,nan) (3.33250,nan) (3.33500,nan) (3.33750,nan) (3.34000,nan) (3.34250,nan) (3.34500,nan) (3.34750,nan) (3.35000,nan) (3.35250,nan) (3.35500,nan) (3.35750,nan) (3.36000,nan) (3.36250,nan) (3.36500,nan) (3.36750,nan) (3.37000,nan) (3.37250,nan) (3.37500,nan) (3.37750,nan) (3.38000,nan) (3.38250,nan) (3.38500,nan) (3.38750,nan) (3.39000,nan) (3.39250,nan) (3.39500,nan) (3.39750,nan) (3.40000,nan) (3.40250,nan) (3.40500,nan) (3.40750,nan) (3.41000,nan) (3.41250,nan) (3.41500,nan) (3.41750,nan) (3.42000,nan) (3.42250,nan) (3.42500,nan) (3.42750,nan) (3.43000,nan) (3.43250,nan) (3.43500,nan) (3.43750,nan) (3.44000,nan) (3.44250,nan) (3.44500,nan) (3.44750,nan) (3.45000,nan) (3.45250,nan) (3.45500,nan) (3.45750,nan) (3.46000,nan) (3.46250,nan) (3.46500,nan) (3.46750,nan) (3.47000,nan) (3.47250,nan) (3.47500,nan) (3.47750,nan) (3.48000,nan) (3.48250,nan) (3.48500,nan) (3.48750,nan) (3.49000,nan) (3.49250,nan) (3.49500,nan) (3.49750,nan) (3.50000,nan) (3.50250,nan) (3.50500,nan) (3.50750,nan) (3.51000,nan) (3.51250,nan) (3.51500,nan) (3.51750,nan) (3.52000,-1.8111666) (3.52250,-1.4920232) (3.52500,-1.3112689) (3.52750,-1.1849022) (3.53000,-1.0878291) (3.53250,-1.0091022) (3.53500,-0.9429519) (3.53750,-0.8859613) (3.54000,-0.8359389) (3.54250,-0.7913947) (3.54500,-0.7512690) (3.54750,-0.7147816) (3.55000,-0.6813410) (3.55250,-0.6504883) (3.55500,-0.6218600) (3.55750,-0.5951634) (3.56000,-0.5701591) (3.56250,-0.5466490) (3.56500,-0.5244672) (3.56750,-0.5034731) (3.57000,-0.4835470) (3.57250,-0.4645855) (3.57500,-0.4464991) (3.57750,-0.4292096) (3.58000,-0.4126483) (3.58250,-0.3967544) (3.58500,-0.3814740) (3.58750,-0.3667589) (3.59000,-0.3525656) (3.59250,-0.3388553) (3.59500,-0.3255925) (3.59750,-0.3127451) (3.60000,-0.3002839) (3.60250,-0.2881820) (3.60500,-0.2764146) (3.60750,-0.2649590) (3.61000,-0.2537941) (3.61250,-0.2429002) (3.61500,-0.2322590) (3.61750,-0.2218534) (3.62000,-0.2116674) (3.62250,-0.2016857) (3.62500,-0.1918940) (3.62750,-0.1822788) (3.63000,-0.1728269) (3.63250,-0.1635261) (3.63500,-0.1543644) (3.63750,-0.1453303) (3.64000,-0.1364127) (3.64250,-0.1276008) (3.64500,-0.1188840) (3.64750,-0.1102520) (3.65000,-0.1016945) (3.65250,-0.0932016) (3.65500,-0.0847631) (3.65750,-0.0763690) (3.66000,-0.0680093) (3.66250,-0.0596738) (3.66500,-0.0513522) (3.66750,-0.0430340) (3.67000,-0.0347085) (3.67250,-0.0263646) (3.67500,-0.0179909) (3.67750,-0.0095755) (3.68000,-0.0011061) (3.68250,0.0074305) (3.68500,0.0160478) (3.68750,0.0247605) (3.69000,0.0335839) (3.69250,0.0425345) (3.69500,0.0516300) (3.69750,0.0608896) (3.70000,0.0703338) (3.70250,0.0799851) (3.70500,0.0898681) (3.70750,0.1000096) (3.71000,0.1104394) (3.71250,0.1211903) (3.71500,0.1322988) (3.71750,0.1438058) (3.72000,0.1557572) (3.72250,0.1682047) (3.72500,0.1812074) (3.72750,0.1948324) (3.73000,0.2091571) (3.73250,0.2242710) (3.73500,0.2402785) (3.73750,0.2573024) (3.74000,0.2754881) (3.74250,0.2950099) (3.74500,0.3160787) (3.74750,0.3389530) (3.75000,0.3639537) (3.75250,0.3914863) (3.75500,0.4220720) (3.75750,0.4563963) (3.76000,0.4953840) (3.76250,0.5403247) (3.76500,0.5930896) (3.76750,0.6565403) (3.77000,0.7353712) (3.77250,0.8380945) (3.77500,0.9827445) (3.77750,1.2200318) (3.78000,1.8716072) (3.78250,nan) (3.78500,nan) (3.78750,nan) (3.79000,nan) (3.79250,nan) (3.79500,nan) (3.79750,nan) (3.80000,nan) (3.80250,nan) (3.80500,nan) (3.80750,nan) (3.81000,nan) (3.81250,nan) (3.81500,nan) (3.81750,nan) (3.82000,nan) (3.82250,nan) (3.82500,nan) (3.82750,nan) (3.83000,nan) (3.83250,nan) (3.83500,nan) (3.83750,nan) (3.84000,nan) (3.84250,nan) (3.84500,-1.5416231) (3.84750,-1.2122777) (3.85000,-1.0372602) (3.85250,-0.9198107) (3.85500,-0.8325903) (3.85750,-0.7639595) (3.86000,-0.7078765) (3.86250,-0.6608038) (3.86500,-0.6204940) (3.86750,-0.5854321) (3.87000,-0.5545497) (3.87250,-0.5270662) (3.87500,-0.5023944) (3.87750,-0.4800827) (3.88000,-0.4597759) (3.88250,-0.4411906) (3.88500,-0.4240971) (3.88750,-0.4083067) (3.89000,-0.3936629) (3.89250,-0.3800343) (3.89500,-0.3673101) (3.89750,-0.3553954) (3.90000,-0.3442091) (3.90250,-0.3336808) (3.90500,-0.3237494) (3.90750,-0.3143614) (3.91000,-0.3054700) (3.91250,-0.2970337) (3.91500,-0.2890156) (3.91750,-0.2813830) (3.92000,-0.2741067) (3.92250,-0.2671605) (3.92500,-0.2605206) (3.92750,-0.2541659) (3.93000,-0.2480771) (3.93250,-0.2422365) (3.93500,-0.2366283) (3.93750,-0.2312380) (3.94000,-0.2260522) (3.94250,-0.2210588) (3.94500,-0.2162466) (3.94750,-0.2116052) (3.95000,-0.2071253) (3.95250,-0.2027979) (3.95500,-0.1986151) (3.95750,-0.1945693) (3.96000,-0.1906535) (3.96250,-0.1868613) (3.96500,-0.1831865) (3.96750,-0.1796236) (3.97000,-0.1761672) (3.97250,-0.1728125) (3.97500,-0.1695549) (3.97750,-0.1663900) (3.98000,-0.1633139) (3.98250,-0.1603226) (3.98500,-0.1574127) (3.98750,-0.1545806) (3.99000,-0.1518234) (3.99250,-0.1491380) (3.99500,-0.1465214) (3.99750,-0.1439712) (4.00000,-0.1414847)};\addlegendentry{X}
\end{axis}\end{tikzpicture}\end{center}

样本数每密度每支2401点, 先检验符号再取log, 在非正值或出视野处以nan分段. 精确边界由解析零点方程定义, 图上采样误差不作为边界证明. 可重现脚本位于\texttt{tools/cold\_wave\_plots.py}.
\end{solution}
\begin{exercise}\label{ex:06-8-2}
原题2. 平面波从真空入射x向密度坡, y/z均匀、$n=n_0x/\ell$, k\ensuremath{\ell}\ensuremath{\gg}1, Bz=B0, $\omega_{pe}(n_0)=\Omega_{ce}$, 只电子响应. (a)在x/\ensuremath{\ell}\ensuremath{\le}4画O传播区; (b)画X传播区; (c)\ensuremath{\omega}<\ensuremath{\Omega}ce时X最大传播深度.
\end{exercise}
\begin{solution}
令u=x/\ensuremath{\ell}、w=\ensuremath{\omega}/\ensuremath{\Omega}ce>0. O:P=1\ensuremath{-}u/w\ensuremath{{}^{2}}, 传播0\ensuremath{\le}u<w\ensuremath{{}^{2}}. X按章\ref{ch:cold}:
$R=1-u/[w(w-1)]$, $L=1-u/[w(w+1)]$, $S=1-u/(w^2-1)$, $n_X^2=RL/S$. w<1只有uL=w(w+1)正截止, 0\ensuremath{\le}u<uL传播, 所以最大深度上确界2\ensuremath{\ell}在w\ensuremath{\to}1\ensuremath{-}. w>1的顺序uR=w(w\ensuremath{-}1)<uH=w\ensuremath{{}^{2}}\ensuremath{-}1<uL=w(w+1), 两传播带为u<uR与uH<u<uL. 真空来波先在uR截止, 第二带需隧穿或模转换才能到达.
\begin{center}\begin{tikzpicture}\begin{axis}[width=.75\textwidth,height=5cm,xlabel={$u=x/\ell$},ylabel={$w=\omega/\Omega_{ce}$},xmin=0,xmax=4,ymin=0,ymax=3,legend style={font=\small,at={(1.02,1)},anchor=north west}]
\addplot[thick,structurecolor,domain=0:4,samples=70]{sqrt(x)};\addlegendentry{O截止}
\addplot[thick,dashed,domain=0:4,samples=70]{(1+sqrt(1+4*x))/2};\addlegendentry{R截止}
\addplot[thick,dotted,domain=0:4,samples=70]{sqrt(1+x)};\addlegendentry{X共振}
\addplot[thick,dashdotted,domain=0:4,samples=70]{(-1+sqrt(1+4*x))/2};\addlegendentry{L截止}
\end{axis}\end{tikzpicture}\end{center}
图线和上述不等式一起确定区域: O在线上方; X对固定w按u的两区间取值. w=1是冷电子回旋奇点, 不能直接在R分母代零当有限区.
\end{solution}
